\ifdefined\XeTeXversion\else\pdfoutput=1\fi
\documentclass[aps,prd,onecolumn,10pt,nofootinbib,floatfix]{revtex4-2}
\usepackage{amsmath,amssymb,mathtools,bm}
\usepackage{array}
\usepackage[dvipsnames]{xcolor}
\usepackage[colorlinks=true,linkcolor=blue!50!black,citecolor=blue!50!black,urlcolor=blue!50!black,hypertexnames=false]{hyperref}
\newcommand{\PI}{Paper~I}
\newcommand{\PII}{Paper~II}
\newcommand{\eqI}[1]{(I.#1)}
\newcommand{\eqII}[1]{(II.#1)}
\newcommand{\Tab}{\mathcal{T}}
\newcommand{\Dop}{\mathcal{D}}
\newcommand{\khat}{\hat k}
\newcommand{\pihat}{\hat\pi}
\newcommand{\dd}{\mathrm{d}}
\newcommand{\ren}{_{\rm ren}}
\newcommand{\reg}{_{\rm reg}}

\hypersetup{pdftitle={Quantum gravity corrections to a massless minimally coupled scalar during inflation revisited: the one-loop self-mass and the late-time mode functions},pdfauthor={Emre Onur Kahya}}
\newcommand{\TOKOUT}{$2.7\times10^{6}$}
\newcommand{\TOKTOT}{about $8\times10^{8}$}
\newcommand{\HOURS}{10}
\newcommand{\NDAYS}{three}
\begin{document}
\renewcommand{\thetable}{\arabic{table}}
\title{Auditing a quantum gravity calculation with an AI model: corrected scalar self-mass
and unchanged late-time mode functions during inflation}

\author{Emre Onur Kahya}
\email{eokahya@itu.edu.tr}
\affiliation{Department of Physics, Istanbul Technical University, Maslak 34469, Istanbul, Turkey}

\begin{abstract}
We present an independent, step-by-step recomputation of the one-loop quantum-gravitational self-mass-squared of a massless,
minimally coupled scalar on de Sitter background in a non-de Sitter-invariant gauge [E.~O.~Kahya and R.~P.~Woodard, Phys.\ Rev.\ D
\textbf{76}, 124005 (2007)], and of the Schwinger--Keldysh analysis of the late-time scalar mode functions based on it [Phys.\ Rev.\ D
\textbf{77}, 084012 (2008)]. Each main-text step was re-derived from its inputs, in general $D$ where the papers keep it
general, and cross-checked by independent code; the appendix of the first paper on extracting derivatives was
audited (worked examples and 30-digit tests of its totals) rather than re-derived. Most results are confirmed, but
(i) the printed self-mass-squared is $-1$ times the one defined by the papers' own effective field equation, with both the
4-point and 3-point parts reversed; (ii) the final table of nonlocal terms needs two corrections (a missing non-analytic term, and
a constant that drops out of the retarded kernel), established by a 30-digit global pointwise test; (iii) two finite errors shift
counterterm coefficients by absorbable amounts. We reproduce Tables~2--5 of the second paper exactly (Table~5 in its own sign
convention; in the definition sign every entry flips, the total remaining zero) and, with the corrected self-mass, redo its late-time
analysis through order $k^{10}$ without its simplifying approximation, whose justification fails term by term. The
non-falling part of the source cancels order by order, and the mode function tends to its tree-order value times a constant: one-loop quantum gravity gives no
significant late-time correction to the scalar mode functions. Limitations: the global tests are numerical, the cancellation is not
proven beyond $k^{10}$, and a possible extra local term is only partially excluded. The work used an AI model (Claude Opus 5.5); notes
and code are public.
\end{abstract}

\maketitle

\section{Introduction}
\label{sec:s0s1s7s8:intro}

\subsection{Physical motivation}

Massless, minimally coupled scalars and gravitons are unique in possessing zero mass without classical conformal invariance.
Simple arguments based on the energy-time uncertainty principle suggest that quantum effects from these particles are strongly
enhanced during primordial inflation~\cite{ref:Woodard2004}. At lowest order inflation rips long-wavelength virtual quanta of these
fields out of the vacuum; this is the origin of the primordial scalar~\cite{ref:Mukhanov1981} and tensor~\cite{ref:Starobinsky1979}
perturbations. At higher orders the sea of long-wavelength scalars and gravitons interacts with itself and with other fields, and loop
corrections can be enhanced by powers of the logarithm of the scale factor~\cite{ref:Weinberg2005}. Gravitons couple through
derivatives, which weakens the enhancement, but their effects need not be small: for massless fermions, one-loop quantum gravity
induces a time-dependent field-strength renormalization that eventually becomes nonperturbatively large~\cite{ref:Miao2006}.

It is therefore natural to ask whether the inflationary production of gravitons has a significant effect on the propagation of a
massless, minimally coupled scalar. The question has phenomenological interest. The inflaton potential of scalar-driven inflation is so
flat that the inflaton mode functions are effectively those of a massless, minimally coupled scalar, and the scalar power spectrum is
routinely computed by setting the scalar-scalar correlator equal to its value for such a scalar on de Sitter
background~\cite{ref:Mukhanov1992,ref:Liddle1993}. A significant one-loop correction to the scalar mode functions, in particular a
secularly growing field-strength renormalization of the kind found for fermions, would invalidate this procedure.

\subsection{What Papers I and II did}

The question was addressed in two papers by R.~P.~Woodard and the present author. In Paper~I~\cite{KW07} the one-loop
self-mass-squared $-iM^2(x;x')$ of a massless, minimally coupled scalar was computed on the conformal coordinate patch of
$D$-dimensional de Sitter space, $ds^2=a^2(\eta)(-d\eta^2+d\vec x\cdot d\vec x)$ with $a=-1/(H\eta)$. The graviton propagator was
taken in the non-de Sitter-invariant gauge \eqI{15}, in which it is a sum of constant tensor factors times three scalar
propagators~\cite{ref:Tsamis1994,ref:Woodard2006}. The calculation used dimensional regularization. The 4-point vertices give a purely
local contribution. The products of two 3-point vertices were split into local, logarithm and ``normal'' parts; the normal parts
were made integrable by extracting derivatives (the Appendix of Paper~I), and their remaining divergences were separated into local
terms. These were removed by fourth-order BPHZ counterterms~\cite{ref:Bogoliubov1957}: two de Sitter-invariant ones and one
noninvariant one, which is needed because the gauge breaks de Sitter invariance. The final result \eqI{167} consists of arbitrary finite counterterm
contributions plus a table of thirteen finite nonlocal terms, Table~X of Paper~I.

In Paper~II~\cite{KW08} this self-mass-squared was used to correct the linearized scalar field equation. To obtain a real and causal
equation the Schwinger--Keldysh formalism~\cite{ref:Schwinger1961,ref:Keldysh1965,ref:Jordan1986,ref:Calzetta1987} was used. Writing
$\Phi=u+\kappa^2\Phi_1$, with $u(\eta,k)$ the tree-order mode function, the one-loop correction $\Phi_1$ obeys a sourced mode
equation, \eqII{29}. Its source was evaluated at late times under a simplification, \eqII{44}: $u(\eta',k)\to u(0,k)$ and
$e^{i\vec k\cdot(\vec x\,'-\vec x)}\to1$, which removes every term of the table containing a spatial Laplacian. Several of the nineteen
remaining contributions (Table~5 of Paper~II) grow individually like $\ln^2a$ or $\ln a$, but their sum vanishes. Paper~II concluded that the plane-wave
mode functions receive no significant one-loop corrections at late times.

\subsection{Why a recomputation}

Both calculations are long. Paper~I runs through about one hundred and eighty numbered equations in the main text, thirty-eight in its
appendix, and twenty tables, and the coefficients of its final table emerge from many cancellations between separately computed
pieces. The late-time conclusion of Paper~II rests on an exact cancellation among individually growing terms, obtained under an
approximation whose validity was argued rather than demonstrated. An independent, step-by-step recomputation is the most reliable way
to establish which parts of these results are correct, and to quantify the effect of any error on the physical conclusion. The present
work provides such a recomputation, with the exceptions and limitations stated in Secs.~\ref{sec:s0s1s7s8:method}
and~\ref{sec:s0s1s7s8:limitations}, and presents the corrected result in one place.

\subsection{Method}
\label{sec:s0s1s7s8:method}

The recomputation was organized as eighteen consecutive tasks, each covering a block of equations of Paper~I or Paper~II
(our Table~\ref{tab:s0s1s7s8:map}). The following rules were applied throughout.
\begin{itemize}
\item Every equation of the main text of Paper~I and of Paper~II was re-derived from its inputs: the Lagrangian, the propagators, and
results established in earlier tasks. The printed target was not used as input; it was compared with only at the end of each step.
The stated exceptions are the following. (a) The gauge-fixed graviton propagator \eqI{15}--\eqI{25} was taken from Paper~I and
Ref.~\cite{ref:Tsamis1994} as input and was not re-derived. (b) The appendix ``Extracting derivatives'' of Paper~I, \eqI{A1}--\eqI{A38}
and Tables~XI--XX, was audited rather than redone: worked examples were checked by hand, Tables~XII and~XIV were transcribed, and the
totals of all ten tables were verified by global pointwise tests (Sec.~\ref{sec:s4b:appendix}). (c) Tables~IV and~V of Paper~I were
computed from these printed appendix totals. In the assembly of Table~X, Table~V enters as printed (apart from one analytic
constant corrected after the global test, Sec.~\ref{sec:s4b:tableV}), in particular its polynomial terms, after its non-analytic terms had been confirmed. Tables~2--5 of Paper~II were never used as input.
\item The spacetime dimension $D$ was kept general wherever Paper~I keeps it general; $D=4$ was set only in terms that are not
multiplied by a pole.
\item Every analytic step was cross-checked by independent code where possible (Python with \texttt{sympy} and \texttt{mpmath}), kept
separate from the derivation: symbolic identities in general $D$ where possible, and numerical tests at random points in 30-digit
arithmetic otherwise. The few steps done by hand only are listed in Sec.~\ref{sec:s0s1s7s8:limitations}.
\item Every discrepancy with the printed papers was classified as typographical, finite and absorbable into the counterterms,
substantive, or a matter of convention.
\end{itemize}
The step-by-step re-derivations and the code checks were carried out with an AI model, Claude Opus~5.5 (Anthropic), used as a
calculation and checking tool. The original working reports were written in Turkish; they were translated into English derivation notes
(nine chapters and a summary chapter of findings), which, together with every script and its expected output, are publicly available~\cite{ref:repo}. The present paper
gives the derivation in condensed form; the notes give the intermediate steps in full.

\subsection{Conventions}
\label{sec:s0s1s7s8:conventions}

Equation $(n)$ of Paper~I is cited as \eqI{n}, equation $(n)$ of Paper~II as \eqII{n}. Tables of Paper~I carry Roman numerals
(Table~X), tables of Paper~II Arabic numerals (Table~5). Equations of the appendix of Paper~I are cited as \eqI{A1}, etc.; equation
and table numbers of Paper~II refer to its arXiv version, arXiv:0710.5282v1. Tables of the present paper are also numbered with Arabic numerals,
but they are always cited as ``our Table''; an unqualified Roman table number refers to Paper~I\@. We follow the notation of Paper~I: $g_{\mu\nu}=a^2(\eta_{\mu\nu}+\kappa
h_{\mu\nu})$ with $\kappa^2=16\pi G$, and the de Sitter length function
\begin{equation}
y(x;x')=aa'H^2\Big[\|\vec x-\vec x\,'\|^2-\big(|\eta-\eta'|-i\delta\big)^2\Big],\qquad a\equiv a(\eta),\quad a'\equiv a(\eta'),\qquad
x\equiv\frac y4 ,
\label{eq:s0s1s7s8:y}
\end{equation}
where the variable $x=y/4$ appears only as the argument of the coefficient functions of the tables. Paper~II uses the same $y$ with
$H=1$.

\paragraph*{Two sign conventions.}
Paper~I defines the self-mass-squared through the linearized effective field equation \eqI{168},
\begin{equation}
\partial_\mu\big(\sqrt{-\hat g}\,\hat g^{\mu\nu}\partial_\nu\Phi\big)(x)-\int d^4x'\,M^2(x;x')\,\Phi(x')=0,
\label{eq:s0s1s7s8:def}
\end{equation}
where $\hat g_{\mu\nu}=a^2\eta_{\mu\nu}$ is the background metric. Paper~I prints only this field equation. Writing the
one-particle-irreducible effective action as $\Gamma=S_0+\Gamma_1$, with $S_0$ the tree action of the background field and $\Gamma_1$ its
quantum part, Eq.~\eqref{eq:s0s1s7s8:def} is $\delta\Gamma/\delta\Phi=0$ with
\begin{equation}
-iM^2(x;x')=\frac{\delta^2\,(i\Gamma_1[\Phi])}{\delta\Phi(x)\,\delta\Phi(x')}\Big|_{\Phi=0}
\label{eq:s0s1s7s8:defM}
\end{equation}
(Sec.~\ref{sec:s2:definition}). We call the sign fixed by Eqs.~\eqref{eq:s0s1s7s8:def}--\eqref{eq:s0s1s7s8:defM} the
\emph{definition sign}. The sign of the self-mass
contributions as printed in Papers~I and~II is called the \emph{paper convention} (or Table-X convention); our
Table~\ref{tab:s2:conventions} lists which printed equations are in which convention. One of our results is
\begin{equation}
\text{definition sign}=-1\times\text{paper convention}.
\label{eq:s0s1s7s8:signs}
\end{equation}
Corrected tables are quoted in the paper convention, so that they can be compared entry by entry with the printed ones, and the
conversion is stated explicitly. Physical statements are made in the definition sign.

\paragraph*{Units.}
The entries of Table~X of Paper~I and of Table~1 of Paper~II are coefficients of $\kappa^2H^4/(4\pi)^4$. The counterterm coefficients
$\alpha_i$ and $\beta_i$ are quoted in units of $H^{D-4}/(4\pi)^{D/2}$. We write $\hat k\equiv k/H$, and Euler's constant as
$\gamma_{\rm E}$ (it is $\gamma$ in Papers~I and~II), to avoid a clash with the operators $\gamma_i$ of Paper~I.

\paragraph*{Counting of $\kappa^2$.}
As in \eqII{23} we write $\Phi=u+\kappa^2\Phi_1$, where $\Phi_1$ is sourced by the one-loop self-mass-squared $M_1^2$ with the factor
$\kappa^2$ stripped off,
\begin{equation}
\Big[\partial_t^2+3H\partial_t+\frac{k^2}{a^2}\Big]\Phi_1=\mathcal R_1\equiv-\frac1{a^4}\int_{\eta_i}^0\!d\eta'\,u(\eta',k)
\int\!d^3x'\,\big\{M^2_{++}(x;x')+M^2_{+-}(x;x')\big\}_1\,e^{i\vec k\cdot(\vec x\,'-\vec x)} .
\label{eq:s0s1s7s8:mode}
\end{equation}
The factor $\kappa^2$ is thus counted once. The source $\mathcal R_1$ is quoted in units of $H^4u(0,k)/16\pi^2$.

\subsection{Organization of the paper}

The paper follows the order of Papers~I and~II, so that it can be read side by side with them. Our Table~\ref{tab:s0s1s7s8:map} maps each
of our sections to the corresponding sections, equations and tables of the printed papers, and to the chapters of the public
derivation notes~\cite{ref:repo}. Within each section we present our derivation as the main line, with explicit pointers (``cf.\
\eqI{n}'') to the printed equations, and compare with the printed values in small tables or footnotes.

\begin{table}[tb]
\caption{Correspondence between the sections of this paper, the printed Papers~I~\cite{KW07} and~II~\cite{KW08}, and the chapters of
the derivation notes~\cite{ref:repo}.}
\label{tab:s0s1s7s8:map}
\renewcommand{\arraystretch}{1.25}
\small
\begin{ruledtabular}
\begin{tabular}{@{}>{\raggedright\arraybackslash}p{0.07\linewidth}>{\raggedright\arraybackslash}p{0.33\linewidth}%
>{\raggedright\arraybackslash}p{0.38\linewidth}>{\raggedright\arraybackslash}p{0.15\linewidth}@{}}
Sec. & Content & Paper I / Paper II & Notes (tasks)\\
\hline
\ref{sec:s2} & Lagrangian, vertices, propagators, definition of $M^2$, derivative rules & I: Sec.~II, \eqI{1}--\eqI{30},
Tables~I--II; derivative rules (text after \eqI{31}); definition \eqI{168}; calibration with \eqI{130}--\eqI{131} & Chs.~1--3
(1--2; parts of 3 and 6--8)\\
\ref{sec:s3} & 4-point contribution, coincidence limits & I: Sec.~III\,A, \eqI{31}--\eqI{44} & Ch.~2 (3--5)\\
\ref{sec:s4a} & 3-point contribution: setup, local and logarithm terms & I: Sec.~III\,B (introduction, B.1, B.2), \eqI{45}--\eqI{67}
& Chs.~3--4 (6--10)\\
\ref{sec:s4b} & 3-point normal terms: Table~III, extraction of derivatives, Tables~IV--V & I: Sec.~III\,B.3, \eqI{68}--\eqI{125},
Tables~III--V; Appendix, \eqI{A1}--\eqI{A38}, Tables~XI--XX & Chs.~5--6 (11--13)\\
\ref{sec:s5} & Counterterms, residual symmetries, divergences, Table~X, renormalized result & I: Sec.~IV, \eqI{126}--\eqI{167},
Tables~VI--X & Ch.~7--8 (14--17)\\
\ref{sec:s6} & Schwinger--Keldysh equation, Tables~2--5, late times beyond \eqII{44}, mode function & II: Secs.~2--3, \eqII{1}--\eqII{45},
Tables~1--5 & Ch.~9 (18)\\
\ref{sec:s0s1s7s8:summary} & Summary of differences, errata, conclusions, limitations & I: Sec.~V; II: Sec.~4 & all\\
\ref{app:s0s1s7s8:verification} & Verification methods & --- & code\\
\end{tabular}
\end{ruledtabular}
\end{table}

\subsection{Summary of results}
\label{sec:s0s1s7s8:results}

Most of the calculation of Paper~I is confirmed: the Feynman rules (apart from misprints), the coincidence limits (apart from the
finite part of \eqI{40}), the split of the
3-point contribution, Table~III, Table~IV (apart from one sign misprint that the paper did not use in its calculation), all non-analytic
terms of Table~V, Tables~VI--IX, and the counterterm analysis, with two invariant counterterms and the single noninvariant one
\eqI{145}. The coefficient $\beta_1$ vanishes, and $\beta_2$ is finite; the pole $-4/(D-4)$ of $\beta_3$ (paper convention) is
confirmed within Paper~I's treatment of the outer operators of the normal terms (Secs.~\ref{sec:s4b:distr} and~\ref{sec:s5:div}).
Tables~2--5 of Paper~II are reproduced exactly (Table~5 in its own sign convention). Three things change in the self-mass-squared.

\emph{The overall sign (S1).} The printed $-iM^2$ is $-1$ times the quantity defined by Paper~I's own effective field equation \eqI{168}
[our Eqs.~\eqref{eq:s0s1s7s8:def}--\eqref{eq:s0s1s7s8:defM}]. The
4-point part acquires the reversed sign through the reading of the derivative rule in \eqI{32}, the 3-point part through an overall
sign slip in the step from \eqI{46} to \eqI{47}--\eqI{49}. Since both parts are reversed, their relative sign, and the local
cancellation between \eqI{44} and \eqI{67}, are consistent. Several independent arguments support this sign
(Sec.~\ref{sec:s3}). The sign of the local part is unphysical, because the finite counterterms
$\Delta\alpha_i$ are arbitrary; the nonlocal part, however, has no counterterm freedom, and in the definition sign it is $-$(corrected
Table~X).

\emph{Two corrections of Table~X (S2, S3).} In the $(aa')^2\nabla^4$ row the contribution $8\ln x-16x\ln x$ of Table~III is missing; the
corrected row is $\tfrac{34}{3}\ln x-40x\ln x+24x^2\ln x-36x^2$. This term is non-analytic and changes the physical nonlocal
self-mass. In the $(aa')^3(a^2+a'^2)H^4$ row the constant is $+18$ instead of $-54$; this term is analytic and drops out of the retarded
kernel identically. Only with both corrections does Table~X agree, pointwise to 30 digits at random points, with the direct sum of
the normal contributions and the nonlocal logarithm term of \eqI{67}.

\emph{Absorbable finite shifts (B1--B6).} Finite errors in \eqI{40} and \eqI{59} shift $\beta_2:-\tfrac{61}{3}\to-\tfrac{49}{3}$ and
$\beta_3:+\tfrac{58}{3}\to+\tfrac{52}{3}$ (paper convention). These shifts are absorbed into the arbitrary finite parts of the
counterterms; the divergence is unchanged.

The renormalized self-mass-squared in the definition sign is therefore
\begin{equation}
-iM^2_{\rm ren}(x;x')=i\kappa^2a^4\Big(\Delta\alpha_1\square^2+\Delta\alpha_2H^2\square+\Delta\alpha_3H^2\frac{\nabla^2}{a^2}\Big)
\delta^4(x-x')-\big[\text{corrected Table X}\big],
\label{eq:s0s1s7s8:Mren}
\end{equation}
where the factors $H^2$, missing in the printed \eqI{167} (erratum C18), are required by dimensions. The corrected Table~X is given in
our Table~\ref{tab:s5:TX}.

\emph{Late times.} None of these corrections changes the late-time physics. We redid the analysis of Paper~II with the corrected
self-mass-squared in the definition sign and without the simplification \eqII{44}, keeping the exact $u(\eta',k)$, the phase
$e^{i\vec k\cdot(\vec x\,'-\vec x)}$, the $k^2/a^2$ part of $\square$ and every row containing $\nabla^2$. The justification given for
\eqII{44} fails term by term (S4): two rows of the table separately grow like $\mp\tfrac{4}{15}\hat k^2a$. Nevertheless the non-falling part
of the source cancels separately at each order $k^0,\dots,k^{10}$, for the printed as well as for the corrected table (verified
exactly through $k^{10}$, not proven to all orders; see Sec.~\ref{sec:s0s1s7s8:limitations}). The correction of the
$(aa')^2\nabla^4$ row changes only falling terms: it removes the $\hat k^4/a$, $\hat k^6/a,\dots$ tails, so that with the corrected
table the source falls like $\ln a/a^2$. The mode
function behaves as
\begin{equation}
\Phi(\eta,k)=u(\eta,k)+\kappa^2\Phi_1(\eta,k)\ \xrightarrow{\ a\to\infty\ }\ u(0,k)\Big\{1+\frac{\kappa^2H^2}{16\pi^2}\Big[\tilde C
+O\Big(\frac{\ln a}{a^2}\Big)\Big]\Big\},
\label{eq:s0s1s7s8:late}
\end{equation}
where $\tilde C$ is a constant field-strength renormalization and the tree-order approach $u(\eta,k)=u(0,k)[1+\hat k^2/(2a^2)+\dots]$
is not displayed. Through order $k^{10}$, and given the kernel and local terms as derived, the main conclusion of Paper~II stands:
one-loop quantum gravity gives no significant correction to the scalar mode functions at late times.

In the classification of Sec.~\ref{sec:s0s1s7s8:summary} there are four substantive findings: the overall sign (S1), the two
corrections of Table~X (S2, S3) and the failure of the justification of \eqII{44} (S4, an argument only, with no effect on the
result). The finite errors and their consequences for $\beta_{2,3}$ and $\alpha_{2,3}$ are the six absorbable items B1--B6. In
addition we list 23 typographical errors, C1--C23, that have no effect on any result. All differences are collected in
Sec.~\ref{sec:s0s1s7s8:summary}, together with the limitations of this work.

\providecommand{\sIIdel}{\delta^D(x-x')}
\providecommand{\sIIdd}{\partial_1\!\cdot\partial_2}
\providecommand{\sIIdp}{\partial\phi\!\cdot\!\partial\phi}
\providecommand{\sIIN}{\mathcal N}
\providecommand{\sIIcode}[1]{\texttt{code/#1}}

\section{The model, propagators and Feynman rules}
\label{sec:s2}

This section collects the ingredients of the computation: the Lagrangian and
its expansion to second order in the graviton field (Sec.~II of \PI,
Eqs.~\eqI{1}--\eqI{8} and Tables~I--II), the gauge-fixed graviton propagator
and the scalar propagators [\eqI{9}--\eqI{30}], the definition of the
self-mass-squared [\eqI{168}], and the rules by which the vertex operators
are turned into kernels. \PI{} states these rules only in words, after
\eqI{31} and before \eqI{45}. The overall sign of the printed self-mass can
be traced back to them, so we state them here as a lemma
(Sec.~\ref{sec:s2:rules}). The spacetime dimension $D$ is general
throughout.

Apart from the sign of the self-mass, every difference from the printed
equations in this part of the calculation is a misprint that did not enter
the paper's computation. The differences are listed in
our Table~\ref{tab:s2:summary} at the end of the section. Full derivations are
in Chs.~1--3 of the public derivation notes, and the scripts are in
\sIIcode{01\_feynman\_rules}, \sIIcode{02\_four\_point} and
\sIIcode{03\_three\_point\_setup} of the repository~\cite{ref:repo}. All
scripts quoted below were run and reproduce their logged output.

\subsection{Lagrangian, background and notation}
\label{sec:s2:lagrangian}

The theory is a massless, minimally coupled scalar coupled to gravity with a
positive cosmological constant~[cf.\ \eqI{1}]:
\begin{equation}
\mathcal L=-\tfrac12\,\partial_\mu\phi\,\partial_\nu\phi\,g^{\mu\nu}\sqrt{-g}
+\frac{1}{16\pi G}\big(R-(D-2)\Lambda\big)\sqrt{-g},
\qquad \Lambda\equiv(D-1)H^2 .
\label{eq:s2:L}
\end{equation}
\PI{} prints the first term as
$-\tfrac12\partial_\mu\phi\,\partial_\mu\phi\,g^{\mu\nu}\sqrt{-g}$. The only
reading without a free index is the one in \eqref{eq:s2:L} (item C1 of
our Table~\ref{tab:s2:summary}). The scalar is a spectator with zero background
value. The background is the conformal patch of $D$-dimensional de Sitter
space, and the graviton field is defined by conformally rescaling the
metric~[cf.\ \eqI{2}, \eqI{3}]:
\begin{equation}
ds^2=a^2\big(-d\eta^2+d\vec x\cdot d\vec x\big),\qquad a(\eta)=-\frac{1}{H\eta},
\qquad
g_{\mu\nu}\equiv a^2\big(\eta_{\mu\nu}+\kappa h_{\mu\nu}\big),\qquad \kappa^2\equiv16\pi G .
\label{eq:s2:background}
\end{equation}
We use the signature $(-,+,\ldots,+)$, so $a>0$ on $-\infty<\eta<0$. The
indices of $h_{\mu\nu}$ are raised and lowered with $\eta_{\mu\nu}$ only, so
$h\equiv\eta^{\mu\nu}h_{\mu\nu}$ and
$h^{\mu\nu}\equiv\eta^{\mu\alpha}\eta^{\nu\beta}h_{\alpha\beta}$. Likewise
$\partial^\mu\equiv\eta^{\mu\nu}\partial_\nu$ and
$\sIIdp\equiv\eta^{\mu\nu}\partial_\mu\phi\,\partial_\nu\phi$. This is not
merely notation: the coefficients of \eqI{6} below hold only in this
convention, since raising with the full metric would change the
$O(\kappa)$ terms. The background quantities carry a hat~[cf.\ \eqI{9}],
\begin{equation}
\hat g_{\mu\nu}=a^2\eta_{\mu\nu},\qquad \hat g^{\mu\nu}=\frac{\eta^{\mu\nu}}{a^2},\qquad
\sqrt{-\hat g}=a^D,\qquad \hat R=D(D-1)H^2 .
\label{eq:s2:hat}
\end{equation}
The gauge used below treats time and space differently, so we also need the
spatial parts of the metric and of the Kronecker delta~[cf.\ \eqI{10}]:
\begin{equation}
\bar\eta_{\mu\nu}\equiv\eta_{\mu\nu}+\delta^0_\mu\delta^0_\nu ,\qquad
\bar\delta^\mu_\nu\equiv\delta^\mu_\nu-\delta^\mu_0\delta^0_\nu .
\label{eq:s2:bar}
\end{equation}
Finally, the scalar d'Alembertian of the background~[cf.\ \eqI{42}],
\begin{equation}
\square\equiv\frac{1}{\sqrt{-\hat g}}\,\partial_\mu\big(\sqrt{-\hat g}\,\hat g^{\mu\nu}\partial_\nu\big)
=\frac{1}{a^D}\,\partial_\mu\big(a^{D-2}\partial^\mu\big),
\label{eq:s2:box}
\end{equation}
is used throughout the rest of the paper.

\subsection{Expansion of the scalar kinetic term}
\label{sec:s2:expansion}

Write the metric as the matrix product $g_{\mu\nu}=a^2\eta_{\mu\rho}(\delta^\rho{}_\nu+\kappa h^\rho{}_\nu)$.
The inverse metric then follows from the Neumann series of
$(\mathbf 1+\kappa\mathsf h)^{-1}$, with $(\mathsf h)^\mu{}_\nu\equiv h^\mu{}_\nu$.
We also obtained it independently by matching coefficients in
$g_{\mu\rho}g^{\rho\nu}=\delta_\mu{}^\nu$. The volume element follows from
$\det(\mathbf 1+\kappa\mathsf h)=\exp\operatorname{Tr}\ln(\mathbf 1+\kappa\mathsf h)$,
which we cross-checked by a direct Leibniz expansion of the determinant. The
results are~[cf.\ \eqI{4}, \eqI{5}]
\begin{align}
g^{\mu\nu}&=\frac{1}{a^2}\Big(\eta^{\mu\nu}-\kappa h^{\mu\nu}+\kappa^2h^\mu{}_\rho h^{\rho\nu}
-\kappa^3h^\mu{}_\rho h^\rho{}_\sigma h^{\sigma\nu}+O(\kappa^4)\Big),
\label{eq:s2:ginv}\\
\sqrt{-g}&=a^D\Big(1+\tfrac12\kappa h+\tfrac18\kappa^2h^2-\tfrac14\kappa^2h^{\rho\sigma}h_{\rho\sigma}+O(\kappa^3)\Big).
\label{eq:s2:vol}
\end{align}
The structure $h^2$ in \eqref{eq:s2:vol} does not come from the logarithm.
It comes only from the term $\tfrac12X^2$ in the expansion of
$\sqrt{\det(\mathbf 1+\kappa\mathsf h)}=e^X$, where
$X\equiv\tfrac12\operatorname{Tr}\ln(\mathbf 1+\kappa\mathsf h)
=\tfrac12\kappa h-\tfrac14\kappa^2h^{\rho\sigma}h_{\rho\sigma}+O(\kappa^3)$.
In the product, the only second-order cross term is
$(\tfrac12h)(-h^{\mu\nu})$. Since $\partial_\mu\phi$ does not depend on the
metric, the kinetic term becomes~[cf.\ \eqI{6}]
\begin{equation}
\begin{aligned}
-\tfrac12\partial_\mu\phi\,\partial_\nu\phi\,g^{\mu\nu}\sqrt{-g}
=-\tfrac12\partial_\mu\phi\,\partial_\nu\phi\,a^{D-2}\Big\{&\eta^{\mu\nu}-\kappa h^{\mu\nu}+\tfrac12\eta^{\mu\nu}\kappa h
+\tfrac18\eta^{\mu\nu}\kappa^2h^2-\tfrac14\eta^{\mu\nu}\kappa^2h^{\rho\sigma}h_{\rho\sigma}\\
&-\tfrac12\kappa^2hh^{\mu\nu}+\kappa^2h^{\mu\rho}h_\rho{}^\nu+O(\kappa^3)\Big\}.
\end{aligned}
\label{eq:s2:eq6}
\end{equation}
The common factor $a^{D-2}$ is $a^{-2}$ from the inverse metric times $a^D$
from the determinant. It equals $1$ at $D=2$, which reflects the Weyl
invariance of the two-dimensional kinetic term. The parts linear and
quadratic in $h$ are
\begin{align}
\kappa\mathcal L^{(1)}&=\kappa a^{D-2}\Big[\tfrac12h^{\mu\nu}\partial_\mu\phi\,\partial_\nu\phi-\tfrac14h\,\sIIdp\Big],
\label{eq:s2:L1}\\
\kappa^2\mathcal L^{(2)}&=\kappa^2a^{D-2}\Big[-\tfrac1{16}h^2\,\sIIdp+\tfrac18h^{\rho\sigma}h_{\rho\sigma}\,\sIIdp
+\tfrac14hh^{\mu\nu}\partial_\mu\phi\,\partial_\nu\phi-\tfrac12h^{\mu\rho}h_\rho{}^\nu\partial_\mu\phi\,\partial_\nu\phi\Big].
\label{eq:s2:L2}
\end{align}

\paragraph*{Checks.}
Two special cases can be computed exactly from the full metric, without
using the expansion. For $h_{\mu\nu}=c\,\eta_{\mu\nu}$ one has
$\sqrt{-g}\,g^{\mu\nu}=a^{D-2}(1+\kappa c)^{(D-2)/2}\eta^{\mu\nu}$, and its
$\kappa^2$ coefficient is $\tfrac18(D-2)(D-4)c^2$. Write the general
second-order bracket as
$A\,h^{\mu\rho}h_\rho{}^\nu+B\,hh^{\mu\nu}+C\,\eta^{\mu\nu}h^2+E\,\eta^{\mu\nu}h^{\rho\sigma}h_{\rho\sigma}$.
Matching powers of $D$ fixes $A=1$, $C=\tfrac18$ and $B+E=-\tfrac34$. For
$h_{\mu\nu}=b\,\delta^0_\mu\delta^0_\nu$ one has
$\sqrt{-g}\,g^{00}=-a^{D-2}(1-\kappa b)^{-1/2}$ and
$\sqrt{-g}\,g^{ij}=a^{D-2}\delta^{ij}(1-\kappa b)^{1/2}$, which gives
$C+E=-\tfrac18$. Together the two cases fix all four coefficients,
$(A,B,C,E)=(1,-\tfrac12,\tfrac18,-\tfrac14)$, in agreement with
\eqref{eq:s2:eq6}. The script \sIIcode{01\_feynman\_rules/check\_eq6.py}
builds the inverse metric and $\sqrt{-g}$ from the exact adjugate and
determinant, expands them, and compares with the bracket of
\eqref{eq:s2:eq6}. The comparison is fully symbolic for $D=2,3$ and uses
random rational $h_{\mu\nu}$ for $D=4,5,6$. The difference vanishes in all
cases.

\emph{Result:} \eqI{4}, \eqI{5} and the bracket of \eqI{6} are exact for
general $D$, with no gauge or trace condition. The printed left-hand side of
\eqI{6} reads $-\tfrac12\partial_\mu\phi\,\partial_\nu g^{\mu\nu}\sqrt{-g}$;
the second $\phi$ is missing. This slip (C2) was noticed only while
preparing the write-up.

\subsection{Vertex operators: corrected Tables I and II}
\label{sec:s2:vertices}

\PI{} represents the interactions as vertex operators acting on the fields.
The two scalar legs are numbered ``1'', ``2'' and the gravitons ``3'',
``4''. A factor $\tfrac12$ is extracted for the two identical scalars, but
not for the two identical gravitons of the 4-point vertices. With the
standard rule that a vertex carries $i$ times the corresponding coefficient
in the action, this defines the normalization
\begin{equation}
\kappa\mathcal L^{(1)}=\frac12\cdot\frac1i\sum_I\Big[V_I^{\alpha\beta}\phi_1\phi_2\Big]_{\phi_1=\phi_2=\phi}h_{\alpha\beta},
\qquad
\kappa^2\mathcal L^{(2)}=\frac12\cdot\frac1i\sum_I\Big[U_I^{\alpha\beta\rho\sigma}\phi_1\phi_2\Big]_{\phi_1=\phi_2=\phi}h_{\alpha\beta}h_{\rho\sigma}.
\label{eq:s2:norm}
\end{equation}
Consider a term $C\,T(\partial\phi,\partial\phi)\,h\cdots$, where $C$ contains
the number, $\kappa^n$ and $a^{D-2}$. By \eqref{eq:s2:norm} its vertex is
$2iC\,T(\partial_1,\partial_2)$. This vertex is the full scalar functional
derivative: for a symmetric $X^{\mu\nu}$, the part of
$C\,X^{\mu\nu}\partial_\mu\phi\,\partial_\nu\phi$ bilinear in
$\phi=\phi_1+\phi_2$ is $2C\,X^{\mu\nu}\partial_\mu\phi_1\partial_\nu\phi_2$.
A 4-point operator $U$ carries only one assignment of the two graviton legs,
and the full vertex is $U[3,4]+U[4,3]$. When the two graviton legs are
joined by the coincident propagator $X_{\alpha\beta\rho\sigma}$, which is
pair symmetric (Sec.~\ref{sec:s2:graviton}), the symmetry factor $\tfrac12$
of the full vertex gives
$\tfrac12(U^{\alpha\beta\rho\sigma}+U^{\rho\sigma\alpha\beta})X_{\alpha\beta\rho\sigma}=U^{\alpha\beta\rho\sigma}X_{\alpha\beta\rho\sigma}$.
This is why \PI{} needs no symmetry factor. Combining $2iC$ with
\eqref{eq:s2:eq6} gives a compact rule: every vertex equals $-i$ times the
coefficient $b_I$ of the corresponding structure in the bracket of \eqI{6},
times $\kappa^na^{D-2}$ and the contracted structure. The resulting vertices
are listed in our Table~\ref{tab:s2:vertices}.

\begin{table}[tb]
\caption{3-point vertex operators $V_I^{\alpha\beta}$, contracted into
$\phi_1\phi_2h_{\alpha\beta}$, and 4-point vertex operators
$U_I^{\alpha\beta\rho\sigma}$, contracted into
$\phi_1\phi_2h_{\alpha\beta}h_{\rho\sigma}$ (cf.\ Tables~I and~II of \PI).
The number $b_I$ is the coefficient of the source structure in the bracket of
\eqref{eq:s2:eq6}, and every vertex equals $-ib_I\kappa^na^{D-2}$ times its
structure. Only $U_2$ differs from the print.}
\label{tab:s2:vertices}
\begin{ruledtabular}
\begin{tabular}{llcll}
Vertex & Source structure in \eqref{eq:s2:eq6} & $b_I$ & Derived (this work) & \PI{} (printed)\\
\hline
$V_1$ & $h^{\mu\nu}$ & $-1$ & $+i\kappa a^{D-2}\partial_1^\alpha\partial_2^\beta$ & same\\
$V_2$ & $\eta^{\mu\nu}h$ & $+\tfrac12$ & $-\tfrac i2\kappa a^{D-2}\eta^{\alpha\beta}\,\sIIdd$ & same\\
$U_1$ & $\eta^{\mu\nu}h^2$ & $+\tfrac18$ & $-\tfrac i8\kappa^2a^{D-2}\eta^{\alpha\beta}\eta^{\rho\sigma}\,\sIIdd$ & same\\
$U_2$ & $\eta^{\mu\nu}h^{\rho\sigma}h_{\rho\sigma}$ & $-\tfrac14$ & $+\tfrac i4\kappa^2a^{D-2}\eta^{\alpha\rho}\eta^{\beta\sigma}\,\sIIdd$ & $-\tfrac i4\kappa^2a^{D-2}\eta^{\alpha\rho}\eta^{\beta\sigma}\,\sIIdd$\\
$U_3$ & $hh^{\mu\nu}$ & $-\tfrac12$ & $+\tfrac i2\kappa^2a^{D-2}\eta^{\alpha\beta}\partial_1^\rho\partial_2^\sigma$ & same\\
$U_4$ & $h^{\mu\rho}h_\rho{}^\nu$ & $+1$ & $-i\kappa^2a^{D-2}\partial_1^\alpha\eta^{\beta\rho}\partial_2^\sigma$ & same\\
\end{tabular}
\end{ruledtabular}
\end{table}

\paragraph*{Eq.~(I.7).}
\PI{} prints
$-\tfrac12\kappa a^{D-2}\partial_\alpha\phi\,\partial_\beta\phi\,h^{\alpha\beta}\Rightarrow V_1^{\alpha\beta}=i\kappa a^{D-2}\partial_1^\alpha\partial_2^\beta$.
The paper's own normalization, which can be read off from \eqI{8}, is
$U_1/C=2i$. With this normalization a coefficient $-\tfrac12\kappa$ would
give $V_1=-i\kappa a^{D-2}\partial_1^\alpha\partial_2^\beta$. The
coefficient in \eqref{eq:s2:L1} is $+\tfrac12\kappa$, so the sign on the
printed left-hand side is a misprint (C3). The printed $V_1$ is correct, and
it is the vertex used in \eqI{46}.

\paragraph*{The sign of $U_2$.}
The printed $-\tfrac i4$ in row~2 of Table~II is a misprint (C5). Two checks
based on the full metric alone, without using \eqref{eq:s2:eq6}, fix all
coefficients uniquely. Take the ansatz
$V=i\kappa a^{D-2}[v_1\,\partial_1^\alpha\partial_2^\beta+v_2\,\eta^{\alpha\beta}\sIIdd]$
and $U=i\kappa^2a^{D-2}\sum_Iu_I\,(\text{structure}_I)$ with the structures
of our Table~\ref{tab:s2:vertices}. For $h_{\mu\nu}=c\,\eta_{\mu\nu}$,
\eqref{eq:s2:norm} must reproduce
$\kappa\mathcal L^{(1)}=-\tfrac14(D-2)\kappa c\,a^{D-2}\sIIdp$ and
$\kappa^2\mathcal L^{(2)}=-\tfrac1{16}(D-2)(D-4)\kappa^2c^2a^{D-2}\sIIdp$.
This requires
\begin{equation}
\tfrac12(v_1+v_2D)=-\tfrac{D-2}{4},\qquad
\tfrac12\big[u_1D^2+(u_2+u_3)D+u_4\big]=-\tfrac{(D-2)(D-4)}{16},
\label{eq:s2:confsys}
\end{equation}
so $v_1=1$, $v_2=-\tfrac12$, $u_1=-\tfrac18$, $u_2+u_3=\tfrac34$ and $u_4=-1$.
For $h_{\mu\nu}=b\,\delta^0_\mu\delta^0_\nu$ one needs $u_1+u_2=\tfrac18$ and
$u_1+u_2+u_3+u_4=-\tfrac38$. Hence $(u_1,u_2,u_3,u_4)=(-\tfrac18,+\tfrac14,+\tfrac12,-1)$
is the unique solution. The printed $u_2=-\tfrac14$ violates
$u_1+u_2=\tfrac18$. With the printed Table~II, the reconstructed
Lagrangian differs from the exact one by
$-\tfrac14\kappa^2a^{D-2}h^{\rho\sigma}h_{\rho\sigma}\,\sIIdp$. Similarly,
the printed sign on the left-hand side of \eqI{7} would correspond to $v_1=-1$
and fail the first condition in \eqref{eq:s2:confsys}. The script
\sIIcode{01\_feynman\_rules/check\_vertices.py} confirms all of this in
three ways:
\begin{itemize}
\item[(A)] it obtains the full vertices by polarization with distinct legs
(symbolic, $D=3,4$);
\item[(B)] it compares the tables with the exact metric ($D=2,\dots,5$);
\item[(C)] it solves the linear system above.
\end{itemize}
The printed \eqI{33} is built with $+\tfrac i4$ (Sec.~\ref{sec:s3:rule}).
The Table~II entry is therefore a misprint only.

\subsection{Gauge fixing and the graviton propagator}
\label{sec:s2:graviton}

The graviton propagator is that of the noninvariant gauge of
Refs.~\cite{ref:Tsamis1994,KW07}. The gauge-fixing term is~[cf.\ \eqI{15}]
\begin{equation}
\mathcal L_{\rm GF}=-\tfrac12a^{D-2}\eta^{\mu\nu}F_\mu F_\nu,\qquad
F_\mu\equiv\eta^{\rho\sigma}\Big(h_{\mu\rho,\sigma}-\tfrac12h_{\rho\sigma,\mu}+(D-2)Ha\,h_{\mu\rho}\delta^0_\sigma\Big).
\label{eq:s2:GF}
\end{equation}
After partial integration, the quadratic part of the gauge-fixed Lagrangian
takes the form $\tfrac12h^{\mu\nu}D_{\mu\nu}{}^{\rho\sigma}h_{\rho\sigma}$,
with~[cf.\ \eqI{16}--\eqI{19}]
\begin{align}
D_{\mu\nu}{}^{\rho\sigma}\equiv{}&\Big\{\tfrac12\bar\delta_\mu{}^{(\rho}\bar\delta_\nu{}^{\sigma)}-\tfrac14\eta_{\mu\nu}\eta^{\rho\sigma}
-\frac{1}{2(D-3)}\delta^0_\mu\delta^0_\nu\delta^\rho_0\delta^\sigma_0\Big\}D_A
+\delta^0_{(\mu}\bar\delta_{\nu)}^{(\rho}\delta^{\sigma)}_0\,D_B
+\frac12\Big(\frac{D-2}{D-3}\Big)\delta^0_\mu\delta^0_\nu\delta^\rho_0\delta^\sigma_0\,D_C,
\label{eq:s2:Dop}\\
D_A\equiv{}&\partial_\mu\big(\sqrt{-\hat g}\,\hat g^{\mu\nu}\partial_\nu\big),\qquad
D_B\equiv D_A-\frac1D\Big(\frac{D-2}{D-1}\Big)\hat R\sqrt{-\hat g},\qquad
D_C\equiv D_A-\frac2D\Big(\frac{D-3}{D-1}\Big)\hat R\sqrt{-\hat g}.
\label{eq:s2:DABC}
\end{align}
By \eqref{eq:s2:hat}, $D_B=D_A-(D-2)H^2a^D$ and $D_C=D_A-2(D-3)H^2a^D$. The
propagator is a sum of constant tensor factors times scalar
propagators~[cf.\ \eqI{20}--\eqI{25}]:
\begin{gather}
i\big[{}_{\mu\nu}\Delta_{\rho\sigma}\big](x;x')=\sum_{I=A,B,C}\big[{}_{\mu\nu}T^I_{\rho\sigma}\big]\,i\Delta_I(x;x'),
\qquad D_I\times i\Delta_I(x;x')=i\sIIdel ,
\label{eq:s2:hprop}\\
\big[{}_{\mu\nu}T^A_{\rho\sigma}\big]=2\bar\eta_{\mu(\rho}\bar\eta_{\sigma)\nu}-\frac{2}{D-3}\bar\eta_{\mu\nu}\bar\eta_{\rho\sigma},\qquad
\big[{}_{\mu\nu}T^B_{\rho\sigma}\big]=-4\delta^0_{(\mu}\bar\eta_{\nu)(\rho}\delta^0_{\sigma)},
\label{eq:s2:TAB}\\
\big[{}_{\mu\nu}T^C_{\rho\sigma}\big]=\frac{2}{(D-2)(D-3)}\big[(D-3)\delta^0_\mu\delta^0_\nu+\bar\eta_{\mu\nu}\big]\big[(D-3)\delta^0_\rho\delta^0_\sigma+\bar\eta_{\rho\sigma}\big],
\label{eq:s2:TC}\\
D_{\mu\nu}{}^{\rho\sigma}\times i\big[{}_{\rho\sigma}\Delta^{\alpha\beta}\big](x;x')=\delta_\mu^{(\alpha}\delta_\nu^{\beta)}\,i\sIIdel .
\label{eq:s2:propeq}
\end{gather}
We take \eqref{eq:s2:GF}--\eqref{eq:s2:propeq} from \PI{} and
Ref.~\cite{ref:Tsamis1994} as input; we did not re-derive the propagator.
The gauge \eqref{eq:s2:GF} breaks de Sitter invariance. Of the de Sitter
transformations it breaks only the spatial special conformal ones. This
property was checked in the counterterm analysis (Sec.~\ref{sec:s5:dS}).

\paragraph*{Properties used later.}
With $u_\mu\equiv\delta^0_\mu$, we have $u^\mu=-\delta^\mu_0$, $u\cdot u=-1$
and $\bar\eta_{\mu\nu}=\eta_{\mu\nu}+u_\mu u_\nu$. The tensor factors are
symmetric within each index pair and under exchange of the pairs. Hence the
graviton propagator obeys
$i[{}_{\alpha\beta}\Delta_{\rho\sigma}](x;x')=i[{}_{\rho\sigma}\Delta_{\alpha\beta}](x';x)$,
and its coincidence limit
$X_{\alpha\beta\rho\sigma}(x)\equiv i[{}_{\alpha\beta}\Delta_{\rho\sigma}](x;x)$
satisfies $X_{\alpha\beta\rho\sigma}=X_{\rho\sigma\alpha\beta}$. In
dimensional regularization, and for general $D$, $X$ is a smooth function of
$x$. The two symmetries together imply that the partial trace
$N_{\lambda\kappa}\equiv X_{\lambda\rho}{}^\rho{}_\kappa$ is symmetric. The
single and double traces needed for the contractions are
\begin{equation}
\begin{gathered}
\eta^{\alpha\beta}T^{A}_{\alpha\beta\mu\nu}=-\frac{4}{D-3}\bar\eta_{\mu\nu},\qquad
\eta^{\alpha\beta}T^{B}_{\alpha\beta\mu\nu}=0,\qquad
\eta^{\alpha\beta}T^{C}_{\alpha\beta\mu\nu}=\frac{4\big[(D-3)u_\mu u_\nu+\bar\eta_{\mu\nu}\big]}{(D-2)(D-3)},\\
\eta^{\alpha\beta}\eta^{\rho\sigma}T^{A}_{\alpha\beta\rho\sigma}=-\frac{4(D-1)}{D-3},\qquad
\eta^{\alpha\beta}\eta^{\rho\sigma}T^{B}_{\alpha\beta\rho\sigma}=0,\qquad
\eta^{\alpha\beta}\eta^{\rho\sigma}T^{C}_{\alpha\beta\rho\sigma}=\frac{8}{(D-2)(D-3)}.
\end{gathered}
\label{eq:s2:traces}
\end{equation}
The pair and within-pair symmetries were verified with explicit component
arrays for $D=4,\dots,7$, by the script
\sIIcode{02\_four\_point/check\_coincidence.py}. The traces enter the
contractions checked there and those checked for $D=4,5,6$ by
\sIIcode{03\_three\_point\_setup/check\_47\_49.py}.

\subsection{Scalar propagators}
\label{sec:s2:scalar}

All propagators are written in terms of the de Sitter length
function $y(x;x')$ of Eq.~\eqref{eq:s0s1s7s8:y}~[cf.\ \eqI{11}]. The A-type propagator obeys the equation of a massless, minimally coupled
scalar, $D_A\,i\Delta_A(x;x')=i\sIIdel$~[cf.\ \eqI{12}]. This equation has
no de Sitter-invariant solution~\cite{ref:Allen1987}. The solution relevant
for cosmology is the one that preserves homogeneity and isotropy, the
``E(3)'' vacuum~\cite{ref:Allen1985}, and its minimal form
is~\cite{ref:Onemli2002,ref:Onemli2004}~[cf.\ \eqI{13}, \eqI{14}]
\begin{equation}
i\Delta_A(x;x')=A(y)+k\ln(aa'),\qquad k\equiv\sIIN\,\frac{\Gamma(D-1)}{\Gamma(\frac D2)},\qquad
\sIIN\equiv\frac{H^{D-2}}{(4\pi)^{D/2}},
\label{eq:s2:DA}
\end{equation}
\begin{equation}
\begin{aligned}
A(y)=\sIIN\bigg\{&\Gamma\Big(\frac D2-1\Big)\Big(\frac4y\Big)^{\frac D2-1}
+\frac{\Gamma(\frac D2+1)}{\frac D2-2}\Big(\frac4y\Big)^{\frac D2-2}
-\pi\cot\Big(\frac{\pi D}{2}\Big)\frac{\Gamma(D-1)}{\Gamma(\frac D2)}\\
&+\sum_{n=1}^\infty\Big[\frac1n\frac{\Gamma(n+D-1)}{\Gamma(n+\frac D2)}\Big(\frac y4\Big)^n
-\frac{1}{n-\frac D2+2}\frac{\Gamma(n+\frac D2+1)}{\Gamma(n+2)}\Big(\frac y4\Big)^{n-\frac D2+2}\Big]\bigg\}.
\end{aligned}
\label{eq:s2:Ay}
\end{equation}
The first coefficient in \eqref{eq:s2:Ay} is $\Gamma(\tfrac D2-1)$.
\PI{} prints $\Gamma(\tfrac D2-1)/(\tfrac D2-1)$ (C4). The B-type and C-type
propagators are de Sitter invariant. Their most singular part is the
conformally coupled propagator~\cite{ref:Birrell1982}~[cf.\
\eqI{26}--\eqI{28}]:
\begin{gather}
i\Delta_{\rm cf}(x;x')=\sIIN\,\Gamma\Big(\frac D2-1\Big)\Big(\frac4y\Big)^{\frac D2-1},
\label{eq:s2:Dcf}\\
i\Delta_B=i\Delta_{\rm cf}-\sIIN\sum_{n=0}^\infty\bigg\{\frac{\Gamma(n+D-2)}{\Gamma(n+\frac D2)}\Big(\frac y4\Big)^n
-\frac{\Gamma(n+\frac D2)}{\Gamma(n+2)}\Big(\frac y4\Big)^{n-\frac D2+2}\bigg\},
\label{eq:s2:DB}\\
i\Delta_C=i\Delta_{\rm cf}+\sIIN\sum_{n=0}^\infty\bigg\{(n+1)\frac{\Gamma(n+D-3)}{\Gamma(n+\frac D2)}\Big(\frac y4\Big)^n
-\Big(n-\frac D2+3\Big)\frac{\Gamma(n+\frac D2-1)}{\Gamma(n+2)}\Big(\frac y4\Big)^{n-\frac D2+2}\bigg\}.
\label{eq:s2:DC}
\end{gather}
Equivalently~\cite{ref:Candelas1975,ref:Dowker1976}~[cf.\ \eqI{29}, \eqI{30}],
$i\Delta_B=\sIIN\,\frac{\Gamma(D-2)\Gamma(1)}{\Gamma(D/2)}\,{}_2F_1\big(D-2,1;\tfrac D2;1-\tfrac y4\big)$
and
$i\Delta_C=\sIIN\,\frac{\Gamma(D-3)\Gamma(2)}{\Gamma(D/2)}\,{}_2F_1\big(D-3,2;\tfrac D2;1-\tfrac y4\big)$.
In \eqref{eq:s2:Ay}--\eqref{eq:s2:DC} the infinite sums vanish at $D=4$, and
each of their terms is a positive power of $y$. Away from coincidence,
\begin{equation}
\square f(y)=H^2\Big[(4y-y^2)f''(y)+D(2-y)f'(y)\Big],
\label{eq:s2:boxy}
\end{equation}
which we verified with explicit coordinates for $D=4,5$
(\sIIcode{03\_three\_point\_setup/check\_52.py}, check (g)). As a check
that the definitions were read correctly, we compared the series
\eqref{eq:s2:DB}--\eqref{eq:s2:DC} with the hypergeometric forms at $D=0.6$
and $y=0.5$. They agree to about $10^{-31}$
(\sIIcode{02\_four\_point/check\_coincidence.py}).

\paragraph*{The first coefficient of (I.14).}
Three independent arguments fix the coefficient $c_0$ of
$\sIIN(4/y)^{D/2-1}$ in $A(y)$ to $\Gamma(\tfrac D2-1)$.
\begin{enumerate}
\item The leading singularity of $A(y)$ is that of $i\Delta_{\rm cf}$,
\eqref{eq:s2:Dcf}. Up to the conformal factor $(aa')^{1-D/2}$, this is the
flat-space Feynman propagator
$\Gamma(\tfrac D2-1)/(4\pi^{D/2})\,(\Delta x^2)^{1-D/2}$, with exactly this
coefficient.
\item Away from coincidence, \eqref{eq:s2:boxy} and the propagator equation
give $(4y-y^2)A''+D(2-y)A'=\text{const}$. The coefficient of $y^{1-D/2}$ must
therefore cancel between the first two terms of \eqref{eq:s2:Ay}. This
requires $c_0=(\tfrac D2-2)c_1/[(\tfrac D2-1)\tfrac D2]$. With the printed
second coefficient $c_1=\Gamma(\tfrac D2+1)/(\tfrac D2-2)$ this gives
$c_0=\Gamma(\tfrac D2-1)$.
\item The mixed second derivative of $i\Delta_A$ has a contact term. The
inner derivatives annihilate $k\ln(aa')$, and~[cf.\ \eqI{52}]
\begin{equation}
\partial_\rho\partial'_\sigma\,i\Delta_A(x;x')=\delta^0_\rho\delta^0_\sigma\,\frac{i}{a^{D-2}}\,\sIIdel
+A''(y)\,\frac{\partial y}{\partial x^\rho}\frac{\partial y}{\partial x'^\sigma}
+A'(y)\,\frac{\partial^2y}{\partial x^\rho\partial x'^\sigma}.
\label{eq:s2:eq52}
\end{equation}
The local coefficient was obtained in three ways: directly from the series
with the $i\delta$ prescription, from the propagator equation \eqI{12}
alone, and from the canonical equal-time commutator. The direct computation
gives $i\,c_0(\tfrac D2-1)a^{2-D}/\Gamma(\tfrac D2)$. The other two routes
do not use the normalization of $A(y)$, and both give $i\,a^{2-D}$.
\end{enumerate}
With the printed $c_0$, the local coefficient in \eqref{eq:s2:eq52} would
carry an extra factor $2/(D-2)$. This factor is $1$ at $D=4$. Multiplied by
the pole of $i\Delta_A(x;x)$, however, it would leave the finite error
$2k/(D-2)$. The printed \eqI{52} has the correct coefficient, so the error
is in the printing only. It does not affect the coincidence limits either,
because the term vanishes as $y\to0$ in dimensional regularization. The
derivation of \eqref{eq:s2:eq52} and its use in the 3-point contribution
are given in Sec.~\ref{sec:s4a:split}
(\sIIcode{03\_three\_point\_setup/check\_52.py}, checks (e) and (f)).

\subsection{The self-mass-squared and the two sign conventions}
\label{sec:s2:definition}

\PI{} defines the self-mass-squared through the effective field equation
for the quantum-corrected mode functions. \PI{} prints~[\eqI{168}]
\begin{equation*}
\partial_\mu\big(\sqrt{-\hat g}\,\hat g^{\mu\nu}\partial_\nu\Phi(x)\big)-\int d^4x'\,M^2_{\rm ren}(x;x')\,\Phi(x')=0 .
\end{equation*}
We adopt this equation as the definition of $M^2$. For the regulated
kernel we use the same form in $D$ dimensions. Write the effective action
of the background $\Phi$ as
\begin{equation}
e^{i\Gamma[\Phi]}=\int[Dh][D\varphi]_{\rm 1PI}\,e^{iS[\Phi+\varphi,h]},\qquad
\Gamma[\Phi]=S_0[\Phi]-\tfrac12\iint\Phi\,M^2\,\Phi+\dots,
\label{eq:s2:Gamma}
\end{equation}
where $S_0=\int-\tfrac12a^{D-2}\partial\Phi\cdot\partial\Phi$ has
$\delta S_0/\delta\Phi=\partial_\mu(\sqrt{-\hat g}\,\hat g^{\mu\nu}\partial_\nu\Phi)$.
Then $\delta\Gamma/\delta\Phi=0$ is \eqI{168}, and
\begin{equation}
-iM^2(x;x')=\frac{\delta^2\big(i\Gamma_1[\Phi]\big)}{\delta\Phi(x)\,\delta\Phi(x')},
\label{eq:s2:defM}
\end{equation}
where $\Gamma_1$ is the quantum part of $\Gamma$. This matches Sec.~IV of
\PI, which refers to the self-mass-squared as a one-particle-irreducible
(1PI) function. A local counterterm action $S_{\rm ct}$
contributes $-iM^2_{\rm ct}=i\,\delta^2S_{\rm ct}/\delta\Phi\,\delta\Phi$.
Section~\ref{sec:s2:rules} shows that this reproduces the paper's own
counterterm vertices \eqI{130}--\eqI{131}.

\paragraph*{Definition sign and paper convention.}
We call the sign fixed by \eqref{eq:s2:defM}, i.e.\ by \eqI{168}, the
\emph{definition sign}. The \emph{paper convention} (or Table-X convention)
is the sign of the self-mass contributions as \PI{} prints them. This
paper establishes that the definition sign is $-1$ times the paper
convention, Eq.~\eqref{eq:s0s1s7s8:signs}.
Our Table~\ref{tab:s2:conventions} lists which printed equations are in which
convention. Both the 4-point and the 3-point contributions are reversed, so
their relative sign in \PI{} is consistent. The local cancellation stated
after \eqI{67} holds in either convention. The sign of the local part of the
renormalized self-mass has no physical meaning, since it is fixed by the
arbitrary finite counterterms $\Delta\alpha_i$. The nonlocal part has no
counterterm freedom. Once \eqI{168} is adopted, its sign is physical: the
physical nonlocal self-mass is $-$(corrected Table~X). We make every
physical statement in the definition sign and state the convention wherever
we quote a printed number. The overall sign reverses the one-loop
correction to the mode function of \PII; as shown in Sec.~\ref{sec:s6},
the non-falling total of that correction is zero at late times, so only
terms that fall off change sign and the late-time conclusion is unchanged.

\begin{table}[tb]
\caption{Sign conventions of the printed equations. The definition sign is
fixed by \eqI{168} and \eqref{eq:s2:defM}; the paper convention is the sign
as printed, and definition sign $=-1\times$ paper convention.}
\label{tab:s2:conventions}
\begin{ruledtabular}
\begin{tabular}{>{\raggedright\arraybackslash}p{0.22\textwidth}>{\raggedright\arraybackslash}p{0.72\textwidth}}
Convention & Printed equations\\
\hline
Paper convention & 4-point: \eqI{32}, \eqI{33}, \eqI{41}, \eqI{44}.
3-point: \eqI{47}--\eqI{49}, \eqI{54}, \eqI{57}--\eqI{59}, \eqI{63}--\eqI{67}, \eqI{89}, Tables~III--X.
Renormalization: \eqI{160}--\eqI{165}, \eqI{167}.
\PII: Table~1, \eqII{10}, Table~5.\\
Definition sign & \eqI{45}, \eqI{46}.\\
Consistent with the definition & counterterm vertices \eqI{130}, \eqI{131}, \eqI{145}, built with the convention-free rule
vertex $=i\,\delta^2S_{\rm ct}/\delta\phi\,\delta\phi$ (Sec.~\ref{sec:s2:rules}).\\
Convention independent & \eqI{50}--\eqI{53}, \eqI{55}--\eqI{56},
\eqI{60}--\eqI{62}, \eqI{68}--\eqI{88}, \eqI{90}--\eqI{103} (identities without an overall coefficient);
\eqI{166}, $\alpha_i=-\beta_i+\Delta\alpha_i$.\\
\end{tabular}
\end{ruledtabular}
\end{table}

\subsection{From vertex operators to kernels: Lemma A and the derivative rules}
\label{sec:s2:rules}

Since $\sIIdel$ depends only on $x-x'$, we have
$\partial'_\mu\sIIdel=-\partial_\mu\sIIdel$. A kernel acts as
$(\mathcal K f)(x)=\int d^Dx'\,\mathcal K(x;x')f(x')$, and two kernels are
equal if they agree on all pairs of test functions $f_1(x)f_2(x')$.

\paragraph*{Lemma A.}
Let $F[\phi]=\int d^Dz\,K^{\mu\nu}(z)\,\partial_\mu\phi\,\partial_\nu\phi$
with $K$ smooth and not necessarily symmetric, and let
$\mathcal K^{\mu\nu}\equiv K^{\mu\nu}+K^{\nu\mu}$. Then
\begin{equation}
\frac{\delta^2F}{\delta\phi(x)\,\delta\phi(x')}
=-\partial_\mu\Big[\mathcal K^{\mu\nu}(x)\,\partial_\nu\sIIdel\Big]
=+\partial_\mu\Big[\mathcal K^{\mu\nu}(x)\,\partial'_\nu\sIIdel\Big].
\label{eq:s2:lemmaA}
\end{equation}
\emph{Proof.} The coefficient of $\epsilon_1\epsilon_2$ in
$F[\epsilon_1f_1+\epsilon_2f_2]$ is
$\int d^Dx\,\mathcal K^{\mu\nu}\partial_\mu f_1\,\partial_\nu f_2$. Apply the
right-hand side of \eqref{eq:s2:lemmaA} to $f_2$, integrate against $f_1$,
and integrate by parts once; the boundary term vanishes by compact support.
The result is the same expression. We also proved the lemma with formal
$\delta$-functions, by integrating by parts in the vertex coordinate $z$.

In words: one derivative is moved, by partial integration, onto the whole
expression and becomes $-\partial$. The other acts on $\sIIdel$ as $+\partial$
\emph{with respect to $x$}, which equals $-\partial'$. A positivity test
fixes the sign. For $D=1$ and $K=c>0$, $F=\int c\,\phi'^2\ge0$, so the
Hessian must be a positive operator. Indeed
$\int f\,[-2(cf')']=\int2cf'^2\ge0$, while the opposite sign gives a
negative operator.

\paragraph*{Calibration with the counterterm vertices.}
\PI{} prints~[\eqI{130}, \eqI{131}]
\begin{equation*}
\tfrac12\alpha_1\kappa^2\square\phi\,\square\phi\,a^D\ \to\ i\alpha_1\kappa^2a^D\square^2\sIIdel,
\qquad
-\tfrac12\alpha_2\kappa^2H^2\partial_\mu\phi\,\partial^\mu\phi\,a^{D-2}\ \to\ i\alpha_2\kappa^2H^2a^D\square\,\sIIdel .
\end{equation*}
For the second term, Lemma~A with
$\mathcal K^{\mu\nu}=-\alpha_2\kappa^2H^2a^{D-2}\eta^{\mu\nu}$ and
\eqref{eq:s2:box} give
\begin{equation}
i\frac{\delta^2S_{\rm ct}}{\delta\Phi\,\delta\Phi}
=-i\,\partial_\mu\big[-\alpha_2\kappa^2H^2a^{D-2}\partial^\mu\sIIdel\big]
=+i\alpha_2\kappa^2H^2a^D\square\,\sIIdel ,
\label{eq:s2:calib131}
\end{equation}
which is identical to \eqI{131}. Because $a^D\square$ is self-adjoint, the
first term gives exactly \eqI{130}. The definition \eqref{eq:s2:defM}
therefore reproduces the paper's own counterterm vertices. Of the two,
the vertex \eqI{131} carries one derivative on each leg, like the 4-point
vertices, and so calibrates the derivative-sign rule directly. The vertex
\eqI{130} carries two derivatives on each leg and does not depend on that
rule.

\paragraph*{4-point contribution.}
At one loop the 4-point vertices contribute
$i\Gamma_{\rm 4pt}=i\int d^Dz\,\langle\kappa^2\mathcal L^{(2)}\rangle$.
Once the two graviton fields are contracted into the coincident propagator
$X_{\alpha\beta\rho\sigma}$, this is a quadratic form of the type of
Lemma~A in $\partial\Phi$. Lemma~A and \eqref{eq:s2:defM} then give the
kernel in the form of \eqI{31},
$-iM^2_{\rm 4pt}=\sum_IU_I^{\alpha\beta\rho\sigma}\,i[{}_{\alpha\beta}\Delta_{\rho\sigma}](x;x)\,\sIIdel$,
with normalization $1$ (the scalar $\tfrac12$ in \eqref{eq:s2:norm}
cancels against the factor $2$ of Lemma~A), provided the derivatives are
read as follows:
\begin{itemize}
\item[(R1)] $\partial_{1\mu}\to-\partial_\mu$, acting on everything to its
right: $a^{D-2}$, $X$ and the factor containing $\delta^D$;
\item[(R2)] $\partial_{2\nu}\to+\partial_\nu$, the derivative with respect to
$x$, acting only on $\sIIdel$; this means $-\partial'_\nu$.
\end{itemize}
The paper's rule for $\partial_1$ is (R1). For $\partial_2$, the
text after \eqI{31} says that ``$\partial_2$'' acts upon $x'$, and \eqI{32}
uses $\partial_2\to+\partial'$. This is the opposite of (R2). Every $U_I$
contains exactly one $\partial_1$ and one $\partial_2$, so the difference is
a common factor $-1$ for the whole 4-point contribution. The explicit
kernel, the resulting \eqI{32}--\eqI{33}, and the independent evidence
for (R2) (the field-equation route, the calibration with \eqI{131}, test
functions, a lattice Hessian and an operator-level check) are given in
Sec.~\ref{sec:s3:rule}.

\paragraph*{3-point contribution.}
For two cubic vertices the two outer legs sit at different vertices, $x$
and $x'$. Each outer-leg derivative produces $\partial^z\delta^D(z-x)$ at
its own vertex, and a partial integration in $z$ turns it into $-\partial$
acting on everything at that vertex; the inner-leg derivatives act only on
$i\Delta_A$. With the overall normalization $1$ (Sec.~\ref{sec:s4a:setup})
this gives~[cf.\ \eqI{45}]
\begin{equation}
-iM^2_{\rm 3pt}(x;x')=\sum_{I=1}^2V_I^{\alpha\beta}(x)\sum_{J=1}^2V_J^{\rho\sigma}(x')\;
i\big[{}_{\alpha\beta}\Delta_{\rho\sigma}\big](x;x')\,i\Delta_A(x;x'),
\label{eq:s2:eq45}
\end{equation}
read with the following rules:
\begin{itemize}
\item[(R1)] $\partial_1\to-\partial$, and likewise (R1$'$)
$\partial'_1\to-\partial'$, each acting on everything at its own vertex;
\item[(R2)] $\partial_2\to+\partial$, and likewise (R2$'$)
$\partial'_2\to+\partial'$, acting only on $i\Delta_A$.
\end{itemize}
This agrees with \eqI{45} and with the rule stated in the text of \PI{}
(``partially integrate any derivative that acts upon an outer leg''). The
effective action behind \eqref{eq:s2:eq45}, its normalization, the
derivation of \eqI{46} and the code checks are given in
Sec.~\ref{sec:s4a:setup}. Our Table~\ref{tab:s2:rules} compares the two rules.

\begin{table}[tb]
\caption{Derivative rules for the outer legs (definition sign). The 3-point
rule of \PI{} is correct; its 4-point rule differs by the treatment of the
second leg.}
\label{tab:s2:rules}
\begin{ruledtabular}
\begin{tabular}{>{\raggedright\arraybackslash}p{0.16\textwidth}>{\raggedright\arraybackslash}p{0.37\textwidth}>{\raggedright\arraybackslash}p{0.39\textwidth}}
 & 3-point, \eqref{eq:s2:eq45} & 4-point, \eqref{eq:s3:kernel}\\
\hline
outer legs & at different vertices, $x$ and $x'$ & at the same vertex $x$, with $\sIIdel$\\
leg 1 & partial integration: $-\partial$, on everything & partial integration: $-\partial$, on everything\\
leg 2 & partial integration: $-\partial'$, on everything & $+\partial$ (w.r.t.\ $x$) on $\sIIdel$, i.e.\ $-\partial'$\\
sign product & $(-)(-)=+$ & $(-)(+)=-$\\
\PI & same as here & \eqI{32}: $\partial_2\to+\partial'$, overall sign reversed\\
\end{tabular}
\end{ruledtabular}
\end{table}

\paragraph*{Origin of the overall sign.}
In the 4-point contribution, the printed result is reversed at its source,
the reading of $\partial_2$ in \eqI{32}. In the 3-point contribution the rule
is correct and \eqI{45}--\eqI{46} agree with the definition sign. However,
the step from \eqI{46} to \eqI{47}--\eqI{49} carries an overall sign slip in
all nine terms. The paper's subsequent 3-point expressions evaluate
left-hand sides that carry the printed signs of \eqI{47}--\eqI{49}
(Secs.~\ref{sec:s4a:graviton}--\ref{sec:s4a:log}). Both parts of the printed $-iM^2$ are therefore $-1$
times the quantity defined by \eqI{168}, and their relative sign is
consistent. Several independent lines of evidence support the definition
sign: Lemma~A (two proofs), the field-equation route, the counterterm vertex
\eqI{131}, test functions, the lattice Hessian and positivity, and the
consistency of the 3-point steps.

\subsection{Summary of this section}
\label{sec:s2:summary}

Everything else in this part of \PI{} is confirmed as printed:
\begin{itemize}
\item \eqI{4}, \eqI{5} and the bracket of \eqI{6} (exact, general $D$);
\item Table~I, \eqI{8}, and rows 1, 3 and 4 of Table~II;
\item the structure and normalization of \eqI{31} and \eqI{45}, and the four
terms of \eqI{46};
\item \eqI{52}, obtained along three routes.
\end{itemize}
The differences are collected in our Table~\ref{tab:s2:summary}. The
labels C$n$ refer to the errata list of
Sec.~\ref{sec:s0s1s7s8:summary} (our Table~\ref{tab:s0s1s7s8:errata}), and
S1 denotes the first substantive finding, the overall sign of the printed
self-mass.

\begin{table}[tb]
\caption{Differences from the printed equations in the part of \PI{}
covered by this section. ``Typo'' means that the printed value did not
enter the paper's computation. S1 is the 4-point origin of the overall
sign, Eq.~\eqref{eq:s0s1s7s8:signs}; it is classified as substantive
because the sign of the nonlocal part of the self-mass is physical. Its
consequences are discussed in the following sections.}
\label{tab:s2:summary}
\begin{ruledtabular}
\begin{tabular}{>{\raggedright\arraybackslash}p{0.04\textwidth}>{\raggedright\arraybackslash}p{0.14\textwidth}>{\raggedright\arraybackslash}p{0.25\textwidth}>{\raggedright\arraybackslash}p{0.25\textwidth}>{\raggedright\arraybackslash}p{0.24\textwidth}}
ID & Where & Printed & Correct & Class / evidence\\
\hline
C1 & \eqI{1} & $-\tfrac12\partial_\mu\phi\,\partial_\mu\phi\,g^{\mu\nu}\sqrt{-g}$ & $-\tfrac12\partial_\mu\phi\,\partial_\nu\phi\,g^{\mu\nu}\sqrt{-g}$ & typo; free index\\
C2 & \eqI{6}, left side & $-\tfrac12\partial_\mu\phi\,\partial_\nu g^{\mu\nu}\sqrt{-g}$ & $-\tfrac12\partial_\mu\phi\,\partial_\nu\phi\,g^{\mu\nu}\sqrt{-g}$ & typo; right side correct\\
C3 & \eqI{7}, left side & $-\tfrac12\kappa a^{D-2}\partial_\alpha\phi\,\partial_\beta\phi\,h^{\alpha\beta}$ & $+\tfrac12\kappa a^{D-2}\partial_\alpha\phi\,\partial_\beta\phi\,h^{\alpha\beta}$ & typo; $V_1$ correct and used in \eqI{46}\\
C4 & \eqI{14}, first coefficient & $\Gamma(\tfrac D2-1)/(\tfrac D2-1)$ & $\Gamma(\tfrac D2-1)$ & typo; required by \eqI{26}, the recursion and \eqI{52}\\
C5 & Table~II, $U_2$ & $-\tfrac i4\kappa^2a^{D-2}\eta^{\alpha\rho}\eta^{\beta\sigma}\sIIdd$ & $+\tfrac i4\kappa^2a^{D-2}\eta^{\alpha\rho}\eta^{\beta\sigma}\sIIdd$ & typo; \eqI{33} uses $+\tfrac i4$\\
S1 & rule of \eqI{32} & $\partial_2\to+\partial'$ on $\sIIdel$ & $\partial_2\to+\partial=-\partial'$ (Lemma~A) & substantive: 4-point origin of the overall sign (physical for the nonlocal part)\\
\end{tabular}
\end{ruledtabular}
\end{table}

\providecommand{\sthreedel}{\delta^D(x-x')}
\providecommand{\sthreedelf}{\delta^4(x-x')}
\providecommand{\sthreeM}{M^2_{\rm 4pt}}
\providecommand{\sthreeK}{\mathcal K}
\providecommand{\sthreec}[1]{i\bigl[#1\bigr](x;x)}
\section{The 4-point contribution}
\label{sec:s3}

This section corresponds to Sec.~III\,A of \PI{} \cite{KW07}, Eqs.~\eqI{31}--\eqI{44}.
The 4-point vertices $U_I$ of Table~II of \PI{}, with the graviton loop closing on itself
at a single point (Fig.~1 of \PI), give a purely local contribution to the self-mass-squared.
Three steps are needed: (i) the kernel generated by the vertices, which fixes how the
derivatives of Table~II have to be read; (ii) the coincidence limits of the three scalar
propagators and of the four contractions of the coincident graviton propagator; (iii) the
assembly into the final local form. The outcome is the following. The structure and
normalization of \eqI{31}, all $D$-dependent coefficients, the coincidence limits
\eqI{34}--\eqI{39} and the algebra \eqI{41}$\to$\eqI{44} are confirmed. There are two
differences. First, the rule by which \PI{} reads the derivative ``$\partial_2$'' introduces a
factor $-1$ relative to the paper's own definition \eqI{168}; this is the 4-point origin of the
overall sign of the printed self-mass. Second, the finite part of \eqI{40} has the wrong sign,
which produces extra finite terms in \eqI{41} and \eqI{44}; these have counterterm form and are
absorbed. Full details are in Ch.~2 of the derivation notes and in the scripts of
\nolinkurl{code/02_four_point/} \cite{ref:repo}.

\paragraph*{Conventions.}
We write $\partial_\mu\equiv\partial/\partial x^\mu$, $\partial'_\mu\equiv\partial/\partial x'^\mu$, and indices
are raised and lowered with $\eta_{\mu\nu}={\rm diag}(-1,1,\dots,1)$. We use conformal
coordinates with $a=-1/(H\eta)$, so $\partial_0a=Ha^2$. The dimension $D$ is kept general
throughout, unless a step says otherwise. We abbreviate
\begin{equation}\label{eq:s3:abbrev}
k\equiv\frac{H^{D-2}}{(4\pi)^{D/2}}\frac{\Gamma(D-1)}{\Gamma(\frac D2)},\qquad
h\equiv\frac{H^2}{16\pi^2},\qquad
\pihat\equiv-\pi\cot\Bigl(\frac{\pi D}{2}\Bigr),
\end{equation}
where $k$ is the constant of \eqI{13}; as $D\to4$, $k\to2h=H^2/8\pi^2$. (The constant $h$ is not to be confused with the graviton
trace $h=\eta^{\mu\nu}h_{\mu\nu}$ of Sec.~\ref{sec:s2}, which does not occur in this section.)
Signs matter in this section. The \emph{definition sign} is the one fixed by \PI's own
definition \eqI{168} through Eq.~\eqref{eq:s2:defM} (Sec.~\ref{sec:s2:definition}); the \emph{paper
convention} is the sign with which \PI{} prints \eqI{32}, \eqI{33}, \eqI{41} and \eqI{44}. All kernels
derived below are in the definition sign, and we show that the paper convention is $-1$ times the
definition sign [Eq.~\eqref{eq:s0s1s7s8:signs}].

\subsection{The kernel of the 4-point vertices and the derivative rule}
\label{sec:s3:rule}

\paragraph*{Vertices.}
The 4-point vertex operators $U_I$, derived from the expansion \eqI{6} in the normalization
\eqref{eq:s2:norm}, are listed in our Table~\ref{tab:s2:vertices}. A factor $\frac12$ is extracted for the
two identical scalars but not for the two identical gravitons, so that no symmetry factor is needed
(Sec.~\ref{sec:s2:vertices}). Table~II of \PI{} prints $U_2$ with $-\frac i4$. The printed \eqI{33}
carries $+\frac i4$ in the corresponding term, so the derived value $+\frac i4$ was the one used in the
calculation and the printed $-\frac i4$ is a typographical error.

\paragraph*{Effective action.}
At order $\kappa^2$ and quadratic order in the background $\Phi$, the 1PI contributions are
(i) a single 4-point vertex with a closed graviton loop, (ii) two 3-point vertices
(Sec.~\ref{sec:s4a}), and (iii) counterterms. A 3-point vertex with a graviton tadpole is not 1PI. For (i),
$i\Gamma_{\rm 4pt}[\Phi]=i\langle S_4[\Phi,h]\rangle$, with
$\langle T\,h_{\alpha\beta}(x)h_{\rho\sigma}(x')\rangle=i[{}_{\alpha\beta}\Delta_{\rho\sigma}](x;x')$.
In dimensional regularization the coincidence limit
\begin{equation}\label{eq:s3:X}
X_{\alpha\beta\rho\sigma}(x)\equiv i\bigl[{}_{\alpha\beta}\Delta_{\rho\sigma}\bigr](x;x)
\end{equation}
is a smooth function of $x$. It is symmetric within each index pair, and it is symmetric under
exchange of the pairs, $X_{\alpha\beta\rho\sigma}=X_{\rho\sigma\alpha\beta}$; the latter can
be read off term by term from the tensor factors \eqI{22}--\eqI{24}. Writing
$U_I^{\alpha\beta\rho\sigma}X_{\alpha\beta\rho\sigma}\equiv\tilde K_I^{\mu\nu}(z)\,\partial_{1\mu}\partial_{2\nu}$,
we obtain
\begin{equation}\label{eq:s3:S4}
\langle S_4\rangle=\int d^Dz\,\frac1{2i}\sum_I\tilde K_I^{\mu\nu}(z)\,\partial_\mu\Phi\,\partial_\nu\Phi .
\end{equation}

\paragraph*{Second functional derivatives.}
We use Lemma~A of Sec.~\ref{sec:s2:rules}, Eq.~\eqref{eq:s2:lemmaA}: for a quadratic form
$F=\int d^Dz\,K^{\mu\nu}\partial_\mu\phi\,\partial_\nu\phi$, one of the two derivatives is integrated by parts
onto the whole expression, where it becomes $-\partial$, while the other acts on $\sthreedel$ as $+\partial$
\emph{with respect to $x$}, i.e.\ as $-\partial'$, because $\partial'_\mu\sthreedel=-\partial_\mu\sthreedel$.
The sign is fixed by a positivity test (Sec.~\ref{sec:s2:rules}).

\paragraph*{The kernel and the reading rule.}
Applying Eq.~\eqref{eq:s2:defM} and Lemma~A to Eq.~\eqref{eq:s3:S4},
\begin{equation}\label{eq:s3:kernel}
-i\sthreeM(x;x')=i\,\frac{\delta^2\langle S_4\rangle}{\delta\Phi(x)\,\delta\Phi(x')}
=-\sum_I\partial_\mu\Bigl[\tilde K^{\mu\nu}_{I,s}(x)\,\partial_\nu\sthreedel\Bigr],
\qquad \tilde K_s\equiv\tfrac12\bigl(\tilde K+\tilde K^T\bigr).
\end{equation}
The scalar factor $\frac12$ of Eq.~\eqref{eq:s2:norm} cancels against the factor 2 from the two
ways of attaching the external legs, so the normalization is 1, and Eq.~\eqref{eq:s3:kernel}
takes the form of \eqI{31}:
\begin{equation}\label{eq:s3:I31}
-i\sthreeM(x;x')=\sum_{I=1}^4U_I^{\alpha\beta\rho\sigma}\;
i\bigl[{}_{\alpha\beta}\Delta_{\rho\sigma}\bigr](x;x)\;\sthreedel .
\end{equation}
This notation is given its meaning by the reading rule that follows from
Eq.~\eqref{eq:s3:kernel}:
\begin{itemize}
\item[(R1)] $\partial_{1\mu}\to-\partial_\mu$, acting on everything to its right ($a^{D-2}$, the
coincident propagator and the factor containing $\delta^D$);
\item[(R2)] $\partial_{2\nu}\to+\partial_\nu$, the derivative with respect to $x$, acting only on
$\sthreedel$; equivalently $\partial_{2\nu}\to-\partial'_\nu$;
\item[(R3)] if $\tilde K_I$ is not symmetric, it is symmetrized. This is not needed here:
$\tilde K_{1,2}$ are proportional to $\eta^{\mu\nu}$, $\tilde K_3$ is symmetric by the
within-pair symmetry, and $\tilde K_4$ is symmetric by the pair-exchange symmetry of
$X_{\alpha\beta\rho\sigma}$.
\end{itemize}
Rule (R1) is the paper's rule that ``$\partial_1$'' is partially integrated back onto the
entire contribution. For ``$\partial_2$'', however, \PI{} states that it acts upon $x'$ and
uses $\partial_2\to+\partial'$ in \eqI{32}. This is exactly the opposite of (R2). Since each
$U_I$ contains exactly one $\partial_1$ and one $\partial_2$, the difference is a common factor
$-1$ for the whole 4-point contribution.

For the first vertex we obtain, in the definition sign,
\begin{equation}\label{eq:s3:I32}
-\tfrac i8\kappa^2a^{D-2}\eta^{\alpha\beta}\eta^{\rho\sigma}\,\partial_1\!\cdot\partial_2
\times i\bigl[{}_{\alpha\beta}\Delta_{\rho\sigma}\bigr](x;x)\,\sthreedel
\ \Rightarrow\ -\tfrac i8\kappa^2\,\partial^\mu\Bigl\{a^{D-2}\,
\sthreec{{}^\alpha_\alpha\Delta^\rho_\rho}\,\partial'_\mu\sthreedel\Bigr\},
\end{equation}
whereas the right-hand side of the printed \eqI{32} carries $+\frac i8$. The contractions for
$U_2$, $U_3$, $U_4$ are $\eta^{\alpha\rho}\eta^{\beta\sigma}X_{\alpha\beta\rho\sigma}
=\sthreec{{}^{\alpha\beta}\Delta_{\alpha\beta}}$,
$\eta^{\alpha\beta}X_{\alpha\beta}{}^{\rho\sigma}=\sthreec{{}^\alpha_\alpha\Delta^{\rho\sigma}}$
and $X^\alpha{}_\rho{}^{\rho\sigma}=\sthreec{{}^\alpha{}_\rho\Delta^{\rho\sigma}}$. Collecting
all four terms (definition sign, general $D$; cf.\ \eqI{33}):
\begin{equation}\label{eq:s3:I33}
\begin{aligned}
-i\sthreeM(x;x')={}&+\tfrac i8\kappa^2\partial^\mu\bigl\{a^{D-2}\,\sthreec{{}^\alpha_\alpha\Delta^\rho_\rho}\,
\partial_\mu\sthreedel\bigr\}
-\tfrac i4\kappa^2\partial^\mu\bigl\{a^{D-2}\,\sthreec{{}^{\alpha\beta}\Delta_{\alpha\beta}}\,
\partial_\mu\sthreedel\bigr\}\\
&-\tfrac i2\kappa^2\partial_\rho\bigl\{a^{D-2}\,\sthreec{{}^\alpha_\alpha\Delta^{\rho\sigma}}\,
\partial_\sigma\sthreedel\bigr\}
+i\kappa^2\partial^\alpha\bigl\{a^{D-2}\,\sthreec{{}_{\alpha\rho}\Delta^{\rho\sigma}}\,
\partial_\sigma\sthreedel\bigr\}.
\end{aligned}
\end{equation}
The printed \eqI{33} has the coefficients $(-\frac i8,+\frac i4,+\frac i2,-i)$. Magnitudes,
tensor contractions and relative signs agree; the overall sign is reversed. In compact form,
\begin{equation}\label{eq:s3:Kcal}
\begin{gathered}
-i\sthreeM=-\partial_\mu\bigl[\sthreeK^{\mu\nu}(x)\,\partial_\nu\sthreedel\bigr],\qquad
\sthreeK^{\mu\nu}=\kappa^2a^{D-2}\Bigl\{-\tfrac i8\eta^{\mu\nu}X_1+\tfrac i4\eta^{\mu\nu}X_2
+\tfrac i2X_3^{\mu\nu}-i\,X_4^{\mu\nu}\Bigr\},\\
X_1=\sthreec{{}^\alpha_\alpha\Delta^\rho_\rho},\quad
X_2=\sthreec{{}^{\alpha\beta}\Delta_{\alpha\beta}},\quad
X_3^{\mu\nu}=\sthreec{{}^\alpha_\alpha\Delta^{\mu\nu}},\quad
X_4^{\mu\nu}=\eta^{\mu\alpha}\,\sthreec{{}_{\alpha\rho}\Delta^{\rho\nu}} .
\end{gathered}
\end{equation}
Physically, $\sthreeK^{\mu\nu}$ is the correction to the kinetic term generated by
$\langle\mathcal L_4\rangle$.

\paragraph*{Independent evidence for the sign.}
The sign of Eq.~\eqref{eq:s3:I33} is supported by several independent arguments.
\begin{enumerate}
\item \emph{Field-equation route.} This route uses only the Euler--Lagrange equation and not
the lemma. From $\Gamma\supset\int\frac1{2i}\sthreeK^{\mu\nu}\partial_\mu\Phi\partial_\nu\Phi$
($\sthreeK$ symmetric),
$\delta\Gamma/\delta\Phi=\partial_\mu(a^{D-2}\partial^\mu\Phi)+i\,\partial_\mu(\sthreeK^{\mu\nu}\partial_\nu\Phi)+\dots$.
Comparison with \eqI{168} gives $\int M^2\Phi=-i\,\partial_\mu(\sthreeK^{\mu\nu}\partial_\nu\Phi)$,
which is Eq.~\eqref{eq:s3:Kcal}.
\item \emph{The paper's own counterterm vertex \eqI{131}.} Like the 4-point vertices, the counterterm
$-\frac12\alpha_2\kappa^2H^2\partial_\mu\phi\,\partial^\mu\phi\,a^{D-2}$ has one derivative on each leg, so it
calibrates the derivative rule. Lemma~A reproduces the printed vertex
$i\alpha_2\kappa^2H^2a^D\square\,\sthreedel$ [Eq.~\eqref{eq:s2:calib131}]. Written as a vertex operator in the
normalization of Eq.~\eqref{eq:s2:norm}, the same term is
$U_{\rm ct}=-i\alpha_2\kappa^2H^2a^{D-2}\partial_1\!\cdot\partial_2$. The rules (R1)--(R2) give
$+i\alpha_2\kappa^2H^2a^D\square\,\sthreedel$, while the paper's 4-point rule
($\partial_2\to+\partial'$) gives $-i\alpha_2\kappa^2H^2a^D\square\,\sthreedel$. Hence \eqI{131} and
the rule of \eqI{32} cannot both hold in a single convention, and \eqI{131} is the one that
agrees with \eqI{168}. The vertex \eqI{130} is also reproduced, but it has two derivatives on
each leg and is insensitive to the rule.
\item \emph{Test functions, lattice Hessian and operator-level check} (script
\mbox{\nolinkurl{check_4pt.py}} \cite{ref:repo}). (a) For $D=2$, a deliberately non-symmetric
polynomial $K^{\mu\nu}$ and two Gaussian-type test functions, the bilinear form is
$\partial^2F/\partial\epsilon_1\partial\epsilon_2=\frac{97\sqrt6\,\pi}{54}e^{-1/2}\approx8.384$;
the kernel \eqref{eq:s2:lemmaA} gives the same value and the opposite-sign kernel gives
$-8.384$. (b) On a periodic one-dimensional lattice with $N=7$ sites and
$S=\sum_zc_z(\phi_{z+1}-\phi_z)^2$, the Hessian is exactly $-\nabla^-(2c\,\nabla^+)$, the lattice
version of $-\partial(2c\,\partial\,\delta)$; for $c=1$ its eigenvalues are
$0,\,1.506,\,4.890,\,7.604$. (c) For $D=3$, a general pair-symmetric $X_{\alpha\beta\rho\sigma}(x)$
built from 21 independent functions and a general $a(\eta)$, the kernel
\eqref{eq:s3:kernel} equals Eq.~\eqref{eq:s3:I33} and is exactly $-1$ times the printed \eqI{33};
the paper's rule combined with $U_2=+\frac i4$ reproduces the printed \eqI{33} exactly, while
with the printed $U_2=-\frac i4$ it does not. (d) The lemma reproduces \eqI{131}; the paper's
4-point rule gives its negative.
\end{enumerate}
The flat-space reduction cannot discriminate between the two signs, because the coincident
flat-space graviton propagator vanishes in dimensional regularization.

\paragraph*{Relation to the 3-point part.}
The paper's derivative rule for the 3-point diagram (``partially integrate any derivative that
acts upon an outer leg'') is correct: each outer-leg derivative becomes $-\partial$ at its own
vertex, as in the lemma. However, as shown in Sec.~\ref{sec:s4a}, the printed \eqI{47}--\eqI{49}
are $-1$ times what follows from the paper's own \eqI{46}. Both parts of the printed result are
therefore reversed. Their relative sign is consistent, the local cancellation between \eqI{44}
and \eqI{67} holds in either convention, and \emph{the printed self-mass is $-1$ times the one
defined by \eqI{168}}. Because the nonlocal part of the self-mass has no counterterm freedom,
this overall sign is physical for the nonlocal part, which is $-$(corrected Table~X). It reverses
the sign of the one-loop mode-function correction of \PII; but, as shown in Sec.~\ref{sec:s6}, the
non-falling total of that correction is zero at late times, so only decaying terms change sign and
the late-time conclusion of \PII{} is unchanged.

\subsection{Coincidence limits}
\label{sec:s3:coinc}

These limits (cf.\ \eqI{34}--\eqI{40}; the paper takes them from Ref.~\cite{ref:Tsamis2006})
are independent of the vertices and of the derivative rule, so the sign issue of
Sec.~\ref{sec:s3:rule} does not enter here.

\paragraph*{Scalar propagators.}
As $x'\to x$ one has $y(x;x')\to0$ and $\ln(aa')\to2\ln a$. The propagators \eqI{13}--\eqI{14}
and \eqI{26}--\eqI{28} are sums of powers $(y/4)^{p(D)}$ with exponents of the form
$1-\frac D2$, $2-\frac D2$, $n\ge1$ and $n-\frac D2+2$ ($n\ge0$). All of them are positive for
${\rm Re}\,D<2$, where the limit exists literally and may be taken term by term (the series
converge uniformly on $0<y\le y_0<4$); the result is then continued analytically in $D$. Only
the terms with $p\equiv0$ survive.\footnote{The first coefficient of $A(y)$ in \eqI{14} is
misprinted (it should be $\Gamma(\frac D2-1)$; erratum C4); since that term vanishes in the coincidence
limit, \eqI{34} is not affected.} This gives (cf.\ \eqI{34}--\eqI{36})
\begin{gather}
\lim_{x'\to x}i\Delta_A(x;x')=k\bigl\{\pihat+2\ln a\bigr\}
=k\Bigl\{-\frac{2}{D-4}+2\ln a+O(D-4)\Bigr\},\label{eq:s3:I34}\\
b\equiv\lim_{x'\to x}i\Delta_B(x;x')=-\frac{k}{D-2}\ \longrightarrow\ -h,\qquad
c\equiv\lim_{x'\to x}i\Delta_C(x;x')=\frac{k}{(D-2)(D-3)}\ \longrightarrow\ +h,\label{eq:s3:I3536}
\end{gather}
where $\pihat=-2/(D-4)+\frac{\pi^2}{6}(D-4)+O\bigl((D-4)^3\bigr)$. The $B$- and $C$-type limits are
finite at $D=4$ and are related by
\begin{equation}\label{eq:s3:cb}
c=-\frac{b}{D-3}.
\end{equation}
As an independent check, $b$ and $c$ follow from Gauss's theorem
${}_2F_1(a_1,a_2;c_1;1)=\Gamma(c_1)\Gamma(c_1-a_1-a_2)/[\Gamma(c_1-a_1)\Gamma(c_1-a_2)]$ applied to
the hypergeometric forms \eqI{29}--\eqI{30} in the region ${\rm Re}(c_1-a_1-a_2)=1-\frac D2>0$;
and at $D=0.6$, where no
continuation is needed, the series approach Eqs.~\eqref{eq:s3:I34}--\eqref{eq:s3:I3536} like
$(y/4)^{0.7}$ (script \mbox{\nolinkurl{check_coincidence.py}}).

\paragraph*{Tensor algebra.}
Let $u_\mu\equiv\delta^0_\mu$, so that $u^\mu=-\delta^\mu_0$, $u\cdot u=-1$,
$\bar\eta_{\mu\nu}=\eta_{\mu\nu}+u_\mu u_\nu$, $\bar\eta_{\mu\nu}u^\nu=0$,
$\eta^{\mu\nu}\bar\eta_{\mu\nu}=\bar\delta^\mu_\mu=D-1$, $u_\alpha u^\sigma=-\delta^0_\alpha\delta^\sigma_0$
and $u^\rho u^\sigma=\delta^\rho_0\delta^\sigma_0$. The tensor factors \eqI{22}--\eqI{24} of the
graviton propagator, with the symmetrizations written out, are
\begin{equation}\label{eq:s3:T}
\begin{aligned}
T^A_{\alpha\beta\rho\sigma}&=\bar\eta_{\alpha\rho}\bar\eta_{\sigma\beta}+\bar\eta_{\alpha\sigma}\bar\eta_{\rho\beta}
-\tfrac{2}{D-3}\bar\eta_{\alpha\beta}\bar\eta_{\rho\sigma},\\
T^B_{\alpha\beta\rho\sigma}&=-\bigl[u_\alpha\bar\eta_{\beta\rho}u_\sigma+u_\beta\bar\eta_{\alpha\rho}u_\sigma
+u_\alpha\bar\eta_{\beta\sigma}u_\rho+u_\beta\bar\eta_{\alpha\sigma}u_\rho\bigr],\\
T^C_{\alpha\beta\rho\sigma}&=\tfrac{2}{(D-2)(D-3)}f_{\alpha\beta}f_{\rho\sigma},\qquad
f_{\alpha\beta}\equiv(D-3)u_\alpha u_\beta+\bar\eta_{\alpha\beta}.
\end{aligned}
\end{equation}
Since these are constant tensors, the coincidence limit acts only on the scalar propagators:
\begin{equation}\label{eq:s3:Xdecomp}
X_{\alpha\beta\rho\sigma}=T^A_{\alpha\beta\rho\sigma}\,i\Delta_A(x;x)+T^B_{\alpha\beta\rho\sigma}\,b
+T^C_{\alpha\beta\rho\sigma}\,c .
\end{equation}
The four contractions needed in Eq.~\eqref{eq:s3:I33} then follow by elementary algebra. For
example, in the double trace the $T^B$ terms vanish because every term contains a $u$ contracted
with a $\bar\eta$, and $\eta^{\alpha\beta}f_{\alpha\beta}=2$. Exactly in $D$ we find
(cf.\ \eqI{37}--\eqI{40})
\begin{align}
\sthreec{{}^\alpha_\alpha\Delta^\rho_\rho}&=-\frac{4(D-1)}{D-3}\,i\Delta_A(x;x)+\frac{8}{(D-2)(D-3)}\,c,
\label{eq:s3:I37}\\
\sthreec{{}^{\alpha\beta}\Delta_{\alpha\beta}}&=\frac{(D-1)(D^2-3D-2)}{D-3}\,i\Delta_A(x;x)
+2(D-1)\,b+\frac{2(D^2-5D+8)}{(D-2)(D-3)}\,c,\label{eq:s3:I38}\\
\sthreec{{}^\alpha_\alpha\Delta^{\rho\sigma}}&=-\frac{4}{D-3}\,\bar\eta^{\rho\sigma}\,i\Delta_A(x;x)
+\frac{4}{(D-2)(D-3)}\bigl[(D-3)\delta^\rho_0\delta^\sigma_0+\bar\eta^{\rho\sigma}\bigr]\,c,\label{eq:s3:I39}\\
\sthreec{{}_{\alpha\rho}\Delta^{\rho\sigma}}&=\frac{D^2-3D-2}{D-3}\,\bar\delta^\sigma_\alpha\,i\Delta_A(x;x)
+\bigl[(D-1)\delta^0_\alpha\delta^\sigma_0+\bar\delta^\sigma_\alpha\bigr]\,b\notag\\
&\quad+\frac{2}{(D-2)(D-3)}\bigl[(D-3)^2\delta^0_\alpha\delta^\sigma_0+\bar\delta^\sigma_\alpha\bigr]\,c.
\label{eq:s3:I40}
\end{align}
Two identities hold by definition and are satisfied identically in $D$: the trace
$\alpha=\sigma$ of Eq.~\eqref{eq:s3:I40} equals Eq.~\eqref{eq:s3:I38} (this uses only the
within-pair symmetry), and $\eta_{\rho\sigma}$ times Eq.~\eqref{eq:s3:I39} equals
Eq.~\eqref{eq:s3:I37}. All four contractions were also verified component by component, with
explicit $D^4$-component arrays for $D=4,5,6,7$ and symbolic $i\Delta_A$, $b$, $c$.

\paragraph*{Finite parts.}
With Eq.~\eqref{eq:s3:cb} the $B$- and $C$-type (finite) parts become, exactly in $D$,
\begin{equation}\label{eq:s3:finiteD}
\begin{aligned}
\text{\eqref{eq:s3:I37}}:&\quad -\frac{8\,b}{(D-2)(D-3)^2},\qquad
\text{\eqref{eq:s3:I38}}:\quad 2b\Bigl[(D-1)-\frac{D^2-5D+8}{(D-2)(D-3)^2}\Bigr],\\
\text{\eqref{eq:s3:I39}}:&\quad -\frac{4\,b}{(D-2)(D-3)^2}\bigl[(D-3)\delta^\rho_0\delta^\sigma_0
+\bar\eta^{\rho\sigma}\bigr],\\
\text{\eqref{eq:s3:I40}}:&\quad \frac{D(D-3)}{D-2}\,b\,\delta^0_\alpha\delta^\sigma_0
+b\Bigl[1-\frac{2}{(D-2)(D-3)^2}\Bigr]\bar\delta^\sigma_\alpha ,
\end{aligned}
\end{equation}
and at $D=4$ ($b\to-h$)
\begin{equation}\label{eq:s3:finite4}
+4h,\qquad -2h,\qquad \bigl[2\delta^\rho_0\delta^\sigma_0+2\bar\eta^{\rho\sigma}\bigr]h,\qquad
-2\,\delta^0_\alpha\delta^\sigma_0\,h ,
\end{equation}
respectively; the coefficient of $\bar\delta^\sigma_\alpha$ in the last one vanishes. Setting
$D=4$ in the finite parts is legitimate as long as they are not multiplied by a pole: $b$ and
$c$ enter Eq.~\eqref{eq:s3:I33} linearly and without a divergent prefactor, so the change is
$O(D-4)$. The $D$-dependent coefficients of the divergent $i\Delta_A(x;x)$, on the other hand,
must be kept exactly, as the paper does.

The comparison with \PI{} is given in our Table~\ref{tab:s3:coinc}. Everything agrees except the
finite part of \eqI{40}, where \PI{} prints $+2\,\delta^0_\alpha\delta^\sigma_0\,h$. The printed
value is inconsistent with the paper's own \eqI{38}: the trace identity, which holds by
definition, would require the trace of the finite part of \eqI{40}, $+2h$ as printed, to equal
the finite part of \eqI{38}, which is $-2h$ as printed. The discrepancy cannot be traced to the
sign of $T^B$ or $T^C$; for example, reversing $T^B$ would make the finite part of \eqI{38}
equal to $+10h$.

\begin{table}[t]
\caption{Coincidence limits, cf.\ \eqI{34}--\eqI{40}. Finite parts are quoted at $D=4$ in
units of $h=H^2/16\pi^2$; the coefficients of $i\Delta_A(x;x)$ are exact in $D$.}
\label{tab:s3:coinc}
\begin{ruledtabular}
\begin{tabular}{llll}
Equation & This work & \PI{} prints & Status\\
\hline
\eqI{34} & $k\{\pihat+2\ln a\}$ & same & agrees\\
\eqI{35}, \eqI{36} & $b=-\frac{k}{D-2}\to-h$, \ $c=\frac{k}{(D-2)(D-3)}\to+h$ & same & agree\\
\eqI{37} & $-\frac{4(D-1)}{D-3}\,i\Delta_A$; \ finite $+4$ & same & agrees\\
\eqI{38} & $\frac{(D-1)(D^2-3D-2)}{D-3}\,i\Delta_A$; \ finite $-2$ & same & agrees\\
\eqI{39} & $-\frac4{D-3}\bar\eta^{\rho\sigma}i\Delta_A$; \ finite $2\delta^\rho_0\delta^\sigma_0+2\bar\eta^{\rho\sigma}$
 & finite $2\delta^\rho_0\delta^\sigma_0+2\bar\eta_{\rho\sigma}$ & index placement only\\
\eqI{40}, $i\Delta_A$ & $\frac{D^2-3D-2}{D-3}\bar\delta^\sigma_\alpha\,i\Delta_A$ & same & agrees\\
\eqI{40}, finite & $-2\,\delta^0_\alpha\delta^\sigma_0$ & $+2\,\delta^0_\alpha\delta^\sigma_0$ & sign error\\
\end{tabular}
\end{ruledtabular}
\end{table}

\subsection{The 4-point result}
\label{sec:s3:result}

\paragraph*{Assembly in general $D$.}
Write $\Lambda\equiv i\Delta_A(x;x)=k(\pihat+2\ln a)$ and split each contraction into its
$\Lambda$ part and its finite part. Inserting Eqs.~\eqref{eq:s3:I37}--\eqref{eq:s3:I40} into
Eq.~\eqref{eq:s3:Kcal} (with $\eta^{\mu\alpha}\bar\delta^\nu_\alpha=\bar\eta^{\mu\nu}$), the
coefficients of $\eta^{\mu\nu}$ and $\bar\eta^{\mu\nu}$ in the $\Lambda$ part are
\begin{equation}\label{eq:s3:coefA}
\tfrac i8\cdot\tfrac{4(D-1)}{D-3}+\tfrac i4\cdot\tfrac{(D-1)(D^2-3D-2)}{D-3}=\tfrac i4D(D-1),\qquad
\tfrac i2\cdot\Bigl(-\tfrac4{D-3}\Bigr)-i\,\tfrac{D^2-3D-2}{D-3}=-iD .
\end{equation}
The factors $1/(D-3)$ of the graviton propagator cancel exactly. The finite parts
\eqref{eq:s3:finiteD} combine, with $\eta^{\mu\nu}=\bar\eta^{\mu\nu}-u^\mu u^\nu$, into a
multiple of $D-4$. Altogether $\sthreeK=\sthreeK_A+\sthreeK_F$ with
\begin{equation}\label{eq:s3:KAKF}
\sthreeK^{\mu\nu}_A=i\kappa^2a^{D-2}\,\Lambda\Bigl[\tfrac14D(D-1)\,\eta^{\mu\nu}-D\,\bar\eta^{\mu\nu}\Bigr],
\qquad
\sthreeK^{\mu\nu}_F=i\kappa^2a^{D-2}\,b\,\frac{D-4}{2(D-3)}\Bigl[(D-2)\,\bar\eta^{\mu\nu}+D\,u^\mu u^\nu\Bigr].
\end{equation}
Thus, with the correct \eqI{40}, the finite $B$- and $C$-type contributions are exactly
$O(D-4)$. For a function $g(\eta)$ of time only one has
$\partial^\mu(g\,\partial_\mu\delta)=-\partial_0(g\,\partial_0\delta)+g\nabla^2\delta$,
$\partial_\mu(g\,\bar\eta^{\mu\nu}\partial_\nu\delta)=g\nabla^2\delta$ and
$\partial_\mu(g\,u^\mu u^\nu\partial_\nu\delta)=\partial_0(g\,\partial_0\delta)$, so the kernel is,
exactly in $D$ and in the definition sign,
\begin{equation}\label{eq:s3:kernelD}
\begin{aligned}
-i\sthreeM={}&-i\kappa^2\Bigl\{\tfrac14D(D-1)\,\partial^\mu\bigl(a^{D-2}\Lambda\,\partial_\mu\sthreedel\bigr)
-D\,a^{D-2}\Lambda\,\nabla^2\sthreedel\Bigr\}\\
&-i\kappa^2\,b\,\frac{D-4}{2(D-3)}\Bigl\{D\,\partial_0\bigl(a^{D-2}\partial_0\sthreedel\bigr)
+(D-2)\,a^{D-2}\nabla^2\sthreedel\Bigr\}.
\end{aligned}
\end{equation}

\paragraph*{The analogue of \eqI{41}.}
In Eq.~\eqref{eq:s3:kernelD} the $k\pihat$ terms carry a pole and are kept exactly in $D$. The
$2k\ln a$ terms are finite, so setting $D=4$ in them ($2k\to H^2/4\pi^2$, $a^{D-2}\to a^2$,
$\delta^D\to\delta^4$) changes them only by $O(D-4)$; the paper does the same. The $b$ term is
explicitly $O(D-4)$, is not multiplied by a pole, and drops out. In the definition sign,
\begin{equation}\label{eq:s3:I41}
\begin{aligned}
-i\sthreeM(x;x')={}&-i\kappa^2k\,\pihat\,\Bigl\{\tfrac14D(D-1)\,\partial^\mu\bigl(a^{D-2}\partial_\mu\sthreedel\bigr)
-D\,a^{D-2}\nabla^2\sthreedel\Bigr\}\\
&-\frac{i\kappa^2H^2}{4\pi^2}\Bigl\{3\,\partial^\mu\bigl(a^2\ln(a)\,\partial_\mu\sthreedelf\bigr)
-4\ln(a)\,a^2\nabla^2\sthreedelf\Bigr\}+O(D-4).
\end{aligned}
\end{equation}
\PI{} prints \eqI{41} as
\begin{equation}\label{eq:s3:I41printed}
\begin{aligned}
-i\sthreeM\big|_{\rm printed}={}&+i\kappa^2k\,\pihat\,\Bigl\{\tfrac14D(D-1)\,\partial^\mu\bigl(a^{D-2}\partial_\mu\sthreedel\bigr)
-D\,a^{D-2}\nabla^2\sthreedel\Bigr\}\\
&+\frac{i\kappa^2H^2}{4\pi^2}\Bigl\{3\,\partial^\mu\bigl(a^2\ln(a)\,\partial_\mu\sthreedelf\bigr)
-4\ln(a)\,a^2\nabla^2\sthreedelf\\
&\hphantom{{}+\frac{i\kappa^2H^2}{4\pi^2}\Bigl\{}
-\partial^\mu\bigl(a^2\partial_\mu\sthreedelf\bigr)+a^2\nabla^2\sthreedelf\Bigr\}+O(D-4),
\end{aligned}
\end{equation}
where we have read the unparenthesized factor $-\pi\cot(\frac\pi2D)$ of the printed equation as
multiplying the prefactor, as \eqI{44} confirms. Apart from the overall sign, the last bracket
is the only difference. Using the paper's \eqI{33} together with the \emph{printed} \eqI{40}, the
finite part at $D=4$ is $4i\kappa^2h\,\partial_0(a^2\partial_0\sthreedelf)
=\frac{i\kappa^2H^2}{4\pi^2}\{-\partial^\mu(a^2\partial_\mu\sthreedelf)+a^2\nabla^2\sthreedelf\}$, which
is exactly this bracket; with the correct \eqI{40} the same place gives zero.

\paragraph*{Extracting the logarithm.}
With \eqI{42}, $\square=a^{-D}\partial^\mu(a^{D-2}\partial_\mu)$, one has
$\partial^\mu(a^{D-2}\partial_\mu\sthreedel)=a^D\square\,\sthreedel$. For the logarithmic terms we
use the distributional identity, valid for any smooth $g(x,x')$,
\begin{equation}\label{eq:s3:lemmaB}
g(x,x')\,\partial_\mu\sthreedel=g(x,x)\,\partial_\mu\sthreedel+\bigl[\partial'_\mu g(x,x')\bigr]_{x'=x}\sthreedel ,
\end{equation}
which follows by testing with $f(x')$ and integrating by parts once. For $g=\ln a(\eta')$ it
gives $\ln(a')\,\partial_\mu\sthreedel=\ln(a)\,\partial_\mu\sthreedel+(\partial_\mu\ln a)\,\sthreedel$.
Combining this with the product rule one finds
$\partial^\mu(a^{D-2}\ln(a)\,\partial_\mu\sthreedel)=\frac12\ln(aa')\,a^D\square\,\sthreedel
-\frac12\,a^D(\square\ln a)\,\sthreedel$, and $\square\ln a=-(D-1)H^2$. This is the general-$D$
form of \eqI{43},
\begin{equation}\label{eq:s3:I43}
\partial^\mu\bigl(a^{D-2}\ln(a)\,\partial_\mu\sthreedel\bigr)
=\tfrac12\ln(aa')\,a^D\square\,\sthreedel+\tfrac12(D-1)H^2a^D\,\sthreedel ,
\end{equation}
which reduces to the printed \eqI{43} at $D=4$. The spatial counterpart has no correction term,
$\ln(a)\,\nabla^2\sthreedel=\frac12\ln(aa')\,\nabla^2\sthreedel$, because $\ln a'$ does not
depend on $\vec x\,'$. The symmetric $\ln(aa')$ is essential: with $\ln a^2$ in its place,
Eq.~\eqref{eq:s3:I43} is false.

\paragraph*{Final 4-point result.}
Inserting Eq.~\eqref{eq:s3:I43} and $\frac{H^2}{4\pi^2}=2k+O(D-4)$ into Eq.~\eqref{eq:s3:I41}
gives the corrected form of \eqI{44}, in the definition sign:
\begin{equation}\label{eq:s3:I44}
\begin{aligned}
-i\sthreeM(x;x')={}&-i\kappa^2k\,\Bigl\{\Bigl[-\tfrac14D(D-1)\pi\cot\bigl(\tfrac\pi2D\bigr)+3\ln(aa')\Bigr]a^D\square\\
&\qquad\quad+\Bigl[D\pi\cot\bigl(\tfrac\pi2D\bigr)-4\ln(aa')\Bigr]a^{D-2}\nabla^2+9H^2a^D+O(D-4)\Bigr\}\sthreedel ,
\end{aligned}
\end{equation}
with $k=\frac{H^{D-2}}{(4\pi)^{D/2}}\frac{\Gamma(D-1)}{\Gamma(D/2)}$, the same prefactor as in the
printed equation. If Eq.~\eqref{eq:s3:I43} is used in general $D$, no finite term needs to be
evaluated at $D=4$:
\begin{equation}\label{eq:s3:I44exact}
\begin{aligned}
-i\sthreeM={}&-i\kappa^2k\Bigl\{\tfrac14D(D-1)\bigl[\pihat+\ln(aa')\bigr]a^D\square
-D\bigl[\pihat+\ln(aa')\bigr]a^{D-2}\nabla^2\\
&\qquad\quad+\tfrac14D(D-1)^2H^2a^D\Bigr\}\sthreedel+O(D-4),
\end{aligned}
\end{equation}
where the $O(D-4)$ is the $b$ term of Eq.~\eqref{eq:s3:kernelD}; at $D=4$,
$\frac14\cdot4\cdot3=3$ and $\frac14\cdot4\cdot9=9$. In the paper convention the 4-point result
is $-1$ times Eq.~\eqref{eq:s3:I44},
\begin{equation}\label{eq:s3:I44paper}
\begin{aligned}
-i\sthreeM\big|_{\rm paper\ conv.}={}&+i\kappa^2k\,\Bigl\{\Bigl[-\tfrac14D(D-1)\pi\cot\bigl(\tfrac\pi2D\bigr)
+3\ln(aa')\Bigr]a^D\square\\
&\qquad\quad+\Bigl[D\pi\cot\bigl(\tfrac\pi2D\bigr)-4\ln(aa')\Bigr]a^{D-2}\nabla^2+9H^2a^D+O(D-4)\Bigr\}\sthreedel ,
\end{aligned}
\end{equation}
which is the printed \eqI{44} with the constant $-2$ removed from the $a^D\square$ bracket and the
constant $+2$ removed from the $a^{D-2}\nabla^2$ bracket. The printed \eqI{44} reads
\begin{equation}\label{eq:s3:I44printed}
\begin{aligned}
-i\sthreeM\big|_{\rm printed}={}&+i\kappa^2k\,\Bigl\{\Bigl[-\tfrac14D(D-1)\pi\cot\bigl(\tfrac\pi2D\bigr)
-2+3\ln(aa')\Bigr]a^D\square\\
&\qquad\quad+\Bigl[D\pi\cot\bigl(\tfrac\pi2D\bigr)+2-4\ln(aa')\Bigr]a^{D-2}\nabla^2+9H^2a^D+O(D-4)\Bigr\}\sthreedel .
\end{aligned}
\end{equation}

\paragraph*{Diagnosis.}
The paper's inputs (the printed \eqI{33} and the printed \eqI{40}) reproduce the printed \eqI{41}
and \eqI{44} exactly, and the printed \eqI{41} gives the printed \eqI{44} exactly through
\eqI{42}--\eqI{43} (script \mbox{\nolinkurl{check_41_44.py}}). The steps
\eqI{33}$\to$\eqI{41}$\to$\eqI{44} therefore contain no new error. Both differences are
inherited: the overall sign from the reading of $\partial_2$, and the $\mp2$ constants (together
with the extra bracket of \eqI{41}) from the finite part of \eqI{40}. Note that
$-2a^D\square+2a^{D-2}\nabla^2=2\,\partial_0(a^{D-2}\partial_0\,\cdot\,)$, which is the form found
for the extra finite part above.

\paragraph*{Physical significance.}
The terms produced by the \eqI{40} slip are harmless. At $D=4$, $k\to H^2/8\pi^2$, the extra part
of the printed \eqI{44} is
\begin{equation}\label{eq:s3:extra}
i\kappa^2k\bigl(-2a^D\square+2a^{D-2}\nabla^2\bigr)\sthreedel
=i\Bigl(-\frac1{4\pi^2}\Bigr)\kappa^2H^2a^D\square\,\sthreedel
+i\Bigl(+\frac1{4\pi^2}\Bigr)\kappa^2H^2a^{D-2}\nabla^2\sthreedel .
\end{equation}
This has exactly the form of the counterterm vertices \eqI{131} and \eqI{145}
[Eqs.~\eqref{eq:s5:invvert} and~\eqref{eq:s5:145}]. It is therefore equivalent to the finite counterterm shifts
$\delta\alpha_2=-4$ and $\delta\alpha_3=+4$ (paper convention), in the units $H^{D-4}/(4\pi)^{D/2}$ of Sec.~\ref{sec:s0s1s7s8:conventions},
whose value at $D=4$ is $1/16\pi^2$; the coefficients in Eq.~\eqref{eq:s3:extra} are $\mp\frac1{4\pi^2}=\mp\frac{4}{16\pi^2}$.
In the printed result these are contributions of $-4$ to $\beta_2$ and $+4$ to $\beta_3$.
Removing them shifts $\beta_2$ by $+4$ and $\beta_3$ by $-4$. Together with a finite arithmetic
slip in the 3-point local term \eqI{59} (Sec.~\ref{sec:s4a}), the corrected values are
$\beta_2=-\frac{49}{3}$ instead of $-\frac{61}{3}$ and
$\beta_3=-\frac4{D-4}+\frac{52}{3}+2\gamma_{\rm E}$ instead of $-\frac4{D-4}+\frac{58}{3}+2\gamma_{\rm E}$
(paper convention); the pole is unchanged. The present section supplies only the first of the two
corrections; Sec.~\ref{sec:s5} collects both. Since the finite parts of
the counterterms are arbitrary, these shifts are absorbed into $\Delta\alpha_2$ and
$\Delta\alpha_3$ and have no physical effect.

The overall sign is a different matter. The terms $3\ln(aa')a^D\square$,
$-4\ln(aa')a^{D-2}\nabla^2$ and $9H^2a^D$ in Eq.~\eqref{eq:s3:I44} are not of counterterm form:
the first two carry a time-dependent $\ln a$, and the last is of mass type, while no mass
counterterm is used. They cancel against the corresponding terms of the local 3-point
contribution \eqI{67} (Sec.~\ref{sec:s4a}). Because the printed 3-point part is reversed as well, this cancellation
holds in either convention, and the overall sign of the 4-point part has no consequence by
itself. What remains is the overall sign of the complete self-mass, discussed at the end of
Sec.~\ref{sec:s3:rule}.

\begin{table}[t]
\caption{Summary of the comparison for \eqI{31}--\eqI{44}. ``Definition sign'' refers to
\eqI{168}; ``paper convention'' is the sign of the printed equations. S1 denotes the overall sign of the printed self-mass
(substantive, because the sign of the nonlocal part is physical); C$n$ and B$n$ refer to our Tables~\ref{tab:s0s1s7s8:errata}
and~\ref{tab:s0s1s7s8:finite}.}
\label{tab:s3:summary}
\begin{ruledtabular}
\begin{tabular}{>{\raggedright\arraybackslash}p{0.29\textwidth}>{\raggedright\arraybackslash}p{0.27\textwidth}>{\raggedright\arraybackslash}p{0.15\textwidth}>{\raggedright\arraybackslash}p{0.21\textwidth}}
Item & This work & \PI{} prints & Classification\\
\hline
\eqI{31}: structure, normalization & normalization 1 & same & agrees\\
Table~II, $U_2$ & $+\frac i4\kappa^2a^{D-2}\eta^{\alpha\rho}\eta^{\beta\sigma}\partial_1\!\cdot\partial_2$
 & $-\frac i4$ & typo C5 ($+\frac i4$ used in \eqI{33})\\
Reading of $\partial_2$ & $+\partial$ on $\delta^D$, i.e.\ $-\partial'$ & $+\partial'$ & substantive (S1): 4-point origin of the overall sign\\
\eqI{32}, \eqI{33} (definition sign) & $(+\frac i8,-\frac i4,-\frac i2,+i)$ & $(-\frac i8,+\frac i4,+\frac i2,-i)$
 & S1: inherited overall sign, $-1$\\
\eqI{34}--\eqI{38} & exact in $D$ & same & agree\\
\eqI{39} & $\bar\eta^{\rho\sigma}$ in finite part & $\bar\eta_{\rho\sigma}$ & typo C6 (index placement)\\
\eqI{40}, finite part & $-2\,\delta^0_\alpha\delta^\sigma_0\,h$ & $+2\,\delta^0_\alpha\delta^\sigma_0\,h$ & absorbable (B1)\\
\eqI{41}, overall sign & $-i\kappa^2k\,\pihat\{\dots\}$, \ $-\frac{i\kappa^2H^2}{4\pi^2}\{\dots\}$ & $+$ & S1: inherited overall sign\\
\eqI{41}, bracket $\{-\partial^\mu(a^2\partial_\mu\delta^4)+a^2\nabla^2\delta^4\}$ & absent ($O(D-4)$) & present
 & absorbable (B1, from \eqI{40})\\
\eqI{42}, \eqI{43} & \eqI{43} also in general $D$ & same & agree\\
\eqI{44}: $\pi\cot$, $3\ln(aa')$, $-4\ln(aa')$, $9H^2a^D$ & same structures & same & agree (overall sign: S1)\\
\eqI{44}: constants $-2$ ($a^D\square$), $+2$ ($a^{D-2}\nabla^2$) & absent & present & absorbable (B1, from \eqI{40})\\
\end{tabular}
\end{ruledtabular}
\end{table}

Our Table~\ref{tab:s3:summary} summarizes the comparison of this section with \PI.
All checks described in this section were run with the scripts \mbox{\nolinkurl{check_4pt.py}},
\mbox{\nolinkurl{check_coincidence.py}} and \mbox{\nolinkurl{check_41_44.py}} in
\nolinkurl{code/02_four_point/} of Ref.~\cite{ref:repo}; their outputs agree with the logged
outputs of the original working sessions.

\providecommand{\sIVaBr}[3]{\big[\!\big[\,#1\,\big|\,#2\,\big]\!\big]_{#3}}
\providecommand{\sIVacode}[1]{\texttt{code/#1}}

\section{The 3-point contribution: setup, local and logarithm terms}
\label{sec:s4a}

This section re-derives the contribution of two 3-point vertices to the one-loop scalar
self-mass-squared, the diagram of Fig.~2 of \PI, following Sec.~III\,B of Ref.~\cite{KW07},
Eqs.~\eqI{45}--\eqI{67}. We derive
(i) the effective-action form \eqI{45}--\eqI{46};
(ii) the nine terms produced by the graviton propagator, \eqI{47}--\eqI{49};
(iii) the split into local, logarithm and normal contributions, \eqI{50}--\eqI{53}, including a
derivation of the distributional identity \eqI{52};
(iv) the local contributions \eqI{54}--\eqI{59}; and
(v) the logarithm contributions \eqI{60}--\eqI{67}, together with their cancellation against the
finite local 4-point terms of \eqI{44}.
The normal contributions, \eqI{68} onward, are treated in Sec.~\ref{sec:s4b}.
The complete derivations and the scripts quoted below are in the public notes and code
\cite{ref:repo} (Chs.~3 and 4 of the notes; folders \sIVacode{03\_three\_point\_setup} and
\sIVacode{04\_three\_point\_local\_log}).

\paragraph*{Sign conventions.}
Every equation derived in this section is in the \emph{definition sign}, fixed by \PI's own definition
\eqI{168} (Sec.~\ref{sec:s2:definition}). The printed \eqI{45}--\eqI{46} agree with it. From \eqI{47} on,
the printed 3-point expressions are in the \emph{paper convention}, which is $-1$ times the definition
sign: the printed \eqI{47}--\eqI{49} are $-1$ times what follows from the printed \eqI{46}
(Sec.~\ref{sec:s4a:graviton}), and the paper evaluates all later 3-point expressions from left-hand
sides that carry these printed signs (Secs.~\ref{sec:s4a:local} and~\ref{sec:s4a:log}). The printed
4-point part is reversed as well (Sec.~\ref{sec:s3:rule}). Hence the printed self-mass is $-1$ times the
one defined by \eqI{168}; the 4-point and 3-point parts are both reversed, so their relative sign,
and the local cancellation stated after \eqI{67}, are consistent.
Equations \eqI{50}--\eqI{53}, \eqI{55}--\eqI{56} and \eqI{60}--\eqI{62} carry no self-mass
coefficient and do not depend on the convention.

\paragraph*{Notation.}
The dimension $D$ is kept general throughout; $D=4$ is set only in finite terms that are not multiplied by a pole.
We use $a=-1/(H\eta)$, so that $\partial_0a=Ha^2$, and the abbreviations
\begin{equation}
k\equiv\frac{H^{D-2}}{(4\pi)^{D/2}}\,\frac{\Gamma(D-1)}{\Gamma(\frac D2)}\;\xrightarrow{D\to4}\;\frac{H^2}{8\pi^2},
\qquad
\pihat\equiv-\pi\cot\Big(\frac{\pi D}{2}\Big)=-\frac{2}{D-4}+O(D-4),
\label{eq:s4a:kpi}
\end{equation}
together with $p\equiv D-2$, $l\equiv\ln a(\eta)$, $l'\equiv\ln a(\eta')$ and $L\equiv\ln(aa')=l+l'$.
The scalar propagator is $i\Delta_A(x;x')=A(y)+k\ln(aa')$ [cf.\ \eqI{13}], with
$y=aa'H^2\{\|\vec x-\vec x'\|^2-(|\eta-\eta'|-i\delta)^2\}$ [cf.\ \eqI{11}], and the covariant
d'Alembertian is $a^D\square\equiv\partial^\mu(a^{D-2}\partial_\mu)$ [cf.\ \eqI{42}].
The coincidence limits derived for the 4-point contribution [cf.\ \eqI{34}--\eqI{36}] are
\begin{equation}
A(x;x)=k\,\pihat,\qquad i\Delta_B(x;x)=-\frac{k}{D-2},\qquad i\Delta_C(x;x)=\frac{k}{(D-2)(D-3)},
\label{eq:s4a:coinc}
\end{equation}
where $A(y)\,\delta^D(x-x')=k\pihat\,\delta^D(x-x')$ is the dimensional-regularization assumption
shared with \PI{} (for $\mathrm{Re}\,D<2$ all powers of $y$ vanish at coincidence; the result is
continued in $D$).

\subsection{The effective action and Eqs.~(I.45)--(I.46)}
\label{sec:s4a:setup}

\paragraph*{Vertices and effective action.}
The cubic interaction is [cf.\ Table~I and \eqI{7}]
\begin{equation}
\kappa\mathcal L^{(1)}=\frac12\cdot\frac1i\sum_{I=1}^2\Big[V_I^{\alpha\beta}\phi_1\phi_2\Big]_{\phi_1=\phi_2=\phi}h_{\alpha\beta},
\qquad
V_1^{\alpha\beta}=i\kappa a^{D-2}\partial_1^\alpha\partial_2^\beta,\qquad
V_2^{\alpha\beta}=-\frac i2\kappa a^{D-2}\eta^{\alpha\beta}\partial_1\!\cdot\partial_2 .
\label{eq:s4a:vertices}
\end{equation}
Table~I of \PI{} prints these vertex operators correctly.\footnote{The left-hand side of \eqI{7} is
printed with $-\tfrac12\kappa a^{D-2}\partial_\alpha\phi\,\partial_\beta\phi\,h^{\alpha\beta}$; the correct sign is
$+\tfrac12\kappa$. The printed $V_1$ and the printed \eqI{46} are built with the correct sign. Had the
printed sign been used, the second and third terms of \eqI{46} would be reversed.}
We write $\phi=\Phi+\varphi$ with background $\Phi$ and quantum field $\varphi$. Because $V_I$ is
symmetric under $1\leftrightarrow2$ once contracted into $h_{\alpha\beta}$, the cross term is
\begin{equation}
S_3^\times=\int d^Dz\;\frac1i\sum_I\big[V_I^{\alpha\beta}(\partial_1,\partial_2)\,\Phi_1\varphi_2\big]h_{\alpha\beta},
\label{eq:s4a:cross}
\end{equation}
with leg~1 the outer leg ($\Phi$) and leg~2 the inner leg ($\varphi$). At second order,
$i\Gamma\supset\frac12\langle(iS_{\rm int})^2\rangle_{\rm c}$. The only 1PI term that is quadratic in
$\Phi$, of order $\kappa^2$ and built from two cubic vertices is $\langle(iS_3^\times)^2\rangle$ with
$\varphi\varphi$ and $hh$ contracted; the diagrams in which the graviton line ends on a scalar loop,
or on the $\langle h\rangle$ tadpole, are one-particle reducible. With
$\langle T\varphi(z)\varphi(z')\rangle=i\Delta_A(z;z')$ and
$\langle Th_{\alpha\beta}(z)h_{\rho\sigma}(z')\rangle=i[{}_{\alpha\beta}\Delta_{\rho\sigma}](z;z')$,
\begin{equation}
i\Gamma_{\rm 3pt}[\Phi]=\frac12\sum_{I,J}\iint d^Dz\,d^Dz'\;V_I^{\alpha\beta}(z)\,V_J^{\rho\sigma}(z')\;
\Phi(z)\,\Phi(z')\;i[{}_{\alpha\beta}\Delta_{\rho\sigma}](z;z')\;i\Delta_A(z;z'),
\label{eq:s4a:Gamma3}
\end{equation}
where $\partial_1$ ($\partial'_1$) acts on $\Phi(z)$ ($\Phi(z')$) and $\partial_2$ ($\partial'_2$) on the first
(second) argument of $i\Delta_A$.

\paragraph*{Normalization and derivative rules.}
The second functional derivative has two assignments, $(z,z')\to(x,x')$ and $(z,z')\to(x',x)$, which
contribute equally because both propagators are symmetric under exchange of their arguments. The
overall factor is
\begin{equation}
\underbrace{\tfrac12}_{\rm cumulant}\times\underbrace{2\times2}_{\text{cross terms}}\times
\underbrace{(2i)^{-2}}_{\text{normalization}}\times\underbrace{i^2}_{e^{iS}}\times
\underbrace{2}_{\text{assignments}}=1 .
\label{eq:s4a:norm}
\end{equation}
For an outer-leg derivative, $\delta/\delta\Phi(x)$ produces $\partial^z_\mu\delta^D(z-x)$, and a partial
integration in $z$ turns it into $-\partial_\mu$ acting on \emph{everything} at the vertex $x$:
$a^{D-2}(x)$, the graviton propagator and the inner derivative of $i\Delta_A$. The same holds at $x'$.
Inner-leg derivatives act only on $i\Delta_A$, with a $+$ sign. The two outer signs give
$(-1)(-1)=+1$, and therefore
\begin{equation}
-iM^2_{\rm 3pt}(x;x')=\sum_{I=1}^2V_I^{\alpha\beta}(x)\sum_{J=1}^2V_J^{\rho\sigma}(x')\;
i\big[{}_{\alpha\beta}\Delta_{\rho\sigma}\big](x;x')\;i\Delta_A(x;x'),
\label{eq:s4a:45}
\end{equation}
with the outer derivatives partially integrated. This is \eqI{45}, together with the rule stated in
\PI{} (``partially integrate any derivative that acts upon an outer leg'').

\paragraph*{Relation to the 4-point rule.}
In the 4-point contribution both outer legs sit at the same vertex. After the first partial
integration no integration variable is left, and the second outer derivative acts on
$\delta^D(x-x')$ with respect to $x$, i.e.\ as $-\partial'$. The rule \eqI{32} reads it as $+\partial'$;
this is the origin of the reversed sign of the printed 4-point part (Sec.~\ref{sec:s3:rule}). The 3-point
rule, with $-\partial$ on each outer leg at its own vertex, is correct. The printed 3-point part is
nonetheless reversed as well, by a separate overall sign slip in the step from \eqI{46} to
\eqI{47}--\eqI{49} (Sec.~\ref{sec:s4a:graviton}).

\paragraph*{Eq.~(I.46).}
Expanding the four $(I,J)$ terms of \eqref{eq:s4a:45} with \eqref{eq:s4a:vertices}: the vertex
coefficients are $(i\kappa)^2=-\kappa^2$, $(-\frac i2\kappa)(i\kappa)=+\frac{\kappa^2}2$ (twice) and
$(-\frac i2\kappa)^2=-\frac{\kappa^2}4$; the outer derivatives give $(-\partial)(-\partial')$ in each case; and
$\eta^{\alpha\beta}$ in $V_2$ produces a trace of the graviton propagator. Thus
\begin{equation}
\begin{aligned}
-iM^2_{\rm 3pt}(x;x')={}&-\kappa^2\partial_\alpha\partial'_\rho\Big\{(aa')^{D-2}\,i[{}^{\alpha\beta}\Delta^{\rho\sigma}]\,\partial_\beta\partial'_\sigma i\Delta_A\Big\}
+\frac{\kappa^2}2\partial^\mu\partial'_\rho\Big\{(aa')^{D-2}\,i[{}^\alpha_{\ \alpha}\Delta^{\rho\sigma}]\,\partial_\mu\partial'_\sigma i\Delta_A\Big\}\\
&+\frac{\kappa^2}2\partial_\alpha\partial'^\nu\Big\{(aa')^{D-2}\,i[{}^{\alpha\beta}\Delta^\rho_{\ \rho}]\,\partial_\beta\partial'_\nu i\Delta_A\Big\}
-\frac{\kappa^2}4\partial^\mu\partial'^\nu\Big\{(aa')^{D-2}\,i[{}^\alpha_{\ \alpha}\Delta^\rho_{\ \rho}]\,\partial_\mu\partial'_\nu i\Delta_A\Big\},
\end{aligned}
\label{eq:s4a:46}
\end{equation}
where all propagators have arguments $(x;x')$. The coefficients do not depend on $D$. Equation
\eqref{eq:s4a:46} agrees with the printed \eqI{46} in every coefficient, sign and index contraction,
and has the same structure and signs as the flat-space 3-point expression \eqI{175} of \PI.

\paragraph*{Checks.}
The script \sIVacode{03\_three\_point\_setup/check\_45\_46.py}~\cite{ref:repo} verifies:
(a) for $D=3$ and a generic tensor $X_{\alpha\beta\rho\sigma}(x;x')$ symmetric within each pair
(36 independent components), $\sum_{IJ}V_IV_JX$ reproduces \eqref{eq:s4a:46}, while $V_1\to-V_1$
(the printed left-hand side of \eqI{7}) reverses exactly the second and third terms;
(b) the factor \eqref{eq:s4a:norm};
(c) in a one-dimensional model with smooth symmetric kernels and test functions, the G\^ateaux second
derivative ($0.779492709611$) is reproduced by the rule $(-\partial)(-\partial')$ on everything, while
a rule with one outer leg reversed gives the opposite value.
A textbook comparison (the cubic theory $\frac g2\phi^2\chi$, whose one-loop self-energy has vertex
$ig$ and symmetry factor~1) confirms the same normalization.

\subsection{Inserting the graviton propagator: Eqs.~(I.47)--(I.49)}
\label{sec:s4a:graviton}

\paragraph*{Inputs.}
The graviton propagator [cf.\ \eqI{20}--\eqI{24}] is
$i[{}_{\alpha\beta}\Delta_{\rho\sigma}]=\sum_{I=A,B,C}T^I_{\alpha\beta\rho\sigma}\,i\Delta_I(x;x')$, with the
constant tensors $T^{A,B,C}$ of Eqs.~\eqref{eq:s2:TAB}--\eqref{eq:s2:TC} and their single and double traces
\eqref{eq:s2:traces}. Because $i\Delta_A$ depends on $\vec x,\vec x'$
only through $\|\vec x-\vec x'\|$,
\begin{equation}
\partial'_i\,i\Delta_A=-\partial_i\,i\Delta_A,\qquad \partial_i\partial'_j\,i\Delta_A=\partial_j\partial'_i\,i\Delta_A .
\label{eq:s4a:hom}
\end{equation}

\paragraph*{Contractions.}
We represent the outer derivatives by formal vectors $P$ (at $x$) and $P'$ (at $x'$), the inner ones
by $Q$ and $Q'$, and write $m\equiv P_0Q_0$, $\bar m\equiv\bar P\!\cdot\!\bar Q\equiv\sum_iP_iQ_i$,
$n\equiv P'_0Q'_0$, $\bar n\equiv\bar P'\!\cdot\!\bar Q'$, so that $P\!\cdot\!Q=-m+\bar m$ and
$P'\!\cdot\!Q'=-n+\bar n$. For the result we use the shorthand
\begin{equation}
\sIVaBr{\partial_a\partial'_c}{\partial_b\partial'_d}{I}\equiv
\partial_a\partial'_c\Big\{(aa')^{D-2}\,i\Delta_I(x;x')\,\partial_b\partial'_d\,i\Delta_A(x;x')\Big\},
\label{eq:s4a:bracket}
\end{equation}
with outer derivatives to the left of the bar and inner ones to the right. The first term of
\eqref{eq:s4a:46} contracts the full tensors; with $(u\!\cdot\!P)(u\!\cdot\!Q')=P_0Q'_0$ for
$u_\mu=\delta^0_\mu$,
\begin{equation}
\begin{aligned}
T^A_{\alpha\beta\rho\sigma}P^\alpha Q^\beta P'^\rho Q'^\sigma&=(\bar P\!\cdot\!\bar P')(\bar Q\!\cdot\!\bar Q')
+(\bar P\!\cdot\!\bar Q')(\bar P'\!\cdot\!\bar Q)-\tfrac{2}{D-3}\,\bar m\,\bar n,\\
T^B_{\alpha\beta\rho\sigma}P^\alpha Q^\beta P'^\rho Q'^\sigma&=-\big[P_0Q'_0(\bar P'\!\cdot\!\bar Q)+Q_0Q'_0(\bar P\!\cdot\!\bar P')
+P_0P'_0(\bar Q\!\cdot\!\bar Q')+Q_0P'_0(\bar P\!\cdot\!\bar Q')\big],\\
T^C_{\alpha\beta\rho\sigma}P^\alpha Q^\beta P'^\rho Q'^\sigma&=\tfrac{2}{(D-2)(D-3)}\,\big[(D-3)m+\bar m\big]\big[(D-3)n+\bar n\big].
\end{aligned}
\label{eq:s4a:c1}
\end{equation}
The second and third terms of \eqref{eq:s4a:46} involve one trace and the fourth term the double trace;
the $T^B$ sector receives contributions from the first term only.

\paragraph*{$AA$ terms.}
Our Table~\ref{tab:s4a:AA} collects the $T^A$ parts of the four terms. In the $\bar m\bar n$ row the factor
$D-3$ cancels exactly, $(2-2-2+D-1)/(D-3)=1$. The total is
\begin{equation}
\kappa^2\Big\{-(\bar P\!\cdot\!\bar P')(\bar Q\!\cdot\!\bar Q')+\tfrac{D-1}{D-3}\,mn-\bar m\,n-m\,\bar n\Big\}
+\kappa^2\Big\{\bar m\,\bar n-(\bar P\!\cdot\!\bar Q')(\bar P'\!\cdot\!\bar Q)\Big\}.
\label{eq:s4a:AAsum}
\end{equation}
The second brace is $\sIVaBr{\partial_i\partial'_j}{\partial_i\partial'_j-\partial_j\partial'_i}{A}$, which vanishes
identically by \eqref{eq:s4a:hom}; \PI{} does not write it either.

\begin{table}[htbp]
\caption{Coefficients of the $AA$ monomials (definition sign), from the $T^A$ parts of the four terms
of Eq.~\eqref{eq:s4a:46}.}
\label{tab:s4a:AA}
\renewcommand{\arraystretch}{1.3}
\begin{ruledtabular}
\begin{tabular}{lccccc}
Monomial & 1st term & 2nd term & 3rd term & 4th term & Total$/\kappa^2$\\
\hline
$\bar m\,\bar n$ & $+\frac2{D-3}$ & $-\frac2{D-3}$ & $-\frac2{D-3}$ & $+\frac{D-1}{D-3}$ & $+1$\\
$m\,\bar n$ & --- & $+\frac2{D-3}$ & --- & $-\frac{D-1}{D-3}$ & $-1$\\
$\bar m\,n$ & --- & --- & $+\frac2{D-3}$ & $-\frac{D-1}{D-3}$ & $-1$\\
$m\,n$ & --- & --- & --- & $+\frac{D-1}{D-3}$ & $+\frac{D-1}{D-3}$\\
$(\bar P\!\cdot\!\bar P')(\bar Q\!\cdot\!\bar Q')$ & $-1$ & --- & --- & --- & $-1$\\
$(\bar P\!\cdot\!\bar Q')(\bar P'\!\cdot\!\bar Q)$ & $-1$ & --- & --- & --- & $-1$\\
\end{tabular}
\end{ruledtabular}
\end{table}

\paragraph*{$AB$ and $AC$ terms.}
Since the traces of $T^B$ vanish, the $AB$ sector is $-\kappa^2$ times the $T^B$ line of
\eqref{eq:s4a:c1}. In the $AC$ sector, with $M_1=(D-3)m+\bar m$, $M_2=P\!\cdot\!Q=-m+\bar m$, the
analogous $N_{1,2}$, and $c_C\equiv2\kappa^2/[(D-2)(D-3)]$, the four terms add up to
\begin{equation}
c_C\big[-M_1N_1+M_2N_1+M_1N_2-M_2N_2\big]=-c_C\,(M_1-M_2)(N_1-N_2)=-c_C\,(D-2)^2\,mn .
\label{eq:s4a:ACfac}
\end{equation}

\paragraph*{Result.}
Mapping the monomials back to derivatives, the nine terms are, in the definition sign,
\begin{subequations}
\label{eq:s4a:47-49}
\begin{align}
AA:\quad&-\kappa^2\,\sIVaBr{\nabla\!\cdot\!\nabla'}{\nabla\!\cdot\!\nabla'}{A}
+\kappa^2\Big(\frac{D-1}{D-3}\Big)\sIVaBr{\partial_0\partial'_0}{\partial_0\partial'_0}{A}
-\kappa^2\,\sIVaBr{\partial_i\partial'_0}{\partial_i\partial'_0}{A}
-\kappa^2\,\sIVaBr{\partial_0\partial'_i}{\partial_0\partial'_i}{A},
\label{eq:s4a:47}\\
AB:\quad&+\kappa^2\,\sIVaBr{\partial_0\partial'_0}{\nabla\!\cdot\!\nabla'}{B}
+\kappa^2\,\sIVaBr{\partial_i\partial'_0}{\partial_0\partial'_i}{B}
+\kappa^2\,\sIVaBr{\partial_0\partial'_i}{\partial_i\partial'_0}{B}
+\kappa^2\,\sIVaBr{\nabla\!\cdot\!\nabla'}{\partial_0\partial'_0}{B},
\label{eq:s4a:48}\\
AC:\quad&-2\kappa^2\Big(\frac{D-2}{D-3}\Big)\sIVaBr{\partial_0\partial'_0}{\partial_0\partial'_0}{C}.
\label{eq:s4a:49}
\end{align}
\end{subequations}
These are \eqI{47}, \eqI{48} and \eqI{49}, respectively. Structures, magnitudes and the $D$-dependent
coefficients $\frac{D-1}{D-3}$ and $\frac{2(D-2)}{D-3}$ agree with \PI{}, but each of the nine printed
terms has the opposite sign (our Table~\ref{tab:s4a:47cmp}). A one-line hand check: the structure
$\sIVaBr{\nabla\!\cdot\!\nabla'}{\nabla\!\cdot\!\nabla'}{A}$ can only come from the first term of
\eqref{eq:s4a:46} (coefficient $-\kappa^2$) and the $\bar\eta\bar\eta$ part of $T^A$ (coefficient $+1$),
so its coefficient is unambiguously $-\kappa^2$; \PI{} prints $+\kappa^2$.
The difference cannot come from the tensors \eqref{eq:s2:TAB}--\eqref{eq:s2:TC}, which reproduce the printed \eqI{37}--\eqI{38}
exactly, nor from the vertices, since \eqref{eq:s4a:46} agrees with the printed \eqI{46}. It is an
overall sign slip in the printed step from \eqI{46} to \eqI{47}--\eqI{49}.

\begin{table}[htbp]
\caption{The nine terms of \eqI{47}--\eqI{49}: derived coefficients (definition sign) and printed
coefficients (paper convention).}
\label{tab:s4a:47cmp}
\renewcommand{\arraystretch}{1.25}
\begin{ruledtabular}
\begin{tabular}{lcc}
Term & Derived, Eq.~\eqref{eq:s4a:47-49} & Printed in \PI\\
\hline
$AA$: $\sIVaBr{\nabla\!\cdot\!\nabla'}{\nabla\!\cdot\!\nabla'}{A}$ & $-\kappa^2$ & $+\kappa^2$\\
$AA$: $\sIVaBr{\partial_0\partial'_0}{\partial_0\partial'_0}{A}$ & $+\kappa^2\frac{D-1}{D-3}$ & $-\kappa^2\frac{D-1}{D-3}$\\
$AA$: $\sIVaBr{\partial_i\partial'_0}{\partial_i\partial'_0}{A}$, $\sIVaBr{\partial_0\partial'_i}{\partial_0\partial'_i}{A}$ & $-\kappa^2$, $-\kappa^2$ & $+\kappa^2$, $+\kappa^2$\\
$AB$: four terms & all $+\kappa^2$ & all $-\kappa^2$\\
$AC$: $\sIVaBr{\partial_0\partial'_0}{\partial_0\partial'_0}{C}$ & $-2\kappa^2\frac{D-2}{D-3}$ & $+2\kappa^2\frac{D-2}{D-3}$\\
antisymmetric $AA$ pair & $0$ by Eq.~\eqref{eq:s4a:hom} & not written\\
\end{tabular}
\end{ruledtabular}
\end{table}

\paragraph*{Checks.}
The script \sIVacode{03\_three\_point\_setup/check\_47\_49.py}~\cite{ref:repo} builds $T^{A,B,C}$ as
explicit component arrays for $D=4,5,6$ and tracks the outer propagator type symbolically. For each of
these $D$ it finds: the route through \eqref{eq:s4a:46} equals the route directly from \eqref{eq:s4a:45}
and Table~I; both equal \eqref{eq:s4a:47-49} plus the antisymmetric pair, which vanishes with
$\partial'_i=-\partial_i$ on $i\Delta_A$; and they equal $-1$ times the printed \eqI{47}--\eqI{49}.
The homogeneity relations \eqref{eq:s4a:hom} are verified for $A(y)+k\ln(aa')$.

\subsection{The three-way split: Eqs.~(I.50)--(I.53)}
\label{sec:s4a:split}

Each of the nine terms has one outer derivative at $x$, one at $x'$, and an inner pair of the form
$\partial_\rho\partial'_\sigma$, i.e.\ the form \eqI{50},
$\kappa^2\partial_\mu\partial'_\nu[(aa')^{D-2}i\Delta_I\,\partial_\rho\partial'_\sigma i\Delta_A]$ with $I=A,B,C$.
By \eqI{13}, \eqI{27} and \eqI{28}, $i\Delta_A=A(y)+k\ln(aa')$, $i\Delta_B=B(y)$ and $i\Delta_C=C(y)$,
which is \eqI{51}. Since each inner pair is mixed,
\begin{equation}
\partial_\rho\partial'_\sigma\,k\ln(aa')=k\,\partial_\rho\partial'_\sigma\big[\ln a(\eta)+\ln a(\eta')\big]=0 ,
\label{eq:s4a:lnkill}
\end{equation}
so the inner derivatives eliminate the de Sitter breaking term, as stated in \PI\@. The only nontrivial
equation of this group is \eqI{52}, derived next.

\paragraph*{Derivation of (I.52): direct computation.}
For $\delta>0$,
$y=aa'H^2\{r^2-\Delta\eta^2+2i\delta|\Delta\eta|+\delta^2\}$ with $r\equiv\|\vec x-\vec x'\|$ and
$\Delta\eta\equiv\eta-\eta'$, so $\mathrm{Im}\,y=2aa'H^2\delta|\Delta\eta|\ge0$, and $y>0$ when
$\Delta\eta=0$. Thus $y$ never reaches zero or the negative real axis, the powers $y^{1-D/2},y^{2-D/2},\dots$
in $A(y)$ are analytic, and the only non-smooth structure is the kink of $|\Delta\eta|$. With
$\partial_0|\Delta\eta|=\mathrm{sgn}(\Delta\eta)$ and $\partial'_0\,\mathrm{sgn}(\Delta\eta)=-2\delta(\Delta\eta)$,
\begin{equation}
\partial_0\partial'_0\big(|\Delta\eta|-i\delta\big)^2=-2+4i\delta\,\delta(\Delta\eta),
\label{eq:s4a:kink}
\end{equation}
while $\partial'_jy=-2aa'H^2(x-x')_j$ contains no $|\Delta\eta|$. The chain rule then gives
\begin{equation}
\partial_\rho\partial'_\sigma A(y)=A''(y)\,\partial_\rho y\,\partial'_\sigma y+A'(y)\big[\partial_\rho\partial'_\sigma y\big]_{\rm reg}
-4i\delta\,\delta^0_\rho\delta^0_\sigma\,a^2H^2A'\big(a^2H^2(r^2+\delta^2)\big)\,\delta(\Delta\eta),
\label{eq:s4a:chain}
\end{equation}
where $\delta(\Delta\eta)$ has been used to set $a'=a$ in the last term. In $D-1$ dimensions
\begin{equation}
\lim_{\delta\to0^+}\delta\,(r^2+\delta^2)^{-D/2}=\frac{\pi^{D/2}}{\Gamma(\frac D2)}\,\delta^{D-1}(\vec x-\vec x'),
\label{eq:s4a:nascent}
\end{equation}
since the total weight is
$\Omega_{D-2}\int_0^\infty ds\,s^{D-2}(1+s^2)^{-D/2}=\pi^{D/2}/\Gamma(\frac D2)$. Writing the leading term of
$A(y)$ as $\mathcal N c_0(4/y)^{D/2-1}$ with $\mathcal N\equiv H^{D-2}/(4\pi)^{D/2}$, and using
$\mathcal N\,4^{D/2}H^{2-D}\pi^{D/2}=1$, the last term of \eqref{eq:s4a:chain} tends to
\begin{equation}
i\,\frac{c_0\,(\frac D2-1)}{\Gamma(\frac D2)}\,a^{2-D}\,\delta^D(x-x').
\label{eq:s4a:lead}
\end{equation}
The subleading terms of $A'$ are at most as singular as $y^{1-D/2}$; their weight is $O(\delta)$ and
they drop out. With $c_0=\Gamma(\frac D2-1)$ (see below), $c_0(\frac D2-1)=\Gamma(\frac D2)$, and
\begin{equation}
\partial_\rho\partial'_\sigma\,i\Delta_A(x;x')=\delta^0_\rho\delta^0_\sigma\,\frac{i}{a^{D-2}}\,\delta^D(x-x')
+A''(y)\,\frac{\partial y}{\partial x^\rho}\frac{\partial y}{\partial x'^\sigma}+A'(y)\,\frac{\partial^2y}{\partial x^\rho\partial x'^\sigma},
\label{eq:s4a:52}
\end{equation}
where the last derivative is the pointwise (regular) one. This is the printed \eqI{52}.

\paragraph*{Independent derivations of the local coefficient.}
(B) \emph{Propagator equation.} For a function of $|\Delta\eta|$ the singular parts obey
$[\partial_0\partial'_0G]_{\rm sing}=-[\partial_0^2G]_{\rm sing}$. In \eqI{12},
$\partial_\mu(a^{D-2}\eta^{\mu\nu}\partial_\nu)i\Delta_A=i\delta^D$, the singular term comes only from
$-a^{D-2}\partial_0^2$, so $[\partial_0\partial'_0i\Delta_A]_{\rm sing}=+i\,a^{2-D}\delta^D$. This route does
not use the normalization of $A(y)$.
(C) \emph{Canonical commutator.} From $T\phi\phi'=\theta(\Delta\eta)\phi\phi'+\theta(-\Delta\eta)\phi'\phi$ and
$\pi=a^{D-2}\dot\phi$, $[\phi,\pi']=i\delta^{D-1}$,
\begin{equation}
\partial_0\partial'_0\langle T\phi(x)\phi(x')\rangle=\langle T\dot\phi(x)\dot\phi(x')\rangle+\frac{i}{a^{D-2}}\,\delta^D(x-x').
\label{eq:s4a:routeC}
\end{equation}
All three routes give the coefficient $\delta^0_\rho\delta^0_\sigma\,i\,a^{2-D}\delta^D(x-x')$.

\paragraph*{First coefficient of (I.14).}
Route~A gives $i\,c_0(\frac D2-1)a^{2-D}/\Gamma(\frac D2)$; routes B and C then force
$c_0=\Gamma(\frac D2)/(\frac D2-1)=\Gamma(\frac D2-1)$, in agreement with the most singular part \eqI{26}
and with the recursion obeyed by the series (Sec.~\ref{sec:s2:scalar}). The printed
$c_0=\Gamma(\frac D2-1)/(\frac D2-1)$ in \eqI{14} is therefore a typographical error (C4): the printed
\eqI{52} carries the correct local coefficient, whereas the printed $c_0$ would multiply it by
$\frac{2}{D-2}$, which is $1$ at $D=4$ but would produce a finite error when multiplied by the pole of
$i\Delta_A(x;x)$.

\paragraph*{The split.}
Inserting \eqref{eq:s4a:52} into \eqI{50} gives three parts:
\begin{equation}
\kappa^2\partial_\mu\partial'_\nu\Big[(aa')^{D-2}\,i\Delta_I\Big(\underbrace{\delta^0_\rho\delta^0_\sigma\,\tfrac{i}{a^{D-2}}\,\delta^D}_{\text{(i) local}}
+\mathcal R_{\rho\sigma}\Big)\Big],
\qquad
i\Delta_I\,\mathcal R_{\rho\sigma}=
\begin{cases}
\underbrace{A(y)\,\mathcal R_{\rho\sigma}}_{\text{(iii) normal}}+\underbrace{k\ln(aa')\,\mathcal R_{\rho\sigma}}_{\text{(ii) logarithm}}, & I=A,\\[2pt]
\underbrace{I(y)\,\mathcal R_{\rho\sigma}}_{\text{(iii) normal}}, & I=B,C,
\end{cases}
\label{eq:s4a:split}
\end{equation}
with $\mathcal R_{\rho\sigma}\equiv A''\partial_\rho y\,\partial'_\sigma y+A'\partial_\rho\partial'_\sigma y$; for $I=B,C$,
$I(y)$ stands for $B(y)$ or $C(y)$ (the propagator label $I$ is unrelated to the integration operator $I[f]$ of
Sec.~\ref{sec:s4b}). Part (iii) is \eqI{53}. Following \PI, the local part for $I=A$ uses only $A(y)$, and the piece
$k\ln(aa')\times\delta^D$ is counted with the logarithm contributions (which then use the full
$\partial\partial'i\Delta_A$). With this assignment the split is complete and free of overlap.
Equations \eqI{50}--\eqI{53} are statements of form without overall coefficients; they do not decide the sign
question of Sec.~\ref{sec:s4a:graviton}.

\paragraph*{Checks.}
The script \sIVacode{03\_three\_point\_setup/check\_52.py}~\cite{ref:repo} verifies
$\partial_0\partial'_0\ln(aa')=0$; the kink identity \eqref{eq:s4a:kink}; the vanishing of the
$\eta^{\rho\sigma}$-contracted regular part in flat space for general $D$; the weight in
\eqref{eq:s4a:nascent} in closed form for $D=3,\dots,6$ and for general $D$, and numerically at $D=4$
($9.650$, $9.847$, $9.867$ for $\delta=10^{-2},10^{-3},10^{-4}$, against $\pi^2=9.8696$); the local
coefficient $i\,a^{2-D}$ with $c_0=\Gamma(\frac D2-1)$ (and $\frac{2i}{D-2}a^{2-D}$ with the printed
$c_0$); the recursion for $c_0$ of Sec.~\ref{sec:s2:scalar}; and $\square f(y)=H^2[(4y-y^2)f''+D(2-y)f']$ for $D=4,5$.

\subsection{Local contributions: Eqs.~(I.54)--(I.59)}
\label{sec:s4a:local}

Only the three terms of \eqref{eq:s4a:47-49} with the inner pair $\partial_0\partial'_0$ give local
contributions: the second term of \eqref{eq:s4a:47}, the fourth term of \eqref{eq:s4a:48}, and
\eqref{eq:s4a:49}. In the definition sign,
\begin{equation}
\begin{aligned}
\mathcal J_{AA}&=+\kappa^2\,\frac{D-1}{D-3}\,\partial_0\partial'_0\Big[(aa')^{D-2}\,i\Delta_A\,\partial_0\partial'_0\,i\Delta_A\Big],\\
\mathcal J_{AB}&=+\kappa^2\,\nabla\!\cdot\!\nabla'\Big[(aa')^{D-2}\,i\Delta_B\,\partial_0\partial'_0\,i\Delta_A\Big],\\
\mathcal J_{AC}&=-2\kappa^2\,\frac{D-2}{D-3}\,\partial_0\partial'_0\Big[(aa')^{D-2}\,i\Delta_C\,\partial_0\partial'_0\,i\Delta_A\Big].
\end{aligned}
\label{eq:s4a:Jterms}
\end{equation}
The text of \PI{} before \eqI{54} attributes the $AA$ local term to ``the 4th term of (47)''; as the
left-hand side of \eqI{54} shows, the second term is meant (typographical, C7).

\paragraph*{Rules.}
(R1) With the local part of \eqref{eq:s4a:52} and the coincidence limits \eqref{eq:s4a:coinc},
\begin{equation}
\begin{gathered}
(aa')^{D-2}\,i\Delta_I(x;x')\,\frac{i}{a^{D-2}}\,\delta^D(x-x')=i\,\Lambda_I\,a^{D-2}\,\delta^D(x-x'),\\
\Lambda_A=k\pihat,\qquad \Lambda_B=-\frac{k}{D-2},\qquad \Lambda_C=\frac{k}{(D-2)(D-3)},
\end{gathered}
\label{eq:s4a:R1}
\end{equation}
where for $I=A$ only $A(y)$ is kept, and derivatives are taken after the coincidence limit.
(R2) Acting on a test function, the primed derivative acts only on the delta function, and with
$a^D\square f=-\partial_0(a^{D-2}\partial_0f)+a^{D-2}\nabla^2f$,
\begin{equation}
\begin{gathered}
\partial_0\partial'_0\big[a^{D-2}\delta^D\big]=-\partial_0\big(a^{D-2}\partial_0\delta^D\big)=\big(a^D\square-a^{D-2}\nabla^2\big)\delta^D,\\
\nabla\!\cdot\!\nabla'\big[a^{D-2}\delta^D\big]=-a^{D-2}\nabla^2\delta^D .
\end{gathered}
\label{eq:s4a:R2}
\end{equation}
The printed \eqI{55}--\eqI{56} is the first relation multiplied by $-1$ and is correct.

\paragraph*{The three local terms.}
Using $(D-1)\Gamma(D-1)=\Gamma(D)$,
\begin{align}
\mathcal J_{AA}^{\rm loc}&=\frac{i\kappa^2H^{D-2}}{(4\pi)^{D/2}}\,\frac{\Gamma(D)}{(D-3)\Gamma(\frac D2)}
\Big(-\pi\cot\big(\tfrac{\pi D}{2}\big)\Big)\big\{+a^D\square-a^{D-2}\nabla^2\big\}\delta^D(x-x'),
\label{eq:s4a:54}\\
\mathcal J_{AB}^{\rm loc}&=+\frac{i\kappa^2k}{D-2}\,a^{D-2}\nabla^2\delta^D(x-x')
=+\frac{i\kappa^2H^2}{16\pi^2}\,a^2\nabla^2\delta^4(x-x')+O(D-4),
\label{eq:s4a:57}\\
\mathcal J_{AC}^{\rm loc}&=-\frac{2i\kappa^2k}{(D-3)^2}\big(a^D\square-a^{D-2}\nabla^2\big)\delta^D(x-x')
=-\frac{i\kappa^2H^2}{4\pi^2}\big\{a^4\square-a^2\nabla^2\big\}\delta^4(x-x')+O(D-4),
\label{eq:s4a:58}
\end{align}
which correspond to \eqI{54}, \eqI{57} and \eqI{58}. In \eqref{eq:s4a:57} and \eqref{eq:s4a:58} the
coefficients are finite and are not multiplied by a pole, so setting $D=4$ there changes only terms
of order $D-4$; the pole coefficient in \eqref{eq:s4a:54} is kept at exact $D$.

\paragraph*{The sum, Eq.~(I.59).}
At exact $D$, in units of $i\kappa^2k$,
\begin{equation}
-iM^2_{\rm 3pt,loc}=i\kappa^2k\,\Big\{\Big[\frac{D-1}{D-3}\,\pihat-\frac{2}{(D-3)^2}\Big]a^D\square
+\Big[-\frac{D-1}{D-3}\,\pihat+\frac{1}{D-2}+\frac{2}{(D-3)^2}\Big]a^{D-2}\nabla^2\Big\}\delta^D(x-x').
\label{eq:s4a:59exact}
\end{equation}
Setting $D=4$ in the finite parts, $\frac{2}{(D-3)^2}\to2$ and $\frac1{D-2}\to\frac12$:
\begin{equation}
\begin{aligned}
-iM^2_{\rm 3pt,loc}(x;x')=\frac{i\kappa^2H^{D-2}}{(4\pi)^{D/2}}\,\frac{\Gamma(D-1)}{\Gamma(\frac D2)}\,\Big\{
&\Big[-\frac{D-1}{D-3}\,\pi\cot\big(\tfrac\pi2D\big)-2\Big]a^D\square\\
&+\Big[\frac{D-1}{D-3}\,\pi\cot\big(\tfrac\pi2D\big)+\frac52\Big]a^{D-2}\nabla^2+O(D-4)\Big\}\delta^D(x-x')
\end{aligned}
\label{eq:s4a:59}
\end{equation}
(definition sign). In the paper convention the two brackets read
$[\frac{D-1}{D-3}\pi\cot(\frac\pi2D)+2]$ and $[-\frac{D-1}{D-3}\pi\cot(\frac\pi2D)-\frac52]$. \PI{} prints
\eqI{59} with $-\frac72$ in place of $-\frac52$; all other entries agree up to the overall sign
(our Table~\ref{tab:s4a:loccmp}).

\begin{table}[htbp]
\caption{Local 3-point contributions. Derived values are in the definition sign, printed values in the
paper convention. Common prefactors are those of Eqs.~\eqref{eq:s4a:54}--\eqref{eq:s4a:59}.}
\label{tab:s4a:loccmp}
\renewcommand{\arraystretch}{1.3}
\begin{ruledtabular}
\begin{tabular}{lll}
Item & Derived & Printed in \PI\\
\hline
\eqI{54} & $\{+a^D\square-a^{D-2}\nabla^2\}\delta^D$ & $\{-a^D\square+a^{D-2}\nabla^2\}\delta^D$\\
\eqI{55}--\eqI{56} & Eq.~\eqref{eq:s4a:R2} & same\\
\eqI{57} & $+\frac{i\kappa^2H^2}{16\pi^2}\,a^2\nabla^2\delta^4$ & $-\frac{i\kappa^2H^2}{16\pi^2}\,a^2\nabla^2\delta^4$\\
\eqI{58} & $-\frac{i\kappa^2H^2}{4\pi^2}\{a^4\square-a^2\nabla^2\}\delta^4$ & $+\frac{i\kappa^2H^2}{4\pi^2}\{a^4\square-a^2\nabla^2\}\delta^4$\\
\eqI{59}, $a^D\square$ & $-\frac{D-1}{D-3}\pi\cot(\frac\pi2D)-2$ & $+\frac{D-1}{D-3}\pi\cot(\frac\pi2D)+2$\\
\eqI{59}, $a^{D-2}\nabla^2$ & $+\frac{D-1}{D-3}\pi\cot(\frac\pi2D)+\frac52$ & $-\frac{D-1}{D-3}\pi\cot(\frac\pi2D)-\frac72$\\
\end{tabular}
\end{ruledtabular}
\end{table}

\paragraph*{Diagnosis.}
(i) \emph{Overall sign.} The left-hand sides of the printed \eqI{54}, \eqI{57} and \eqI{58} carry the
signs of the printed \eqI{47}--\eqI{49}, and their right-hand sides are correct evaluations of those
left-hand sides. The sign slip of the step \eqI{46}$\to$\eqI{47}--\eqI{49} was therefore carried into
the rest of the 3-point calculation, which establishes the convention statement made at the beginning
of this section.
(ii) \emph{Arithmetic of \eqI{59}.} The paper's own \eqI{57} and \eqI{58} give
$-\frac12-2=-\frac52$ for the finite part of the $a^{D-2}\nabla^2$ coefficient (paper convention).
The printed \eqI{59} minus the sum of the paper's own terms is
\begin{equation}
-\frac{i\kappa^2H^2}{8\pi^2}\,a^2\nabla^2\delta^4(x-x'),
\label{eq:s4a:59diff}
\end{equation}
which has the form of the noninvariant counterterm vertex \eqI{145},
$i\alpha_3\kappa^2H^2a^{D-2}\nabla^2\delta^D(x-x')$. Since $-\frac1{8\pi^2}=-\frac{2}{16\pi^2}$, it is equivalent to a finite
counterterm shift $\delta\alpha_3=-2$ in the units $H^{D-4}/(4\pi)^{D/2}$ of Sec.~\ref{sec:s0s1s7s8:conventions} (paper convention).
This finite error is absorbed by the arbitrary finite part of the counterterm. Its only trace in the renormalized
result is a shift of the finite part of $\beta_3$ by $+2$ (paper convention); together with the
correction of \eqI{40} in the 4-point part, this changes $\beta_3$ from
$-\frac{4}{D-4}+\frac{58}{3}+2\gamma_{\rm E}$ to $-\frac{4}{D-4}+\frac{52}{3}+2\gamma_{\rm E}$ (paper convention, units of
$H^{D-4}/(4\pi)^{D/2}$), with the pole unchanged; see Sec.~\ref{sec:s5:beta}.

\paragraph*{Checks.}
The script \sIVacode{04\_three\_point\_local\_log/check\_54\_59.py}~\cite{ref:repo} checks:
(A) the local polynomial extracted directly from \eqref{eq:s4a:45} and Table~I, without the
intermediate step \eqref{eq:s4a:47-49}, is
$\frac{D-1}{D-3}P_0P'_0\,d_A+(\bar P\!\cdot\!\bar P')\,d_B-\frac{2(D-2)}{D-3}P_0P'_0\,d_C$ for $D=4,5,6$,
i.e.\ exactly the local part of \eqref{eq:s4a:Jterms};
(B) the rules \eqref{eq:s4a:R2} and \eqI{55}--\eqI{56} on a test function, with symbolic $D$ in the
powers of $a$;
(C) the sum \eqref{eq:s4a:59exact} with finite parts $-2$ and $+\frac52$, the paper's own terms adding up
to $-\frac52$ rather than the printed $-\frac72$, and the printed \eqI{54}, \eqI{57}, \eqI{58} being
$-1$ times the derived ones term by term (\eqI{54} at symbolic $D$; \eqI{57}--\eqI{58}, printed only at
$D=4$, through their $D=4$ coefficients).

\subsection{Logarithm contributions: Eqs.~(I.60)--(I.67)}
\label{sec:s4a:log}

The logarithm contributions come from the $AA$ terms \eqref{eq:s4a:47} with the outer $i\Delta_A$
replaced by $k\ln(aa')$, and the \emph{full} $\partial\partial'i\Delta_A$ inside (including the local part
of \eqref{eq:s4a:52}). With $\mathcal P\equiv(aa')^{D-2}\,k\,L$, they are
$-\kappa^2\mathcal V_1+\kappa^2\frac{D-1}{D-3}\mathcal W-\kappa^2\mathcal V_3-\kappa^2\mathcal V_4$, where
\begin{equation}
\begin{aligned}
\mathcal V_1&\equiv\nabla\!\cdot\!\nabla'\big[\mathcal P\,\nabla\!\cdot\!\nabla' i\Delta_A\big],&
\mathcal W&\equiv\partial_0\partial'_0\big[\mathcal P\,\partial_0\partial'_0 i\Delta_A\big],\\
\mathcal V_3&\equiv\partial_i\partial'_0\big[\mathcal P\,\partial_i\partial'_0 i\Delta_A\big],&
\mathcal V_4&\equiv\partial_0\partial'_i\big[\mathcal P\,\partial_0\partial'_i i\Delta_A\big].
\end{aligned}
\label{eq:s4a:logterms}
\end{equation}
The printed \eqI{63}--\eqI{66} have the opposite signs, $+,-,+,+$, inherited from the printed \eqI{47}.

\paragraph*{Eqs.~(I.60)--(I.62).}
Since $\sqrt{-g}\,g^{\mu\nu}=a^{D-2}\eta^{\mu\nu}$ and $\nabla^2[k\ln(aa')]=0$, the propagator equation \eqI{12}
gives
\begin{equation}
\partial_0\big(a^{D-2}\partial_0\,i\Delta_A\big)=-i\delta^D(x-x')+a^{D-2}\nabla^2A(y),
\qquad
\partial'_0\big(a'^{D-2}\partial'_0\,i\Delta_A\big)=-i\delta^D(x-x')+a'^{D-2}\nabla^2A(y),
\label{eq:s4a:60}
\end{equation}
which are \eqI{60} and \eqI{61}. With $\partial_0(a^p\partial_0l)=(D-1)H^2a^D$ and $A=i\Delta_A-kL$, applying
$\nabla^2$ gives the corollary
\begin{equation}
\partial_0\big(a^p\partial_0\nabla^2A\big)=-i\nabla^2\delta^D+a^p\nabla^4A,
\qquad
\partial'_0\big(a'^p\partial'_0\nabla^2A\big)=-i\nabla^2\delta^D+a'^p\nabla^4A .
\label{eq:s4a:corol}
\end{equation}
For $D=4+\epsilon$, the first three terms of \eqI{14} (with $c_0=\Gamma(\frac D2-1)$) expand as
$4/y$, $\frac4\epsilon+3-2\gamma_{\rm E}+2\ln\frac4y$ and $-\frac4\epsilon-4+2\gamma_{\rm E}$, all times
$\mathcal N\to H^2/16\pi^2$; the poles cancel and the $n$-th series term vanishes identically at $D=4$.
Hence
\begin{equation}
A(y)=\frac{H^2}{16\pi^2}\Big\{\frac4y-2\ln\frac y4-1+O(D-4)\Big\},
\label{eq:s4a:62}
\end{equation}
which is \eqI{62}. (The printed $c_0$ would differ only at order $D-4$ here.) As a check,
$(4y-y^2)A''+4(2-y)A'=3k$, i.e.\ $\square A=(D-1)kH^2$ at $D=4$.

\paragraph*{Distributional tools.}
For a function $f$ of time only,
$f(\eta')\partial_0\delta^D=f(\eta)\partial_0\delta^D+f'(\eta)\delta^D$ and
$f(\eta')\partial_0^2\delta^D=f\partial_0^2\delta^D+2f'\partial_0\delta^D+f''\delta^D$. As kernels,
$\partial'_0[F(\eta')\partial'_0\delta^D]\equiv\partial_0[F(\eta)\partial_0\delta^D]$ (both act on a test function
$g(x')$ as $\partial_0(F\partial_0g)$), and $f(\eta,\eta')\nabla^2\delta^D=f(\eta,\eta)\nabla^2\delta^D$; in
particular $L\nabla^2\delta^D=2l\nabla^2\delta^D$. From these and \eqref{eq:s4a:R2} follows the lemma
\begin{equation}
-2i\,\partial_0\big[a^pl\,\partial_0\delta^D\big]=i\,L\,\big(a^D\square-a^{D-2}\nabla^2\big)\delta^D+i(D-1)H^2a^D\delta^D ,
\label{eq:s4a:lemma}
\end{equation}
obtained from $L\,\partial_0(a^p\partial_0\delta^D)=2\,\partial_0[a^pl\,\partial_0\delta^D]+(D-1)H^2a^D\delta^D$, which in turn follows
from the first identity with $f=l(\eta')$, $f'=Ha$, $f''=H^2a^2$.

\paragraph*{The four terms at general $D$.}
(1) By \eqref{eq:s4a:hom}, $\nabla\!\cdot\!\nabla'i\Delta_A=-\nabla^2A$, and $\mathcal P$ is constant in space:
\begin{equation}
\mathcal V_1=k\,(aa')^{D-2}L\,\nabla^4A .
\label{eq:s4a:V1}
\end{equation}
(2) For $\mathcal W$ split $L=l+l'$. In the $l$ part, $a^pl$ depends only on $\eta$ and
$a'^p\partial_0\partial'_0i\Delta_A=\partial_0(a'^p\partial'_0i\Delta_A)$, so \eqref{eq:s4a:60} gives
$\partial_0[a^pl\,\partial_0(-i\delta^D+a'^p\nabla^2A)]$; the $l'$ part is its mirror image under
$x\leftrightarrow x'$. The delta-function pieces combine through the kernel symmetry and
\eqref{eq:s4a:lemma}; the $A$ pieces follow from \eqref{eq:s4a:corol} with $\partial_0l=Ha$:
\begin{equation}
\frac{\mathcal W}{k}=iL\big(a^D\square-2a^{D-2}\nabla^2\big)\delta^D+i(D-1)H^2a^D\delta^D
+(aa')^{D-2}L\,\nabla^4A+H(aa')^{D-2}\big(a\partial_0+a'\partial'_0\big)\nabla^2A .
\label{eq:s4a:W}
\end{equation}
(3)--(4) In $\mathcal V_3$ the outer $\partial_i$ passes through $\mathcal P$ and meets the inner $\partial_i$, so
$\mathcal V_3=k\,\partial'_0[(aa')^{D-2}L\,\partial'_0\nabla^2A]$; the primed corollary and
$\partial'_0l'=Ha'$ give
\begin{equation}
\frac{\mathcal V_3}{k}=-ia^{D-2}L\,\nabla^2\delta^D+(aa')^{D-2}L\,\nabla^4A+H(aa')^{D-2}a'\partial'_0\nabla^2A,
\label{eq:s4a:V3}
\end{equation}
and $\mathcal V_4$ is its mirror image, with $a'\partial'_0\to a\partial_0$ in the last term.
At $D=4$ the four contributions are
\begin{align}
-\kappa^2\mathcal V_1&=-\frac{\kappa^2H^2}{8\pi^2}\,L\,(aa')^2\nabla^4A,
\label{eq:s4a:63}\\
3\kappa^2\mathcal W&=\frac{3i\kappa^2H^2}{8\pi^2}\,L\,\{a^4\square-2a^2\nabla^2\}\delta^4+\frac{3\kappa^2H^2}{8\pi^2}\,L\,(aa')^2\nabla^4A
\notag\\
&\quad+\frac{9i\kappa^2H^4}{8\pi^2}\,a^4\delta^4+\frac{3\kappa^2H^3}{8\pi^2}\,(aa')^2\big(a\partial_0+a'\partial'_0\big)\nabla^2A,
\label{eq:s4a:64}\\
-\kappa^2\mathcal V_3&=+\frac{i\kappa^2H^2}{8\pi^2}\,L\,a^2\nabla^2\delta^4-\frac{\kappa^2H^2}{8\pi^2}\,L\,(aa')^2\nabla^4A
-\frac{\kappa^2H^3}{8\pi^2}\,(aa')^2a'\partial'_0\nabla^2A,
\label{eq:s4a:65}
\end{align}
and $-\kappa^2\mathcal V_4$ equals \eqref{eq:s4a:65} with $a'\partial'_0\to a\partial_0$; these correspond to
\eqI{63}--\eqI{66}, each of which is printed with the opposite overall sign.

\paragraph*{The sum.}
Our Table~\ref{tab:s4a:logsum} collects the coefficients at general $D$. Altogether,
\begin{equation}
\begin{aligned}
-iM^2_{\rm 3pt,log}=\kappa^2k\,\Big\{&\ln(aa')\Big[\frac{D-1}{D-3}\,a^D\square-\frac{4}{D-3}\,a^{D-2}\nabla^2\Big]i\delta^D
+\frac{(D-1)^2}{D-3}\,H^2a^D\,i\delta^D\\
&-\frac{2(D-4)}{D-3}\,(aa')^{D-2}\ln(aa')\,\nabla^4A
+\frac{2}{D-3}\,H(aa')^{D-2}\big(a\partial_0+a'\partial'_0\big)\nabla^2A\Big\}.
\end{aligned}
\label{eq:s4a:67D}
\end{equation}
Each term is either $\delta^D$-local with a finite coefficient, or a smooth function times derivatives
of a \emph{single} propagator. Since $A\sim y^{1-D/2}$ is locally integrable, its derivatives are
well-defined distributions without poles near $D=4$ (unlike the products of two propagators met in the
normal contributions). Thus the $\nabla^4A$ term, whose coefficient is $O(D-4)$, vanishes, and setting
$D=4$ elsewhere changes only terms of order $D-4$. This confirms the statement of \PI{} that all
logarithm contributions are finite. The result is
\begin{equation}
\begin{aligned}
-iM^2_{\rm 3pt,log}(x;x')=\frac{\kappa^2H^2}{8\pi^2}\Big\{&\ln(aa')\big[3a^4\square-4a^2\nabla^2\big]i\delta^4(x-x')
+9H^2a^4\,i\delta^4(x-x')\\
&+2H(aa')^2\big(a\partial_0+a'\partial'_0\big)\nabla^2A(y)+O(D-4)\Big\}
\end{aligned}
\label{eq:s4a:67}
\end{equation}
(definition sign). \PI{} prints \eqI{67} with the bracket $\{\ln(aa')[-3a^4\square+4a^2\nabla^2]i\delta^4-9H^2a^4i\delta^4
-2H(aa')^2(a\partial_0+a'\partial'_0)\nabla^2A\}$, i.e.\ exactly $-1$ times \eqref{eq:s4a:67}; the cancellation of
the $\nabla^4A$ terms and the coefficients $3,4,9,2$ agree. No new error occurs in \eqI{60}--\eqI{67}:
the printed \eqI{67} is the correct sum of the printed \eqI{63}--\eqI{66}, whose left-hand sides carry
the signs of the printed \eqI{47}. The nonlocal term of \eqref{eq:s4a:67} is processed together with
the normal contributions in Sec.~\ref{sec:s4b:tableIII}, where it yields \eqI{89} and the
corresponding rows of Table~III.

\begin{table}[htbp]
\caption{Coefficients of the logarithm contributions in units of $\kappa^2k$ (definition sign), from
Eqs.~\eqref{eq:s4a:V1}--\eqref{eq:s4a:V3}. In the last row $-\mathcal V_4$ contributes to the $a\partial_0$
term and $-\mathcal V_3$ to its mirror $a'\partial'_0$, so each receives $\frac{D-1}{D-3}-1$.}
\label{tab:s4a:logsum}
\renewcommand{\arraystretch}{1.4}
\begin{ruledtabular}
\begin{tabular}{lcccccc}
Structure & $-\mathcal V_1$ & $\frac{D-1}{D-3}\mathcal W$ & $-\mathcal V_3$ & $-\mathcal V_4$ & Total & $D=4$\\
\hline
$L\,a^D\square\,i\delta^D$ & --- & $\frac{D-1}{D-3}$ & --- & --- & $\frac{D-1}{D-3}$ & $3$\\
$L\,a^{D-2}\nabla^2\,i\delta^D$ & --- & $-\frac{2(D-1)}{D-3}$ & $+1$ & $+1$ & $-\frac{4}{D-3}$ & $-4$\\
$H^2a^D\,i\delta^D$ & --- & $\frac{(D-1)^2}{D-3}$ & --- & --- & $\frac{(D-1)^2}{D-3}$ & $9$\\
$(aa')^{D-2}L\,\nabla^4A$ & $-1$ & $\frac{D-1}{D-3}$ & $-1$ & $-1$ & $-\frac{2(D-4)}{D-3}$ & $0$\\
$H(aa')^{D-2}a\partial_0\nabla^2A$ (and mirror) & --- & $\frac{D-1}{D-3}$ & $-1$ & $-1$ & $\frac{2}{D-3}$ & $2$\\
\end{tabular}
\end{ruledtabular}
\end{table}

\paragraph*{Cancellation against the local 4-point terms.}
The corrected 4-point result at general $D$ [Eq.~\eqref{eq:s3:I44exact}; definition sign] is
\begin{equation}
-iM^2_{\rm 4pt}=-i\kappa^2k\Big\{\tfrac14D(D-1)\big[\pihat+\ln(aa')\big]a^D\square-D\big[\pihat+\ln(aa')\big]a^{D-2}\nabla^2
+\tfrac14D(D-1)^2H^2a^D\Big\}\delta^D+O(D-4).
\label{eq:s4a:44D}
\end{equation}
Adding the local terms of \eqref{eq:s4a:67D} to its $\ln(aa')$ and $H^2$ terms, and using
$\frac{1}{D-3}-\frac D4=-\frac{(D-4)(D+1)}{4(D-3)}$ and $D-\frac{4}{D-3}=\frac{(D-4)(D+1)}{D-3}$, gives
\begin{equation}
-i\kappa^2k\,\frac{(D-4)(D+1)}{4(D-3)}\Big\{(D-1)\ln(aa')\,a^D\square-4\ln(aa')\,a^{D-2}\nabla^2+(D-1)^2H^2a^D\Big\}\delta^D
\;\xrightarrow{D\to4}\;0 .
\label{eq:s4a:cancel}
\end{equation}
The remainder is not multiplied by a pole, so it vanishes at $D=4$: the terms
$3\ln(aa')a^D\square$, $-4\ln(aa')a^{D-2}\nabla^2$ and $9H^2a^D$, which are not of counterterm form, cancel
exactly between \eqI{44} and \eqI{67}. With the opposite relative sign they would add instead. Since
both printed parts are reversed, the statement of \PI{} after \eqI{67} holds in the paper convention
as well. (The divergent $\pihat$ terms of \eqI{44} do not take part in this cancellation. The constants
$\mp2$ in the brackets of the printed \eqI{44}, which originate from the finite-part error in \eqI{40},
are absent from the corrected 4-point result; they are of counterterm form.)

\paragraph*{Checks.}
The script \sIVacode{04\_three\_point\_local\_log/check\_60\_67.py}~\cite{ref:repo} checks:
(a) the expansion \eqref{eq:s4a:62} from \eqI{14}, with the printed $c_0$ differing only at $O(D-4)$;
(b) the differential equation satisfied by \eqref{eq:s4a:62};
(c) the lemma \eqref{eq:s4a:lemma} and the time-only rule for $L\nabla^2\delta^D$ on a test function at
general $D$;
(d) a pointwise check of the nonlocal parts of $\mathcal V_1,\mathcal W,\mathcal V_3,\mathcal V_4$ at $D=4$,
by brute-force differentiation of $A(y)+k\ln(aa')$ with the actual $A(y)$ of \eqref{eq:s4a:62}, at three
random non-coincident points with 30-digit precision (exact agreement; this check is numerical, at
sample points, and at $D=4$ only);
(e) the general-$D$ coefficients of our Table~\ref{tab:s4a:logsum} and the remainder \eqref{eq:s4a:cancel};
(f) the printed \eqI{63}--\eqI{66} are $-1$ times the derived terms, and their sum is the printed \eqI{67}.

\subsection{Status of Eqs.~(I.45)--(I.67)}
\label{sec:s4a:status}

Our Table~\ref{tab:s4a:status} summarizes the comparison. Apart from the overall sign, the only item in
this part that changes a number entering the renormalized result is the $-\frac72\to-\frac52$ slip in
\eqI{59}, and it is absorbable (B2). The overall sign rests on the paper's own definition \eqI{168}; once
that definition is adopted it is physical for the nonlocal part, which has no counterterm freedom: the
physical nonlocal self-mass is $-1$ times the (corrected) Table~X\@. In \PII~\cite{KW08} this flips every entry of Table~5 and the sign of the
one-loop particular solution; since the non-falling total is zero, only decaying terms change sign, and
the late-time conclusion of \PII{} is unchanged (Sec.~\ref{sec:s6}).

\begin{table}[htbp]
\caption{Status of the equations of this part. ``S1 (inherited sign)'': the printed expression is $-1$ times the
definition sign, inherited from the step \eqI{46}$\to$\eqI{47}--\eqI{49}; this is the 3-point part of the overall sign S1,
classified as substantive because the sign of the nonlocal part of the self-mass is physical. C$n$ and B$n$ refer to our
Tables~\ref{tab:s0s1s7s8:errata} and~\ref{tab:s0s1s7s8:finite}.}
\label{tab:s4a:status}
\renewcommand{\arraystretch}{1.25}
\begin{ruledtabular}
\begin{tabular}{>{\raggedright\arraybackslash}p{0.2\textwidth}>{\raggedright\arraybackslash}p{0.5\textwidth}>{\raggedright\arraybackslash}p{0.2\textwidth}}
Equations & Finding & Class\\
\hline
\eqI{7}, left side & $-\frac12\kappa\to+\frac12\kappa$; \eqI{46} uses the correct vertex & typo (C3)\\
\eqI{45}--\eqI{46} & reproduced exactly (definition sign) & confirmed\\
\eqI{47}--\eqI{49} & structures, magnitudes, $D$-dependence agree; overall sign reversed in all nine terms & S1 (origin of the 3-point sign)\\
\eqI{14}, first coefficient & $\Gamma(\frac D2-1)/(\frac D2-1)\to\Gamma(\frac D2-1)$; \eqI{52} not affected & typo (C4)\\
\eqI{50}--\eqI{53} & confirmed; \eqI{52} derived along three independent routes & confirmed\\
text before \eqI{54} & ``4th term of (47)'' $\to$ ``2nd term of (47)'' & typo (C7)\\
\eqI{54}, \eqI{57}, \eqI{58} & reproduced up to the overall sign & S1 (inherited sign)\\
\eqI{55}--\eqI{56} & reproduced & confirmed\\
\eqI{59} & overall sign; finite constant $-\frac72\to-\frac52$ (paper convention), equivalent to $\delta\alpha_3=-2$ in units of $H^{D-4}/(4\pi)^{D/2}$ & S1 (inherited sign); absorbable (B2)\\
\eqI{60}--\eqI{62} & reproduced & confirmed\\
\eqI{63}--\eqI{67} & reproduced up to the overall sign; no new error & S1 (inherited sign)\\
cancellation after \eqI{67} & confirmed at general $D$, Eq.~\eqref{eq:s4a:cancel} & confirmed\\
\end{tabular}
\end{ruledtabular}
\end{table}

\providecommand{\sivbdd}{\nabla\!\cdot\!\nabla'}
\providecommand{\sivbAl}{A_{\rm lead}}
\providecommand{\sivbAh}{\hat A}
\providecommand{\sivbO}{O}
\providecommand{\sivbF}{\mathcal{F}}

\section{The 3-point contribution: normal terms}
\label{sec:s4b}

This section corresponds to Sec.~III\,B.3 of \PI{} (``Normal contributions''), Eqs.~\eqI{68}--\eqI{125} and Tables~III--V,
together with its Appendix ``Extracting derivatives'', Eqs.~\eqI{A1}--\eqI{A38} and Tables~XI--XX. The setup, the local and the
logarithm terms of the 3-point contribution were treated in Sec.~\ref{sec:s4a}.

\subsection{Normal contributions: extracting derivatives, Eqs.~(I.68)--(I.80)}
\label{sec:s4b:extract}

The \emph{normal} contributions are the 3-point terms in which a pair of inner derivatives acts on one
propagator function $A(y)$ and the result multiplies a second propagator function $f(y)\in\{A,B,C\}$, after the
local and logarithm pieces of \eqI{52} have been split off, cf.\ \eqI{53}. \PI~\cite{KW07} reduces them by
extracting the derivatives \emph{generically}, so that only functions of $y$ remain, acted upon by ten
external operators. In this and the following subsections we follow this reduction through Tables~III, IV and V
of \PI\@. The detailed derivations and the audit are in Chapters~5 and 6 of the derivation notes, and the scripts quoted below are
in \texttt{code/05\_normal\_terms\_tableIII/} and \texttt{code/06\_tables\_IV\_V/} of the
repository~\cite{ref:repo}.

\paragraph*{Notation.}
We write $\Delta x^\mu\equiv x^\mu-x'^\mu$, $\Delta\eta\equiv\eta-\eta'$, and $\Delta x_\sigma\equiv\eta_{\sigma\nu}\Delta x^\nu$, so that
$\Delta x_0=-\Delta\eta$. The length function is
$y=aa'H^2\{\|\Delta\vec x\|^2-(|\Delta\eta|-i\delta)^2\}$, cf.\ \eqI{11}, with $a=-1/(H\eta)$ and $\partial_0a=Ha^2$.
Primes on functions of $y$ denote $d/dy$; $\partial$ and $\partial'$ act on $x$ and $x'$. The constants are
\begin{equation}
\mathcal N\equiv\frac{H^{D-2}}{(4\pi)^{D/2}},\qquad
k\equiv\mathcal N\frac{\Gamma(D-1)}{\Gamma(\frac D2)}\xrightarrow{D\to4}\frac{H^2}{8\pi^2},\qquad
\pihat\equiv-\pi\cot\Big(\frac{\pi D}{2}\Big)=-\frac{2}{D-4}+\sivbO(D-4),\qquad
h\equiv\frac{H^2}{16\pi^2}.
\label{eq:s4b:constants}
\end{equation}
The leading term of $A(y)$ is $\sivbAl\equiv\mathcal N\,\Gamma(\frac D2-1)\,(4/y)^{\frac D2-1}$. Here the first coefficient
of \eqI{14} is read as $\Gamma(\frac D2-1)$ (the printed value is a misprint, C4; Sec.~\ref{sec:s2:scalar}). At $D=4$, $A=h\{4/y-2\ln(y/4)-1\}$, cf.\ \eqI{62}, and $B=C=4h/y$, because every
series term of \eqI{27}--\eqI{28} vanishes there. Table numbers without qualification (Table~III, IV, \dots) refer to
\PI; tables of the present paper are referred to as ``our Table''.

\paragraph*{Conventions.}
The identities \eqI{68}--\eqI{88} and \eqI{90}--\eqI{103} are exact in general $D$ and do not depend on the overall
sign of the self-mass. Only \eqI{89}, Table~III, and the tables assembled from the parts $P_1,\dots,P_7$ defined
below (Tables~IV and V) carry that sign. As shown in Secs.~\ref{sec:s4a:graviton}--\ref{sec:s4a:log}, the printed 3-point terms
\eqI{47}--\eqI{49} are $-1$ times those that follow from the definition \eqI{168}; the 4-point part is reversed in
the same way, so the relative sign of the two parts is consistent. We give \eqI{89} and Table~III in both
conventions. Tables~IV and V are quoted in the paper convention, which allows an entry-by-entry comparison; in the
definition sign every entry is multiplied by $-1$.

\paragraph*{Indefinite integrals, cf.\ Eq.~(I.68).}
\begin{equation}
I[f](y)\equiv\int^y\!dy'\,f(y'),\qquad \frac{d}{dy}I[f]=f,\qquad I^n\equiv I\circ\cdots\circ I .
\label{eq:s4b:Idef}
\end{equation}
The integral $I[f]$ is defined up to a constant and $I^2[f]$ up to $c_1y+c_0$. In each identity below this ambiguity cancels
when the constants are chosen consistently; e.g.\ in \eqref{eq:s4b:I70} the shift
$I^2[fA'']\to I^2[fA'']+c_1y$ is compensated by $I[f'A']\to I[f'A']-c_1$.

\paragraph*{Chain rule and delta-function content, cf.\ Eq.~(I.69).}
For smooth $F$ the chain rule holds pointwise,
\begin{equation}
\partial_\rho\partial'_\sigma F(y)=F''(y)\,\frac{\partial y}{\partial x^\rho}\frac{\partial y}{\partial x'^\sigma}
+F'(y)\,\frac{\partial^2y}{\partial x^\rho\partial x'^\sigma}.
\label{eq:s4b:chain}
\end{equation}
\PI{} states that two derivatives acting on a product of two propagators can never produce a delta function. We
find that this holds for the \emph{naive} (pointwise) derivative. At fixed $\delta>0$ the kink of $|\Delta\eta|$
produces in addition the singular term
$\delta^0_\rho\delta^0_\sigma\,(-4i\delta\,a^2H^2)\,F'(y)|_{\Delta\eta=0}\,\delta(\Delta\eta)$, where
$y|_{\Delta\eta=0}=a^2H^2(r^2+\delta^2)$ and $r=\|\Delta\vec x\|$. For a component $F\supset y^s$ its spatial integral
scales as $\delta\int d^{D-1}x\,(r^2+\delta^2)^{s-1}\propto\delta^{D+2s-2}$. This vanishes as $\delta\to0$ (by analytic
continuation in $D$) unless $s=1-\frac D2$, in which case a finite $\delta^D$ term survives. Hence
\begin{equation}
\big[\partial_\rho\partial'_\sigma F\big]_{\rm distr}=\big[\partial_\rho\partial'_\sigma F\big]_{\rm naive}
+\delta^0_\rho\delta^0_\sigma\,\gamma_F\,\frac{i}{a^{D-2}}\,\delta^D(x-x'),
\qquad F\supset\gamma_F\,\sivbAl ,
\label{eq:s4b:distr}
\end{equation}
normalized so that $F=\sivbAl$ reproduces the local coefficient $i\,a^{2-D}$ of \eqI{52}.

\paragraph*{The basic identity, cf.\ Eq.~(I.70).}
With $G=I^2[fA'']$ (so $G''=fA''$) in \eqref{eq:s4b:chain} and the integration by parts
$I[fA'']=fA'-I[f'A']$,
\begin{equation}
f(y)\Big\{A''(y)\frac{\partial y}{\partial x^\rho}\frac{\partial y}{\partial x'^\sigma}
+A'(y)\frac{\partial^2y}{\partial x^\rho\partial x'^\sigma}\Big\}
=\partial_\rho\partial'_\sigma I^2[fA''](y)+\frac{\partial^2y}{\partial x^\rho\partial x'^\sigma}\,I[f'A'](y),
\label{eq:s4b:I70}
\end{equation}
identical to the printed \eqI{70}.

\paragraph*{Delta content and the naive-derivative rule.}
Write $f=c_f+(\text{powers of }y)$, where $c_f$ is the constant part of $f$: $c_f=k\pihat$ for $f=A$, and
$c_f=b,\,c$, the coincidence limits of $i\Delta_B$ and $i\Delta_C$, for $f=B,\,C$. In general $D$ the $D$-dependent
powers do not mix with the integer powers, so the $y^{1-D/2}$ component of $I^2[fA'']$ comes only from
$c_f\,\sivbAl''$:
\begin{equation}
I^2[fA'']\supset c_f\,\sivbAl\;(+\text{linear in }y)\qquad\Longrightarrow\qquad\gamma_{I^2[fA'']}=c_f .
\label{eq:s4b:gamma}
\end{equation}
The function $I[f'A']$ has no $y^{-D/2}$ component, so its product with the $\sivbO(\delta)$ singular part of
$\partial_0\partial'_0y$ gives no delta function. Consequently \eqref{eq:s4b:I70} holds both with naive derivatives on
both sides and with distributional derivatives on both sides: the distributional left-hand side gains
$f\,i\,a^{2-D}\delta^D=c_f\,i\,a^{2-D}\delta^D$ (in dimensional regularization $f(y)\delta^D=f(0)\delta^D=c_f\delta^D$), and
by \eqref{eq:s4b:distr}--\eqref{eq:s4b:gamma} the right-hand side gains the same term. This delta content is exactly the local contribution \eqI{54}--\eqI{58} computed
separately in Sec.~\ref{sec:s4a:local}; for $f=A$ it is \eqI{54}, with its $k\pihat$ pole. The inner derivatives of
the normal contributions must therefore be taken \emph{naively}, as the footnote to \eqI{75} states; otherwise the
local term would be counted twice. The corresponding question for the \emph{outer} operators is taken up in
Sec.~\ref{sec:s4b:distr}.

\paragraph*{Mixed derivatives of $y$, cf.\ Eqs.~(I.71)--(I.73).}
For the regular part ($\delta\to0$),
\begin{equation}
\partial_0y=Ha\,y+2aa'H^2\Delta x_0,\qquad
\partial'_0y=Ha'\,y-2aa'H^2\Delta x_0,\qquad
\partial_iy=2aa'H^2\Delta x_i=-\partial'_iy ,
\label{eq:s4b:dy}
\end{equation}
and one more derivative gives the single expression
\begin{equation}
\frac{\partial^2y}{\partial x^\rho\partial x'^\sigma}
=H^2aa'\big\{y\,\delta^0_\rho\delta^0_\sigma-2a\,\delta^0_\rho H\Delta x_\sigma+2a'H\Delta x_\rho\,\delta^0_\sigma
-2\eta_{\rho\sigma}\big\},
\label{eq:s4b:I71}
\end{equation}
which is \eqI{71}, provided $\Delta x_\sigma$ is read with a lowered index (\PI{} does not state this explicitly). With
$H(a-a')=H^2aa'\Delta\eta$ [Eq.~\eqref{eq:s4b:aux} below] the components are
\begin{equation}
\begin{aligned}
&\frac{\partial^2y}{\partial x^0\partial x'^0}=H^2aa'\big[2-y+2aa'H^2\|\Delta\vec x\|^2\big],\qquad
\frac{\partial^2y}{\partial x^0\partial x'^j}=-2H^3a^2a'\Delta x_j,\\
&\frac{\partial^2y}{\partial x^i\partial x'^0}=2H^3aa'^2\Delta x_i,\qquad
\frac{\partial^2y}{\partial x^i\partial x'^j}=-2H^2aa'\eta_{ij},
\end{aligned}
\label{eq:s4b:I72}
\end{equation}
in agreement with \eqI{72}--\eqI{73}.

\paragraph*{Spatial identity, cf.\ Eq.~(I.74).}
From \eqref{eq:s4b:chain} and \eqref{eq:s4b:dy},
$\sivbdd\,I^2[f]=-4(aa'H^2)^2\|\Delta\vec x\|^2f-2(D-1)aa'H^2\,I[f]$. Since $aa'$ is constant in space,
\begin{equation}
aa'H^2\|\Delta\vec x\|^2f(y)=-\tfrac12(D-1)\,I[f](y)-\frac{\sivbdd}{4aa'H^2}\,I^2[f](y),
\label{eq:s4b:I74}
\end{equation}
which is \eqI{74}.

\paragraph*{Component identities, cf.\ Eqs.~(I.75)--(I.78).}
These are the components of \eqref{eq:s4b:I70} with \eqref{eq:s4b:I72}. For $\rho=\sigma=0$ the piece
$2H^2aa'\cdot aa'H^2\|\Delta\vec x\|^2\,I[f'A']$ is converted with \eqref{eq:s4b:I74} applied to $I[f'A']$. The result is
\begin{align}
f\,\partial_0\partial'_0A&=\partial_0\partial'_0I^2[fA'']-\tfrac12\sivbdd\,I^3[f'A']
+H^2aa'\big\{(2-y)\,I[f'A']-(D-1)\,I^2[f'A']\big\},\label{eq:s4b:I75}\\
f\,\partial_0\partial'_jA&=\partial_0\partial'_jI^2[fA'']+Ha\,\partial'_jI^2[f'A'],\label{eq:s4b:I76}\\
f\,\partial_i\partial'_0A&=\partial_i\partial'_0I^2[fA'']+Ha'\,\partial_iI^2[f'A'],\label{eq:s4b:I77}\\
f\,\partial_i\partial'_jA&=\partial_i\partial'_jI^2[fA'']-2H^2aa'\eta_{ij}\,I[f'A'],\label{eq:s4b:I78}
\end{align}
identical to \eqI{75}--\eqI{78}. On the left the second derivatives are naive (footnote to \eqI{75}).

\paragraph*{The first $A$ term, cf.\ Eqs.~(I.79)--(I.80).}
Take $f=A$ and the trace of \eqref{eq:s4b:I78}. Passing the outer $\sivbdd$ through the spatially constant
$(aa')^{D-2}$ and using $\sivbdd F(y)=-\nabla^2F(y)$,
\begin{equation}
\sivbdd\big[(aa')^{D-2}A\,\sivbdd A\big]=(aa')^{D-2}\nabla^4I^2[AA''](y)+2(D-1)H^2(aa')^{D-1}\nabla^2I[A'^2](y),
\label{eq:s4b:I80}
\end{equation}
in agreement with \eqI{79}--\eqI{80}. In \eqI{47} this bracket carries the coefficient $+\kappa^2$ as printed, and
$-\kappa^2$ in the definition sign.

Since $y^s$ is locally integrable in $D$ dimensions if and only if $s>-\frac D2$, the powers present are
\begin{equation}
I^2[AA'']:\ y^{2-D},\,y^{3-D},\,y^{4-D},\,y^{1-\frac D2},\,y^{2-\frac D2},\dots;\qquad
I[A'^2]:\ y^{1-D},\,y^{2-D},\,y^{3-D},\,y^{1-\frac D2},\,y^{2-\frac D2},\dots.
\end{equation}
At $D=4$ only $y^{2-D}$ (marginal) in $I^2[AA'']$, and $y^{1-D}$, $y^{2-D}$ in $I[A'^2]$, fail to be integrable. This
confirms the divergence count after \eqI{80}. One qualification: in $I^2[AA'']$ the coefficients of $y^{1-D/2}$ and
$y^{3-D}$ each carry a pole, but their sum is finite at $D=4$ and gives a logarithm, as in \eqI{62}; ``setting $D=4$
in the higher terms'' means taking these two terms together.

\paragraph*{Code check.}
The script \texttt{check\_68\_80.py} uses $f=y^\alpha$ and $A=y^\beta$ with generic exponents (so no integration constants
arise), explicit conformal coordinates and the exact $y(x;x')$. It confirms \eqref{eq:s4b:I70}--\eqref{eq:s4b:I80} at
$D=4$ and $D=5$ at random rational points to 25 digits, verifies \eqref{eq:s4b:gamma} symbolically for a model $f$,
$A$, and generates the power lists above with symbolic $D$. No discrepancy with \PI{} was found in
\eqI{68}--\eqI{80}.

\subsection{Time derivatives, Eq.~(I.89) and Table III, Eqs.~(I.81)--(I.93)}
\label{sec:s4b:tableIII}

\paragraph*{First derivatives, cf.\ Eqs.~(I.81)--(I.83).}
Because $y$ depends on the spatial coordinates only through $\Delta\vec x$, $\partial'_if(y)=-\partial_if(y)$ \eqI{81}.
From $a=-1/(H\eta)$,
\begin{equation}
H\Delta\eta=\frac{a-a'}{aa'},\qquad a'H\Delta\eta=1-\frac{a'}{a},\qquad aH\Delta\eta=\frac{a}{a'}-1,
\label{eq:s4b:aux}
\end{equation}
and therefore
\begin{equation}
\frac{\partial y}{\partial x^0}=Ha\Big(y-2+2\frac{a'}{a}\Big),\qquad
\frac{\partial y}{\partial x'^0}=Ha'\Big(y-2+2\frac{a}{a'}\Big),
\label{eq:s4b:I82}
\end{equation}
which are \eqI{82}--\eqI{83}.

\paragraph*{Three identities, cf.\ Eqs.~(I.84)--(I.88).}
Adding the two time derivatives,
\begin{equation}
(\partial_0+\partial'_0)f(y)=H(a+a')\,y\,f'(y).
\label{eq:s4b:I84}
\end{equation}
The weighted combination gives $(a'\partial_0+a\partial'_0)f=2Haa'\cdot aa'H^2\|\Delta\vec x\|^2f'$ \eqI{85}, where
\eqref{eq:s4b:aux} was used in the form $2H\Delta\eta\,(a-a')=2H^2aa'\Delta\eta^2$. Applying \eqref{eq:s4b:I74} to $f'$
(the integration constants drop out because $\sivbdd y=-2(D-1)aa'H^2$) and $\sivbdd\to-\nabla^2$,
\begin{equation}
(a'\partial_0+a\partial'_0)f(y)=-(D-1)Haa'f(y)+\frac{\nabla^2}{2H}\,I[f](y).
\label{eq:s4b:I86}
\end{equation}
Finally, $a\partial_0+a'\partial'_0=(a+a')(\partial_0+\partial'_0)-(a'\partial_0+a\partial'_0)$ \eqI{87} gives
\begin{equation}
(a\partial_0+a'\partial'_0)f(y)=H(a+a')^2y\,f'(y)+(D-1)Haa'f(y)-\frac{\nabla^2}{2H}\,I[f](y).
\label{eq:s4b:I88}
\end{equation}
Equations~\eqref{eq:s4b:I84}--\eqref{eq:s4b:I88} agree with \eqI{84}--\eqI{88}. They hold in general $D$ and do not depend on
the integration constant.

\paragraph*{The nonlocal logarithm term, cf.\ Eq.~(I.89).}
In the definition sign, the re-derived \eqI{67} contains the nonlocal term
\begin{equation}
-iM^2_{\rm 3pt}\supset\kappa^2k\,\frac{2}{D-3}\,H(aa')^{D-2}\big(a\partial_0+a'\partial'_0\big)\nabla^2A(y)
\;\xrightarrow{D\to4}\;\frac{\kappa^2H^2}{8\pi^2}\,2H(aa')^2\big(a\partial_0+a'\partial'_0\big)\nabla^2A(y),
\label{eq:s4b:log67}
\end{equation}
which \PI{} prints with $-2H$ in place of $2H$ (paper convention). Since $a\partial_0+a'\partial'_0$ commutes with
$\nabla^2$, Eq.~\eqref{eq:s4b:I88} gives, in general $D$,
\begin{equation}
\kappa^2k\,\frac{2}{D-3}\,H(aa')^{D-2}\,\nabla^2\Big[H(a+a')^2yA'+(D-1)Haa'A-\frac{\nabla^2}{2H}\,I[A]\Big],
\label{eq:s4b:I89D}
\end{equation}
and at $D=4$ (definition sign)
\begin{equation}
\frac{\kappa^2H^2}{8\pi^2}\,2H(aa')^2\big(a\partial_0+a'\partial'_0\big)\nabla^2A
=\frac{\kappa^2H^2}{16\pi^2}\Big\{12(aa')^3H^2\nabla^2A+4H^2(aa')^2(a+a')^2\nabla^2(yA')-2(aa')^2\nabla^4I[A]\Big\}.
\label{eq:s4b:I89}
\end{equation}
\PI{} prints, in the paper convention,
\begin{equation*}
\frac{\kappa^2H^2}{8\pi^2}\,(-2H)(aa')^2\big(a\partial_0+a'\partial'_0\big)\nabla^2A
=\frac{\kappa^2H^2}{16\pi^2}\Big\{-12(aa')^3H^2\nabla^2A-4(aa')^2(a+a')^2\nabla^2(yA')+2(aa')^2\nabla^4I[A]\Big\}.
\end{equation*}
Apart from the inherited overall sign, the second term lacks a factor $H^2$: as printed it has dimension $L^{-2}$
times that of $A$, whereas the other two terms have $L^{-4}$. Table~III carries the $H^2$, so this is a misprint (C8). A
numerical check at test points with $H\neq1$ confirms that \eqI{89} holds with $H^2$ restored and fails without it.

\paragraph*{Table III.}
With $A=h\,\sivbAh$, $\sivbAh\equiv4/y-2\ln(y/4)-1$, the prefactor becomes
$\frac{\kappa^2H^2}{16\pi^2}\,h=\frac{\kappa^2H^4}{(4\pi)^4}$, and the coefficients are read off from the brace of
\eqref{eq:s4b:I89} with $A\to\sivbAh$ and $x\equiv y/4$. Constants are dropped under $\nabla^2$ and $\nabla^4$, and since
$\nabla^2y=2(D-1)aa'H^2$ is constant in space, $\nabla^4y=0$ (the ``terms which eventually vanish such as
$\nabla^4y$'' of \PI). Thus $12\sivbAh\to\frac{12}{x}-24\ln x$; $y\sivbAh'=-\frac4y-2$, so $4y\sivbAh'\to-\frac4x$; and
$I[\sivbAh]=4\ln y-2y\ln\frac y4+y$, so $-2I[\sivbAh]\to-8\ln x+16x\ln x$. The result is our Table~\ref{tab:s4b:III}.
The same computation applied to the printed \eqI{89}, with $H^2$ restored, reproduces the printed Table~III exactly.

\begin{table}[tb]
\caption{Table III of \PI{} (nonlocal logarithm contributions), coefficients of $\kappa^2H^4/(4\pi)^4$, $x\equiv y/4$.
The derived table in the paper convention coincides with the printed one; in the definition sign every entry is
reversed (the inherited overall sign, not a new error).}
\label{tab:s4b:III}
\renewcommand{\arraystretch}{1.3}
\begin{ruledtabular}
\begin{tabular}{@{}llll@{}}
External operator & Derived, definition sign & Derived, paper convention & Printed (Table~III)\\
\hline
$(aa')^3H^2\nabla^2$ & $+\frac{12}{x}-24\ln x$ & $-\frac{12}{x}+24\ln x$ & $-\frac{12}{x}+24\ln x$\\
$(aa')^2(a+a')^2H^2\nabla^2$ & $-\frac4x$ & $+\frac4x$ & $+\frac4x$\\
$(aa')^2\nabla^4$ & $-8\ln x+16x\ln x$ & $+8\ln x-16x\ln x$ & $+8\ln x-16x\ln x$\\
\end{tabular}
\end{ruledtabular}
\end{table}

The functions $1/x$, $\ln x$ and $x\ln x$ are locally integrable at $D=4$, and the external operators contain only
spatial derivatives, so they produce no $\delta^D$ term: Table~III is purely nonlocal, as its caption states.
\PI{} states that Table~X includes Table~III\@. The first two rows are indeed added to the corresponding rows of
Table~X, but the third row, $8\ln x-16x\ln x$ (paper convention), is missing from the $(aa')^2\nabla^4$ row of
Table~X\@. This omission is non-analytic and changes the nonlocal self-mass; it is treated together with Table~X in
Sec.~\ref{sec:s5:TX}.

\paragraph*{Covariant d'Alembertian, cf.\ Eqs.~(I.90)--(I.93).}
From $a^D\square=\partial_\mu(a^{D-2}\eta^{\mu\nu}\partial_\nu)$ \eqI{42},
\begin{equation}
\square=-\frac{1}{a^2}\partial_0^2-\frac{(D-2)H}{a}\partial_0+\frac{1}{a^2}\nabla^2 .
\label{eq:s4b:I90}
\end{equation}
For the mixed $A$ term with inner $\partial_0\partial'_i$ (the fourth term of \eqI{47}), Eq.~\eqref{eq:s4b:I76} with $f=A$,
summed over $i$, and moving $\partial'_i$ outside gives, with $G\equiv I^2[AA'']$ and $K\equiv I^2[A'^2]$,
\begin{equation}
\partial_0\partial'_i\big[(aa')^{D-2}A\,\partial_0\partial'_iA\big]
=\nabla'^2\partial_0\big[(aa')^{D-2}\partial_0G\big]+H\nabla'^2\partial_0\big[a^{D-1}a'^{D-2}K\big]
\label{eq:s4b:I91}
\end{equation}
and, with \eqI{42} and $\nabla'^2=\nabla^2$,
\begin{equation}
\partial_0\partial'_i\big[(aa')^{D-2}A\,\partial_0\partial'_iA\big]
=-a^Da'^{D-2}\nabla^2\square G+(aa')^{D-2}\nabla^4G+Ha^{D-1}a'^{D-2}\nabla^2\partial_0K
+(D-1)H^2a^Da'^{D-2}\nabla^2K .
\label{eq:s4b:I92}
\end{equation}
On functions of $y$ alone, $\eta^{\mu\nu}\partial_\mu y\,\partial_\nu y/a^2=H^2(4y-y^2)$ and $\square y=DH^2(2-y)$, so
\begin{equation}
\square f(y)=H^2\big[(4y-y^2)f''(y)+D(2-y)f'(y)\big]=\square'f(y),
\label{eq:s4b:I93}
\end{equation}
the equality with $\square'$ following from $y(x;x')=y(x';x)$. Equations~\eqref{eq:s4b:I90}--\eqref{eq:s4b:I93} agree with
\eqI{90}--\eqI{93}.

\paragraph*{Distributional meaning of Eqs.~(I.92)--(I.93).}
The inner pair $\partial_0\partial'_i$ in \eqref{eq:s4b:I91} is mixed and contains no delta function even for the full
$i\Delta_A$, so the step \eqref{eq:s4b:I91}$\to$\eqref{eq:s4b:I92} is exact in the distributional sense and the
$\square$ in \eqref{eq:s4b:I92} is a \emph{distributional} operator. Since $G\supset c_f\sivbAl$ with $c_f=k\pihat$
[Eq.~\eqref{eq:s4b:gamma}] and $a^D\square\sivbAl\supset i\delta^D$, it contains the local term
\begin{equation}
-a^Da'^{D-2}\nabla^2\square I^2[AA'']\supset-ik\pihat\,a^{D-2}\nabla^2\delta^D(x-x'),
\label{eq:s4b:loc}
\end{equation}
which is not part of \eqI{54}--\eqI{59} (those contain only the inner $(0,0)$ terms). By contrast, \eqref{eq:s4b:I93}
is a pointwise identity and never produces $\delta^D$.\footnote{Pointwise only the leading power cancels for
$f=y^{1-D/2}$; one finds $\square\,y^{1-D/2}=\tfrac{D(D-2)}{4}H^2y^{1-D/2}$, a locally integrable function.} If
$\square I^2[AA'']$ in \eqref{eq:s4b:I92} is evaluated with \eqref{eq:s4b:I93}, the term \eqref{eq:s4b:loc} is lost.
Conversely, if a naive inner $(0,0)$ derivative, as in \eqref{eq:s4b:I75}, is recombined with outer derivatives into a
distributional $\square$, its local part is counted a second time. The consistent rule is that outer operators are
distributional while the local piece $c_f\,i\,a^{2-D}\delta^D$, already counted in \eqI{54}--\eqI{58}, is removed from
inner $(0,0)$ derivatives. Whether the reduction of \PI{} respects this rule is discussed in
Sec.~\ref{sec:s4b:distr}.

\paragraph*{Code check.}
The script \texttt{check\_81\_93.py} verifies \eqI{81}--\eqI{88} symbolically or at random points (25 digits), \eqI{90} in
general $D$, and \eqI{91}--\eqI{93} at $D=4,5$ by brute-force differentiation. With the true $A(y)$ of \eqI{62}, at three
test points with $H\neq1$, the derived nonlocal term of \eqI{67} plus the printed Table~III vanishes exactly (so
derived $=-$printed), \eqI{89} with $H^2$ restored equals $-1$ times the derived term, and the printed \eqI{89}
without $H^2$ fails at every point.

\subsection{Reduction to ten external operators, Eqs.~(I.94)--(I.115)}
\label{sec:s4b:operators}

\paragraph*{The second mixed $A$ term, cf.\ Eqs.~(I.94)--(I.96).}
Equation~\eqref{eq:s4b:I77} with $f=A$, summed over $i$, gives
\begin{align}
\partial_i\partial'_0\big[(aa')^{D-2}A\,\partial_i\partial'_0A\big]
&=\nabla^2\partial'_0\big[(aa')^{D-2}\partial'_0G\big]+H\nabla^2\partial'_0\big[a^{D-2}a'^{D-1}K\big]
\label{eq:s4b:I94}\\
&=-a^{D-2}a'^D\nabla^2\square'G+(aa')^{D-2}\nabla^4G+Ha^{D-2}a'^{D-1}\nabla^2\partial'_0K
+(D-1)H^2a^{D-2}a'^D\nabla^2K .
\label{eq:s4b:I95}
\end{align}
With $\square'G=\square G$ and Eq.~\eqref{eq:s4b:I86} for $f=K$,
$a\partial'_0K=-a'\partial_0K-(D-1)Haa'K+\frac{\nabla^2}{2H}I[K]$, the last term of \eqref{eq:s4b:I95} cancels and
\begin{equation}
\partial_i\partial'_0\big[(aa')^{D-2}A\,\partial_i\partial'_0A\big]
=-a^{D-2}a'^D\nabla^2\square G+(aa')^{D-2}\nabla^4G-Ha^{D-3}a'^D\nabla^2\partial_0K
+\tfrac12a^{D-3}a'^{D-1}\nabla^4I^3[A'^2].
\label{eq:s4b:I96}
\end{equation}
Equations~\eqref{eq:s4b:I94}--\eqref{eq:s4b:I96} agree with \eqI{94}--\eqI{96}. Adding \eqref{eq:s4b:I92}, the sum $P_2$ of the
two mixed $A$ terms takes the symmetric form [operators defined in Eq.~\eqref{eq:s4b:ops} below]
\begin{align}
P_2&=(-\delta+2\zeta)\,G+(D-1)\,\epsilon_2K+H(aa')^{D-2}\big(a\partial_0+a'\partial'_0\big)\nabla^2K\nonumber\\
&=(-\delta+2\zeta)\,I^2[AA'']+(D-1)(\epsilon_1+\epsilon_2)\,I^2[A'^2]+\epsilon_3\,yI[A'^2]-\tfrac12\zeta\,I^3[A'^2],
\label{eq:s4b:A9}
\end{align}
where \eqref{eq:s4b:I88} was used for $f=K$; this agrees with \eqI{A8}--\eqI{A9}.

\paragraph*{The ubiquitous reduction, cf.\ Eqs.~(I.97)--(I.101).}
Equation~\eqref{eq:s4b:I90} and its primed counterpart give
\begin{equation}
\partial_0\big[(aa')^{D-2}\partial_0f(y)\big]=-a^Da'^{D-2}\square f+(aa')^{D-2}\nabla^2f,\qquad
\partial'_0\big[(aa')^{D-2}\partial'_0f(y)\big]=-a^{D-2}a'^D\square f+(aa')^{D-2}\nabla^2f .
\label{eq:s4b:I97}
\end{equation}
Writing $\partial_0\partial'_0=\frac12(\partial_0+\partial'_0)^2-\frac12\partial_0^2-\frac12\partial_0'^2$ \eqI{99} inside
\begin{equation}
\partial_0\partial'_0\big[(aa')^pf\big]=(aa')^p\partial_0\partial'_0f+pH(aa')^p(a'\partial_0+a\partial'_0)f+p^2H^2(aa')^{p+1}f,
\qquad p\equiv D-2,
\label{eq:s4b:I100}
\end{equation}
and using $(\partial_0+\partial'_0)^2f=H^2(a^2+a'^2)yf'+H^2(a+a')^2(yf'+y^2f'')$,
$-\frac12(\partial_0^2+\partial_0'^2)f=\frac12(a^2+a'^2)\square f+\frac12pH(a\partial_0+a'\partial'_0)f-\nabla^2f$, and
\eqref{eq:s4b:I86}, \eqref{eq:s4b:I88}, we obtain
\begin{equation}
\begin{split}
\partial_0\partial'_0\big[(aa')^{D-2}f(y)\big]={}&\tfrac12(aa')^{D-2}(a^2+a'^2)\big[\square f+H^2yf'\big]
-(aa')^{D-2}\nabla^2f+\tfrac14(D-2)(aa')^{D-2}\nabla^2I[f]\\
&+\tfrac12(D-2)(D-3)H^2(aa')^{D-1}f+\tfrac12H^2(a+a')^2(aa')^{D-2}\big[(D-1)yf'+y^2f''\big],
\end{split}
\label{eq:s4b:I101}
\end{equation}
as printed in \eqI{101}.

\paragraph*{The two $B$ terms, cf.\ Eqs.~(I.102)--(I.103).}
Expanding $\partial_i(\partial_0+\partial'_0)[X(\partial_0+\partial'_0)\partial'_iA]$, $X\equiv(aa')^{D-2}B$, and using
$\partial'_i\to-\partial_i$ on functions of $y$,
\begin{equation}
\begin{split}
&\partial_i\partial'_0\big[(aa')^{D-2}B\,\partial_0\partial'_iA\big]+\partial_0\partial'_i\big[(aa')^{D-2}B\,\partial_i\partial'_0A\big]\\
&\qquad=\partial_i(\partial_0+\partial'_0)\big[(aa')^{D-2}B\,(\partial_0+\partial'_0)\partial'_iA\big]
-\partial_i\partial_0\big[(aa')^{D-2}B\,\partial_0\partial'_iA\big]-\partial_i\partial'_0\big[(aa')^{D-2}B\,\partial'_0\partial'_iA\big],
\end{split}
\label{eq:s4b:I102}
\end{equation}
as in \eqI{102}. The first piece is reduced with \eqref{eq:s4b:I84} and
$B\,\partial'_i(yA')=\partial'_iI[A'B+yA''B]$; the two subtracted pieces with \eqref{eq:s4b:I76}--\eqref{eq:s4b:I77} for
$f=B$, \eqref{eq:s4b:I97} and \eqref{eq:s4b:I88}. Collecting terms,
\begin{equation}
\begin{split}
={}&-(a^2+a'^2)(aa')^{D-2}\nabla^2\big\{\square I^2[A''B]+H^2I[A'B+yA''B]\big\}
+(aa')^{D-2}\nabla^4\big\{2I^2[A''B]-\tfrac12I^3[A'B']\big\}\\
&-H^2(a+a')^2(aa')^{D-2}\nabla^2\big\{(D-2)I[A'B+yA''B]+yA'B+y^2A''B-yI[A'B']\big\}\\
&+(D-1)H^2(a^2+aa'+a'^2)(aa')^{D-2}\nabla^2I^2[A'B'],
\end{split}
\label{eq:s4b:I103}
\end{equation}
which agrees term by term with \eqI{103}.

\paragraph*{External operators, cf.\ Eqs.~(I.104)--(I.113).}
The reduced normal terms are functions of $y$ acted upon by
\begin{equation}
\begin{aligned}
\alpha&\equiv(aa')^D\square^2, & \beta&\equiv(aa')^{D-1}(a^2+a'^2)H^2\square,\\
\gamma_1&\equiv(aa')^DH^4, & \gamma_2&\equiv(aa')^{D-1}(a^2+a'^2)H^4,\\
\gamma_3&\equiv(aa')^{D-1}(a+a')^2H^4=2\gamma_1+\gamma_2, & \delta&\equiv(aa')^{D-2}(a^2+a'^2)\nabla^2\square,\\
\epsilon_1&\equiv(aa')^{D-1}H^2\nabla^2, & \epsilon_2&\equiv(aa')^{D-2}(a^2+a'^2)H^2\nabla^2,\\
\epsilon_3&\equiv(aa')^{D-2}(a+a')^2H^2\nabla^2=2\epsilon_1+\epsilon_2, & \zeta&\equiv(aa')^{D-2}\nabla^4 .
\end{aligned}
\label{eq:s4b:ops}
\end{equation}
Under $\eta\to\lambda\eta$, $\vec x\to\lambda\vec x$ one has $a\to a/\lambda$ and $\nabla^2\to\nabla^2/\lambda^2$, while $H$,
$y$ and $\square$ are invariant; all ten operators have weight $\lambda^{-2D}$. The header of the appendix Table~XII
of \PI{} prints $\beta=(aa')^{D-2}(a^2+a'^2)H^2\square$, of weight $\lambda^{-2D+2}$; the correct factor is
$(aa')^{D-1}$, as in \eqI{105}, which the global check below also confirms.

\paragraph*{Differences and products, cf.\ Eqs.~(I.114)--(I.115).}
\PI{} introduces $\Delta B\equiv B-A$ and $\Delta C\equiv2\frac{D-2}{D-3}(C-A)$. The normalization of $\Delta C$ is chosen
because the $\partial_0\partial'_0$ terms of \eqI{47} combine as
$-\frac{D-1}{D-3}A+\frac{2(D-2)}{D-3}C=A+\Delta C$. The appendix totals then contain only the eight products
\begin{equation}
A'^2,\quad AA'',\quad A'B',\quad A''B,\quad A'\Delta B',\quad A''\Delta B,\quad A'\Delta C',\quad A''\Delta C .
\label{eq:s4b:products}
\end{equation}

\subsection{Expansions and the \texorpdfstring{$\alpha$}{alpha} example, Eqs.~(I.116)--(I.125)}
\label{sec:s4b:expansions}

With $z\equiv4/y$ and $p\equiv D-2$, only the two leading terms of each propagator function reach the powers $z^D$
and $z^{D-1}$ of the products:
\begin{equation}
A\simeq\mathcal N\big[c_0z^{p/2}+a_1z^{p/2-1}\big],\quad
B\simeq\mathcal N\big[c_0z^{p/2}+b_1z^{p/2-1}\big],\quad
C\simeq\mathcal N\big[c_0z^{p/2}+c_1z^{p/2-1}\big],
\end{equation}
with $c_0=\Gamma(\frac D2-1)$, $a_1=\Gamma(\frac D2+1)/(\frac D2-2)$, $b_1=\Gamma(\frac D2)$ and
$c_1=-(3-\frac D2)\Gamma(\frac D2-1)$. Using $\frac{d}{dy}z^q=-\frac q4z^{q+1}$ one obtains the coefficients of
our Table~\ref{tab:s4b:products}, in the unit $\frac{\Gamma^2(D/2)}{16}\frac{H^{2D-4}}{(4\pi)^D}$ of \eqI{116}--\eqI{123}. For
example, with $q\equiv\frac D2-1$ the coefficient of $z^{D-1}$ in $A'B'$ is $(q-1)(a_1+b_1)/\Gamma(q+1)=(q+1)+(q-1)=D-2$.

\begin{table}[tb]
\caption{Leading coefficients of the eight products \eqref{eq:s4b:products} in units of
$\frac{\Gamma^2(D/2)}{16}\frac{H^{2D-4}}{(4\pi)^D}$, and their $D=4$ closed forms (in units of $h^2/16$), with
$L\equiv\ln(y/4)$ and $z=4/y$. All entries agree with \eqI{116}--\eqI{123}. $^{*}$The $z^{D-1}$ coefficient carries a pole;
it combines with the constant ($\pihat$) term into $z^3L$ at $D=4$. These products enter Table~IV only through
$z^D$, because $I^2[z^{D-1}]\sim z^{D-3}$ is not divergent.}
\label{tab:s4b:products}
\renewcommand{\arraystretch}{1.3}
\begin{ruledtabular}
\begin{tabular}{@{}lccl@{}}
Product & $z^D$ (exact $D$) & $z^{D-1}$ (exact $D$) & $D=4$ closed form $(\times h^2/16)$\\
\hline
$A'^2$ & $1$ & $D$ & $z^4+4z^3+4z^2$\\
$AA''$ & $\frac{D}{D-2}$ & $\frac{D(D-2)}{D-4}{}^{*}$ & $2z^4-4z^3L-4z^2L-2z^2$\\
$A'B'$ & $1$ & $D-2$ & $z^4+2z^3$\\
$A''B$ & $\frac{D}{D-2}$ & $D$ & $2z^4+2z^3$\\
$A'\Delta B'$ & $0$ & $-2$ & $-2z^3-4z^2$\\
$A''\Delta B$ & $0$ & $-\frac{2D}{D-4}{}^{*}$ & $4z^3L+2z^3+4z^2L+2z^2$\\
$A'\Delta C'$ & $0$ & $-8$ & $-8z^3-16z^2$\\
$A''\Delta C$ & $0$ & $-\frac{8D}{D-4}{}^{*}$ & $16z^3L+8z^3+16z^2L+8z^2$\\
\end{tabular}
\end{ruledtabular}
\end{table}

The product $A'^2$ appears in the appendix totals only as $yI[\cdot]$, $y^2\,\cdot$ and $I^2[\cdot]$, where $z^{D-1}\to z^{D-3}$ is
not divergent; the $4z^3$ of \eqI{116} is the $D=4$ value of $Dz^{D-1}$. The products $A'B'$, $A'\Delta B'$ and
$A'\Delta C'$ appear with a single integration, so their exact-$D$ coefficients of $z^{D-1}$ are needed, and \PI{}
gives them exactly. For $A''B$ the exact-$D$ coefficient of $z^{D-1}$ is $D$ while the $D=4$ form has $2z^3$; there is
no conflict, since the constant term of $B$ times the leading term of $A''\propto z^{D/2+1}$ is a separate power in
general $D$ that coincides with $z^3$ at $D=4$, where it contributes $-2$. Only the $D=4$ value is needed, because
$A''B$ is always shielded by at least two powers of $y$.

\paragraph*{The $\alpha$ example, cf.\ Eqs.~(I.124)--(I.125).}
The total of the appendix Table~XI (operator $\alpha$) is $I^2[AA'']+I^2[\Delta C\,A'']$. With
our Table~\ref{tab:s4b:products},
\begin{equation}
\kappa^2\big\{I^2[AA'']+I^2[A''\Delta C]\big\}=\frac{\Gamma^2(\frac D2)}{16}\frac{\kappa^2H^{2D-4}}{(4\pi)^D}\,
I^2\Big[\frac{D}{D-2}z^D+12z^3L+8z^3+12z^2L+6z^2+\sivbO\Big(\frac{D-4}{y^3}\Big)\Big],
\label{eq:s4b:I124}
\end{equation}
as printed. With $I^2[z^D]=\frac{16}{(D-1)(D-2)}z^{D-2}$ for the exact-$D$ term, and two integrations at $D=4$ for the
rest,
\begin{equation}
=\frac{\kappa^2H^{2D-4}}{(4\pi)^D}\Gamma^2\Big(\frac D2\Big)\Big\{\frac{D}{(D-1)(D-2)^2}\Big(\frac4y\Big)^{D-2}
+6\Big(\frac4y\Big)\ln\frac y4+13\Big(\frac4y\Big)-6\ln^2\frac y4-18\ln\frac y4+\sivbO\Big(\frac{D-4}{y}\Big)\Big\},
\label{eq:s4b:I125}
\end{equation}
in agreement with \eqI{125} up to the constant $(24\ln2-36)\ln2$, which $\alpha=(aa')^D\square^2$ annihilates.

\paragraph*{Code check.}
The script \texttt{check\_94\_125.py} confirms \eqI{94}--\eqI{103} in explicit coordinates at $D=4$ and $D=5$ with
$A=y^\beta$, $B=y^\gamma$, $f=y^\alpha$ (generic rational exponents, 30-digit arithmetic), \eqI{99}, \eqI{108},
\eqI{112}, the exact-$D$ leading coefficients and $D=4$ closed forms of our Table~\ref{tab:s4b:products}, and
\eqI{124}--\eqI{125} (exact-$D$ term at $D=5,6,7,\frac92$; finite part up to the constant above).

\subsection{The appendix decomposition and the scope of its audit}
\label{sec:s4b:appendix}

\PI{} splits the normal contribution into seven parts, \eqI{A1}--\eqI{A7} of its appendix ``Extracting derivatives'':
\begin{equation}
\begin{aligned}
P_1&\equiv\sivbdd\big[(aa')^{D-2}A\,\sivbdd A\big],\qquad
P_2\equiv\partial_i\partial'_0\big[(aa')^{D-2}A\,\partial_i\partial'_0A\big]+\partial_0\partial'_i\big[(aa')^{D-2}A\,\partial_0\partial'_iA\big],\\
P_3&\equiv\partial_0\partial'_0\big[(aa')^{D-2}A\,\partial_0\partial'_0A\big],\qquad
P_4\equiv-\partial_0\partial'_0\big[(aa')^{D-2}B\,\sivbdd A\big],\\
P_5&\equiv-\partial_i\partial'_0\big[(aa')^{D-2}B\,\partial_0\partial'_iA\big]-\partial_0\partial'_i\big[(aa')^{D-2}B\,\partial_i\partial'_0A\big],\\
P_6&\equiv-\sivbdd\big[(aa')^{D-2}B\,\partial_0\partial'_0A\big],\qquad
P_7\equiv\partial_0\partial'_0\big[(aa')^{D-2}\Delta C\,\partial_0\partial'_0A\big],
\end{aligned}
\label{eq:s4b:Pk}
\end{equation}
with naive inner derivatives and the factors $\kappa^2$ suppressed. Their sum is the normal part of the
\emph{printed} \eqI{47}--\eqI{49} (paper convention), because
$-\frac{D-1}{D-3}P_3+\frac{2(D-2)}{D-3}\partial_0\partial'_0[(aa')^{D-2}C\,\partial_0\partial'_0A]=P_3+P_7$; in the
definition sign the normal part is $-\sum_kP_k$. The appendix uses \eqI{A17}--\eqI{A18}, i.e.\ \eqref{eq:s4b:I75} and
\eqref{eq:s4b:I78}; the generic reductions \eqI{A35}--\eqI{A38}; and \eqI{A9}, \eqI{A15} for $P_2$ and $P_5$. Each pair (part,
operator) gives one row of one of the Tables~XI--XX. As an example of a generic reduction, \eqI{A36} reads
\begin{equation}
\partial_0\partial'_0\big[(aa')^{D-1}H^2f(y)\big]=\Big\{\frac\beta2+\frac{(D-1)(D-2)}{2}\gamma_1+\frac{\gamma_2}{2}\,y\frac{\partial}{\partial y}
+\frac{\gamma_3}{2}\Big[(D-1)y\frac{\partial}{\partial y}+y^2\frac{\partial^2}{\partial y^2}\Big]
+\epsilon_1\Big[-1+\frac D4I\Big]\Big\}f(y).
\label{eq:s4b:A36}
\end{equation}
For instance, the part $P_{3c}$ contains $f=-(D-1)K-yI[A'^2]$ under
$\partial_0\partial'_0[(aa')^{D-1}H^2\,\cdot\,]$, so by \eqref{eq:s4b:A36} its $\gamma_3$ row is
$\frac12[(D-1)yf'+y^2f'']=-\frac12D(D-1)\,yI[A'^2]-D\,y^2A'^2-\frac12y^3(A'^2)'$, which is the $P_{3c}$ row of
Table~XV\@. Similarly, the rows $P_{3b}$ and $P_{4b}$ of the $\beta$ table give
$I[A'^2]+(D-1)I[A'B']=DI[A'B']-I[A'\Delta B']$.

\paragraph*{The ten totals.}
Denoting by $\sivbF_{\rm op}$ the total of the appendix table for operator op (Tables~XI--XX, as printed after the
five corrections of our Table~\ref{tab:s4b:apptypos}), the content of the appendix is
\begin{equation}
\sum_{k=1}^7P_k=\alpha\,\sivbF_\alpha+\beta\,\sivbF_\beta+\sum_{i=1}^3\gamma_i\,\sivbF_{\gamma_i}+\delta\,\sivbF_\delta
+\sum_{i=1}^3\epsilon_i\,\sivbF_{\epsilon_i}+\zeta\,\sivbF_\zeta ,
\label{eq:s4b:decomp}
\end{equation}
with
\begin{align}
\sivbF_\alpha={}&I^2[AA'']+I^2[\Delta C\,A''],\nonumber\\
\sivbF_\beta={}&DI[A'B']-I[A'\Delta B']-\tfrac12yI[A'^2]-\tfrac{D-1}{2}I^2[A'^2]+I[A'\Delta C']-\tfrac12yI[A'\Delta C']
-\tfrac{D-1}{2}I^2[A'\Delta C'],\nonumber\\
\sivbF_{\gamma_1}={}&(D-1)(D-2)\,\sivbF_\beta,\nonumber\\
\sivbF_{\gamma_2}={}&DyA'B'-yA'\Delta B'-\tfrac12y^2A'^2-\tfrac D2yI[A'^2]+yA'\Delta C'-\tfrac D2yI[A'\Delta C']
-\tfrac12y^2A'\Delta C',\nonumber\\
\sivbF_{\gamma_3}={}&(D-1)D\,yA'B'+Dy^2(A'B')'-(D-1)yA'\Delta B'-y^2(A'\Delta B')'-Dy^2A'^2-\tfrac12y^3(A'^2)'\nonumber\\
&-\tfrac12(D-1)D\,yI[A'^2]+(D-1)yA'\Delta C'+y^2(A'\Delta C')'-Dy^2A'\Delta C'-\tfrac12y^3(A'\Delta C')'\nonumber\\
&-\tfrac12(D-1)D\,yI[A'\Delta C'],\nonumber\\
\sivbF_\delta={}&2I^2[\Delta B\,A'']+\tfrac14I^3[A'^2]-I^2[\Delta C\,A'']+\tfrac14I^3[A'\Delta C'],\label{eq:s4b:totals}\\
\sivbF_{\epsilon_1}={}&-2(D-2)I[A'\Delta B']-yI[A'\Delta B']-\tfrac{5D-4}{2}I^2[A'\Delta B']+\tfrac{D^2}{2}I^2[A'B']
-\tfrac D4yI^2[A'^2]+(D-2)^2I^2[BA'']\nonumber\\
&-\tfrac34(D-2)I^3[A'^2]-2I[A'\Delta C']+yI[A'\Delta C']+\tfrac{3D-2}{2}I^2[A'\Delta C']\nonumber\\
&-\tfrac D4yI^2[A'\Delta C']-\tfrac34(D-2)I^3[A'\Delta C'],\nonumber\\
\sivbF_{\epsilon_2}={}&2yI[BA'']+I^2[A'B']-(D-1)I^2[A'\Delta B']+\tfrac14yI^2[A'^2]+\tfrac14yI^2[A'\Delta C'],\nonumber\\
\sivbF_{\epsilon_3}={}&2y^2BA''+2(D-1)yI[BA'']+(D-2)I^2[A'B']+yI[A'^2]+\tfrac{D-1}{4}yI^2[A'^2]+\tfrac14y^2I[A'^2]\nonumber\\
&+\tfrac{D-1}{4}yI^2[A'\Delta C']+\tfrac14y^2I[A'\Delta C'],\nonumber\\
\sivbF_\zeta={}&-4I^2[\Delta B\,A'']+I^3[A'\Delta B']+\tfrac{D-2}{8}I^4[A'^2]+I^2[\Delta C\,A'']-\tfrac12I^3[A'\Delta C']
+\tfrac{D-2}{8}I^4[A'\Delta C'].\nonumber
\end{align}

\paragraph*{Global pointwise check.}
The ten operators were applied to the totals \eqref{eq:s4b:totals} and compared pointwise with the direct sum of
$P_1,\dots,P_7$ (naive inner derivatives), (i) with generic powers $A=y^{-7/3}$, $B=y^{-5/4}$, $\Delta C=y^{2/7}$ at
$D=5$ and $D=4$ (two random spacelike points each), and (ii) with the actual $D=4$ propagators (three random points),
where the constants of the first integrals must be linked, $I[fA'']=fA'-I[f'A']$ (scripts \texttt{check\_global.py},
\texttt{check\_tables45.py}). In every case the difference vanishes to the 30-digit working precision, once the five
misprints C9--C13 of our Table~\ref{tab:s4b:apptypos} are corrected (C23 in that table is a wording slip that does not enter the
totals).

\begin{table}[tb]
\caption{Typographical errors in the appendix of \PI{} (numbering of our Table~\ref{tab:s0s1s7s8:errata}). None of them affects any
downstream result (Tables~IV--V are consistent with the corrected totals).}
\label{tab:s4b:apptypos}
\renewcommand{\arraystretch}{1.3}
\begin{ruledtabular}
\begin{tabular}{@{}l>{\raggedright\arraybackslash}p{0.14\textwidth}>{\raggedright\arraybackslash}p{0.22\textwidth}%
>{\raggedright\arraybackslash}p{0.40\textwidth}>{\raggedright\arraybackslash}p{0.13\textwidth}@{}}
& Location & Printed & Correct & Reason\\
\hline
C9 & Table XII header & $\beta=(aa')^{D-2}(a^2+a'^2)H^2\square$ & $(aa')^{D-1}(a^2+a'^2)H^2\square$, as in \eqI{105} & weight counting; global check\\
C10 & \eqI{A18} & $-2(D-1)aa'H^2\,I[f'A]$ & $-2(D-1)aa'H^2\,I[f'A']$ & trace of \eqI{78}\\
C11 & Table XIII, row $P_{7c}$ & repeats part of the total & $-\frac12(D-1)(D-2)\,yI[A'\Delta C']-\frac12(D-1)^2(D-2)\,I^2[A'\Delta C']$ & printed total correct\\
C12 & Table XV, row $P_{4b}$ & $(D-1)\,y'(A'B')'$ & $(D-1)\,y^2(A'B')'$ & from \eqI{A36}\\
C13 & Table XVIII, total & $\frac14\,yI^2[A'^2\Delta C']$ & $\frac14\,yI^2[A'\Delta C']$ & row $P_{7d}$\\
C23 & text before \eqI{A31} & ``\ldots\ to $P_1$ and $P_6$'' & ``\ldots\ to $P_1$ and $P_4$'' & subparts $P_{1a},\dots,P_{4a},\dots$\\
\end{tabular}
\end{ruledtabular}
\end{table}

\paragraph*{Scope of the audit.}
The appendix ``Extracting derivatives'' was \emph{audited}, not redone from scratch row by row. What is
established is: (i) the worked examples above (the $P_{3c}$ row of the $\gamma_3$ table and the $\beta$ and $\gamma_1$
totals) were checked by hand; (ii) the definitions \eqI{A1}--\eqI{A7}, the reduction \eqI{A8}--\eqI{A9} and Tables~XII ($\beta$) and XIV
($\gamma_2$) were transcribed directly from the printed appendix; (iii) the remaining tables (XI, XIII, XV--XX) were used
through their totals, which were coded line by line and verified by the global pointwise test. This test is
numerical, at a few random sample points with 30-digit arithmetic, and it checks only the sum over operators; it is not
a symbolic identity. An earlier working report stated that every row had been reproduced by hand; that was an
overstatement. Finally, Tables~IV and V below were computed from the appendix totals after these totals had been
checked, so for this step the printed appendix (not the target tables) is an input.

\subsection{Table IV in general \texorpdfstring{$D$}{D}}
\label{sec:s4b:tableIV}

The coefficients of $z^{D-1}$ and $z^{D-2}$ in each total follow from our Table~\ref{tab:s4b:products} by the rules
\begin{equation}
I[z^q]=-\frac{4\,z^{q-1}}{q-1},\qquad y\,z^q=4z^{q-1},\qquad \frac{d}{dy}z^q=-\frac q4z^{q+1}.
\label{eq:s4b:rules}
\end{equation}
The unit of Table~IV is $\Gamma^2(\frac D2)\kappa^2H^{2D-4}/(4\pi)^D$, so the factor $\frac1{16}$ of the products is
divided out at the end.

\paragraph*{Example: $\beta$.}
In $\sivbF_\beta$ the terms reaching $z^{D-1}$ or $z^{D-2}$ are $DI[A'B']\to-\frac{4D}{D-1}z^{D-1}-4Dz^{D-2}$,
$-I[A'\Delta B']\to-\frac{8}{D-2}z^{D-2}$, $-\frac12yI[A'^2]\to+\frac{8}{D-1}z^{D-2}$,
$-\frac{D-1}{2}I^2[A'^2]\to-\frac{8}{D-2}z^{D-2}$ and $I[A'\Delta C']\to+\frac{32}{D-2}z^{D-2}$. Dividing by 16,
\begin{equation}
z^{D-1}:\ -\frac{D}{4(D-1)},\qquad
z^{D-2}:\ \frac1{16}\Big[-4D+\frac{16}{D-2}+\frac{8}{D-1}\Big]=-\frac{D^3-3D^2-4D+8}{4(D-1)(D-2)},
\end{equation}
as printed. Since $\sivbF_{\gamma_1}=(D-1)(D-2)\sivbF_\beta$, the $\gamma_1$ row follows.

\paragraph*{Example: $\gamma_2$.}
In $\sivbF_{\gamma_2}$ the contributions are $DyA'B'\to4Dz^{D-1}+4D(D-2)z^{D-2}$, $-yA'\Delta B'\to+8z^{D-2}$,
$-\frac12y^2A'^2\to-8z^{D-2}$, $-\frac D2yI[A'^2]\to+\frac{8D}{D-1}z^{D-2}$ and $yA'\Delta C'\to-32z^{D-2}$; the last two
terms of $\sivbF_{\gamma_2}$ start at $z^{D-3}$. The first column is $4D/16=D/4$, as printed, and the second is
\begin{equation}
\frac1{16}\Big[4D(D-2)+\frac{8D}{D-1}-32\Big]=\frac{D(D-1)(D-2)+2D-8(D-1)}{4(D-1)}=+\frac{D^3-3D^2-4D+8}{4(D-1)} .
\label{eq:s4b:gamma2}
\end{equation}
\PI{} prints $-\frac{D^3-3D^2-4D+8}{4(D-1)}$. The other rows were obtained mechanically with the same rules
(\texttt{check\_table4.py}); the $\alpha$ row also follows from \eqref{eq:s4b:I125}. The result is
our Table~\ref{tab:s4b:IV}: 19 of the 20 entries agree with \PI{} in general $D$. The printed sign of the $\gamma_2$
entry is a misprint (C15): \PI's own Tables~VI--IX use the $+$ sign, and shift symmetry independently requires it (with
the printed sign a forbidden term $-\frac{16}{3}$ would remain in Table~IX at $D=4$); this is discussed in
Sec.~\ref{sec:s5:div}.

\begin{table}[tb]
\caption{Table IV of \PI{} (divergent normal contributions) derived in general $D$, in units of
$\Gamma^2(\frac D2)\,\kappa^2H^{2D-4}/(4\pi)^D$, paper convention (signs of $P_1,\dots,P_7$). The last column compares
with the printed table, column by column.}
\label{tab:s4b:IV}
\renewcommand{\arraystretch}{1.45}
\begin{ruledtabular}
\begin{tabular}{@{}lccl@{}}
Op. & Coef.\ of $(\frac4y)^{D-1}$ & Coef.\ of $(\frac4y)^{D-2}$ & Printed\\
\hline
$\alpha$ & $0$ & $\frac{D}{(D-1)(D-2)^2}$ & agrees, agrees\\
$\beta$ & $-\frac{D}{4(D-1)}$ & $-\frac{D^3-3D^2-4D+8}{4(D-1)(D-2)}$ & agrees, agrees\\
$\gamma_1$ & $-\frac{D(D-2)}{4}$ & $-\frac{D^3-3D^2-4D+8}{4}$ & agrees, agrees\\
$\gamma_2$ & $\frac D4$ & $+\frac{D^3-3D^2-4D+8}{4(D-1)}$ & agrees; prints $-\frac{D^3-3D^2-4D+8}{4(D-1)}$\\
$\gamma_3$ & $-\frac D4$ & $0$ & agrees, agrees\\
$\delta$ & $0$ & $0$ & agrees, agrees\\
$\epsilon_1$ & $0$ & $\frac{D^2-6D+4}{2(D-1)(D-2)}$ & agrees, agrees\\
$\epsilon_2$ & $0$ & $\frac{1-2D}{(D-1)(D-2)}$ & agrees, agrees\\
$\epsilon_3$ & $0$ & $0$ & agrees, agrees\\
$\zeta$ & $0$ & $0$ & agrees, agrees\\
\end{tabular}
\end{ruledtabular}
\end{table}

\subsection{Table V at \texorpdfstring{$D=4$}{D=4} and the missing \texorpdfstring{$72\gamma_2$}{72 gamma2}}
\label{sec:s4b:tableV}

The totals \eqref{eq:s4b:totals} were evaluated exactly at $D=4$ with $A=h\{4/y-2\ln(y/4)-1\}$ and $B=C=4h/y$ (linked
first integrals, higher integrals as \texttt{sympy} antiderivatives, $\ln y=\ln x+\ln4$). The resulting functions are listed in
our Table~\ref{tab:s4b:D4totals}. Their $1/x^3$ and $1/x^2$ terms are Table~IV at $D=4$ (with the derived $\gamma_2$ entry
$+\frac23$); the rest is to be compared with Table~V.

\begin{table}[tb]
\caption{Derived operator functions at $D=4$ (coefficients of $\kappa^2H^4/(4\pi)^4$, paper convention, $x=y/4$), with
the integration constants as produced by \texttt{check\_tables45b.py}. Polynomial terms (constants, $x$, $x^2$ and the
$\ln2$ terms) depend on the integration-constant convention.}
\label{tab:s4b:D4totals}
\renewcommand{\arraystretch}{1.35}
\begin{ruledtabular}
\begin{tabular}{@{}ll@{}}
Op. & Derived $D=4$ function\\
\hline
$\alpha$ & $\frac{1}{3x^2}+\frac{13}{x}+\frac{6\ln x}{x}-6\ln^2x-18\ln x+24\ln^22-36\ln2$\\
$\beta$ & $-\frac{1}{3x^3}-\frac{1}{3x^2}+\frac5x-18\ln x-6-36\ln2$\\
$\gamma_1$ & $-\frac{2}{x^3}-\frac{2}{x^2}+\frac{30}{x}-108\ln x-36-216\ln2$\\
$\gamma_2$ & $\frac{1}{x^3}+\frac{2}{3x^2}-\frac5x-18$\\
$\gamma_3$ & $-\frac{1}{x^3}-\frac5x-36$\\
$\delta$ & $-\frac{4\ln x}{x}-\frac{49}{6x}+4\ln^2x+10\ln x+12x\ln x-12x+24x\ln2-16\ln^22+20\ln2$\\
$\epsilon_1$ & $-\frac{1}{3x^2}-\frac{26}{3x}+60\ln x-120x\ln x+72x+20+120\ln2-240x\ln2$\\
$\epsilon_2$ & $-\frac{7}{6x^2}+\frac{13}{6x}-12\ln x+12x\ln x-2-24\ln2+24x\ln2$\\
$\epsilon_3$ & $-\frac{11}{6x}+36x\ln x+12x-8+72x\ln2$\\
$\zeta$ & $\frac{10}{3}\ln x-24x\ln x+24x^2\ln x-36x^2+24x+48x^2\ln2-48x\ln2+\frac{20}{3}\ln2$\\
\end{tabular}
\end{ruledtabular}
\end{table}

\paragraph*{Comparison operator by operator.}
Our Table~\ref{tab:s4b:Vdiff} lists the derived function minus [printed Table~IV at $D=4$ $+$ printed Table~V]. All
logarithmic and inverse-power terms ($1/x$, $\ln x/x$, $\ln x$, $\ln^2x$, $x\ln x$, $x^2\ln x$) agree. Apart from
the $\gamma_2$ entry of Table~IV, the differences are polynomials. To see their effect we need, at $D=4$,
\begin{equation}
\nabla^2x=\tfrac32aa'H^2,\qquad \nabla^4x=0,\qquad \nabla^4x^2=\tfrac{15}{2}(aa'H^2)^2,\qquad
\square x=H^2(2-4x)\ \Rightarrow\ \nabla^2\square x=-6aa'H^4,
\label{eq:s4b:polyops}
\end{equation}
the last from \eqref{eq:s4b:I93}. Then:
(i) constants are annihilated by $\square$, $\nabla^2$ and $\nabla^4$;
(ii) all $\ln2$ terms cancel globally, using $\gamma_3=2\gamma_1+\gamma_2$ and $\epsilon_3=2\epsilon_1+\epsilon_2$: in units
of $\ln2$, the $\gamma_1$ structure receives $-216$ (from $\gamma_1$), $-360$ (from $\epsilon_1$), $+360$ (from $\zeta$) and
$+216$ (from $\epsilon_3$), and the $\gamma_2$ structure receives $-144$ (from $\delta$), $+36$ (from $\epsilon_2$) and
$+108$ (from $\epsilon_3$);
(iii) what remains is the $-12x$ of the $\delta$ row,
\begin{equation}
\delta[-12x]=(aa')^2(a^2+a'^2)(-12)\,\nabla^2\square x=72\,(aa')^3(a^2+a'^2)H^4=72\,\gamma_2[1].
\label{eq:s4b:72g2}
\end{equation}
The choice of integration constants does not change the sum over operators (every consistent convention gives exact
identities), so this remainder is convention independent.

\begin{table}[tb]
\caption{Per-operator differences at $D=4$, derived minus [printed Table~IV($D=4$) $+$ printed Table~V], in units of
$\kappa^2H^4/(4\pi)^4$, paper convention (\texttt{check\_tables45.py}).}
\label{tab:s4b:Vdiff}
\renewcommand{\arraystretch}{1.35}
\begin{ruledtabular}
\begin{tabular}{@{}lll@{}}
Op. & Difference & Effect\\
\hline
$\alpha$ & $24\ln^22-36\ln2$ & constant, annihilated by $\square^2$\\
$\beta$ & $-36\ln2-6$ & constant, annihilated by $\square$\\
$\gamma_1$ & $-216\ln2$ & convention; cancels against $\epsilon_1$, $\zeta$, $\epsilon_3$\\
$\gamma_2$ & $\frac{4}{3x^2}$ & sign of the Table~IV entry, Eq.~\eqref{eq:s4b:gamma2}\\
$\gamma_3$ & $0$ & ---\\
$\delta$ & $-12x+24x\ln2-16\ln^22+20\ln2$ & $\ln2$ part: convention; $-12x$ gives $+72\gamma_2[1]$\\
$\epsilon_1$ & $-240x\ln2+20+120\ln2$ & convention\\
$\epsilon_2$ & $24x\ln2-24\ln2-2$ & convention\\
$\epsilon_3$ & $72x\ln2-8$ & convention\\
$\zeta$ & $48x^2\ln2-48x\ln2+24x+\frac{20}{3}\ln2$ & $\nabla^4x=0$; $x^2\ln2$: convention\\
\end{tabular}
\end{ruledtabular}
\end{table}

\paragraph*{Global test.}
The direct sum $\sum_kP_k$, with the actual $D=4$ propagators, was compared at three random spacelike points with
several versions of Tables~IV$+$V acted upon by the ten operators (our Table~\ref{tab:s4b:hyp}). Only the version with both
corrections agrees, to the 30-digit working precision.

\begin{table}[tb]
\caption{Residual $\sum_kP_k-\sum_{\rm op}{\rm op}\,[\text{Table IV}+\text{Table V}]$ at three random spacelike points
(\texttt{check\_tables45b.py}); ``0'' means zero to 30 digits.}
\label{tab:s4b:hyp}
\renewcommand{\arraystretch}{1.3}
\begin{ruledtabular}
\begin{tabular}{@{}ll@{}}
Version of Tables IV$+$V & Residual at the three points\\
\hline
as printed & $0.0356,\ 17.05,\ 0.607$\\
$\gamma_2$ sign of Table~IV corrected only & $0.0351,\ 16.97,\ 0.602$\\
$-12x$ added to the $\delta$ row only & $5\times10^{-4},\ 0.074,\ 0.0047$\\
$\gamma_2$ sign corrected $+$ $-12x$ added to the $\delta$ row & $0,\ 0,\ 0$\\
$\gamma_2$ sign corrected $+$ $\gamma_2$ constant $-18\to+54$ (equivalent) & $0,\ 0,\ 0$\\
\end{tabular}
\end{ruledtabular}
\end{table}

\paragraph*{Corrected Tables IV and V.}
Thus the printed Tables~IV$+$V reproduce the direct sum only with two corrections: the sign of the second column of
the $\gamma_2$ row of Table~IV, and a finite term equivalent to $72\,\gamma_2$ acting on a constant. The latter is most
naturally read as a term $-12x$ dropped from the $\delta$ row of Table~V (\PI{} may have assumed that
$\nabla^2\square x$ vanishes); equivalently the $\gamma_2$ row becomes $-\frac5x+54$. Where the term belongs is an
interpretation: the tests fix only the sum. The corrected Table~V is given in our Table~\ref{tab:s4b:V}; in the
definition sign all entries of Tables~IV and V are multiplied by $-1$.

\begin{table}[tb]
\caption{Table V of \PI{} (finite normal contributions), corrected, coefficients of $\kappa^2H^4/(4\pi)^4$, $x=y/4$,
paper convention. The only change with respect to the printed table is the term $-12x$ in the $\delta$ row
(equivalently, the $\gamma_2$ row $-\frac5x+54$ instead of the printed $-\frac5x-18$). Together with the corrected
$\gamma_2$ entry of Table~IV, $(\frac D4,\ +\frac{D^3-3D^2-4D+8}{4(D-1)})$, this reproduces the direct sum pointwise.}
\label{tab:s4b:V}
\renewcommand{\arraystretch}{1.35}
\begin{ruledtabular}
\begin{tabular}{@{}ll@{}}
Op. & Coefficient of $\kappa^2H^4/(4\pi)^4$\\
\hline
$\alpha$ & $\frac{6\ln x}{x}+\frac{13}{x}-6\ln^2x-18\ln x$\\
$\beta$ & $\frac5x-18\ln x$\\
$\gamma_1$ & $\frac{30}{x}-108\ln x-36$\\
$\gamma_2$ & $-\frac5x-18$\\
$\gamma_3$ & $-\frac5x-36$\\
$\delta$ & $-\frac{4\ln x}{x}-\frac{49}{6x}+4\ln^2x+10\ln x+12x\ln x\ \bm{-12x}$\\
$\epsilon_1$ & $-\frac{26}{3x}+60\ln x-120x\ln x+72x$\\
$\epsilon_2$ & $\frac{13}{6x}-12\ln x+12x\ln x$\\
$\epsilon_3$ & $-\frac{11}{6x}+36x\ln x+12x$\\
$\zeta$ & $\frac{10}{3}\ln x-24x\ln x+24x^2\ln x-36x^2$\\
\end{tabular}
\end{ruledtabular}
\end{table}

\paragraph*{Physical significance of $72\gamma_2$.}
The term $72\gamma_2[1]$ (S3) is independent of $y$, i.e.\ analytic. In the Schwinger--Keldysh formalism of
\PII~\cite{KW08} it is the same in
the $++$ and $+-$ kernels and therefore drops out of the retarded self-mass-squared identically; the same holds for all
polynomial (integration-constant) ambiguities of Table~V\@. It reappears in Sec.~\ref{sec:s5:TX} as the constant
of the row $(aa')^3(a^2+a'^2)H^4$ of Table~X, $-54\to+18$ (paper convention). There, the global pointwise test of
Table~X against the direct sum of the normal contributions succeeds only when both this correction and the restoration
of the missing Table~III row in the $(aa')^2\nabla^4$ row are applied; neither correction alone suffices. The
$72\gamma_2$ correction has no effect on any retarded quantity, and hence none on the late-time mode function.

\subsection{Distributional content of the external operators}
\label{sec:s4b:distr}

By Eq.~\eqref{eq:s4b:gamma}, the function multiplying $\alpha$, $I^2[(A+\Delta C)A'']$, contains the component
$c_{A+\Delta C}\,\sivbAl$ with
\begin{equation}
c_{A+\Delta C}=k\pihat+\frac{2(D-2)}{D-3}\big(c-k\pihat\big)=-\frac{D-1}{D-3}\,k\pihat+\frac{2(D-2)}{D-3}\,c ,\qquad
c=\frac{k}{(D-2)(D-3)} ,
\end{equation}
which is exactly the combination in the local contribution \eqI{54}$+$\eqI{58}. The generic reduction \eqI{A35} is
correct pointwise. However, if the $\square$ in $\alpha=(aa')^D\square^2$ and in $\delta$ is applied outside the
integral as a distribution, as \PI{} states, then $\square\sivbAl\supset i\delta^D/a^D$ would reproduce
\eqI{54}$+$\eqI{58} a second time; in the $D=4$ tables this content sits in the $1/x$ and $\ln x/x$ terms of $\alpha$,
$\beta$ and $\delta$. The local term \eqref{eq:s4b:loc} raises the complementary question.

This point is only partially resolved. In the counterterm analysis (Sec.~\ref{sec:s5:div}) we find that in \PI{}
divergent local terms arise only from Table~IV, while in the finite Table~V functions the $c_f\sivbAl$ component has
combined with the pole-carrying $z^{D-3}$ terms into finite logarithms. Any double counting then has exactly the form
$c(D)\,(a^D\square-a^{D-2}\nabla^2)\delta^D(x-x')$, with a constant $c(D)$, which may contain a pole, and no $\ln a$; this is the form
of the counterterms \eqI{131} and \eqI{145}. Both the pole and the finite part of $c(D)$ are absorbed into $\alpha_2$ and $\alpha_3$, and
the renormalized result changes only through the arbitrary $\Delta\alpha_i$. Statements about the poles and finite parts of $\beta_2$ and
$\beta_3$ (Sec.~\ref{sec:s5:beta}) therefore hold within \PI's treatment of the outer operators. (The term \eqref{eq:s4b:loc} is likewise of the form $a^{D-2}\nabla^2\delta^D$ of
\eqI{145}.) It has not been shown that the outer-operator reduction \eqI{A35}--\eqI{A38} produces no \emph{other} local
difference at coincidence, for example a finite $H^4a^D\delta^D$; that would require a full distributional analysis.
Partial support comes from the vanishing $a^D\square^2$ total of Table~VIII and the finiteness of Table~IX.

\subsection{Summary of this section}
\label{sec:s4b:summary}

Our Table~\ref{tab:s4b:summary} summarizes the comparison of Eqs.~\eqI{68}--\eqI{125} and Tables~III--V with \PI.

\begin{table}[tb]
\caption{Status of Eqs.~\eqI{68}--\eqI{125} and Tables~III--V of \PI. ``Inherited sign'': the printed quantity is $-1$
times the definition-sign result, the overall sign established for \eqI{47}--\eqI{49}.}
\label{tab:s4b:summary}
\renewcommand{\arraystretch}{1.25}
\begin{ruledtabular}
\begin{tabular}{@{}>{\raggedright\arraybackslash}p{0.22\textwidth}>{\raggedright\arraybackslash}p{0.47\textwidth}%
>{\raggedright\arraybackslash}p{0.24\textwidth}@{}}
Item & Result & Status\\
\hline
\eqI{68}--\eqI{80} & derived; \eqI{69} holds for naive derivatives; divergence count confirmed & agrees\\
\eqI{81}--\eqI{88}, \eqI{90}--\eqI{93} & derived in general $D$; $\square$ in \eqI{92} distributional, \eqI{93} pointwise & agrees\\
\eqI{89} & second term lacks $H^2$; overall sign & misprint (C8); inherited sign\\
Table III & all three rows reproduced & agrees (inherited sign)\\
\eqI{94}--\eqI{115}, \eqI{A8}--\eqI{A9} & derived; weight counting consistent & agrees\\
\eqI{116}--\eqI{125} & leading coefficients and $D=4$ forms; \eqI{125} up to a constant annihilated by $\alpha$ & agrees\\
Appendix, Tables XI--XX & totals verified by the global pointwise test (audit, not row-by-row re-derivation) & misprints C9--C13; wording slip C23\\
Table IV & 19/20 entries; $\gamma_2$, second column, sign reversed & misprint (C15)\\
Table V & all non-analytic terms agree; finite term $72\gamma_2$ missing globally & S3: analytic; no retarded effect\\
Outer $\square$ as a distribution & possible double counting of a local term of counterterm form & partially resolved\\
\end{tabular}
\end{ruledtabular}
\end{table}

Within Eqs.~\eqI{68}--\eqI{125} the misprints in \eqI{89} (C8), in Table~IV (C15) and in the appendix (C9--C13, C23) have no effect on the
printed results downstream, and the overall sign is the one inherited from \eqI{47}--\eqI{49}. Only the $72\gamma_2$
term (S3) changes a quantity used later (the constant of one row of Table~X); it is analytic and has no retarded effect.
The substantive consequence of this part, S2, the omission of the third row of Table~III from the $(aa')^2\nabla^4$ row
of Table~X, which is non-analytic and changes the nonlocal self-mass (but not the late-time mode function), arises
only at the assembly of Table~X.

\providecommand{\sfiveRes}{\operatorname{Res}}
\providecommand{\sfiveLie}{\pounds}
\providecommand{\sfiveLGF}{\mathcal{L}_{\rm GF}}
\providecommand{\sfiveLy}{\mathcal{L}_y}
\providecommand{\sfiveCd}{C_{\delta}}
\providecommand{\sfivegE}{\gamma_{\rm E}}
\providecommand{\sfivekb}{\bar k}

\section{Renormalization}
\label{sec:s5}

This section corresponds to Sec.~IV of \PI~\cite{KW07}, Eqs.~\eqI{126}--\eqI{167} and Tables~VI--X.
It has three parts.
First we determine the counterterms that the symmetries allow: two de Sitter-invariant ones and, because the gauge-fixing functional \eqI{15} breaks de Sitter invariance, exactly one noninvariant one (Secs.~\ref{sec:s5:inv}--\ref{sec:s5:noninv}).
Second, we separate the divergences of the normal 3-point contributions of Table~IV into local terms and finite nonlocal terms (Sec.~\ref{sec:s5:div}, Tables~VI--IX).
Third, we collect all local terms into the coefficients $\beta_1,\beta_2,\beta_3$ of the three counterterm structures, collect all finite nonlocal terms into Table~X, and renormalize (Secs.~\ref{sec:s5:beta}--\ref{sec:s5:final}).
Every step is written out in Chapters~7 and~8 of the derivation notes; the scripts are in \texttt{code/07\_desitter\_counterterms/} and \texttt{code/08\_renormalization\_tableX/} of Ref.~\cite{ref:repo}.

In brief, the counterterm analysis \eqI{126}--\eqI{145}, the separation of divergences \eqI{146}--\eqI{159} with Tables~VI--IX, the vanishing of $\beta_1$, and the pole of $\beta_3$ are confirmed.
The finite parts of $\beta_2$ and $\beta_3$ shift because of the two finite local errors in \eqI{40} and \eqI{59} found in Secs.~\ref{sec:s3} and~\ref{sec:s4a}; the arbitrary finite counterterms absorb these shifts.
Table~X needs two corrections: a non-analytic contribution is missing from the $(aa')^2\nabla^4$ row, and the constant of the $(aa')^3(a^2+a'^2)H^4$ row changes.
Only with both corrections does Table~X agree with a direct pointwise sum of the normal contributions.
In the definition sign the physical nonlocal part of the renormalized self-mass-squared is $-$(corrected Table~X).

\paragraph*{Conventions for this section.}
Tables~VI--X and the coefficients $\beta_i$ are first given in the \emph{paper convention}, the sign of the printed \eqI{47}--\eqI{49} and Tables~III--X, so that they can be compared entry by entry with \PI.
Values in the \emph{definition sign}, fixed by \eqI{168}, are $-1$ times these (Sec.~\ref{sec:s0s1s7s8:conventions}).
Counterterm vertices are computed with the rule ``vertex $=i\,\delta^2S_{\rm ct}/\delta\phi(x)\,\delta\phi(x')$''.
This rule reproduces the printed \eqI{130}, \eqI{131} and \eqI{145}, and it is the rule consistent with \eqI{168}.
The overall sign of the printed self-mass-squared therefore affects only the fitted values of the counterterm coefficients $\alpha_i$, not the structure of the counterterms.
All consistency checks on Tables~VI--IX below are linear and homogeneous, so they hold in either convention.
Following \PI, the letters $\alpha,\beta,\gamma_i,\delta,\epsilon_i,\zeta$ denote the external operators of the normal contributions, while $\alpha_i$ and $\beta_i$ denote counterterm coefficients.
We write Euler's constant as $\sfivegE$ (it is $\gamma$ in \PI) to avoid a clash with the operators $\gamma_i$.
A bare table number (``Table~VI'') refers to \PI; tables of the present paper are cited as ``our Table''.
The coefficients $\beta_i$ and $\alpha_i$ are in units of $H^{D-4}/(4\pi)^{D/2}$; the entries of Tables~VII and~X are coefficients of $\kappa^2H^4/(4\pi)^4$, functions of $x\equiv y/4$, with $y(x;x')$ the de Sitter length function \eqI{11}.
We also use
\begin{equation}
z\equiv\frac4y,\qquad
k\equiv\frac{H^{D-2}}{(4\pi)^{D/2}}\,\frac{\Gamma(D-1)}{\Gamma(\frac D2)}\equiv\frac{H^{D-2}}{(4\pi)^{D/2}}\,\sfivekb,\qquad
\pihat\equiv-\pi\cot\Big(\frac{\pi D}{2}\Big),\qquad
L\equiv\ln(aa') ,
\label{eq:s5:notation}
\end{equation}
so that $\sfivekb\to2$ at $D=4$.

\subsection{Invariant counterterms}
\label{sec:s5:inv}

This subsection corresponds to \eqI{126}--\eqI{131}.
A one-loop counterterm for the scalar self-mass-squared is proportional to $\kappa^2\phi\phi$.
Since $[\kappa^2]=-2$ and $[\phi]=1$ in $D=4$, four further units of mass are needed, in the form $m^4$, $m^2\partial^2$ or $\partial^4$ (a curvature counts as two derivatives).
This is the superficial degree of divergence four stated in \PI: the vertices $\kappa^nh^n\partial\phi\partial\phi$ carry two derivatives, like the $\kappa^nh^n\partial h\partial h$ vertices of pure gravity.
The candidate invariants are $\kappa^2m^4\phi^2\sqrt{-g}$ \eqI{126}; none with one or three derivatives; $\partial\phi\cdot\partial\phi$ and $\phi^2R$ with two derivatives \eqI{127} (after $\phi\square\phi\simeq-\partial\phi\cdot\partial\phi$, where $\simeq$ denotes equality up to a total derivative); and, with four derivatives,
\begin{equation}
(\square\phi)^2,\quad R\,\partial\phi\cdot\partial\phi,\quad R^{\mu\nu}\partial_\mu\phi\,\partial_\nu\phi,\quad \phi^2R^2,\quad \phi^2R_{\mu\nu}R^{\mu\nu},\quad \phi^2R_{\mu\nu\rho\sigma}R^{\mu\nu\rho\sigma},\quad \phi^2\square R .
\label{eq:s5:four}
\end{equation}
To reach this list we used $\phi\square^2\phi\simeq(\square\phi)^2$, $\phi\,\partial\phi\cdot\partial R\simeq-\frac12\phi^2\square R$, the contracted Bianchi identity, and the commutator identity
$\phi_{;\mu\nu}\phi^{;\mu\nu}=(\square\phi)^2-R_{\mu\nu}\phi^{;\mu}\phi^{;\nu}+\nabla_\mu\big(\phi^{;\nu}\phi_{;\nu}{}^{;\mu}-\phi^{;\mu}\square\phi\big)$.%
\footnote{The list \eqI{128} does not contain the last two terms of Eq.~\eqref{eq:s5:four} (erratum C14). Both carry an undifferentiated $\phi$ and are removed by the shift symmetry, so \eqI{129} is unaffected.}
Three arguments shorten the list.
(i)~The Lagrangian \eqI{1} is invariant under $\phi\to\phi+\text{const}$, and dimensional regularization preserves this symmetry; every term with an undifferentiated $\phi$ is excluded.
(ii)~The scalar is massless and mass is renormalized multiplicatively in dimensional regularization, so no term carries $m$.
(iii)~On the background, $\hat R_{\mu\nu}=(D-1)H^2\hat g_{\mu\nu}$, so $R\,\partial\phi\cdot\partial\phi$ and $R^{\mu\nu}\partial_\mu\phi\,\partial_\nu\phi$ become degenerate.
With $\sqrt{-\hat g}=a^D$ and $\sqrt{-\hat g}\,\hat g^{\mu\nu}=a^{D-2}\eta^{\mu\nu}$, two invariant counterterms remain (cf.\ \eqI{129}),
\begin{equation}
\tfrac12\alpha_1\kappa^2\,\square\phi\,\square\phi\,a^D\qquad\text{and}\qquad-\tfrac12\alpha_2\kappa^2H^2\,\partial_\mu\phi\,\partial_\nu\phi\,\eta^{\mu\nu}a^{D-2}.
\label{eq:s5:inv}
\end{equation}
Lemma~A of Sec.~\ref{sec:s2:rules}, Eq.~\eqref{eq:s2:lemmaA}, together with $a^D\square=\partial_\mu(a^{D-2}\eta^{\mu\nu}\partial_\nu)$, gives the vertices (cf.\ \eqI{130}--\eqI{131})
\begin{equation}
\tfrac12\alpha_1\kappa^2\square\phi\,\square\phi\,a^D\;\longrightarrow\;i\alpha_1\kappa^2a^D\square^2\delta^D(x-x'),\qquad
-\tfrac12\alpha_2\kappa^2H^2\partial\phi\cdot\partial\phi\,a^{D-2}\;\longrightarrow\;i\alpha_2\kappa^2H^2a^D\square\,\delta^D(x-x'),
\label{eq:s5:invvert}
\end{equation}
exactly as printed.
Had the gauge condition been de Sitter invariant, these two counterterms would absorb every divergence.
It is not, so we must determine which part of the de Sitter group the gauge respects.

\subsection{The de Sitter group and the residual symmetries of the gauge}
\label{sec:s5:dS}

This subsection corresponds to \eqI{132}--\eqI{135} and the text that follows them.
\PI\ states that the $\frac12D(D+1)$ de Sitter transformations split into four classes, and that the gauge \eqI{15} breaks only the fourth, with a reference to the literature~\cite{ref:Miao2006}.
We prove both statements directly.

\paragraph*{Killing vectors.}
With $\hat g_{\mu\nu}=a^2\eta_{\mu\nu}$, $a=-1/(H\eta)$, $\eta<0$ and $\xi_\mu\equiv\eta_{\mu\nu}\xi^\nu$ (index lowered with the flat metric), the Killing equation reads
\begin{equation}
\sfiveLie_\xi\hat g_{\mu\nu}=a^2\Big[\partial_\mu\xi_\nu+\partial_\nu\xi_\mu-\frac{2\xi^0}{\eta}\,\eta_{\mu\nu}\Big]=0 .
\label{eq:s5:killing}
\end{equation}
Its components are $\partial_0\xi^0=\xi^0/\eta$, $\partial_0\xi^i=\partial_i\xi^0$ and $\partial_i\xi^j+\partial_j\xi^i=2(\xi^0/\eta)\delta_{ij}$.
The first gives $\xi^0=\eta f(\vec x)$ and the second $\xi^i=\frac12\eta^2\partial_if+g^i(\vec x)$.
In the third, the coefficient of $\eta^2$ must vanish separately, so $\partial_i\partial_jf=0$ and $f=\lambda+2\vec b\cdot\vec x$.
What remains, $\partial_ig^j+\partial_jg^i=2f\delta_{ij}$, is the conformal Killing equation of $\mathbb R^{D-1}$.
Its general solution gives
\begin{equation}
\xi^0=\eta\,(\lambda+2\vec b\cdot\vec x),\qquad
\xi^i=c^i+\omega^i{}_jx^j+\lambda x^i+2(\vec b\cdot\vec x)\,x^i-(x\cdot x)\,b^i,\qquad x\cdot x\equiv\|\vec x\|^2-\eta^2 ,
\label{eq:s5:killsol}
\end{equation}
with $\omega_{ij}=-\omega_{ji}$.
This is the general solution, so there are no other Killing vectors.
The four parameter sets generate the four classes of \PI\ (our Table~\ref{tab:s5:classes}), and their number is
\begin{equation}
(D-1)+\tfrac12(D-1)(D-2)+1+(D-1)=\tfrac12D(D+1)=\dim SO(D,1).
\label{eq:s5:count}
\end{equation}
As a cross-check, Eq.~\eqref{eq:s5:killing} says that $\xi$ is a conformal Killing vector of Minkowski space with conformal factor $\xi^0/\eta$.
Of the $\frac12(D+1)(D+2)$ flat conformal generators this condition removes time translation, the $D-1$ boosts and the temporal special conformal transformation, leaving $\frac12D(D+1)$.
The flows of $c$, $\omega$ and $\lambda$ give \eqI{132}--\eqI{134}, with the dilatation factor $s=e^\lambda$ (called $k$ in \eqI{134}; it is unrelated to the propagator constant $k$ of Eq.~\eqref{eq:s5:notation}). The factor must be positive, since $s<0$ would map the patch $\eta<0$ to $\eta>0$.
A special conformal transformation is an inversion, followed by a translation by $\theta^\mu=(0,\vec\theta)$, followed by an inversion:
\begin{equation}
x'^\mu=\frac{x^\mu-\theta^\mu\,x\cdot x}{\Omega},\qquad \Omega\equiv1-2\,\vec\theta\cdot\vec x+\|\vec\theta\|^2\,x\cdot x,\qquad
dx'\cdot dx'=\frac{dx\cdot dx}{\Omega^2},\qquad \eta'=\frac{\eta}{\Omega} .
\label{eq:s5:sct}
\end{equation}
Hence $ds^2=dx\cdot dx/(H^2\eta^2)$ is invariant.
This is \eqI{135}, in which $x\cdot x$ must be read as the Minkowski product $\|\vec x\|^2-\eta^2$.
The generators close into the conformal algebra of $\mathbb R^{D-1}$, isomorphic to $\mathfrak{so}(D,1)$, with
$[\mathcal D,P_i]=-P_i$, $[\mathcal D,K_i]=K_i$ and $[K_i,P_j]=-2(\delta_{ij}\mathcal D+M_{ij})$.

\begin{table}[tb]
\caption{The four classes of de Sitter transformations in conformal coordinates, from the general Killing vector \eqref{eq:s5:killsol}. The last column states whether the gauge-fixing term \eqI{15} is invariant (Eq.~\eqref{eq:s5:xiu}).}
\label{tab:s5:classes}
\renewcommand{\arraystretch}{1.2}
\begin{ruledtabular}
\begin{tabular}{llcll}
Parameter & Generator & Number & \PI & $\sfiveLGF$\\
\hline
$c^i$ & $P_i=\partial_i$ & $D-1$ & (i) spatial translations, \eqI{132} & invariant\\
$\omega_{ij}$ & $M_{ij}=x_i\partial_j-x_j\partial_i$ & $\tfrac12(D-1)(D-2)$ & (ii) rotations, \eqI{133} & invariant\\
$\lambda$ & $\mathcal D=\eta\partial_0+x^i\partial_i$ & $1$ & (iii) dilatation, \eqI{134} & invariant\\
$b^i$ & $K_i=2x_i(\eta\partial_0+x^j\partial_j)-(x\cdot x)\partial_i$ & $D-1$ & (iv) spatial special conformal, \eqI{135} & broken\\
\end{tabular}
\end{ruledtabular}
\end{table}

\paragraph*{The gauge condition.}
The gauge-fixing term \eqI{15} is
$\sfiveLGF=-\frac12a^{D-2}\eta^{\mu\nu}F_\mu F_\nu=-\frac12\sqrt{-\hat g}\,\hat g^{\mu\nu}F_\mu F_\nu$, with
$F_\mu=\eta^{\rho\sigma}\big(h_{\mu\rho,\sigma}-\frac12h_{\rho\sigma,\mu}+(D-2)Ha\,h_{\mu\rho}\delta^0_\sigma\big)$.
Let $\psi_{\mu\nu}\equiv a^2h_{\mu\nu}$, which transforms as a tensor under background isometries, and let $u^\mu\equiv\delta^\mu_0/a$ be the unit normal of the surfaces $\eta=\text{const}$ (in this subsection only; it differs from the vector $u_\mu=\delta^0_\mu$, $u^\mu=-\delta^\mu_0$ of Secs.~\ref{sec:s2}--\ref{sec:s4b}).
Computing the background covariant derivatives with $\hat\Gamma^\lambda_{\mu\nu}=Ha(\delta^\lambda_\mu\delta^0_\nu+\delta^\lambda_\nu\delta^0_\mu+\eta_{\mu\nu}\delta^\lambda_0)$, we find
\begin{equation}
F_\mu=G_\mu+E_\mu,\qquad
G_\mu=\hat g^{\rho\sigma}\big(\hat\nabla_\sigma\psi_{\mu\rho}-\tfrac12\hat\nabla_\mu\psi_{\rho\sigma}\big),\qquad
E_\mu=H\big(2\psi_{\mu\nu}u^\nu-u_\mu\psi\big),
\label{eq:s5:FGE}
\end{equation}
with $\psi\equiv\hat g^{\rho\sigma}\psi_{\rho\sigma}$.
The vector $G_\mu$ is the de Donder vector of the background and is covariant under every isometry; the whole de Sitter-breaking part $E_\mu$ enters through $u^\mu$, that is, through the equal-time slicing.
For a Killing vector $\xi$, $\delta F_\mu=\sfiveLie_\xi F_\mu+\Delta_\mu$, and since $\nabla\cdot\xi=0$,
\begin{equation}
\delta\sfiveLGF=\partial_\lambda\big(\xi^\lambda\sfiveLGF\big)-\sqrt{-\hat g}\,\hat g^{\mu\nu}F_\mu\Delta_\nu,\qquad
\Delta_\mu=-H\big(2\psi_{\mu\nu}w^\nu-w_\mu\psi\big),\qquad w\equiv\sfiveLie_\xi u=[\xi,u].
\label{eq:s5:deltaL}
\end{equation}
An isometry is therefore a symmetry of the gauge-fixing term if and only if it preserves $u$.
With $u^\mu=-H\eta\,\delta^\mu_0$ and Eq.~\eqref{eq:s5:killsol},
\begin{equation}
[\xi,u]^\mu=-H\big(\xi^0\delta^\mu_0-\eta\,\partial_0\xi^\mu\big)=\big(0,\;2H\eta^2\,b^i\big),
\label{eq:s5:xiu}
\end{equation}
because the only $\eta$-dependence of $\xi^i$ is in $-(x\cdot x)b^i$.
Classes (i)--(iii) are symmetries of $\sfiveLGF$, and class (iv) is broken.
The explicit breaking is
\begin{equation}
\Delta_\mu=-\big(4\,h_{\mu i}b^i-2\,\delta^i_\mu b_i\,h^\rho{}_\rho\big),
\label{eq:s5:Delta}
\end{equation}
with the graviton trace $h^\rho{}_\rho=\eta^{\rho\sigma}h_{\rho\sigma}$, independent of $H$.
The breaking is real: on the gauge surface $F_\mu=0$ one has $\delta F_\mu=\Delta_\mu\neq0$, and the change of the action is not a boundary term, because its Euler--Lagrange derivatives with respect to $h_{00}$ and $h_{12}$ are nonzero (checked symbolically in $D=3$).
Geometrically, (i)--(iii) map the slices $\eta=\text{const}$ into themselves, whereas (iv) bends them, since $\eta'=\eta/\Omega$ depends on $\vec x$.%
\footnote{The term $k\ln(aa')$ of the scalar propagator shifts by a constant under a dilatation. This comes from the infrared state of the massless minimally coupled scalar, not from the gauge functional; a constant drops out under any derivative and does not affect the ultraviolet divergences.}

\subsection{Noninvariant counterterms}
\label{sec:s5:noninv}

This subsection corresponds to \eqI{136}--\eqI{145}.
We assume, as \PI\ does, that dimensional regularization preserves the symmetries of the gauge-fixed action, so that the divergences and the counterterms respect (i)--(iii).
Translations forbid explicit $x^i$ dependence and rotations require contracted spatial indices.
Under a dilatation $d^Dx\to s^Dd^Dx$, $a\to a/s$ and $\partial_\mu\to s^{-1}\partial_\mu$, so the invariant building blocks are $a^{-1}\partial_\mu$ and $H$, with an overall measure factor $a^D$.
With the shift symmetry and the dimension counting, every allowed counterterm density has the form
\begin{equation}
\mathcal L_{\rm ct}=\kappa^2a^D\,Q\big(a^{-1}\partial_0,\;a^{-1}\partial_i,\;H\big)[\varphi,\varphi],
\label{eq:s5:ctform}
\end{equation}
where $Q$ is a rotation-invariant polynomial of total weight four, bilinear in the two fields, with each field differentiated at least once.
(Following \PI, the scalar is written $\varphi$ in this subsection.)
We reduce all such densities to the basis
\begin{equation}
\begin{aligned}
B_1&\equiv a^D(\square\varphi)^2, &
B_2&\equiv H^2a^{D-2}\eta^{\mu\nu}\partial_\mu\varphi\,\partial_\nu\varphi, &
B_3&\equiv a^{D-2}\,\square\varphi\,\nabla^2\varphi,\\
B_4&\equiv a^{D-4}(\nabla^2\varphi)^2, &
B_5&\equiv H^2a^{D-2}\,\nabla\varphi\cdot\nabla\varphi, &&
\end{aligned}
\label{eq:s5:basis}
\end{equation}
of which $B_1$ and $B_2$ are the invariant counterterms \eqref{eq:s5:inv}.
The reduction tools are the identities (cf.\ \eqI{136}--\eqI{143}; $\partial_0a=Ha^2$)
\begin{align}
\frac{1}{a^2}\partial_0^2&=-\square-(D-2)\frac{H}{a}\partial_0+\frac{\nabla^2}{a^2},\qquad
\frac{(\partial_0\varphi)^2}{a^2}=-\partial_\mu\varphi\,\partial_\nu\varphi\,g^{\mu\nu}+\frac{1}{a^2}\nabla\varphi\cdot\nabla\varphi,
\label{eq:s5:136}\\
Ha^{D-1}\partial_0\varphi\,\square\varphi&\simeq-Ha^{D-3}\partial_\mu\partial_0\varphi\,\partial_\nu\varphi\,\eta^{\mu\nu}-H^2a^{D-2}(\partial_0\varphi)^2
\simeq\tfrac12(D-1)B_2-B_5,
\label{eq:s5:138}\\
Ha^{D-3}\partial_0\varphi\,\nabla^2\varphi&\simeq-\tfrac12Ha^{D-3}\partial_0(\nabla\varphi\cdot\nabla\varphi)\simeq\tfrac12(D-3)B_5 .
\label{eq:s5:141}
\end{align}
In each step marked $\simeq$ the discarded total derivative was written out explicitly; for example, the first step of Eq.~\eqref{eq:s5:138} discards $\partial_\mu\big(Ha^{D-3}\partial_0\varphi\,\partial^\mu\varphi\big)$.
The second identity of Eq.~\eqref{eq:s5:136} gives $H^2a^{D-2}(\partial_0\varphi)^2=-B_2+B_5$.
Our Table~\ref{tab:s5:reductions} lists the resulting general-$D$ reductions of eleven candidate densities.

\begin{table}[tb]
\caption{Reduction of candidate counterterm densities (each multiplied by $\kappa^2$) to the basis \eqref{eq:s5:basis}, in general $D$, up to total derivatives. All fifteen candidates of the form \eqref{eq:s5:ctform}, including these eleven, were also reduced by computing their Euler--Lagrange operators at $D=4,5,6$; all agree.}
\label{tab:s5:reductions}
\renewcommand{\arraystretch}{1.25}
\begin{ruledtabular}
\begin{tabular}{ll}
Candidate & Reduction ($\simeq$)\\
\hline
$a^D(a^{-2}\partial_0^2\varphi)^2$ & $B_1-2B_3+B_4+(D-2)(B_2-B_5)$\\
$a^D(a^{-2}\partial_0^2\varphi)(a^{-2}\nabla^2\varphi)$ & $-B_3+B_4-\tfrac12(D-2)(D-3)B_5$\\
$a^{D-4}(\nabla^2\varphi)^2$, $a^{D-4}(\partial_i\partial_j\varphi)^2$ & $B_4$\\
$a^{D-4}(\partial_0\partial_i\varphi)^2$ & $-B_3+B_4-(D-3)B_5$\\
$a^{D-4}\partial_0\varphi\,\partial_0^3\varphi$ & $-B_1+2B_3-B_4-\big[(D-2)+\tfrac12(D-3)(D-4)\big](B_2-B_5)$\\
$Ha^{D-3}\partial_0\varphi\,\partial_0^2\varphi$ & $\tfrac12(D-3)(B_2-B_5)$\\
$Ha^{D-3}\partial_0\varphi\,\nabla^2\varphi$ & $\tfrac12(D-3)B_5$\\
$Ha^{D-3}\partial_i\varphi\,\partial_0\partial_i\varphi$ & $-\tfrac12(D-3)B_5$\\
$Ha^{D-1}\partial_0\varphi\,\square\varphi$ & $\tfrac12(D-1)B_2-B_5$\\
$H^2a^{D-2}(\partial_0\varphi)^2$ & $-B_2+B_5$\\
\end{tabular}
\end{ruledtabular}
\end{table}

Every allowed density lies in $\mathrm{span}\{B_1,\dots,B_5\}$, and the five are independent, because their Euler--Lagrange operators are linearly independent.
Removing the invariant $B_1$ and $B_2$ leaves exactly three noninvariant counterterms (cf.\ \eqI{144}),
\begin{equation}
\tfrac12\kappa^2a^{D-2}\,\square\varphi\,\nabla^2\varphi,\qquad
\tfrac12\kappa^2a^{D-4}\,\nabla^2\varphi\,\nabla^2\varphi,\qquad
-\tfrac12\kappa^2H^2a^{D-2}\,\nabla\varphi\cdot\nabla\varphi ,
\label{eq:s5:144}
\end{equation}
as printed.
General-$D$ formulas were derived for the eleven densities of our Table~\ref{tab:s5:reductions}; the identities \eqref{eq:s5:136}--\eqref{eq:s5:141} were verified with symbolic $D$, while the completeness of the basis (fifteen candidates) and the independence of $B_1,\dots,B_5$ were verified by code at $D=4,5,6$ only.

\paragraph*{Flat-space limit.}
For $H\to0$ at fixed comoving time, $a\to1$ and $E_\mu\to0$ in Eq.~\eqref{eq:s5:FGE}: $F_\mu$ becomes the flat de Donder vector and the gauge-fixed theory becomes Poincar\'e invariant.
Assuming, as \PI\ does, that the ultraviolet divergences depend polynomially on $H$, the coefficients of noninvariant counterterms must vanish in this limit.
The first two terms of Eq.~\eqref{eq:s5:144} have dimensionless coefficients and tend to non-Poincar\'e-invariant densities, so they are excluded; $B_5\propto H^2$ survives.
The single noninvariant counterterm and its vertex are therefore (cf.\ \eqI{145})
\begin{equation}
-\tfrac12\alpha_3\kappa^2H^2a^D\,\frac{\nabla\varphi}{a}\cdot\frac{\nabla\varphi}{a}\;\longrightarrow\;i\alpha_3\kappa^2H^2a^{D-2}\,\nabla^2\delta^D(x-x'),
\label{eq:s5:145}
\end{equation}
exactly as printed.
This also confirms the two statements after \eqI{145}.
No local term $Ha^{D-3}\nabla^2\partial_0\delta^D$ can arise: the only candidate, $\int c(\eta)\,\partial_0\varphi\,\nabla^2\varphi$, has the Euler--Lagrange operator $-c'(\eta)\nabla^2\varphi$, in which the time derivative has dropped out.
No divergence $H^3a^{D-1}\partial_0\delta^D$ can arise either, because it would require an undifferentiated field, which the shift symmetry forbids.
The checks of this subsection and the previous one are in \texttt{check\_126\_135.py} and \texttt{check\_136\_145.py}~\cite{ref:repo}.

\subsection{Divergences of the normal contributions}
\label{sec:s5:div}

This subsection corresponds to \eqI{146}--\eqI{159} and Tables~VI--IX.
The 4-point contribution and the local 3-point contribution are already local.
The normal 3-point contributions of Table~IV contain the powers
\begin{equation}
z^{D-1}=\Big(\frac4y\Big)^{D-1}\qquad\text{and}\qquad z^{D-2}=\Big(\frac4y\Big)^{D-2},
\label{eq:s5:146}
\end{equation}
which are not integrable at $D=4$ (cf.\ \eqI{146}): $y^s$ is locally integrable only for $s>-\frac D2$, so the first diverges quadratically and the second logarithmically.
We write $c_1$ and $c_2$ for the coefficients of $z^{D-1}$ and $z^{D-2}$ in Table~IV, in units of $\kappa^2H^{2D-4}\Gamma^2(\frac D2)/(4\pi)^D$, for each external operator (\eqI{104}--\eqI{113}; general-$D$ forms in our Table~\ref{tab:s5:T6}).
For the $\gamma_2$ operator we use $c_2=+\frac{q}{4(D-1)}$, $q\equiv D^3-3D^2-4D+8$, as derived in Sec.~\ref{sec:s4b:tableIV}; Table~IV prints $-\frac{q}{4(D-1)}$, a misprint (C15) that the paper did not carry into its later tables (see below).
Following the partial-integration technique used in \PI~\cite{ref:Onemli2002}, we extract d'Alembertians until the powers become integrable.

\paragraph*{The key identity.}
Let $\sfiveLy[f]\equiv(4y-y^2)f''+D(2-y)f'$, so that the pointwise part of \eqI{93} is $\square f(y)=H^2\sfiveLy[f]$.
In distributional derivatives a delta function arises only from the component of $f$ proportional to $y^{1-D/2}$; all other powers give zero by analytic continuation.
The coefficient follows from the propagator equation $a^D\square\,i\Delta_A=i\delta^D(x-x')$ and the leading term $\frac{H^{D-2}}{(4\pi)^{D/2}}\Gamma(\frac D2-1)z^{D/2-1}$ of $i\Delta_A$.
The result is (cf.\ \eqI{147})
\begin{equation}
\square f(y)=H^2\sfiveLy[f]+\sfiveRes\Big[y^{\frac D2-2}f\Big]\,\frac{4\pi^{D/2}H^{2-D}}{\Gamma(\frac D2-1)}\,\frac{i\,\delta^D(x-x')}{\sqrt{-\hat g}},
\label{eq:s5:147}
\end{equation}
where $\sfiveRes$ denotes the coefficient of $1/y$, and $\sqrt{-\hat g}=a^D$ is the background determinant (written $\sqrt{-g}$ in \PI).
For powers, $\sfiveLy[z^q]=q\big[\{4(q+1)-2D\}/y+(D-q-1)\big]z^q$, with $z^q/y=\frac14z^{q+1}$.
For $q=D-2$ and $q=D-3$ the residue vanishes for generic $D$; for $q=\frac D2-1$ the first bracket vanishes and the residue is $4^{D/2-1}$.
This gives (cf.\ \eqI{148}--\eqI{150})
\begin{align}
z^{D-1}&=\frac{2}{(D-2)^2}\frac{\square}{H^2}z^{D-2}-\frac{2}{D-2}\,z^{D-2},\qquad
z^{D-2}=\frac{2}{(D-3)(D-4)}\frac{\square}{H^2}z^{D-3}-\frac{4}{D-4}\,z^{D-3},
\label{eq:s5:148}\\
0&=\frac{\square}{H^2}z^{\frac D2-1}-\frac D2\Big(\frac D2-1\Big)z^{\frac D2-1}-\sfiveCd\,\frac{i\,\delta^D(x-x')}{a^D},\qquad
\sfiveCd\equiv\frac{(4\pi)^{D/2}H^{-D}}{\Gamma(\frac D2-1)} .
\label{eq:s5:150}
\end{align}

\paragraph*{Segregating the pole.}
Subtracting $\frac{2}{(D-3)(D-4)}$ times Eq.~\eqref{eq:s5:150} from the second identity of Eq.~\eqref{eq:s5:148} gives (cf.\ \eqI{151})
\begin{equation}
z^{D-2}=\frac{2}{(D-3)(D-4)}\Big\{\sfiveCd\,\frac{i\,\delta^D}{a^D}+\frac{\square}{H^2}\big[z^{D-3}-z^{\frac D2-1}\big]\Big\}
-\frac{4}{D-4}\Big\{z^{D-3}-\frac{D(D-2)}{8(D-3)}\,z^{\frac D2-1}\Big\}.
\label{eq:s5:151}
\end{equation}
At $D=4$, $z^{D-3}$ and $z^{D/2-1}$ are the same function, so both braces are $O(D-4)$ and the pole sits entirely on the delta function.
Expanding in $D-4$ gives (cf.\ \eqI{152}, \eqI{153})
\begin{align}
z^{D-2}={}&\frac{iH^{-D}(4\pi)^{D/2}}{(D-3)(D-4)\Gamma(\frac D2)}\,(D-2)\,\frac{\delta^D(x-x')}{a^D}
-\frac{\square}{H^2}\Big\{\frac4y\ln\frac y4\Big\}+2\,\frac4y\ln\frac y4-\frac4y+O(D-4),
\label{eq:s5:152}\\
z^{D-1}={}&\frac{iH^{-D}(4\pi)^{D/2}}{(D-3)(D-4)\Gamma(\frac D2)}\Big\{\frac{2}{D-2}\frac{\square}{H^2}-2\Big\}\frac{\delta^D(x-x')}{a^D}
-\frac12\frac{\square^2}{H^4}\Big\{\frac4y\ln\frac y4\Big\}\notag\\
&+\frac{\square}{H^2}\Big\{2\,\frac4y\ln\frac y4-\frac12\,\frac4y\Big\}-2\,\frac4y\ln\frac y4+\frac4y+O(D-4),
\label{eq:s5:153}
\end{align}
in agreement with \PI.
These are equalities of distributions: $[z^{D-3}-z^{D/2-1}]/(D-4)$ converges to $\frac12z\ln z$ in $L^1_{\rm loc}$, and the outer d'Alembertians act distributionally.
The flat Euclidean analogue of Eqs.~\eqref{eq:s5:148}--\eqref{eq:s5:151}, $r^{-2(D-2)}=[\nabla^2(r^{-2(D-3)}-r^{2-D})-(D-2)\Omega_D\delta^D]/[2(D-3)(D-4)]$ with $\Omega_D=2\pi^{D/2}/\Gamma(\frac D2)$, was checked exactly in $D$ against the test function $e^{-r^2}$.

\paragraph*{Table VI: local divergent terms.}
Keeping the local parts of Eqs.~\eqref{eq:s5:152}--\eqref{eq:s5:153}, each operator contributes
\begin{equation}
c_1\Big\{\frac{2}{D-2}\frac{\square}{H^2}-2\Big\}+c_2\,(D-2)
\qquad\text{times}\qquad
\frac{i\kappa^2H^{D-4}}{(4\pi)^{D/2}}\,\frac{\Gamma(\frac D2)}{(D-3)(D-4)}\,\frac{\delta^D(x-x')}{a^D},
\label{eq:s5:T6rule}
\end{equation}
after the redundant operators have been eliminated with $\gamma_3=2\gamma_1+\gamma_2$ and $\epsilon_3=2\epsilon_1+\epsilon_2$, which follow from $(a+a')^2=a^2+a'^2+2aa'$.
Our Table~\ref{tab:s5:T6} gives the result; it agrees with the printed Table~VI row by row.
The printed $\gamma_2$ row, $+\frac{(D-2)q}{4(D-1)}$, follows from $c_2=+\frac{q}{4(D-1)}$; the sign printed in Table~IV would give the opposite sign.
So \PI\ used the correct sign in its calculation, and the sign in Table~IV is only a misprint.

\begin{table}[tb]
\caption{Table~VI of \PI, derived from Table~IV (paper convention), with $q\equiv D^3-3D^2-4D+8$; $(c_1,c_2)$ are the Table~IV coefficients after eliminating $\gamma_3$ and $\epsilon_3$. The operators $\delta$ and $\zeta$ have $c_1=c_2=0$. Units as in Eq.~\eqref{eq:s5:T6rule}. Every row agrees with the printed Table~VI.}
\label{tab:s5:T6}
\renewcommand{\arraystretch}{1.5}
\begin{ruledtabular}
\begin{tabular}{lll}
Operator & $(c_1,c_2)$ & Table~VI coefficient\\
\hline
$\alpha=(aa')^D\square^2$ & $\big(0,\frac{D}{(D-1)(D-2)^2}\big)$ & $\frac{D}{(D-1)(D-2)}$\\
$\beta=(aa')^{D-1}(a^2+a'^2)H^2\square$ & $\big(-\frac{D}{4(D-1)},-\frac{q}{4(D-1)(D-2)}\big)$ & $-\frac{D}{2(D-1)(D-2)}\frac{\square}{H^2}-\frac{(D+2)(D-4)}{4}$\\
$\gamma_1=(aa')^DH^4$ & $\big(-\frac{D^2}{4},-\frac q4\big)$ & $-\frac{D^2}{2(D-2)}\frac{\square}{H^2}-\frac{D^4-5D^3+16D-16}{4}$\\
$\gamma_2=(aa')^{D-1}(a^2+a'^2)H^4$ & $\big(0,+\frac{q}{4(D-1)}\big)$ & $+\frac{(D-2)q}{4(D-1)}$\\
$\epsilon_1=(aa')^{D-1}H^2\nabla^2$ & $\big(0,\frac{D^2-6D+4}{2(D-1)(D-2)}\big)$ & $\frac{D^2-6D+4}{2(D-1)}$\\
$\epsilon_2=(aa')^{D-2}(a^2+a'^2)H^2\nabla^2$ & $\big(0,\frac{1-2D}{(D-1)(D-2)}\big)$ & $\frac{1-2D}{D-1}$\\
\end{tabular}
\end{ruledtabular}
\end{table}

\paragraph*{Table VII: finite nonlocal terms.}
At $D=4$, with $x=y/4$, the nonlocal parts of Eqs.~\eqref{eq:s5:152} and~\eqref{eq:s5:153} are
\begin{equation}
N_{2}=-\frac{\square}{H^2}\Big[\frac{\ln x}{x}\Big]+\frac{2\ln x}{x}-\frac1x,\qquad
N_{1}=-\frac12\frac{\square^2}{H^4}\Big[\frac{\ln x}{x}\Big]+\frac{\square}{H^2}\Big[\frac{2\ln x}{x}-\frac1{2x}\Big]-\frac{2\ln x}{x}+\frac1x ,
\label{eq:s5:N}
\end{equation}
and each operator contributes $c_1N_1+c_2N_2$ with its $D=4$ values of $(c_1,c_2)$.
The results (our Table~\ref{tab:s5:T7}) agree with the printed Table~VII.
The header of Table~VII reads ``Coefficient of $i\kappa^2H^4/(4\pi)^4$''; the factor $i$ is a misprint (C16), since the units of Table~IV carry no $i$, the nonlocal parts of \eqI{152}--\eqI{153} are real, and the header of Table~X, which contains Table~VII, has no $i$.

\begin{table}[tb]
\caption{Table~VII of \PI, derived (paper convention); coefficients of $\kappa^2H^4/(4\pi)^4$, $x=y/4$. All rows agree with the printed table. In the $\beta$ row the terms without $\square$ cancel.}
\label{tab:s5:T7}
\renewcommand{\arraystretch}{1.5}
\begin{ruledtabular}
\begin{tabular}{lcl}
Operator & $(c_1,c_2)_{D=4}$ & Finite nonlocal contribution\\
\hline
$\alpha$ & $(0,\tfrac13)$ & $\frac{\square}{H^2}\big[-\frac{\ln x}{3x}\big]+\frac{2\ln x}{3x}-\frac1{3x}$\\
$\beta$ & $(-\tfrac13,-\tfrac13)$ & $\frac{\square^2}{H^4}\big[\frac{\ln x}{6x}\big]+\frac{\square}{H^2}\big[-\frac{\ln x}{3x}+\frac1{6x}\big]$\\
$\gamma_1$ & $(-4,-2)$ & $\frac{\square^2}{H^4}\big[\frac{2\ln x}{x}\big]+\frac{\square}{H^2}\big[-\frac{6\ln x}{x}+\frac2x\big]+\frac{4\ln x}{x}-\frac2x$\\
$\gamma_2$ & $(0,\tfrac23)$ & $\frac{\square}{H^2}\big[-\frac{2\ln x}{3x}\big]+\frac{4\ln x}{3x}-\frac2{3x}$\\
$\epsilon_1$ & $(0,-\tfrac13)$ & $\frac{\square}{H^2}\big[\frac{\ln x}{3x}\big]-\frac{2\ln x}{3x}+\frac1{3x}$\\
$\epsilon_2$ & $(0,-\tfrac76)$ & $\frac{\square}{H^2}\big[\frac{7\ln x}{6x}\big]-\frac{7\ln x}{3x}+\frac7{6x}$\\
\end{tabular}
\end{ruledtabular}
\end{table}

\paragraph*{Converting $a'$ into $a$.}
To compare the local terms with the counterterms, the factors $a'$ of the external operators are converted into $a$ with the delta function, and all scale factors are moved to the left.
Acting on a test function,
$\int d^Dx'\,a^p(x)\,a'^q\,\square^n\big[a^{-D}\delta^D(x-x')\big]\phi(x')=a^p\square^n\big[a^{q-D}\phi\big](x)$,
so $a'^q$ becomes $a^q$ \emph{inside} the d'Alembertians.
For $\alpha$, $\gamma_1$, $\gamma_2$, $\epsilon_1$, $\epsilon_2$ and $\gamma_1\square/H^2$ this is immediate (for example $\gamma_2\to2a^DH^4$ and $\epsilon_2\to2a^{D-2}H^2\nabla^2$).
For the two operators that contain $(a^2+a'^2)$ together with $\square$ we find, using $\square a=-DH^2a$ and $\square(1/a)=(D-2)H^2/a$ (cf.\ \eqI{154}--\eqI{155}),
\begin{align}
(aa')^{D-1}(a^2+a'^2)\square^2\Big[\frac{\delta^D(x-x')}{a^D}\Big]&=\big[2a^D\square^2-12H^2a^D\square+8a^{D-2}H^2\nabla^2+2(D^2-2D+2)H^4a^D\big]\delta^D(x-x'),
\label{eq:s5:154}\\
(aa')^{D-1}(a^2+a'^2)H^2\square\Big[\frac{\delta^D(x-x')}{a^D}\Big]&=2a^D\big(H^2\square-H^4\big)\delta^D(x-x'),
\label{eq:s5:155}
\end{align}
as printed; both were verified symbolically for general $D$ on an arbitrary test function.
The term $8a^{D-2}H^2\nabla^2$ has the form of the noninvariant counterterm \eqref{eq:s5:145}.

\paragraph*{Table VIII and $\beta_1=0$.}
Converting each row of Table~VI in this way gives our Table~\ref{tab:s5:T8}, which agrees with the printed Table~VIII row by row and in the totals.
The total of the $a^D\square^2$ column vanishes.
\PI\ explains this by the vanishing of this counterterm in flat space, and we confirmed that statement independently.
For $H\to0$ the $a^D\square^2$ coefficient must equal the divergence of the flat-space one-loop self-energy.
In flat space the 4-point contribution vanishes in dimensional regularization, and in the 3-point contribution the vertex $V^{\alpha\beta}\propto p^{(\alpha}q^{\beta)}-\frac12\eta^{\alpha\beta}p\cdot q$ (with $p$ the external and $q$ the loop momentum) contracted with the de Donder tensor $P=\frac12(\eta\eta+\eta\eta)-\frac1{D-2}\eta\eta$ gives
$V\!\cdot\!P\!\cdot\!V=\frac12p^2q^2$, so the self-energy is proportional to $\frac12p^2\int d^Dq/(p-q)^2=0$, a scaleless integral.

\begin{table}[tb]
\caption{Table~VIII of \PI, derived (paper convention): local normal contributions to the three counterterm structures. All entries are multiplied by $\frac{i\kappa^2H^{D-4}}{(4\pi)^{D/2}}\frac{\Gamma(D/2)}{(D-3)(D-4)}$. Rows and totals agree with the printed table.}
\label{tab:s5:T8}
\renewcommand{\arraystretch}{1.5}
\begin{ruledtabular}
\begin{tabular}{lccc}
From & $a^D\square^2\delta^D$ & $a^DH^2\square\delta^D$ & $a^{D-2}H^2\nabla^2\delta^D$\\
\hline
$\alpha$ & $\frac{D}{(D-1)(D-2)}$ & $0$ & $0$\\
$\beta\,\square/H^2$ (Eq.~\eqref{eq:s5:154}) & $-\frac{D}{(D-1)(D-2)}$ & $\frac{6D}{(D-1)(D-2)}$ & $-\frac{4D}{(D-1)(D-2)}$\\
$\beta$ (Eq.~\eqref{eq:s5:155}) & $0$ & $-\frac{(D+2)(D-4)}{2}$ & $0$\\
$\gamma_1\,\square/H^2$ & $0$ & $-\frac{D^2}{2(D-2)}$ & $0$\\
$\epsilon_1$ & $0$ & $0$ & $\frac{D^2-6D+4}{2(D-1)}$\\
$\epsilon_2$ & $0$ & $0$ & $\frac{2-4D}{D-1}$\\
\hline
Total & $0$ & $\frac{(D-4)(-D^3+D-4)}{2(D-1)(D-2)}$ & $\frac{D^3-16D^2+28D-16}{2(D-1)(D-2)}$\\
\end{tabular}
\end{ruledtabular}
\end{table}

\paragraph*{Table IX and shift symmetry.}
The terms proportional to $a^DH^4\delta^D$ (our Table~\ref{tab:s5:T9}) agree with the printed Table~IX.
A counterterm $H^4a^D\phi^2$ is forbidden by the shift symmetry, so their total must be finite.
It is: the $D=4$ values of the four rows are $-\frac{20}{3}+0+4+\frac83=0$, so the factor $(D-4)$ of the total cancels the pole of the units.
With the $\gamma_2$ sign printed in Table~IV the sum would be $-\frac{20}{3}+4-\frac83=-\frac{16}{3}$, a forbidden divergence.
This confirms the $\gamma_2$ sign independently of its direct computation.

\begin{table}[tb]
\caption{Table~IX of \PI, derived (paper convention): coefficients of $a^DH^4\delta^D(x-x')$, in the units of our Table~\ref{tab:s5:T8}; $q\equiv D^3-3D^2-4D+8$. Rows and total agree with the printed table.}
\label{tab:s5:T9}
\renewcommand{\arraystretch}{1.5}
\begin{ruledtabular}
\begin{tabular}{lll}
From & Coefficient & Origin\\
\hline
$\beta\,\square/H^2$ & $-\frac{D(D^2-2D+2)}{(D-1)(D-2)}$ & $H^4$ term of Eq.~\eqref{eq:s5:154}\\
$\beta$ & $\frac{(D+2)(D-4)}{2}$ & $H^4$ term of Eq.~\eqref{eq:s5:155}\\
$\gamma_1$ & $-\frac{D^4-5D^3+16D-16}{4}$ & $\gamma_1\to a^DH^4$\\
$\gamma_2$ & $\frac{(D-2)q}{2(D-1)}$ & $\gamma_2\to2a^DH^4$\\
\hline
Total & $\frac{(D-4)(-D^4+5D^3-16D+16)}{4(D-2)}$ & finite\\
\end{tabular}
\end{ruledtabular}
\end{table}

Since the total is finite we may take $D=4$ (cf.\ \eqI{157}--\eqI{158}):
\begin{equation}
\text{Table IX}=\frac{i\kappa^2H^{D-4}}{(4\pi)^{D/2}}\,\frac{(-D^4+5D^3-16D+16)\,\Gamma(\frac D2)}{4(D-2)(D-3)}\,a^DH^4\delta^D(x-x')
\;\xrightarrow{D=4}\;\frac{\kappa^2H^4}{8\pi^2}\,ia^4\delta^4(x-x') .
\label{eq:s5:157}
\end{equation}
This finite local term is rewritten in nonlocal form.
The propagator equation, with $i\Delta_A=A+k\ln(aa')$ and $\square\ln a=-(D-1)H^2$, gives (cf.\ \eqI{156})
\begin{equation}
ia^D\delta^D(x-x')=(aa')^D\big\{\square A(y)-(D-1)kH^2\big\},
\label{eq:s5:156}
\end{equation}
where $a^D$ may be replaced by $(aa')^D$ because the bracket is supported at coincidence, and $\square A$ is distributional.
At $D=4$, $k=H^2/8\pi^2$ and $A=\frac{H^2}{16\pi^2}\{\frac4y-2\ln\frac y4-1\}$, so (cf.\ \eqI{159})
\begin{equation}
\frac{\kappa^2H^4}{8\pi^2}\,ia^4\delta^4(x-x')=\frac{\kappa^2H^4}{(4\pi)^4}\Big\{(aa')^4H^2\square\Big[\frac2x-4\ln x\Big]-12\,(aa')^4H^4\Big\}.
\label{eq:s5:159}
\end{equation}
Equations \eqref{eq:s5:146}--\eqref{eq:s5:159} and Tables~VI--IX were checked by \texttt{check\_146\_IX.py}~\cite{ref:repo}.

\paragraph*{A possible extra local term.}
The question raised in Sec.~\ref{sec:s4b:distr}, whether the outer d'Alembertians of the operators $\alpha$ and $\delta$ count the local term \eqI{54}$+$\eqI{58} a second time, bears on this subsection.
In \PI\ the divergent local terms come only from Table~IV through Eqs.~\eqref{eq:s5:152}--\eqref{eq:s5:153}; in the finite functions of Table~V the relevant component has combined with the pole-carrying $z^{D-3}$ terms into finite logarithms.
Any double counting would then be a constant $c(D)$, which may contain a pole, times $(a^D\square-a^{D-2}\nabla^2)\delta^D(x-x')$, with no $\ln a$, of the form of the counterterms \eqref{eq:s5:invvert} and \eqref{eq:s5:145}; its pole and its finite part would be absorbed into $\alpha_2$ and $\alpha_3$.
Whether the outer-operator reduction produces other local differences at coincidence has not been settled (Sec.~\ref{sec:s4b:distr}); the point is only partially resolved.
In particular, a finite local term of mass type, $H^4a^D\delta^D(x-x')$, would not be of counterterm form; its consequences for the late-time analysis are discussed in Secs.~\ref{sec:s6:late} and~\ref{sec:s0s1s7s8:limitations}.

\subsection{\texorpdfstring{Local terms: the coefficients $\beta_i$}{Local terms: the beta coefficients}}
\label{sec:s5:beta}

This subsection corresponds to \eqI{160}--\eqI{165}.
The regulated self-mass-squared is the sum of the 4-point contribution and the local, logarithm and normal 3-point contributions.
Our Table~\ref{tab:s5:parts} shows where each piece goes.
The terms $3L\,a^D\square$, $-4L\,a^{D-2}\nabla^2$ and $9H^2a^D$ (times $i\delta^D$) of \eqI{44} are not of counterterm form.
They cancel against the local part of \eqI{67} for general $D$, up to an $O(D-4)$ remainder (Sec.~\ref{sec:s4a}); both carry the reversed sign in \PI, so the cancellation holds in either convention.

\begin{table}[tb]
\caption{Destination of the parts of the regulated self-mass-squared.}
\label{tab:s5:parts}
\renewcommand{\arraystretch}{1.25}
\begin{ruledtabular}
\begin{tabular}{p{0.24\linewidth}p{0.38\linewidth}p{0.3\linewidth}}
Part & Local ($\to\beta_i$) & Finite nonlocal ($\to$ Table~X)\\
\hline
4-point, \eqI{44} & $\pihat$ terms; constants $\mp2$ from \eqI{40}; $3L$, $-4L$, $9H^2$ terms (cancel) & ---\\
3-point local, \eqI{59} & $\pihat$ terms and finite constants & ---\\
3-point logarithm, \eqI{67} & $L$ and $9H^2$ terms (cancel) & $2H(aa')^2(a\partial_0+a'\partial'_0)\nabla^2A$ $\to$ Table~III\\
Normal, Table~IV & Table~VI $\to$ Table~VIII; Table~IX $\to$ Eq.~\eqref{eq:s5:159} & Table~VII; Eq.~\eqref{eq:s5:159}\\
Normal, Table~V & --- & Table~V\\
\end{tabular}
\end{ruledtabular}
\end{table}

What remains has counterterm form (paper convention; cf.\ \eqI{160}):
\begin{equation}
-iM^2\reg(x;x')=i\kappa^2a^D\Big(\beta_1\square^2+\beta_2H^2\square+\beta_3H^2\frac{\nabla^2}{a^2}\Big)\delta^D(x-x')+\text{Table X}+O(D-4).
\label{eq:s5:160}
\end{equation}
\PI\ prints $\beta_2\square+\beta_3\nabla^2/a^2$ without the factors $H^2$ (C17).%
\footnote{The $\beta_i$ of \eqI{161}--\eqI{165} all have the dimension $H^{D-4}/(4\pi)^{D/2}$, the counterterm vertices \eqI{131} and \eqI{145} carry $H^2$, and \eqI{166} requires $\alpha_i=-\beta_i+\Delta\alpha_i$. The factors $H^2$ are therefore required; the same misprint occurs in \eqI{167}, \eqI{170}, \eqII{10} and \eqII{34} (C18, C22, C20, C21).}

\paragraph*{$\beta_1$.}
The structure $\square^2$ occurs only in the first column of Table~VIII, whose total vanishes. Hence (cf.\ \eqI{161})
\begin{equation}
\beta_1=0 .
\label{eq:s5:161}
\end{equation}

\paragraph*{$\beta_2$ and $\beta_3$ from the printed inputs.}
The coefficient of $i\kappa^2H^2a^D\square\delta^D$ receives $k[\frac14D(D-1)\pihat-2]$ from \eqI{44}, $k[-\frac{D-1}{D-3}\pihat+2]$ from \eqI{59}, and the second column of Table~VIII.
The coefficient of $i\kappa^2H^2a^{D-2}\nabla^2\delta^D$ receives $k[-D\pihat+2]$ from \eqI{44}, $k[\frac{D-1}{D-3}\pihat-\frac72]$ from \eqI{59}, and the third column of Table~VIII.
In both cases the $\pihat$ terms combine, and the constants give $0$ and $-\frac32k$ respectively; the latter is written by its $D=4$ value, which is legitimate because it multiplies no pole.
With $(D-1)\Gamma(D-1)=\Gamma(D)$ and $\Gamma(\frac D2)/(D-2)=\frac12\Gamma(\frac D2-1)$ we obtain (cf.\ \eqI{162}--\eqI{165})
\begin{align}
\beta_2&=\frac{H^{D-4}}{(4\pi)^{D/2}}\Big\{\frac{(-D^3+D-4)\,\Gamma(\frac D2-1)}{4(D-1)(D-3)}-\frac{(D+1)(D-4)\,\Gamma(D)\,\pi\cot(\frac{\pi D}2)}{4(D-3)\,\Gamma(\frac D2)}\Big\}
=\frac{H^{D-4}}{(4\pi)^{D/2}}\Big\{-\frac{61}{3}+O(D-4)\Big\},
\label{eq:s5:162}\\
\beta_3&=\frac{H^{D-4}}{(4\pi)^{D/2}}\Big\{\frac{(D^3-16D^2+28D-16)\,\Gamma(\frac D2)}{2(D-1)(D-2)(D-3)(D-4)}+\frac{(D^2-4D+1)\,\Gamma(D-1)\,\pi\cot(\frac{\pi D}2)}{(D-3)\,\Gamma(\frac D2)}-3\Big\}
\notag\\
&=\frac{H^{D-4}}{(4\pi)^{D/2}}\Big\{-\frac{4}{D-4}+\frac{58}{3}+2\sfivegE+O(D-4)\Big\},
\label{eq:s5:164}
\end{align}
where $(D-4)\pi\cot(\frac{\pi D}{2})\to2$ and $\sfivegE$ enters through $\psi(2)=1-\sfivegE$ and $\psi(3)=\frac32-\sfivegE$.
These reproduce the printed \eqI{161}--\eqI{165} exactly (\texttt{check\_160\_167.py}~\cite{ref:repo}), so \PI\ did use the printed \eqI{44} and \eqI{59}.
Within \PI's treatment of the outer operators of the normal terms, the coefficient $\beta_2$ is finite and the only divergence is in the de Sitter-breaking coefficient $\beta_3$, consistent with Sec.~\ref{sec:s5:noninv}.
A double counting of the kind discussed at the end of Sec.~\ref{sec:s5:div}, if present, would shift poles and finite parts of $\beta_2$ and $\beta_3$ by counterterm-form amounts, absorbed into $\alpha_2$ and $\alpha_3$.

\paragraph*{Corrected $\beta_2$ and $\beta_3$.}
Two finite errors in these inputs were found earlier (B1 and B2 of our Table~\ref{tab:s0s1s7s8:finite}).
With the correct \eqI{40} (Sec.~\ref{sec:s3}), the constants $\mp2$ are absent from \eqI{44}; this adds $+2\sfivekb\to+4$ to $\beta_2$ and $-2\sfivekb\to-4$ to $\beta_3$.
The constant of the $\nabla^2$ term of \eqI{59} is $-\frac52$, the sum of the paper's own \eqI{57} and \eqI{58}, not $-\frac72$ (Sec.~\ref{sec:s4a}); this adds $+\sfivekb\to+2$ to $\beta_3$.
Equivalently, the bracket of Eq.~\eqref{eq:s5:162} acquires $+2\sfivekb$, and the constant $-3$ in Eq.~\eqref{eq:s5:164} becomes $-5$.
The corrected values are compared with the printed ones in our Table~\ref{tab:s5:beta}.
Only the finite parts change; they are absorbed into the arbitrary finite counterterms, and the pole is unchanged by these two corrections.
The same shifts apply to the counterterm coefficients $\alpha_2$ and $\alpha_3$ quoted in \eqII{8}--\eqII{9} (B5, B6).

\begin{table}[tb]
\caption{The counterterm coefficients $\beta_i$ of \eqI{161}--\eqI{165}, in units of $H^{D-4}/(4\pi)^{D/2}$, up to $O(D-4)$. The shifts of $\beta_2$ and $\beta_3$ are B3 and B4.}
\label{tab:s5:beta}
\renewcommand{\arraystretch}{1.45}
\begin{ruledtabular}
\begin{tabular}{lccc}
& Printed (paper convention) & Corrected (paper convention) & Corrected (definition sign)\\
\hline
$\beta_1$ & $0$ & $0$ & $0$\\
$\beta_2$ & $-\frac{61}{3}$ & $-\frac{49}{3}$ & $+\frac{49}{3}$\\
$\beta_3$ & $-\frac{4}{D-4}+\frac{58}{3}+2\sfivegE$ & $-\frac{4}{D-4}+\frac{52}{3}+2\sfivegE$ & $+\frac{4}{D-4}-\frac{52}{3}-2\sfivegE$\\
\end{tabular}
\end{ruledtabular}
\end{table}

\subsection{Table X}
\label{sec:s5:TX}

This subsection corresponds to Table~X and the text before \eqI{160}.
Table~X collects every finite nonlocal term: Table~III (the nonlocal logarithm term of \eqI{67}, Sec.~\ref{sec:s4b:tableIII}), Table~V (the finite normal terms, as printed after the audit of Sec.~\ref{sec:s4b:appendix}), Table~VII and Eq.~\eqref{eq:s5:159}.
The redundant operators are eliminated with $\gamma_3=2\gamma_1+\gamma_2$ and $\epsilon_3=2\epsilon_1+\epsilon_2$; the row $(aa')^2(a+a')^2H^2\nabla^2$ of Table~III is $\epsilon_3$.
Combined operators of Table~VII are read at $D=4$, for example $\beta\,\square/H^2=(aa')^3(a^2+a'^2)\square^2$ and $\gamma_1\square^2/H^4=(aa')^4\square^2$.
In all operators the scale factors stand to the left of the derivatives, which act on the coefficient functions.
Our Table~\ref{tab:s5:assembly} lists the sources of each row.

\begin{table}[tb]
\caption{Sources of the rows of Table~X of \PI\ (paper convention, coefficients of $\kappa^2H^4/(4\pi)^4$). ``VII~$\alpha$'' is the $\alpha$ row of Table~VII, ``V~$\alpha$'' the $\alpha$ row of Table~V, and so on. The last column compares the total of the listed sources with the printed Table~X.}
\label{tab:s5:assembly}
\renewcommand{\arraystretch}{1.3}
\begin{ruledtabular}
\begin{tabular}{lp{0.53\linewidth}l}
Row & Sources & vs.\ printed\\
\hline
$(aa')^4\square^3/H^2$ & VII $\alpha$ & equal\\
$(aa')^4\square^2$ & VII $\alpha$; VII $\gamma_1$ ($\square^2/H^4$ part); V $\alpha$ & equal\\
$(aa')^4H^2\square$ & VII $\gamma_1$ ($\square/H^2$ part); Eq.~\eqref{eq:s5:159}: $\frac2x-4\ln x$ & equal\\
$(aa')^4H^4$ & VII $\gamma_1$; V $\gamma_1$; $2\times$V $\gamma_3$; Eq.~\eqref{eq:s5:159}: $-12$ & equal\\
$(aa')^3(a^2+a'^2)\square^3/H^2$ & VII $\beta$ & equal\\
$(aa')^3(a^2+a'^2)\square^2$ & VII $\beta$ & equal\\
$(aa')^3(a^2+a'^2)H^2\square$ & VII $\gamma_2$ ($\square/H^2$ part); V $\beta$ & equal\\
$(aa')^3(a^2+a'^2)H^4$ & VII $\gamma_2$; V $\gamma_2$; V $\gamma_3$ & equal to printed $-54$; see text\\
$(aa')^3H^2\nabla^2$ & VII $\epsilon_1$; V $\epsilon_1$; $2\times$V $\epsilon_3$; III: $-\frac{12}x+24\ln x$; $2\times$III $\epsilon_3$: $\frac8x$ & equal\\
$(aa')^3\nabla^2\square$ & VII $\epsilon_1$ ($\square/H^2$ part) & equal\\
$(aa')^2(a^2+a'^2)H^2\nabla^2$ & VII $\epsilon_2$; V $\epsilon_2$; V $\epsilon_3$; III $\epsilon_3$: $\frac4x$ & equal\\
$(aa')^2(a^2+a'^2)\nabla^2\square$ & VII $\epsilon_2$ ($\square/H^2$ part); V $\delta$ & equal (printed twice)\\
$(aa')^2\nabla^4$ & V $\zeta$; III $\zeta$: $8\ln x-16x\ln x$ & III omitted in print\\
\end{tabular}
\end{ruledtabular}
\end{table}

Twelve of the thirteen rows follow exactly from the inputs; the row $(aa')^2(a^2+a'^2)\nabla^2\square$ is printed twice in Table~X, a misprint (C19; Table~1 of \PII\ lists it once).
\PI\ added the first two rows of Table~III to the $\epsilon_1$ and $\epsilon_2$ rows; for example $84\ln x=60\ln x+24\ln x$ and $-\frac{16}{x}=(\frac13-\frac{26}3-\frac{11}3-12+8)\frac1x$.
It omitted the third row, which belongs to $\zeta$, although the text states that Table~X includes Table~III (S2).
\PI\ prints
\begin{equation}
(aa')^2\nabla^4:\qquad
\begin{aligned}
&\tfrac{10}{3}\ln x-24x\ln x+24x^2\ln x-36x^2 &&\text{(printed)},\\
&\tfrac{34}{3}\ln x-40x\ln x+24x^2\ln x-36x^2 &&\text{(corrected)}.
\end{aligned}
\label{eq:s5:zeta}
\end{equation}
The difference, $8\ln x-16x\ln x$, is non-analytic.

The row $(aa')^3(a^2+a'^2)H^4$ agrees with the print when Table~V is used as printed, but Table~V lacks a finite term $72\gamma_2$ (paper convention; S3), found globally in the audit of Tables~IV--V (Sec.~\ref{sec:s4b:tableV}).
In Table~X this changes the constant of that row from $-54$ to $+18$.
Equivalently, one may add $-12x$ to the $\delta$ row of Table~V, since $\nabla^2\square x=-6aa'H^4$ gives $\delta[-12x]=72\gamma_2[1]$, or make the $\gamma_2$ row of Table~V $-\frac5x+54$.
The pointwise test below fixes only the sum, so which row the term belongs to is a matter of interpretation.
The corrected Table~X is given in our Table~\ref{tab:s5:TX}, in both conventions; we denote its content by
\begin{equation}
\Tab(x;x')\equiv\frac{\kappa^2H^4}{(4\pi)^4}\sum_{\mathcal O}\mathcal O\big[c_{\mathcal O}(y/4)\big],\qquad y\equiv y(x;x'),
\label{eq:s5:Tdef}
\end{equation}
where the sum runs over the thirteen operators $\mathcal O$ of our Table~\ref{tab:s5:TX} and $c_{\mathcal O}$ are the paper-convention coefficients, written in the table as functions of $x\equiv y/4$.

\begin{table}[tb]
\caption{Corrected Table~X of \PI: all finite nonlocal contributions, coefficients of $\kappa^2H^4/(4\pi)^4$, $x=y/4$. The rows are numbered from 0 as in the derivation notes and the code; these row numbers are used in Sec.~\ref{sec:s6}. The paper-convention column is to be compared with the printed Table~X and with Table~1 of \PII{} (which lists the same rows in the same order); the definition-sign column is $-1$ times it. The two changed entries (S3 and S2) are marked with an asterisk: printed $(aa')^3(a^2+a'^2)H^4$ constant $-54$, printed $(aa')^2\nabla^4$ row $\frac{10}{3}\ln x-24x\ln x+24x^2\ln x-36x^2$. The row $(aa')^2(a^2+a'^2)\nabla^2\square$ appears once.}
\label{tab:s5:TX}
\renewcommand{\arraystretch}{1.5}
\begin{ruledtabular}
\begin{tabular}{rlll}
\# & Operator $\mathcal O$ & Paper convention, $c_{\mathcal O}$ & Definition sign, $-c_{\mathcal O}$\\
\hline
0 & $(aa')^4\square^3/H^2$ & $-\frac{\ln x}{3x}$ & $\frac{\ln x}{3x}$\\
1 & $(aa')^4\square^2$ & $\frac{26\ln x}{3x}+\frac{38}{3x}-6\ln^2x-18\ln x$ & $-\frac{26\ln x}{3x}-\frac{38}{3x}+6\ln^2x+18\ln x$\\
2 & $(aa')^4H^2\square$ & $-\frac{6\ln x}{x}+\frac4x-4\ln x$ & $\frac{6\ln x}{x}-\frac4x+4\ln x$\\
3 & $(aa')^4H^4$ & $\frac{4\ln x}{x}+\frac{18}{x}-120-108\ln x$ & $-\frac{4\ln x}{x}-\frac{18}{x}+120+108\ln x$\\
4 & $(aa')^3(a^2+a'^2)\square^3/H^2$ & $\frac{\ln x}{6x}$ & $-\frac{\ln x}{6x}$\\
5 & $(aa')^3(a^2+a'^2)\square^2$ & $-\frac{\ln x}{3x}+\frac1{6x}$ & $\frac{\ln x}{3x}-\frac1{6x}$\\
6 & $(aa')^3(a^2+a'^2)H^2\square$ & $-\frac{2\ln x}{3x}+\frac5x-18\ln x$ & $\frac{2\ln x}{3x}-\frac5x+18\ln x$\\
7 & $(aa')^3(a^2+a'^2)H^4$ & $\frac{4\ln x}{3x}-\frac{32}{3x}+18^{\,*}$ & $-\frac{4\ln x}{3x}+\frac{32}{3x}-18^{\,*}$\\
8 & $(aa')^3H^2\nabla^2$ & $-\frac{2\ln x}{3x}-\frac{16}{x}+84\ln x-48x\ln x+96x$ & $\frac{2\ln x}{3x}+\frac{16}{x}-84\ln x+48x\ln x-96x$\\
9 & $(aa')^3\nabla^2\square$ & $\frac{\ln x}{3x}$ & $-\frac{\ln x}{3x}$\\
10 & $(aa')^2(a^2+a'^2)H^2\nabla^2$ & $-\frac{7\ln x}{3x}+\frac{11}{2x}-12\ln x+48x\ln x+12x$ & $\frac{7\ln x}{3x}-\frac{11}{2x}+12\ln x-48x\ln x-12x$\\
11 & $(aa')^2(a^2+a'^2)\nabla^2\square$ & $-\frac{17\ln x}{6x}-\frac{49}{6x}+4\ln^2x+10\ln x+12x\ln x$ & $\frac{17\ln x}{6x}+\frac{49}{6x}-4\ln^2x-10\ln x-12x\ln x$\\
12 & $(aa')^2\nabla^4$ & $\big(\frac{34}{3}\ln x-40x\ln x\big)^{*}+24x^2\ln x-36x^2$ & $\big(-\frac{34}{3}\ln x+40x\ln x\big)^{*}-24x^2\ln x+36x^2$\\
\end{tabular}
\end{ruledtabular}
\end{table}

\paragraph*{Global pointwise test.}
The two corrections were established by a test that does not depend on how Table~X was assembled.
At $D=4$ and $x\neq x'$ it compares two quantities, both in the paper convention.
The left-hand side is the direct sum $\sum_kP_k$ of the normal 3-point contributions plus the nonlocal term $-\frac{\kappa^2H^2}{8\pi^2}\,2H(aa')^2(a\partial_0+a'\partial'_0)\nabla^2A$ of \eqI{67}, evaluated with naive (pointwise) inner derivatives and the actual $D=4$ propagator functions $A=h\{4/y-2\ln(y/4)-1\}$, $B=C=4h/y$, $h\equiv H^2/16\pi^2$; it does not use the outer-operator reduction of the appendix of \PI.
The right-hand side is the thirteen operators of Table~X applied pointwise to their coefficient functions.
Equation~\eqref{eq:s5:159} does not enter, since its two terms cancel pointwise ($\sfiveLy[8/y-4\ln(y/4)]=12$), and at $D=4$ the nonlocal parts of Eqs.~\eqref{eq:s5:152}--\eqref{eq:s5:153} reproduce $(4/y)^2$ and $(4/y)^3$ pointwise.
The evaluation uses three random rational spacelike pairs of points (seed 29) and 30-digit arithmetic (\texttt{check\_tableX\_global.py}~\cite{ref:repo}).
Our Table~\ref{tab:s5:global} shows the residuals.
With both corrections the agreement is exact to 30 digits; neither correction alone is enough.
The test is numerical, at sample points, and checks only the sum of all rows; it is not a symbolic identity.

\begin{table}[tb]
\caption{Global pointwise test of Table~X of \PI: left-hand side minus the operator sum at three random spacelike pairs. For scale, $|{\rm LHS}|=7.10$, $0.107$ and $8.60$.}
\label{tab:s5:global}
\renewcommand{\arraystretch}{1.3}
\begin{ruledtabular}
\begin{tabular}{lccc}
Version of Table~X & Point 1 & Point 2 & Point 3\\
\hline
As printed & $-11.52$ & $64.14$ & $-3.526$\\
With the $(aa')^2\nabla^4$ correction only & $1.182$ & $65.28$ & $0.0161$\\
With the $(aa')^2\nabla^4$ and $72\gamma_2$ corrections & $4\times10^{-31}$ & $0$ & $-3\times10^{-31}$\\
\end{tabular}
\end{ruledtabular}
\end{table}

\paragraph*{Physical significance of the two corrections.}
The $72\gamma_2$ term is independent of $y$ (analytic). In the Schwinger--Keldysh formalism it is identical in the $++$ and $+-$ kernels, so it drops out of the retarded self-mass-squared identically and has no effect on any retarded quantity.
The term $(aa')^2\nabla^4(8\ln x-16x\ln x)$ is non-analytic through $\ln x$; it enters the retarded kernel and changes the physical nonlocal self-mass-squared.
Its effect on the late-time mode function is analyzed in Sec.~\ref{sec:s6}: it contributes only terms that fall off at late times (it removes the $\hat k^4/a$, $\hat k^6/a,\dots$ tails that the printed row leaves in the source).

\subsection{The renormalized self-mass-squared}
\label{sec:s5:final}

This subsection corresponds to \eqI{166}--\eqI{167}.
The counterterm vertices \eqref{eq:s5:invvert} and \eqref{eq:s5:145} add $i\kappa^2a^D[\alpha_1\square^2+\alpha_2H^2\square+\alpha_3H^2\nabla^2/a^2]\delta^D$ to $-iM^2\reg$.
The choice (cf.\ \eqI{166})
\begin{equation}
\alpha_i=-\beta_i+\Delta\alpha_i
\label{eq:s5:166}
\end{equation}
removes the divergence and leaves arbitrary finite constants $\Delta\alpha_i$; this relation holds in either convention.
Because $\alpha_ia^D$ cancels $\beta_ia^D$ exactly, the expansion $a^D=a^4[1+(D-4)\ln a+\dots]$ leaves no $\ln a$ behind, and we may set $D=4$.
In the paper convention, with the factors $H^2$ restored and the corrected Table~X, this is (cf.\ \eqI{167})
\begin{equation}
-iM^2\ren(x;x')\Big|_{\text{paper conv.}}=i\kappa^2a^4\Big(\Delta\alpha_1\square^2+\Delta\alpha_2H^2\square+\Delta\alpha_3H^2\frac{\nabla^2}{a^2}\Big)\delta^4(x-x')+\Tab(x;x').
\label{eq:s5:167p}
\end{equation}

As established in Secs.~\ref{sec:s3} and~\ref{sec:s4a}, the printed self-mass-squared is $-1$ times the one defined by \PI's own \eqI{168}.
Both the 4-point part and the 3-point part are reversed, so their relative sign, and the local cancellations used above, are consistent.
In the definition sign the regulated result is Eq.~\eqref{eq:s5:160} with $\beta_i\to-\beta_i$ and Table~X $\to-$Table~X; with the corrections of Sec.~\ref{sec:s5:beta},
\begin{equation}
\beta_1=0,\qquad
\beta_2=\frac{H^{D-4}}{(4\pi)^{D/2}}\Big\{\frac{49}{3}+O(D-4)\Big\},\qquad
\beta_3=\frac{H^{D-4}}{(4\pi)^{D/2}}\Big\{\frac{4}{D-4}-\frac{52}{3}-2\sfivegE+O(D-4)\Big\}\qquad\text{(definition sign)},
\label{eq:s5:betadef}
\end{equation}
and Eq.~\eqref{eq:s5:166} again removes the divergence.
The final result of this section is
\begin{equation}
\boxed{\;
\begin{aligned}
-iM^2\ren(x;x')={}&i\kappa^2a^4\Big(\Delta\alpha_1\square^2+\Delta\alpha_2H^2\square+\Delta\alpha_3H^2\frac{\nabla^2}{a^2}\Big)\delta^4(x-x')\\
&-\frac{\kappa^2H^4}{(4\pi)^4}\sum_{\mathcal O}\mathcal O\big[c_{\mathcal O}(y/4)\big]\qquad\text{(definition sign)}
\end{aligned}\;}
\label{eq:s5:final}
\end{equation}
where the sum runs over the thirteen operators of our Table~\ref{tab:s5:TX}, $c_{\mathcal O}$ are its paper-convention coefficients (so the nonlocal part is $-\Tab$, the definition-sign column), $y\equiv y(x;x')$, and the $\Delta\alpha_i$ are arbitrary finite constants.

Since the $\Delta\alpha_i$ are arbitrary, the sign of the local part has no physical meaning.
The nonlocal part has no counterterm freedom, and its sign is physical: it is opposite to the printed one.
In the late-time analysis of \PII\ this flips every entry of Table~5 (whose total remains zero) and the sign of the particular solution $\Phi_1$.
Section~\ref{sec:s6} shows that the non-falling total of the source is zero at late times, so only terms that fall off change sign, and the late-time conclusion of \PII\ is unchanged.

\paragraph*{Comparison with Sec.~IV of \PI.}
Equations \eqI{126}--\eqI{159}, Tables~VI--IX and the structure of \eqI{160} and \eqI{167} are confirmed, with the misprints noted above: the omission of two (shift-forbidden) terms from the list \eqI{128} (C14), the sign of the $\gamma_2$ entry of Table~IV (C15), the factor $i$ in the header of Table~VII (C16), the missing factors $H^2$ in \eqI{160} and \eqI{167} (C17, C18), and the duplicated row of Table~X (C19)\@. We also noted a notational remark without effect: the dilatation factor ($k$ in \eqI{134}) must be positive.
The coefficients \eqI{161}--\eqI{165} are reproduced exactly from the printed inputs; with the corrected \eqI{40} and \eqI{59} (B1, B2), the finite parts of \eqI{163} and \eqI{165} change (B3, B4), $-\frac{61}{3}\to-\frac{49}{3}$ and $\frac{58}{3}\to\frac{52}{3}$ (paper convention), as in our Table~\ref{tab:s5:beta}, with the pole unchanged; these shifts are absorbable. The relation \eqI{166} is confirmed.
Table~X needs the two corrections of Sec.~\ref{sec:s5:TX} (S2, S3), and the overall sign of the self-mass-squared is reversed with respect to the definition \eqI{168} (S1).

\paragraph*{Limitations.}
The coefficient functions of Table~V enter as printed; the appendix of \PI\ from which they come was audited through worked examples and the global pointwise test, not redone row by row (Sec.~\ref{sec:s4b:appendix}).
The possible extra local term of Sec.~\ref{sec:s4b:distr} is only partially resolved; any double counting of the kind identified has counterterm form.
The global tests are numerical, at sample points, with 30-digit precision; the completeness of the counterterm basis was checked at $D=4,5,6$, and the breaking of the gauge symmetry by class (iv) at the level of the action in $D=3$.
These points, and the assumptions shared with \PI, are collected in Sec.~\ref{sec:s0s1s7s8:limitations}.

\providecommand{\sixR}{\mathcal{R}}
\providecommand{\sixDk}{\mathcal{D}_k}
\providecommand{\sixO}{O}
\providecommand{\sixunit}{\frac{H^4u(0,k)}{16\pi^2}}

\section{Schwinger--Keldysh effective field equation and late-time mode functions}
\label{sec:s6}

In this section the renormalized self-mass-squared obtained above, i.e.\ the corrected Table~X with the overall sign
fixed by \PI's own definition \eqI{168}, is used as the input of the Schwinger--Keldysh (SK) analysis of
\PII~\cite{KW08}. We follow the order of \PII, Secs.~2--3: the input table and the counterterms
(cf.\ \eqII{7}--\eqII{12} and Table~1 of \PII), the four SK polarities and the effective field equation
(cf.\ \eqII{15}--\eqII{21}), the perturbative mode equation and its local part (cf.\ \eqII{22}--\eqII{34}),
the SK differences and the integrals of Tables~2--4 (cf.\ \eqII{35}--\eqII{43}), and the Table~5 cancellation
within the simplification \eqII{44}. We then drop \eqII{44}, keep the full dependence on $k$, and solve for the
late-time one-loop mode function. Tables~2--5 of \PII{} were not used as input; they were compared with our results
only at the end. The full derivation is Ch.~9 of the derivation notes, and the scripts quoted below are in
\texttt{code/09\_schwinger\_keldysh/} of Ref.~\cite{ref:repo}.

\paragraph*{Conventions.}
Two signs appear. The \emph{paper convention}, which in this section we call the \emph{Table-X convention}, is the sign in
which \PI{} prints Table~X and \PII{} prints Table~1, Eq.~\eqII{10} and Table~5. The \emph{definition sign} is the sign
fixed by \eqI{168} [Eq.~\eqref{eq:s0s1s7s8:def}; Sec.~\ref{sec:s2:definition}]; \PII{} restates the same equation as \eqII{1}.
The definition sign is $-1$ times the Table-X convention, and physical statements are made in the definition sign.
As in \PII, everything refers to the renormalized self-mass-squared in the unregulated limit $D=4$; general $D$
appears only in the counterterm coefficients. We use
\begin{equation}
x\equiv\frac y4,\qquad \Delta\eta\equiv\eta-\eta',\qquad r\equiv\|\vec x\,'-\vec x\|,\qquad
X_0\equiv\tfrac14aa'H^2\Delta\eta^2,\qquad \khat\equiv\frac kH,
\label{eq:s6:notation}
\end{equation}
set $H=1$ in intermediate steps and restore it by dimensional analysis. We count $\kappa^2$ once: $\Phi=u+\kappa^2\Phi_1$
(cf.\ \eqII{23}), and $\Phi_1$ is sourced by the one-loop kernel with $\kappa^2$ stripped off, Eq.~\eqref{eq:s6:mode}
below. (The original working report of this step wrote an explicit $\kappa^2$ also in the source and in the particular solution,
i.e.\ it counted $\kappa^2$ twice. Ch.~9 of the derivation notes~\cite{ref:repo} and the present paper count it once. No conclusion
depends on this.)
Table~X/Table~1 entries are coefficients of $\kappa^2H^4/(4\pi)^4$; the source $\sixR_1$ and the entries of Table~5 are
quoted in units of $H^4u(0,k)/16\pi^2$.

\subsection{Input: the renormalized self-mass-squared}
\label{sec:s6:input}

\PII{} prints in its Table~1 the nonlocal contributions to $-iM^2(x;x')$ of \PI{} Table~X, with the row
$(aa')^2(a^2+a'^2)\nabla^2\square$ listed once (\PI{} prints it twice, erratum C19). The corrected table is our Table~\ref{tab:s5:TX}
(paper-convention column); its rows are listed in the same order as in Table~1 of \PII,
and we use its row numbers 0--12 throughout this section. It differs from the printed one in two entries, both
established by the global pointwise test of Table~X in Sec.~\ref{sec:s5:TX}: in row~12 (the $\zeta$ row
$(aa')^2\nabla^4$) the Table~III contribution $8\ln x-16x\ln x$ is restored (S2, non-analytic), and in row~7 the constant
is $+18$ instead of $-54$ (S3, the $72\gamma_2$ correction, analytic).

The length function is $y(x;x')=aa'H^2\{\|\vec x-\vec x\,'\|^2-(|\eta-\eta'|-i\delta)^2\}$; Eq.~\eqII{12} is its
$H=1$ form. We denote by $\Tab(y)$ the sum of all rows of our Table~\ref{tab:s5:TX} (Table-X convention). Scale factors
stand to the left of the operators, and $\square$, $\nabla^2$ act only on $f(y)$. This ordering is an assumption;
it is the ordering under which the pointwise test of Table~X agrees to 30 digits, and the factor $a^{4-K}$ in \eqII{45}
implies the same ordering.

\paragraph*{Counterterms (cf.\ \eqII{7}--\eqII{9}).}
\PII{} prints $\alpha_1=0+\Delta\alpha_1$,
$\alpha_2=\frac{H^{D-4}}{(4\pi)^{D/2}}\{\frac{61}{3}\}+\Delta\alpha_2$ and
$\alpha_3=\frac{H^{D-4}}{(4\pi)^{D/2}}\{\frac4{D-4}-\frac{58}3-2\gamma_{\rm E}\}+\Delta\alpha_3$.
With $\alpha_i=-\beta_i+\Delta\alpha_i$, \eqI{166}, and the corrected $\beta_{2,3}$ of Sec.~\ref{sec:s5:beta} we obtain
\begin{equation}
\alpha_1=0+\Delta\alpha_1,\qquad
\alpha_2=\frac{H^{D-4}}{(4\pi)^{D/2}}\Big\{\pm\frac{49}{3}\Big\}+\Delta\alpha_2,\qquad
\alpha_3=\frac{H^{D-4}}{(4\pi)^{D/2}}\Big\{\pm\Big[\frac{4}{D-4}-\frac{52}{3}-2\gamma_{\rm E}\Big]\Big\}+\Delta\alpha_3,
\label{eq:s6:alpha}
\end{equation}
with the upper sign in the Table-X convention and the lower sign in the definition sign. The pole is unchanged; only the
finite parts $\frac{61}3\to\frac{49}3$ and $-\frac{58}3\to-\frac{52}3$ (Table-X convention) move, and they are absorbed into
the arbitrary $\Delta\alpha_{2,3}$ (B5, B6). They enter none of the results below.

\paragraph*{Renormalized self-mass (cf.\ \eqII{10}).}
\PII{} writes $-iM^2=i\kappa^2a^4(\Delta\alpha_1\square^2+\Delta\alpha_2\square+\Delta\alpha_3\nabla^2/a^2)\delta^4(x-x')
+(\text{Table 1})$, without the factors $H^2$ next to $\Delta\alpha_2$ and $\Delta\alpha_3$ that dimensions require
(erratum C20; the same omission as in \eqI{160}, \eqI{167} and \eqI{170}, errata C17, C18 and C22; it has no effect). In the definition sign, with these
factors shown,
\begin{equation}
-iM^2_{++}(x;x')=i\kappa^2a^4\Big(\Delta\alpha_1\square^2+\Delta\alpha_2H^2\square+\Delta\alpha_3H^2\frac{\nabla^2}{a^2}\Big)
\delta^4(x-x')-\Tab(y_{++}).
\label{eq:s6:Mpp}
\end{equation}
The minus sign in front of $\Tab$ is the overall sign correction: the physical nonlocal part is $-$(corrected Table~X).
The sign of the local part is immaterial, since the $\Delta\alpha_i$ are arbitrary.

\subsection{Schwinger--Keldysh polarities and the effective field equation}
\label{sec:s6:sk}

In the SK formalism~\cite{ref:Schwinger1961,ref:Keldysh1965} each propagator has four polarities, obtained by replacing
$y$ in the Feynman propagator by (cf.\ \eqII{15}--\eqII{18})
\begin{equation}
y_{\pm\pm'}=aa'H^2\big[r^2-(\ldots)^2\big],\qquad
(\ldots)=|\Delta\eta|-i\delta\ (++),\quad \Delta\eta+i\delta\ (+-),\quad \Delta\eta-i\delta\ (-+),\quad |\Delta\eta|+i\delta\ (--).
\label{eq:s6:ypol}
\end{equation}
The $+$ vertex is the in-out vertex and the $-$ vertex is its conjugate; all lines at a vertex carry its polarity.
At one loop $-iM^2_{++}$ is the in-out self-mass-squared, Eq.~\eqref{eq:s6:Mpp}. The $+-$ self-mass follows in three steps:
(i) the vertex at $x'$ has $-$ polarity, and since the vertex functionals are real its factor is
$(i\,\delta^nS)^*=-i\,\delta^nS$, an overall factor $-1$;
(ii) every line from $x$ to $x'$ becomes $i\Delta_{+-}$, i.e.\ $y\to y_{+-}$, while the external operators (derivatives
and scale factors) are unchanged;
(iii) there are no $\delta^4(x-x')$ terms, because the 4-point vertex needs coincident legs and the
$\partial_0|\Delta\eta|$ terms of the time-ordered propagator do not occur in $y_{+-}$. Hence
\begin{equation}
-iM^2_{+-}(x;x')=+\Tab(y_{+-}),\qquad -iM^2_{--}=\big(-iM^2_{++}\big)^*,\qquad -iM^2_{-+}=\big(-iM^2_{+-}\big)^*,
\label{eq:s6:Mpm}
\end{equation}
where the last two relations are \eqII{20}, evaluated with $\Tab(y)^*=\Tab(y^*)$ (real coefficients and scale factors),
$y_{++}^*=y_{--}$ and $y_{+-}^*=y_{-+}$. This is the rule stated in \PII{} below \eqII{19} (drop the $\delta$-function
terms, reverse the sign, replace $y$ by $y_{+-}$), written in the definition sign. Since the operators of
our Table~\ref{tab:s5:TX} are symmetric under $x\leftrightarrow x'$ and $y_{-+}(x;x')=y_{+-}(x';x)$, one has
$M^2_{-+}(x';x)=M^2_{+-}(x;x')$.

Varying $\Gamma[\varphi_+,\varphi_-]$ of \eqII{19} with respect to $\varphi_+(x)$, using the symmetry of $M^2_{++}$ and
the last relation, and setting $\varphi_\pm=\Phi$ gives \eqII{21},
\begin{equation}
a^4\square\Phi(x)-\int_{\eta_i}^0d\eta'\!\int d^3x'\,\big\{M^2_{++}(x;x')+M^2_{+-}(x;x')\big\}\Phi(x')=0,
\qquad \eta_i=-\frac1H,
\label{eq:s6:effeq}
\end{equation}
where $\delta S/\delta\Phi=\partial_\mu(a^2\partial^\mu\Phi)=a^4\square\Phi$. The variation with respect to $\varphi_-$
gives the complex conjugate equation. From Eqs.~\eqref{eq:s6:Mpp} and \eqref{eq:s6:Mpm}, in the definition sign,
\begin{equation}
M^2_{++}+M^2_{+-}=-\kappa^2a^4\Big(\Delta\alpha_1\square^2+\Delta\alpha_2H^2\square+\Delta\alpha_3H^2\frac{\nabla^2}{a^2}\Big)
\delta^4(x-x')+i\big[\Tab(y_{+-})-\Tab(y_{++})\big].
\label{eq:s6:retarded}
\end{equation}
For $\eta'>\eta$, $(\Delta\eta+i\delta)^2=(|\Delta\eta|-i\delta)^2$, so $y_{+-}=y_{++}$ and the nonlocal part vanishes.
For $\eta'<\eta$, $y_{+-}=y_{++}^*$ and the nonlocal kernel is $2\,\mathrm{Im}\,\Tab(y_{++})$, which is real. For spacelike
separation $y$ is real and positive for both polarities and the difference vanishes. The kernel is therefore real and
causal, with support on and inside the past light cone, as stated in \PII{} below \eqII{21}.

\subsection{Perturbative mode equation and the local term}
\label{sec:s6:mode}

With $M^2_{++}+M^2_{+-}=\sum_{\ell\ge1}\kappa^{2\ell}M^2_\ell$ and $\Phi(x;\vec k)=\sum_\ell\kappa^{2\ell}\Phi_\ell(\eta,k)
e^{i\vec k\cdot\vec x}$ (cf.\ \eqII{22}--\eqII{23}), the tree-order solution is (cf.\ \eqII{24}, \eqII{28})
\begin{equation}
\Phi_0=u(\eta,k)=\frac{H}{\sqrt{2k^3}}\Big(1-\frac{ik}{aH}\Big)e^{ik/(aH)}
=\frac{H}{\sqrt{2k^3}}\Big\{1+\frac{\khat^2}{2a^2}+\frac{i\khat^3}{3a^3}+\sixO\Big(\frac{\khat^4}{a^4}\Big)\Big\}.
\label{eq:s6:u}
\end{equation}
On a plane wave $a^4\square=-a^4[\partial_t^2+3H\partial_t+k^2/a^2]$, with $t$ comoving time. The first correction,
$\Phi=u+\kappa^2\Phi_1$, therefore obeys (cf.\ \eqII{29})
\begin{equation}
\Big[\partial_t^2+3H\partial_t+\frac{k^2}{a^2}\Big]\Phi_1=\sixR_1(\eta,k)\equiv
-\frac1{a^4}\int_{\eta_i}^0\!d\eta'\,u(\eta',k)\!\int\!d^3x'\,M^2_1(x;x')\,e^{i\vec k\cdot(\vec x\,'-\vec x)},
\label{eq:s6:mode}
\end{equation}
where $M^2_1$ is the one-loop part of Eq.~\eqref{eq:s6:retarded} with $\kappa^2$ stripped off.

\paragraph*{Late-time inversion (cf.\ \eqII{30}--\eqII{33}).}
With $\theta\equiv a\partial_a$, $\partial_t=H\theta$ and $P'\equiv dP/d\ln a$,
\begin{equation}
(\theta^2+3\theta)\big[a^nP(\ln a)\big]=a^n\big[(n^2+3n)P+(2n+3)P'+P''\big].
\label{eq:s6:theta}
\end{equation}
For $n=0$, $3P'+P''=R$ gives $P=\frac16\ln^2a-\frac19\ln a$ for $R=\ln a$ and $P=\frac13\ln a$ for $R=1$ (the leading terms of
\eqII{30}, \eqII{31}). For $n=-1$, $-2P+P'+P''=R$ gives $P=-\frac12\ln a-\frac14$ for $R=\ln a$ and $P=-\frac12$ for $R=1$
(\eqII{32}, \eqII{33}, which keep only the leading terms and the arbitrary constant). For $n=-2$,
$-2P-P'+P''=\alpha\ln a+\beta$ gives $P=-\frac\alpha2\ln a+\frac\alpha4-\frac\beta2$. The case $n=-3$ is resonant with
the homogeneous solution: $-3P'+P''=R$ gives $P=-\frac R3\ln a$ for constant $R$. As argued in \PII, Sec.~3.2, a
constant in $\Phi_1$ is a field-strength renormalization, and terms that fall off relative to $u$ cannot be separated from
corrections to the free Bunch--Davies initial state. Only contributions to $\sixR_1$ that grow or tend to a nonzero constant
can produce a significant late-time effect.

\paragraph*{Local part (cf.\ \eqII{34}).}
Inserting the local part of Eq.~\eqref{eq:s6:retarded} into Eq.~\eqref{eq:s6:mode} gives
$[\Delta\alpha_1\square^2+\Delta\alpha_2H^2\square+\Delta\alpha_3H^2\nabla^2/a^2]$ acting on $u(\eta,k)e^{i\vec k\cdot\vec x}$.
The d'Alembertian annihilates the tree-order solution, so
\begin{equation}
\sixR_{1,\rm local}=-\Delta\alpha_3\frac{H^2k^2}{a^2}\,u(\eta,k)\ \longrightarrow\ 0 .
\label{eq:s6:local}
\end{equation}
The corrected finite parts of $\beta_{2,3}$ only redefine $\Delta\alpha_{2,3}$ and leave no trace in the mode function.

\subsection{SK differences and the integrals of Tables 2--4}
\label{sec:s6:tab234}

\paragraph*{Identities (cf.\ \eqII{35}--\eqII{36}).}
At $D=4$, $\square f(y)=H^2[(4y-y^2)f''+4(2-y)f']$. For $f=\ln(y/4)$ this gives $\square f/H^2=4/y-3$, and for
$f=\frac12L^2-L$ with $L=\ln(y/4)$ it gives $\square f/H^2=\frac4yL-3L+2$. Hence
\begin{equation}
\frac4y=\frac{\square}{H^2}\ln\frac y4+3,\qquad
\frac4y\ln\frac y4=\frac{\square}{H^2}\Big\{\frac12\ln^2\frac y4-\ln\frac y4\Big\}+3\ln\frac y4-2,
\label{eq:s6:3536}
\end{equation}
in agreement with \PII\@. Since $\ln y$ and $\ln^2y$ are integrable, no $\delta^4$ remainder arises, and $\square_x$ commutes
with $\int d^4x'$ over the fixed region $[\eta_i,0]\times\mathbb R^3$.

\paragraph*{SK differences (cf.\ \eqII{39}--\eqII{41}).}
Inside the past light cone ($\Delta\eta>r$), $y_{++}=-|y|+i0$ and $y_{+-}=-|y|-i0$, so on the principal branch
$\ln y_{+\pm}=\ln|y|\pm i\pi$. For $\Delta f\equiv f(y_{++}/4)-f(y_{+-}/4)$ and $y_r\equiv aa'H^2(r^2-\Delta\eta^2)$,
\begin{equation}
\begin{aligned}
\Delta\ln x&=2\pi i\,\theta(\Delta\eta-r),&\qquad
\Delta\ln^2x&=4\pi i\,\theta(\Delta\eta-r)\ln|x|,\\
\Delta\big(x^j\ln x\big)&=2\pi i\,x^j\,\theta(\Delta\eta-r),&\qquad
\Delta\Big(\frac1x\Big)&=4\Big[\frac1{y_r+i0}-\frac1{y_r-i0}\Big]=-8\pi i\,\delta(y_r),
\end{aligned}
\label{eq:s6:Delta}
\end{equation}
with $x\to-\frac14aa'H^2(\Delta\eta^2-r^2)$ in the polynomial factors, \eqII{41}, and $|x|=X_0(1-z^2)$, $z\equiv r/\Delta\eta$ (in
this section $z$ denotes $r/\Delta\eta$, not the variable $4/y$ of Secs.~\ref{sec:s4b} and~\ref{sec:s5}).
Analytic functions (constants, $x$, $x^2$) are identical for both polarities, so $\Delta=0$ for them. Consequently the
$72\gamma_2$ correction ($-54\to+18$ in row~7), and the terms $-120$, $96x$, $12x$ and $-36x^2$, drop out of the retarded
kernel identically. The equivalent form of the $72\gamma_2$ correction as an extra $-12x$ in the $\delta$ row of \PI{}
Table~V drops out for the same reason.

\paragraph*{Master formulas.}
\PII{} lists in Tables~2--4 the integrals
\begin{equation}
S_K[f]\equiv-\frac{i}{16\pi^2}\int d^4x'\,a'^K\,\Delta f\qquad(H=1),\qquad K=3,4,5,\quad
f=\frac1x,\ \frac{\ln x}x,\ \ln x,\ \ln^2x .
\label{eq:s6:SK}
\end{equation}
The entries for $K=3,4$ were taken there from Ref.~\cite{ref:Prokopec2007}. We derived all of them independently. With
$d\eta'=da'/a'^2$, $\Delta\eta=1/a'-1/a$, $a(\eta_i)=1$, $\int d^3r\,\theta(\Delta\eta-r)=\frac{4\pi}{3}\Delta\eta^3$
and $B_j\equiv\int_0^1z^2(1-z^2)^jdz$ ($B_0=\frac13$, $B_1=\frac2{15}$, $B_2=\frac8{105}$),
Eq.~\eqref{eq:s6:Delta} gives
\begin{equation}
\begin{aligned}
S_K[\ln x]&=\frac16\int_1^a\!da'\,a'^{K-2}\Delta\eta^3,\qquad
S_K[x^j\ln x]=\frac{(-1)^jB_j}{2}\int_1^a\!da'\,a'^{K-2}X_0^j\,\Delta\eta^3,\\
S_K[\ln^2x]&=\int_1^a\!da'\,a'^{K-2}\Delta\eta^3\Big[\frac13\ln X_0+C_0\Big],\qquad
C_0\equiv\int_0^1\!z^2\ln(1-z^2)\,dz=\frac23\ln2-\frac89,\\
S_K\Big[\frac1x\Big]&=-\frac1a\int_1^a\!da'\,a'^{K-3}\Delta\eta .
\end{aligned}
\label{eq:s6:master}
\end{equation}
The last line uses $\int d^3r\,\delta\big(aa'(r^2-\Delta\eta^2)\big)=2\pi\Delta\eta/(aa')$. It evaluates the $1/x$ terms
directly on the light cone, independently of \eqII{35}; the results agree term by term with $\Dop\,S_K[\ln x]$, where
\begin{equation}
\Dop\equiv-\big(a^2\partial_a^2+4a\,\partial_a\big),\qquad \sixDk\equiv\Dop-\frac{\khat^2}{a^2},
\qquad \Dop a^n=-n(n+3)a^n,
\label{eq:s6:Dop}
\end{equation}
is $\square/H^2$ acting on a function of $\eta$ alone; $\sixDk$ is the operator of \eqII{37}.

With $a'=as$ one has $\Delta\eta=(1-s)/(as)$ and $X_0=(1-s)^2/(4s)$, independent of $a$, so that
\begin{equation}
\int_1^a\!da'\,a'^{Q-2}\Delta\eta^NX_0^jG(X_0)=a^{Q-1-N}\int_{1/a}^1\!ds\,s^{Q-2-N}(1-s)^NX_0^jG(X_0),\qquad
\ln X_0=2\ln(1-s)-\ln s-2\ln2 .
\label{eq:s6:scaling}
\end{equation}
Everything reduces to $\int_\epsilon^1s^pds$, $\int_\epsilon^1s^p\ln s\,ds$ and $\int_\epsilon^1s^p\ln(1-s)\,ds$
($\epsilon=1/a$). In the last one the terms $s^n/n$ of $-\ln(1-s)$ with $n\le q\equiv\max(0,-p-1)$ are kept as
divergent or logarithmic integrals, and the remainder leaves the constant
\begin{equation}
C_p=-\frac{H_{p+1}}{p+1}\ (p\ge0),\qquad C_{-1}=-\frac{\pi^2}{6},\qquad C_p=-\frac{H_q}{q}\ (p\le-2),
\label{eq:s6:Cp}
\end{equation}
with $H_n$ the harmonic numbers. The only source of $\pi^2$ in Tables~2--5 is $C_{-1}=-\pi^2/6$. Two examples:
\begin{equation}
\begin{aligned}
S_4[\ln x]&=\frac16\Big[\ln a-3\frac{a-1}a+3\frac{a^2-1}{2a^2}-\frac{a^3-1}{3a^3}\Big]
=\frac16\ln a-\frac{11}{36}+\frac1{2a}-\frac1{4a^2}+\frac1{18a^3},\\
S_4[\ln^2x]&=\frac23\Big(-\frac{\pi^2}6+\frac{49}{36}\Big)-\frac13\Big(-\frac{\ln^2a}2+\frac{85}{36}\Big)
-\frac89\Big(\ln a-\frac{11}6\Big)+\sixO\Big(\frac{\ln a}{a}\Big)\\
&=\frac16\ln^2a-\frac89\ln a+\frac74-\frac{\pi^2}9+\sixO\Big(\frac{\ln a}{a}\Big),
\end{aligned}
\label{eq:s6:examples}
\end{equation}
where in the second line $\frac13\ln X_0+C_0=\frac23\ln(1-s)-\frac13\ln s-\frac89$ (the $\ln2$ cancels) and the three
brackets are the $\ln(1-s)$, $\ln s$ and constant integrals of $(1-s)^3/s$.

\begin{table}[htbp]
\caption{Tables~2--4 of \PII: $S_K[f]$ of Eq.~\eqref{eq:s6:SK} ($H=1$). The third column is both the printed entry and
our result, which agree exactly through the order kept in \PII; the last column is the next order, which \PII{} leaves
as $\sixO(\cdot)$.}
\label{tab:s6:tab234}
\renewcommand{\arraystretch}{1.35}
\begin{ruledtabular}
\begin{tabular}{c c l l}
$K$ & $f$ & Printed $=$ derived & Next order (this work)\\
\hline
3 & $\frac1x$ & $-\frac{\ln a}a+\frac1a$ & $-\frac1{a^2}$\\
3 & $\frac{\ln x}x$ & $-\frac{\ln^2a}{2a}+\frac{2\ln a}a-\frac3a+\frac{\pi^2}{3a}$ & $-\frac{\ln a+1}{a^2}$\\
3 & $\ln x$ & $\frac16-\frac{\ln a}{2a}+\frac1{4a}$ & $-\frac1{2a^2}$\\
3 & $\ln^2x$ & $\frac13\ln a-\frac{11}9-\frac{\ln^2a}{2a}+\frac{2\ln a}a-\frac9{4a}+\frac{\pi^2}{3a}$ & $-\frac{\ln a}{a^2}$\\
\hline
4 & $\frac1x$ & $-\frac12$ & $\frac1a$\\
4 & $\frac{\ln x}x$ & $\frac34$ & $\frac{\ln a-1}a$\\
4 & $\ln x$ & $\frac16\ln a-\frac{11}{36}$ & $\frac1{2a}$\\
4 & $\ln^2x$ & $\frac16\ln^2a-\frac89\ln a+\frac74-\frac{\pi^2}9$ & $\frac{\ln a-1}a$\\
\hline
5 & $\frac1x$ & $-\frac16a$ & $\frac1{2a}$\\
5 & $\frac{\ln x}x$ & $\frac{17}{36}a$ & $\frac{2\ln a-3}{4a}$\\
5 & $\ln x$ & $\frac1{24}a-\frac16$ & $\frac1{4a}$\\
5 & $\ln^2x$ & $-\frac{13}{144}a-\frac13\ln a+\frac59$ & $\frac{2\ln a-3}{4a}$\\
\end{tabular}
\end{ruledtabular}
\end{table}

\paragraph*{Result.}
Equations~\eqref{eq:s6:master}--\eqref{eq:s6:Cp} were implemented in exact rational arithmetic
(\texttt{engine.py}, \texttt{tables234.py}~\cite{ref:repo}). All 12 entries of Tables~2--4 agree exactly with \PII{}
through the order kept there (our Table~\ref{tab:s6:tab234}). As an independent test, seven integrals of the type
\eqref{eq:s6:master} (including $\ln^2x$, $x\ln x$, $x^2\ln x$, an $r^2$ weight and a shifted power of $a'$) were evaluated
by direct 30-digit quadrature at $a=50$ and $a=400$ (\texttt{numcheck.py}); the largest difference from the exact
expansion is $1.0\times10^{-19}$.

\subsection{Table 5 within the simplification (II.44), definition sign}
\label{sec:s6:tab5}

\PII{} replaces $u(\eta',k)\to u(0,k)$ and $e^{i\vec k\cdot\Delta\vec x}\to1$ in Eq.~\eqref{eq:s6:mode}, which is the
simplification \eqII{44}. Rows containing $\nabla^2$ then vanish, because $\int d^3x'\,F(\|\vec x-\vec x\,'\|)$ does
not depend on $\vec x$. For a row term $c\,a^pa'^q(\square/H^2)^n f$ of our Table~\ref{tab:s5:TX}, using
$(4\pi)^4=(16\pi^2)^2$ and Eq.~\eqref{eq:s6:retarded} in the Table-X convention,
\begin{equation}
\sixR_1\ni-\frac{u(0,k)}{a^4}\,\frac{ic}{(4\pi)^4}\,a^p\square^n\!\int\!d^4x'\,a'^q\,\Delta f
=\sixunit\times\Big\{c\,a^{p-4}\,\Dop^nS_q[f]\Big\}\qquad(\text{Table-X convention}),
\label{eq:s6:rule}
\end{equation}
and in the definition sign the entry is $-c\,a^{p-4}\Dop^nS_q[f]$. The factor $(aa')^3(a^2+a'^2)$ is the sum of
$(p,q)=(5,3)$ and $(3,5)$; since $p+q=8$, $a^{p-4}=a^{4-q}$, which is the factor $a^{4-K}$ of \eqII{45}. Three examples
(Table-X convention), using $\Dop\ln^2a=-2-6\ln a$, $\Dop\ln a=-3$ and Eq.~\eqref{eq:s6:examples}:
\begin{equation}
\begin{aligned}
(aa')^4\square^2\times(-6\ln^2x):&\quad -6\,\Dop^2S_4[\ln^2x]=-6\,\Dop\Big(\frac73-\ln a\Big)=-18,\\
(aa')^4H^4\times(-108\ln x):&\quad -108\Big(\frac{\ln a}6-\frac{11}{36}\Big)=33-18\ln a,\\
(aa')^3(a^2+a'^2)H^4\times\Big(-\frac{32}{3x}\Big):&\quad
-\frac{32}3\Big[a\,S_3\Big[\frac1x\Big]+\frac1aS_5\Big[\frac1x\Big]\Big]=-\frac{80}9+\frac{32}3\ln a .
\end{aligned}
\label{eq:s6:t5ex}
\end{equation}
Since $\Dop$ preserves powers of $a$, truncating before or after acting with $\Dop$ is consistent.

\begin{table}[htbp]
\caption{Table~5 of \PII{} (Table-X convention, as printed) and the same quantity in the definition sign (this work),
in units of $H^4u(0,k)/16\pi^2$; non-falling part, within \eqII{44}. Rows of our Table~\ref{tab:s5:TX} not shown
contribute zero.}
\label{tab:s6:tab5}
\renewcommand{\arraystretch}{1.3}
\begin{ruledtabular}
\begin{tabular}{l l l}
External operator $\times f(x)$ & \PII{} (Table-X convention) & This work (definition sign)\\
\hline
$(aa')^4\square^3/H^2\times\big(-\frac{\ln x}{3x}\big)$ & $0$ & $0$\\
$(aa')^4\square^2\times\frac{26\ln x}{3x}$ & $0$ & $0$\\
$(aa')^4\square^2\times\frac{38}{3x}$ & $0$ & $0$\\
$(aa')^4\square^2\times(-6\ln^2x)$ & $-18$ & $18$\\
$(aa')^4\square^2\times(-18\ln x)$ & $0$ & $0$\\
$(aa')^4H^2\square\times\big(-\frac{6\ln x}{x}\big)$ & $0$ & $0$\\
$(aa')^4H^2\square\times\frac4x$ & $0$ & $0$\\
$(aa')^4H^2\square\times(-4\ln x)$ & $2$ & $-2$\\
$(aa')^4H^4\times\frac{4\ln x}{x}$ & $3$ & $-3$\\
$(aa')^4H^4\times\frac{18}x$ & $-9$ & $9$\\
$(aa')^4H^4\times(-108\ln x)$ & $33-18\ln a$ & $-33+18\ln a$\\
$(aa')^3(a^2{+}a'^2)\square^3/H^2\times\frac{\ln x}{6x}$ & $\frac{-325+12\pi^2}{27}+\frac{14}3\ln a-\frac23\ln^2a$ &
$\frac{325-12\pi^2}{27}-\frac{14}3\ln a+\frac23\ln^2a$\\
$(aa')^3(a^2{+}a'^2)\square^2\times\big(-\frac{\ln x}{3x}\big)$ & $\frac{85-12\pi^2}{27}-4\ln a+\frac23\ln^2a$ &
$\frac{-85+12\pi^2}{27}+4\ln a-\frac23\ln^2a$\\
$(aa')^3(a^2{+}a'^2)\square^2\times\frac1{6x}$ & $\frac89-\frac23\ln a$ & $-\frac89+\frac23\ln a$\\
$(aa')^3(a^2{+}a'^2)H^2\square\times\big(-\frac{2\ln x}{3x}\big)$ & $\frac{160-12\pi^2}{27}-\frac{10}3\ln a+\frac23\ln^2a$ &
$\frac{-160+12\pi^2}{27}+\frac{10}3\ln a-\frac23\ln^2a$\\
$(aa')^3(a^2{+}a'^2)H^2\square\times\frac5x$ & $\frac{55}3-10\ln a$ & $-\frac{55}3+10\ln a$\\
$(aa')^3(a^2{+}a'^2)H^2\square\times(-18\ln x)$ & $-15+18\ln a$ & $15-18\ln a$\\
$(aa')^3(a^2{+}a'^2)H^4\times\frac{4\ln x}{3x}$ & $\frac{-91+12\pi^2}{27}+\frac83\ln a-\frac23\ln^2a$ &
$\frac{91-12\pi^2}{27}-\frac83\ln a+\frac23\ln^2a$\\
$(aa')^3(a^2{+}a'^2)H^4\times\big(-\frac{32}{3x}\big)$ & $-\frac{80}9+\frac{32}3\ln a$ & $\frac{80}9-\frac{32}3\ln a$\\
$(aa')^3(a^2{+}a'^2)H^4\times18$ (corrected constant) & (not listed) & $0$ (analytic)\\
rows 8--12 (with $\nabla^2$, including $\zeta$) & (not listed) & $0$ (under \eqII{44})\\
\hline
Total & $0$ & $0$\\
\end{tabular}
\end{ruledtabular}
\end{table}

The script \texttt{table5.py}~\cite{ref:repo} reproduces all 19 printed entries exactly in the Table-X convention. In the
definition sign every entry is $-1$ times the printed one (our Table~\ref{tab:s6:tab5}), and the total is again zero.
Individual entries grow as fast as $\frac23\ln^2a$; the $\ln^2a$, $\pi^2$, $\ln a$ and constant parts cancel separately:
\begin{equation}
\begin{aligned}
\ln^2a:&\quad\tfrac23(1-1-1+1)=0,\qquad \pi^2:\quad\tfrac{12}{27}(-1+1+1-1)=0,\\
\ln a:&\quad 18-\tfrac{14}3+4+\tfrac23+\tfrac{10}3+10-18-\tfrac83-\tfrac{32}3=0,\\
1:&\quad 18-2-3+9-33+\tfrac{325-85-24-160-495+405+91+240}{27}=-11+11=0
\end{aligned}
\label{eq:s6:t5sums}
\end{equation}
(definition sign). Neither the $\zeta$ correction (a $\nabla^4$ row) nor the $72\gamma_2$ correction (analytic) enters
our Table~\ref{tab:s6:tab5}. A heuristic reason to expect the $k=0$ cancellation, which we state as an assumption, is that
the scalar couples to gravity only through derivatives, so that diagrammatically $M^2(x;x')=\partial_\mu\partial'_\nu
K^{\mu\nu}(x;x')$ and $\int d^4x'\,M^2_{\rm ret}(x;x')\cdot1$ leaves only terms from the initial surface $\eta'=\eta_i$,
which decay. The vanishing total is then a consistency test of the rows without $\nabla^2$; it does not test the
$\nabla^2$ rows, hence not the $\zeta$ row.

\subsection{Beyond (II.44): the full \texorpdfstring{$k$}{k} dependence}
\label{sec:s6:beyond}

\paragraph*{The justification of \eqII{44} fails term by term (S4).}
\PII{} argues that the deviations $u(\eta',k)-u(0,k)$ and $e^{i\vec k\cdot\Delta\vec x}-1$ bring at least two factors of
$1/a'$ or $1/a$, so that they, and with them all $\nabla^2$ rows, only produce falling terms. Two effects spoil this
argument for individual rows: positive powers $x^j\propto(aa')^j$ and the prefactors $a^p$ multiply the contribution of
the early-time region by powers of $a$; and at early $\eta'$, $\ln|x|\simeq\ln[a/(4a')]$ contains $\ln a$. For example,
the $(p,q)=(4,2)$ parts of rows~10 and~11 contain $x\ln x$, and from Eq.~\eqref{eq:s6:master} with $X_0=\frac14aa'\Delta\eta^2$,
\begin{equation}
S_2[x\ln x]=-\frac12\cdot\frac2{15}\cdot\frac a4\int_1^a\!da'\,a'\Big(\frac1{a'}-\frac1a\Big)^5
=-\frac a{60}\Big(\frac13-\frac5{2a}+\ldots\Big)=-\frac a{180}+\frac1{24}+\ldots
\label{eq:s6:S2xlnx}
\end{equation}
With $\nabla^2\to-\khat^2$, $u\to u(0,k)$ and $\Dop a=-4a$, in the Table-X convention,
\begin{equation}
\text{row 10:}\ \ 48(-\khat^2)\Big(-\frac a{180}\Big)=+\frac4{15}\khat^2a,\qquad
\text{row 11:}\ \ 12(-\khat^2)\,\Dop\Big(-\frac a{180}\Big)=-\frac4{15}\khat^2a .
\label{eq:s6:row1011}
\end{equation}
In the definition sign rows~10 and~11 thus grow like $\mp\frac4{15}\khat^2a$. Either one alone would give
$\Phi_1\propto a$; they cancel only against each other. Likewise the $k^2$ corrections of the rows without $\nabla^2$
give $\sixO(1)$ constants. A complete expansion in $k$ is therefore required.

\paragraph*{Systematic expansion in $k$.}
We keep the exact $u(\eta',k)$, the phase, the $k^2/a^2$ part of $\square$ and all $\nabla^2$ rows. With $w\equiv\khat/a'$,
\begin{equation}
\frac{u(\eta',k)}{u(0,k)}=(1-iw)e^{iw}=1+\frac{w^2}2+\frac{iw^3}3-\frac{w^4}8+\ldots,\qquad
\big\langle e^{i\vec k\cdot\vec r}\big\rangle_{\rm angles}=\frac{\sin kr}{kr}=\sum_{m\ge0}\frac{(-1)^m(kr)^{2m}}{(2m+1)!},
\label{eq:s6:kser}
\end{equation}
together with $\square/H^2\to\sixDk$, $\nabla^2\to-k^2$, $1/x\to\sixDk\ln x$ and
$\ln x/x\to\sixDk(\frac12\ln^2x-\ln x)+3\ln x$ (Eq.~\eqref{eq:s6:3536}; the constants drop out of $\Delta$). Each term
reduces to integrals of the type \eqref{eq:s6:master}, with the power of $a'$ shifted by $-n$ for the $w^n$ term and an
extra weight $r^{2m}$ (integrals of $z^{2+2m}$). The integrand is an entire function of $k$ and the region is compact
($\eta'\in[\eta_i,\eta]$, $r\le\Delta\eta$), so the expansion converges at every fixed $a$ and for every $k$. Convergence at fixed
$a$ does not, however, control the behavior of the remainder as $a\to\infty$: individual rows contain terms growing like $a$ at every
even order computed (for example $\pm\frac{8a}{441}\khat^6$ in rows~10 and~11, see below), and for $\khat\gtrsim1$, i.e.\ for modes inside the
horizon at $\eta_i$, a truncated series need not approximate the source. The script \texttt{kexp.py}~\cite{ref:repo} expands every
row exactly through $k^{10}$, and the statements below hold order by order through $k^{10}$.

As a hand check, consider row~3 at order $\khat^2$ (Table-X convention). The $\khat^2$ part of $S_4$ comes from
$\frac{w^2}2$ and from $-\frac{k^2r^2}6$ (radial weight $\frac15$ instead of $\frac13$). At leading order
$\int_1^a da'\,\Delta\eta^3\to\frac12$ and $\int_1^a da'\,a'^2\Delta\eta^5\to\frac12$, so
\begin{equation}
G^{(2)}_4[\ln x]=\frac1{12}\int\!da'\,\Delta\eta^3-\frac1{60}\int\!da'\,a'^2\Delta\eta^5\to\frac1{24}-\frac1{120}=\frac1{30},
\qquad G^{(2)}_4[\ln^2x]\ni\frac{\ln a}{15},
\label{eq:s6:G2}
\end{equation}
where the $\ln a$ comes from $\ln X_0\simeq\ln a$ at early times. Then $-108\ln x\to-\frac{18}5$,
$\frac{4\ln x}x\to4[\Dop(\frac{\ln a}{30})+\frac3{30}]=0$ and $\frac{18}x\to0$: row~3 gives $-\frac{18}5$ at order
$\khat^2$ ($+\frac{18}5$ in the definition sign), as does the code. The early-time $\ln a$ in $\ln|x|$ is essential here.

\begin{table}[htbp]
\caption{Non-falling part of $\sixR_1$ beyond \eqII{44}, row by row (rows of our Table~\ref{tab:s5:TX}), in the
definition sign and in units of $H^4u(0,k)/16\pi^2$. The $\khat^0$ column collects, row by row, the last column of our Table~\ref{tab:s6:tab5}; the order $\khat^1$
vanishes in every row. Rows 0 and 9 give zero at these orders.}
\label{tab:s6:t5pp}
\renewcommand{\arraystretch}{1.4}
\begin{ruledtabular}
\begin{tabular}{r l l l l}
\# & $\khat^0$ & $\khat^2$ & $\khat^3$ & $\khat^4$\\
\hline
1 & $18$ & $0$ & $0$ & $0$\\
2 & $-2$ & $0$ & $0$ & $0$\\
3 & $18\ln a-27$ & $\frac{18}5$ & $2i$ & $-\frac{27}{35}$\\
4 & $\frac23\ln^2a-\frac{14}3\ln a+\frac{325-12\pi^2}{27}$ & $\frac29\ln a-\frac{22}{27}$ & $\frac{2i(6\ln a-23)}{81}$ & $\frac{239}{900}-\frac{\ln a}{15}$\\
5 & $-\frac23\ln^2a+\frac{14}3\ln a+\frac{12\pi^2-109}{27}$ & $\frac{22}{27}-\frac29\ln a$ & $\frac{2i(23-6\ln a)}{81}$ & $\frac{\ln a}{15}-\frac{239}{900}$\\
6 & $-\frac23\ln^2a-\frac{14}3\ln a+\frac{12\pi^2-250}{27}$ & $-\frac29\ln a-\frac{20}{27}$ & $-\frac{2i(6\ln a+19)}{81}$ & $\frac{\ln a}{15}+\frac{181}{900}$\\
7 & $\frac23\ln^2a-\frac{40}3\ln a+\frac{331-12\pi^2}{27}$ & $\frac29\ln a-\frac{61}{27}$ & $\frac{4i(3\ln a-31)}{81}$ & $\frac{629}{900}-\frac{\ln a}{15}$\\
8 & $0$ & $\frac25$ & $0$ & $\frac3{35}$\\
10 & $0$ & $1-\frac{4a}{15}$ & $0$ & $\frac15-\frac{12a}{175}$\\
11 & $0$ & $\frac{4a}{15}-2$ & $0$ & $\frac{12a}{175}-\frac25$\\
12 & $0$ & $0$ & $0$ & $-\frac1{70}$\\
\hline
Total & $0$ & $0$ & $0$ & $0$\\
\end{tabular}
\end{ruledtabular}
\end{table}

\paragraph*{Cancellation order by order.}
Our Table~\ref{tab:s6:t5pp} lists the non-falling part of each row through $\khat^4$ (\texttt{kexp.py}, \texttt{extract.py}).
At order $\khat^2$ the terms $\propto a$ cancel between rows~10 and~11, the $\ln a$ terms among rows~4--7, and the
constants add to $\frac{18}5+\frac25-\frac{81}{27}+1-2=0$. At order $\khat^4$ rows~3 and~8 give $-\frac{24}{35}$, rows~4--7
give $+\frac9{10}$, rows~10--11 give $-\frac15$ and row~12 gives $-\frac1{70}$, with total $\frac{-48+63-14-1}{70}=0$.
Without the $\zeta$ row the $\khat^4$ column would not cancel. Its $-\frac1{70}$ comes from the term $24x^2\ln x$, which
is the same in the printed and in the corrected row, so the $\khat^4$ cancellation does not depend on the $\zeta$ correction: with
$S_2[x^2\ln x]=\frac12\cdot\frac8{105}\cdot\frac{a^2}{16}\int da'\,a'^2\Delta\eta^7$ and
$\int_1^a da'\,a'^2\Delta\eta^7\to\frac14$, one finds $24\,a^{-2}\khat^4\,S_2[x^2\ln x]\to\frac{\khat^4}{70}$ in the
Table-X convention. The code finds that the non-falling total vanishes separately at every order
$k^0,k^1,\dots,k^{10}$, for the corrected table and for the printed one. At orders $k^6$ to $k^{10}$ the terms
$\propto a$ of rows~10 and~11 (e.g.\ $\pm\frac{8a}{441}\khat^6$) again cancel each other. The series truncated at $k^8$ was
compared with direct quadrature using the exact $u(\eta',k)$ and $\sin kr/kr$ at $a=200$ (\texttt{numcheck\_k.py}): the
differences are $10^{-10}$--$2\times10^{-8}$ at $\khat=0.3$ and $5\times10^{-6}$--$9.3\times10^{-4}$ at $\khat=1$, the size
of the expected $k^9$ truncation error. The conclusion of \eqII{44} therefore holds through $k^{10}$, although its
justification does not hold row by row.

\paragraph*{Contribution of the $\zeta$ correction.}
The correction is $\delta\Tab=(aa')^2\nabla^4(8\ln x-16x\ln x)$. By hand, in the Table-X convention, with
$S_2[\ln x]=\frac16(\frac12-\frac3a+\ldots)$ and Eq.~\eqref{eq:s6:S2xlnx},
$8a^{-2}\khat^4S_2[\ln x]=\frac23\khat^4/a^2+\ldots$ and
$-16a^{-2}\khat^4S_2[x\ln x]=\frac4{45}\khat^4/a-\frac23\khat^4/a^2+\ldots$; the exact result (\texttt{analysis2.py},
\texttt{analysis3.py}) is
\begin{equation}
\delta\sixR_1=\sixunit\Big[\khat^4\Big(\frac4{45a}-\frac4{3a^3}\Big)+\khat^6\Big(\frac4{175a}-\frac2{15a^3}\Big)+\ldots\Big]
\qquad(\text{Table-X convention}).
\label{eq:s6:zeta}
\end{equation}
Only falling terms appear. This also holds for every $k$, without the $k$ expansion: since
$|u(\eta',k)/u(0,k)|=\sqrt{1+k^2\eta'^2}\le\sqrt{1+\khat^2}$, $|e^{i\vec k\cdot\vec r}|=1$ and $|x|\le\frac14aa'\Delta\eta^2$,
\begin{equation}
\big|\delta\sixR_{1,\ln x}\big|\le C\frac{\khat^4}{a^2}\int_1^a\!da'\Big(\frac1{a'}-\frac1a\Big)^3\le C'\frac{\khat^4}{a^2},\qquad
\big|\delta\sixR_{1,x\ln x}\big|\le C\frac{\khat^4}{a}\int_1^a\!da'\,a'\Big(\frac1{a'}-\frac1a\Big)^5\le C'\frac{\khat^4}{a},
\label{eq:s6:bound}
\end{equation}
where the constants contain the factor $\sqrt{1+\khat^2}$ and the unit $H^4u(0,k)/16\pi^2$. The $\zeta$ correction
therefore falls off at late times for every mode, independently of the $k$ expansion.

\subsection{Late-time source and the one-loop mode function}
\label{sec:s6:late}

Summing all rows of the corrected table (definition sign, through $k^{10}$; \texttt{extract.py}) gives
\begin{equation}
\sixR_1=\sixunit\bigg\{\frac{1}{a^2}\Big[12-\Big(8\ln a+\frac{103}3\Big)\khat^2-8i\khat^3+\frac83\khat^4
+\frac{32i}{45}\khat^5-\frac2{15}\khat^6+\ldots\Big]-\frac{12i\khat^3}{a^3}+\sixO\Big(\frac{\ln a}{a^4}\Big)\bigg\}
+\sixR_{1,\rm local}.
\label{eq:s6:Rlate}
\end{equation}
There is no $a^{-1}$ term. The next $a^{-2}$ coefficients are $-\frac{16i}{1575}\khat^7-\frac4{945}\khat^8
-\frac{16i}{6615}\khat^9+\frac{11}{14175}\khat^{10}$, and there is no other $a^{-3}$ term through $k^{10}$. With the
printed $\zeta$ row (definition sign) $\sixR_1$ would contain the additional terms
\begin{equation}
\sixunit\,\frac1a\Big[\frac4{45}\khat^4+\frac4{175}\khat^6+\frac{2i}{135}\khat^7-\frac8{1323}\khat^8
-\frac{i}{525}\khat^9+\frac4{8019}\khat^{10}+\ldots\Big],
\label{eq:s6:Rprinted}
\end{equation}
plus $a^{-3}$ terms. The correction thus changes the fall-off of the right-hand side from $\khat^4/a$ to $\ln a/a^2$;
we record this as an observation, not as an independent proof of the correction. Both behaviors are consistent with the
statement of \PII{} after Table~5 that the right-hand side falls at least like $1/a$ times powers of $\ln a$. The $a^{-2}$
coefficients depend on the initial time and on $\khat$ (e.g.\ the oscillatory $\khat^3,\khat^5$ terms), i.e.\ they are
sensitive to the initial state and, in the sense of \PII, Sec.~3.2, cannot be separated from initial-state corrections.

\paragraph*{Solving for $\Phi_1$.}
In Eq.~\eqref{eq:s6:theta} the $a^{-2}$ terms of Eq.~\eqref{eq:s6:Rlate} have $\alpha=-8\khat^2$ and
$\beta=12-\frac{103}3\khat^2-8i\khat^3+\ldots$, so $P=4\khat^2\ln a-6+\frac{91}6\khat^2+4i\khat^3+\ldots$; the resonant
$a^{-3}$ term gives $P=4i\khat^3\ln a$. The $k^2/a^2$ term on the left-hand side changes the $a^{-2}$ solution only at
order $a^{-4}$. Dividing by $H^2$ (from $\partial_t^2+3H\partial_t=H^2(\theta^2+3\theta)$) we obtain
\begin{equation}
\Phi_1=C\,u(\eta,k)+\frac{H^2u(0,k)}{16\pi^2}\bigg\{\frac{1}{a^2}\Big[4\khat^2\ln a-6+\frac{91}6\khat^2+4i\khat^3+\ldots\Big]
+\frac{4i\khat^3\ln a}{a^3}+\ldots\bigg\},
\label{eq:s6:Phi1}
\end{equation}
where the local source \eqref{eq:s6:local}, itself of order $\khat^2/a^2$, only adds a $\Delta\alpha_3$-dependent
$\khat^2/a^2$ term inside the braces. Here $C\,u(\eta,k)$ stands for the homogeneous part, which the
late-time analysis does not fix; we write $C\equiv\frac{H^2}{16\pi^2}\tilde C$. In general the initial conditions at $\eta_i$ fix a
combination $c_1u+c_2u^*$ of both homogeneous solutions; since $u(0,k)=H/\sqrt{2k^3}$ is real, this combination tends to
$(c_1+c_2)\,u(0,k)[1+\sixO(a^{-2})]$ and is absorbed into $\tilde C$. Therefore
\begin{equation}
\boxed{\;\Phi(\eta,k)=u(\eta,k)+\kappa^2\Phi_1(\eta,k)\ \xrightarrow{\ a\to\infty\ }\
u(0,k)\Big\{1+\frac{\kappa^2H^2}{16\pi^2}\Big[\tilde C+\sixO\Big(\frac{\ln a}{a^2}\Big)\Big]\Big\}\;}
\label{eq:s6:final}
\end{equation}
where $\tilde C$ is a field-strength renormalization constant, and the tree-order approach of $u(\eta,k)$ to $u(0,k)$,
Eq.~\eqref{eq:s6:u}, is not displayed. Through order $k^{10}$, and given the kernel and the local terms as derived in
Secs.~\ref{sec:s3}--\ref{sec:s5}, there is no growing ($a$, $\ln a$, $\ln^2a$) or drifting term: one-loop quantum gravity gives
no significant correction to the mode function of the massless minimally coupled scalar at late times. This confirms the
conclusion of \PII{} with the corrected self-mass, in the definition sign, and without the simplification \eqII{44}. Two
qualifications are stated at the end of this section: the cancellation is not proven beyond $k^{10}$, and a possible finite local
term from the outer-operator reduction of \PI{} has not been excluded.

\begin{table}[htbp]
\caption{Effect of the corrections of \PI{} on the SK kernel and on the late-time mode function.}
\label{tab:s6:effect}
\renewcommand{\arraystretch}{1.25}
\begin{ruledtabular}
\begin{tabular}{>{\raggedright\arraybackslash}p{0.27\textwidth} >{\raggedright\arraybackslash}p{0.30\textwidth} >{\raggedright\arraybackslash}p{0.36\textwidth}}
Correction & SK kernel / Table~5 & Late-time mode function\\
\hline
S1: overall sign (definition $=-$printed) & every Table~5 entry flips; total still $0$ &
the particular solution $\Phi_1$ changes sign; since the non-falling total is zero, only falling terms change sign; no late-time effect\\
S2: $\zeta$ row, $+8\ln x-16x\ln x$ ($\nabla^4$, non-analytic) & absent under \eqII{44}; in the full calculation only falling
terms & none ($\le C\khat^4/a$ for every $k$); removes the $\khat^{4},\khat^6,\ldots$ tails $\propto1/a$\\
S3: $72\gamma_2$, $-54\to+18$ (analytic) & $\Delta f\equiv0$, drops out of the retarded kernel & identically none\\
B1--B6: finite parts of $\beta_2,\beta_3$ and $\alpha_2,\alpha_3$ & absorbed into $\Delta\alpha_{2,3}$ & none; $\sixR_{1,\rm local}=-\Delta\alpha_3H^2k^2u/a^2\to0$\\
C19: row printed twice in Table~X & already single in \PII{} Table~1 & none\\
\end{tabular}
\end{ruledtabular}
\end{table}

\begin{table}[htbp]
\caption{Statements of \PII{} compared with this work.}
\label{tab:s6:comp}
\renewcommand{\arraystretch}{1.25}
\begin{ruledtabular}
\begin{tabular}{>{\raggedright\arraybackslash}p{0.34\textwidth} >{\raggedright\arraybackslash}p{0.60\textwidth}}
\PII & This work\\
\hline
\eqII{7}--\eqII{9}: $\alpha_2\ni\frac{61}3$, $\alpha_3\ni\frac4{D-4}-\frac{58}3-2\gamma_{\rm E}$ &
$\frac{49}3$ and $\frac4{D-4}-\frac{52}3-2\gamma_{\rm E}$ (Table-X convention; B5, B6); pole unchanged; finite parts absorbed\\
Table~1 & two entries corrected (rows 7 and 12, our Table~\ref{tab:s5:TX}; S3 and S2)\\
\eqII{15}--\eqII{21}: SK rules, real and causal kernel & derived; Eqs.~\eqref{eq:s6:Mpm}--\eqref{eq:s6:retarded} in the
definition sign\\
\eqII{34}: only $\Delta\alpha_3$ contributes locally & confirmed ($H^2$ restored, C21)\\
\eqII{35}--\eqII{41} & confirmed; $1/x$ also evaluated directly on the light cone\\
Tables~2--4 (12 entries) & identical through the order kept; next order obtained\\
Table~5 (19 entries), total $0$ & identical in the Table-X convention; $-1$ times each entry in the definition sign;
total $0$\\
\eqII{44} and its justification & S4: justification fails term by term (rows 10, 11 grow like $\mp\frac4{15}\khat^2a$, definition sign);
the conclusion holds, verified through $k^{10}$\\
Right-hand side falls at least like $\ln^na/a$ & corrected table: $\ln a/a^2$; printed table: $\khat^4/a$\\
No significant one-loop late-time correction & confirmed, Eq.~\eqref{eq:s6:final}\\
\end{tabular}
\end{ruledtabular}
\end{table}

\paragraph*{Status and limitations.}
Our Tables~\ref{tab:s6:effect} and~\ref{tab:s6:comp} summarize the effect of the corrections of \PI{} and the comparison
with \PII\@. The cancellation of the non-falling part beyond \eqII{44} is exact through $k^{10}$ but has not been proven
to all orders in $k$. The convergence of the $k$ series at fixed $a$ gives no control over the late-time behavior of the terms
beyond $k^{10}$, since individual rows carry growing terms $\propto a$ at every even order computed; in particular, for $\khat\gtrsim1$ the
truncated series does not establish the late-time form of the source. The $a^{-2}$ coefficients
of $\sixR_1$ and $\Phi_1$ are sensitive to the initial state. The explanation of the $k=0$ cancellation by the
derivative structure of $M^2$ is a heuristic. The operator ordering of our Table~\ref{tab:s5:TX} is assumed; it is the
one verified by the pointwise test of Table~X\@. The direct light-cone evaluation of $S_K[1/x]$ was done by hand only;
the code uses \eqII{35} and agrees with it. Finally, this section takes the renormalized self-mass of
Secs.~\ref{sec:s3}--\ref{sec:s5} as given, together with its stated limitations: in particular, the possible extra local term from the operator
reduction in the appendix of \PI{} is only partially resolved (a double counting would have counterterm form, but other
local differences at coincidence have not been excluded by a full distributional analysis), and the global checks of
the input table are numerical, at sample points. A finite local term of mass type, such as $H^4a^4\delta^4(x-x')$ in $-iM^2$, would
give a source $\sixR_1$ that tends to a constant, and by \eqII{31} a secular $\Phi_1\propto\frac13\ln a$; this is exactly the kind
of effect excluded above, so the late-time conclusion holds given the local terms as derived. Heuristically, such a term would
conflict with the shift symmetry that underlies the $k=0$ cancellation, but this has not been demonstrated for the outer-operator
reduction.

\section{Summary of differences and conclusions}
\label{sec:s0s1s7s8:summary}

\subsection{Main differences with the printed papers}

Our Table~\ref{tab:s0s1s7s8:main} summarizes, item by item, how our results compare with Papers~I and~II and whether each difference
matters physically. The first row is the physical question that motivated both papers; its answer does not change.

\begin{table*}[tbp]
\caption{Main results of this work compared with the printed Papers~I~\cite{KW07} and~II~\cite{KW08}. Signs of table entries and of
$\beta_i$ are in the paper convention unless stated otherwise; the definition sign is $-1$ times the paper convention,
Eq.~\eqref{eq:s0s1s7s8:signs}. Table~X entries are coefficients of $\kappa^2H^4/(4\pi)^4$, $\beta_i$ are in units of
$H^{D-4}/(4\pi)^{D/2}$, and $x=y/4$. The class of each difference is given in parentheses in the first column: substantive findings
S1--S4, absorbable finite errors B1--B6, and typographical errors C1--C23.}
\label{tab:s0s1s7s8:main}
\renewcommand{\arraystretch}{1.3}
\small
\begin{ruledtabular}
\begin{tabular}{@{}>{\raggedright\arraybackslash}p{0.15\linewidth}>{\raggedright\arraybackslash}p{0.22\linewidth}%
>{\raggedright\arraybackslash}p{0.33\linewidth}>{\raggedright\arraybackslash}p{0.24\linewidth}@{}}
Item & Papers I/II (printed) & This work & Does it matter physically?\\
\hline
\textbf{Late-time scalar mode function} &
\textbf{No significant one-loop correction} at late times; the right-hand side of \eqII{29} falls at least like $\ln^na/a$
(Paper~II, Secs.~3.2 and~3.4). &
\textbf{Confirmed} with the corrected self-mass, in the definition sign and without the simplification \eqII{44}, through order
$k^{10}$:
$\displaystyle\Phi\to u(0,k)\Big\{1+\frac{\kappa^2H^2}{16\pi^2}\Big[\tilde C+O\Big(\frac{\ln a}{a^2}\Big)\Big]\Big\}$
(tree-order $O(\hat k^2/a^2)$ terms of $u(\eta,k)$ not displayed); limitations in Sec.~\ref{sec:s0s1s7s8:limitations}. &
\textbf{The main conclusion stands: no significant one-loop correction at late times.}\\
\hline
Overall sign of $-iM^2$ (S1, substantive) &
$-iM^2_{\rm ren}=$ local $+$ Table~X, \eqI{167}. &
The printed self-mass is $-1$ times the one defined by \eqI{168}. The 4-point and 3-point parts are both reversed, so their relative
sign is consistent. Physical nonlocal part: $-$(corrected Table~X). &
Every Table~5 entry flips and the particular solution $\Phi_1$ changes sign; since the non-falling total is zero, only decaying
terms change sign. No late-time effect.\\
\hline
Table~X, row $(aa')^2\nabla^4$ (S2, substantive) &
$\tfrac{10}{3}\ln x-24x\ln x+24x^2\ln x-36x^2$ &
$\tfrac{34}{3}\ln x-40x\ln x+24x^2\ln x-36x^2$: the Table~III term $8\ln x-16x\ln x$ was left out. Corrected Table~X:
our Table~\ref{tab:s5:TX}. &
Non-analytic: changes the nonlocal self-mass. In the mode equation it adds only falling terms ($\le C\hat k^4/a$); it removes the
$\hat k^4/a$, $\hat k^6/a,\dots$ tails, so the source falls like $\ln a/a^2$. The $k^4$ cancellation beyond \eqII{44} uses the
row's $24x^2\ln x$ term, which the correction leaves unchanged.\\
\hline
Table~X, row $(aa')^3(a^2+a'^2)H^4$ (S3, substantive, analytic) &
constant $-54$ &
constant $+18$ (equivalently, a term $-12x$ in the $\delta$ row of Table~V). Only with this and the previous correction does Table~X
pass the global pointwise test (30 digits). &
Analytic: drops out of the retarded kernel identically. No effect.\\
\hline
Finite local terms, $\beta_2$, $\beta_3$ (B1--B6, absorbable) &
$\beta_2=-\tfrac{61}{3}$, $\beta_3=-\tfrac{4}{D-4}+\tfrac{58}{3}+2\gamma_{\rm E}$ &
$\beta_2=-\tfrac{49}{3}$, $\beta_3=-\tfrac{4}{D-4}+\tfrac{52}{3}+2\gamma_{\rm E}$, from finite errors in \eqI{40} and \eqI{59}
(our Table~\ref{tab:s0s1s7s8:finite}). Pole unchanged; $\beta_1=0$ confirmed. &
No: absorbed into the arbitrary finite counterterms $\Delta\alpha_2$, $\Delta\alpha_3$.\\
\hline
Divergences and counterterms &
Two invariant counterterms and one noninvariant counterterm \eqI{145}; only $\beta_3$ diverges. &
Confirmed, including the de Sitter group analysis and the completeness of the noninvariant counterterms (checked at $D=4,5,6$);
the divergence structure holds within \PI's treatment of the outer operators (Sec.~\ref{sec:s4b:distr}). &
---\\
\hline
Paper~II, Tables~2--4 &
12 integrals &
All 12 reproduced exactly in closed form; next order obtained. &
---\\
\hline
Paper~II, Table~5 &
19 entries, total $0$ &
Reproduced exactly in the paper convention; in the definition sign every entry flips; total $0$. &
Confirms the cancellation at $k=0$.\\
\hline
Simplification \eqII{44} (S4, argument) &
Dropping $u(\eta')-u(0)$, $e^{i\vec k\cdot\Delta\vec x}-1$ and the $\nabla^2$ rows is justified by a fall-off argument. &
The justification fails term by term: two rows grow like $\mp\tfrac{4}{15}\hat k^2a$. The full non-falling part cancels at each order
$k^0,\dots,k^{10}$, for the corrected and for the printed table. &
No change of the conclusion.\\
\hline
Misprints (typographical) &
--- &
23 typographical errors C1--C23 (our Table~\ref{tab:s0s1s7s8:errata}). &
None.\\
\end{tabular}
\end{ruledtabular}
\end{table*}

The finite errors that shift $\beta_2$ and $\beta_3$, and their consequences B3--B6, are listed in our
Table~\ref{tab:s0s1s7s8:finite}. The printed \eqI{161}--\eqI{165}
are reproduced exactly from the printed \eqI{44} and \eqI{59}; the shifts come only from correcting these two inputs. In the
definition sign all $\beta_i$ change sign, while $\alpha_i=-\beta_i+\Delta\alpha_i$, \eqI{166}, holds in either convention.

\begin{table}[tb]
\caption{Finite errors of Paper~I that are absorbed into the counterterms (paper convention; $h\equiv H^2/16\pi^2$; $\beta_i$ and
$\alpha_i$ in units of $H^{D-4}/(4\pi)^{D/2}$). The corresponding counterterm coefficients of Paper~II, \eqII{8}--\eqII{9}, change
accordingly.}
\label{tab:s0s1s7s8:finite}
\renewcommand{\arraystretch}{1.3}
\small
\begin{ruledtabular}
\begin{tabular}{@{}l>{\raggedright\arraybackslash}p{0.19\linewidth}>{\raggedright\arraybackslash}p{0.21\linewidth}%
>{\raggedright\arraybackslash}p{0.21\linewidth}>{\raggedright\arraybackslash}p{0.27\linewidth}@{}}
& Location & Printed & Corrected & Effect\\
\hline
B1 & \eqI{40}, finite part & $+2\,\delta^0_\alpha\delta^\sigma_0\,h$ & $-2\,\delta^0_\alpha\delta^\sigma_0\,h$ &
the printed value contradicts \eqI{38} through a trace identity; $\beta_2\to\beta_2+4$, $\beta_3\to\beta_3-4$\\
B2 & \eqI{59}, constant of the $a^{D-2}\nabla^2$ term & $-\tfrac72$ & $-\tfrac52$, the sum of the paper's own \eqI{57} and \eqI{58} &
$\beta_3\to\beta_3+2$\\
B3 & $\beta_2$, \eqI{163} & $-\tfrac{61}{3}$ & $-\tfrac{49}{3}$ & finite; absorbed into $\Delta\alpha_2$\\
B4 & $\beta_3$, \eqI{165} & $-\tfrac{4}{D-4}+\tfrac{58}{3}+2\gamma_{\rm E}$ & $-\tfrac{4}{D-4}+\tfrac{52}{3}+2\gamma_{\rm E}$ &
pole unchanged; finite part absorbed into $\Delta\alpha_3$\\
B5 & $\alpha_2$, \eqII{8} & $\tfrac{61}{3}+\Delta\alpha_2$ & $\tfrac{49}{3}+\Delta\alpha_2$ & enters no result\\
B6 & $\alpha_3$, \eqII{9} & $\tfrac{4}{D-4}-\tfrac{58}{3}-2\gamma_{\rm E}+\Delta\alpha_3$ & $\tfrac{4}{D-4}-\tfrac{52}{3}-2\gamma_{\rm E}+\Delta\alpha_3$ &
enters no result (the local source $-\Delta\alpha_3H^2k^2u/a^2$ depends only on the arbitrary $\Delta\alpha_3$ and falls off)\\
\end{tabular}
\end{ruledtabular}
\end{table}

\subsection{Errata}

Our Table~\ref{tab:s0s1s7s8:errata} lists the 23 typographical errors C1--C23. None of them affects any result: in each case the
printed calculation used, or is consistent with, the correct expression, and the two terms missing from \eqI{128} (C14) are excluded by
the shift symmetry. The labels C1--C23, S1--S4 and B1--B6 are those of the summary of findings in the public
repository~\cite{ref:repo}. Three of the typographical errors (C2, C20 and C21) were first noticed while the English derivation notes
were prepared. Two more, the missing factors $H^2$ in \eqI{170} (C22, the same omission as in C17, C18 and C20) and the slip
``$P_1$ and $P_6$'' for ``$P_1$ and $P_4$'' before \eqI{A31} (C23), were noticed while the write-up was refereed and were confirmed
on the printed pages of Paper~I; they are not discussed in the derivation-note chapters. The list of misprints is therefore not
claimed to be exhaustive.

The derivation notes also record notational remarks, which are not counted as errors. In \eqI{6} the indices of $h^\nu{}_\rho$ are
not staggered, which is unambiguous because $h_{\mu\nu}$ is symmetric. \eqI{71} requires $\Delta x_\sigma=\eta_{\sigma\nu}\Delta
x^\nu$, so that $\Delta x_0=-\Delta\eta$. The statement at \eqI{69} that no delta function arises holds for naive (pointwise)
derivatives; the distributional derivative adds exactly the local term counted separately in \eqI{54}--\eqI{58}, so the inner
derivatives must be naive, as the footnote to \eqI{75} says. The $\square$ in \eqI{92} is distributional, while \eqI{93} is pointwise.
The dilatation factor in \eqI{134} must be positive, and in \eqI{135} $x\cdot x$ is the Minkowski product $\|\vec x\|^2-\eta^2$.
Finally, \eqI{125} and the polynomial and $\ln2$ terms of Table~V depend on the choice of integration constants, which the external
operators annihilate or which cancel globally, except for S3.

\begin{table*}[tbp]
\caption{Typographical errors in Papers~I and~II\@. None of them affects any result. Paper~I's Table~X and Paper~II's Table~1 are in the
paper convention.}
\label{tab:s0s1s7s8:errata}
\renewcommand{\arraystretch}{1.25}
\footnotesize
\begin{ruledtabular}
\begin{tabular}{@{}l>{\raggedright\arraybackslash}p{0.19\linewidth}>{\raggedright\arraybackslash}p{0.30\linewidth}%
>{\raggedright\arraybackslash}p{0.43\linewidth}@{}}
& Location & Printed & Correct\\
\hline
C1 & \eqI{1}, first term & $-\tfrac12\partial_\mu\phi\,\partial_\mu\phi\,g^{\mu\nu}\sqrt{-g}$ &
$-\tfrac12\partial_\mu\phi\,\partial_\nu\phi\,g^{\mu\nu}\sqrt{-g}$\\
C2 & \eqI{6}, left-hand side & $-\tfrac12\partial_\mu\phi\,\partial_\nu g^{\mu\nu}\sqrt{-g}$ &
$-\tfrac12\partial_\mu\phi\,\partial_\nu\phi\,g^{\mu\nu}\sqrt{-g}$\\
C3 & \eqI{7}, left-hand side & $-\tfrac12\kappa a^{D-2}\partial_\alpha\phi\,\partial_\beta\phi\,h^{\alpha\beta}$ &
$+\tfrac12\kappa a^{D-2}\partial_\alpha\phi\,\partial_\beta\phi\,h^{\alpha\beta}$; the vertex $V_1=i\kappa a^{D-2}\partial_1^\alpha
\partial_2^\beta$ is correct\\
C4 & \eqI{14}, first coefficient of $A(y)$ & $\Gamma(\tfrac D2-1)/(\tfrac D2-1)$ &
$\Gamma(\tfrac D2)/(\tfrac D2-1)=\Gamma(\tfrac D2-1)$, as used in \eqI{52}\\
C5 & Table~II, $U_2$ & $-\tfrac i4\kappa^2a^{D-2}\eta^{\alpha\rho}\eta^{\beta\sigma}\partial_1\!\cdot\partial_2$ &
$+\tfrac i4\kappa^2a^{D-2}\eta^{\alpha\rho}\eta^{\beta\sigma}\partial_1\!\cdot\partial_2$, as used in \eqI{33}\\
C6 & \eqI{39}, finite part & $[2\delta^\rho_0\delta^\sigma_0+2\bar\eta_{\rho\sigma}]H^2/16\pi^2$ &
$[2\delta^\rho_0\delta^\sigma_0+2\bar\eta^{\rho\sigma}]H^2/16\pi^2$ (index placement)\\
C7 & text before \eqI{54} & ``the 4th term of (47)'' & ``the 2nd term of (47)''\\
C8 & \eqI{89}, second term & $-4(aa')^2(a+a')^2\nabla^2(yA')$ & $-4H^2(aa')^2(a+a')^2\nabla^2(yA')$, as in Table~III\\
C9 & Table~XII, header & $\beta=(aa')^{D-2}(a^2+a'^2)H^2\square$ & $\beta=(aa')^{D-1}(a^2+a'^2)H^2\square$, as in \eqI{105}\\
C10 & Appendix, \eqI{A18} & $-2(D-1)aa'H^2\,I[f'A]$ & $-2(D-1)aa'H^2\,I[f'A']$\\
C11 & Table~XIII, row $P_{7c}$ & repeats part of the total &
$-\tfrac12(D-1)(D-2)\,y\,I[A'\Delta C']-\tfrac12(D-1)^2(D-2)\,I^2[A'\Delta C']$\\
C12 & Table~XV, row $P_{4b}$ & $(D-1)\,y'(A'B')'$ & $(D-1)\,y^2(A'B')'$\\
C13 & Table~XVIII, total & $\tfrac14\,y\,I^2[A'^2\Delta C']$ & $\tfrac14\,y\,I^2[A'\Delta C']$\\
C14 & \eqI{128}, list of four-derivative invariants & five terms &
the reduced list also contains $\phi^2R_{\mu\nu\rho\sigma}R^{\mu\nu\rho\sigma}$ and $\phi^2\square R$; both are excluded by the shift
symmetry, so \eqI{129} is unaffected\\
C15 & Table~IV, $\gamma_2$ row, second column & $-\frac{D^3-3D^2-4D+8}{4(D-1)}$ &
$+\frac{D^3-3D^2-4D+8}{4(D-1)}$, as used in Tables~VI--IX and required by shift symmetry\\
C16 & Table~VII, header & coefficient of $i\kappa^2H^4/(4\pi)^4$ & coefficient of $\kappa^2H^4/(4\pi)^4$\\
C17 & \eqI{160} & $\beta_2\square+\beta_3\nabla^2/a^2$ & $\beta_2H^2\square+\beta_3H^2\nabla^2/a^2$\\
C18 & \eqI{167} & $\Delta\alpha_2\square+\Delta\alpha_3\nabla^2/a^2$ & $\Delta\alpha_2H^2\square+\Delta\alpha_3H^2\nabla^2/a^2$\\
C19 & Table~X & row $(aa')^2(a^2+a'^2)\nabla^2\square$ printed twice & printed once, as in Table~1 of Paper~II\\
C20 & \eqII{10} & $\Delta\alpha_2\square+\Delta\alpha_3\nabla^2/a^2$ & $\Delta\alpha_2H^2\square+\Delta\alpha_3H^2\nabla^2/a^2$\\
C21 & \eqII{34} & $\Delta\alpha_3$ term without $H^2$ & with $H^2$\\
C22 & \eqI{170} & $\Delta\alpha_2\square+\Delta\alpha_3\nabla^2/a^2$; result $-\kappa^2a^4\Delta\alpha_3(k^2/a^2)\Phi_0$ &
$\Delta\alpha_2H^2\square+\Delta\alpha_3H^2\nabla^2/a^2$; result $-\kappa^2a^4\Delta\alpha_3H^2(k^2/a^2)\Phi_0$\\
C23 & Appendix, text before \eqI{A31} & ``Applying the second identity (A18) to $P_1$ and $P_6$'' &
``\ldots\ to $P_1$ and $P_4$'' (the subparts that follow are $P_{1a},P_{1b},P_{4a},\dots$)\\
\end{tabular}
\end{ruledtabular}
\end{table*}

\subsection{Conclusions}

We have recomputed, step by step and independently, the one-loop quantum-gravitational self-mass-squared of a massless, minimally
coupled scalar on de Sitter background in the gauge of Paper~I, and the late-time analysis of the effective mode equation of Paper~II.
The structure of the calculation is confirmed: the Feynman rules, the coincidence limits (apart from the finite part of \eqI{40}), the
decomposition of the 3-point
contribution, the extraction of derivatives (audited through worked examples and a global pointwise test), the separation of
divergences, the counterterm analysis and the integral tables of
Paper~II\@. The differences we found fall into three classes.

First, an overall sign: the printed self-mass-squared is $-1$ times the one defined by the effective field
equation \eqI{168} that both papers use. The error arises separately in the 4-point and in the 3-point part, so the two parts have the
correct relative sign and the local cancellations between them hold. The sign of the local part is irrelevant, because the finite
counterterms are arbitrary, but the nonlocal part has no counterterm freedom, so its sign is physical. Second, two corrections to the
final table of nonlocal terms: a missing non-analytic contribution in the $(aa')^2\nabla^4$ row, which changes the physical nonlocal
self-mass, and an analytic constant in the $(aa')^3(a^2+a'^2)H^4$ row, which cannot affect any retarded quantity. Third, finite errors
in two local terms, which shift the finite parts of $\beta_2$ and $\beta_3$ and are absorbed into the arbitrary finite counterterms.
Within Paper~I's treatment of the outer operators of the normal terms, the divergence structure, a single de Sitter-breaking pole in
$\beta_3$, is unchanged (see point~2 of Sec.~\ref{sec:s0s1s7s8:limitations}). In the labels of our Table~\ref{tab:s0s1s7s8:main}, the
first two classes are the substantive findings S1 and S2--S3, and the third consists of the absorbable errors B1--B6; the fourth
substantive finding, S4, concerns the argument for \eqII{44} and is discussed below. In addition there are 23 typographical errors, C1--C23.

None of these corrections changes the late-time physics, within the limitations listed below. With the corrected self-mass-squared in
the definition sign, and without the simplification \eqII{44}, the non-falling part of the source of the one-loop mode equation cancels
separately at each order in $k$ through $k^{10}$; this is verified exactly through $k^{10}$, not proven to all orders. The overall sign
reverses the particular solution $\Phi_1$, but since the non-falling total vanishes only decaying terms change sign. The correction of the $(aa')^2\nabla^4$ row contributes only decaying terms; it removes the $\hat k^4/a$, $\hat k^6/a,\dots$
tails that the printed row would leave, so that the source falls like $\ln a/a^2$. We record this faster fall-off as an observation, not as an
independent confirmation of the correction. The mode function approaches its tree-order late-time value times a constant,
Eq.~\eqref{eq:s0s1s7s8:late}, the constant being a field-strength renormalization. Through $k^{10}$, and given the kernel and local
terms as derived, there is no term that grows like $a$, $\ln a$ or $\ln^2a$. One-loop quantum gravity therefore gives no significant
correction to the mode functions of a massless, minimally coupled
scalar at late times, as concluded in Paper~II\@. The expectation expressed in Paper~II, that the same holds for the inflaton of
scalar-driven inflation, is not affected by our corrections.

The analysis beyond \eqII{44} also clarifies how the cancellation works. The argument given for \eqII{44} does not hold row by row:
the rows $(aa')^2(a^2+a'^2)H^2\nabla^2$ and $(aa')^2(a^2+a'^2)\nabla^2\square$ each grow like $a$ and cancel only against each other,
and at order $k^4$ the $(aa')^2\nabla^4$ row is needed for the cancellation. Its contribution at that order comes from the term
$24x^2\ln x$, which is the same in the printed and in the corrected row; the correction $8\ln x-16x\ln x$ changes only falling
terms. A heuristic explanation of the cancellation at $k=0$ is
that the scalar couples to gravity only through its derivatives, so that the self-mass-squared has the form $\partial_\mu\partial'_\nu
K^{\mu\nu}(x;x')$ and annihilates a constant mode up to initial-surface terms that decay. This is consistent with the observation of
Paper~II that the tree-order mode function approaches a constant whose derivatives vanish.

Finally, a remark on method. The recomputation was carried out with an AI model as the working tool, under the discipline described in
Sec.~\ref{sec:s0s1s7s8:method}: apart from the exceptions stated there, every step was derived from its inputs without using the
target, and every analytic result was checked by independent code where possible. Every tabulated or numerical result is backed by a
script in the public repository~\cite{ref:repo} that reproduces the logged output (Appendix~\ref{app:s0s1s7s8:verification}); the few
steps done by hand only are listed in Sec.~\ref{sec:s0s1s7s8:limitations} (points~1 and~5). The global pointwise tests proved particularly effective: they test
sums of many separately derived pieces against a reference that does not depend on the intermediate manipulations. The test of
Tables~IV--V located the $72\gamma_2$ inconsistency, and the test of Table~X showed that both corrections of Table~X are needed and
that together they are sufficient at the sample points tested.

\subsection{Limitations and open points}
\label{sec:s0s1s7s8:limitations}

We list explicitly what has not been established.
\begin{enumerate}
\item \emph{The appendix of Paper~I was audited, not redone.} For the appendix ``Extracting derivatives,'' worked examples (for
instance the $P_{3c}$ row of the $\gamma_3$ operator and the totals of the $\beta$ and $\gamma_1$ tables) were checked by hand, and the
totals of all ten tables were verified by the global pointwise test, with generic power functions at $D=4$ and $D=5$ and with the
actual $D=4$ propagators. The tables were not rederived row by row from scratch. Tables~IV and~V were computed from the printed
appendix totals after these had been checked, and the polynomial terms of Table~V enter Table~X as printed. In the same spirit, the
gauge-fixed graviton propagator \eqI{15}--\eqI{25} was taken from Paper~I and Ref.~\cite{ref:Tsamis1994} as input and was not
re-derived.
\item \emph{A possible additional local term is only partially resolved.} If the outer d'Alembertian in the operators $\alpha$ and
$\delta$ acts as a distribution, the local term \eqI{54}$+$\eqI{58} could be counted twice. We showed that any such double counting
has exactly the counterterm form $c(D)\,(a^D\square-a^{D-2}\nabla^2)\delta^D(x-x')$, with a constant $c(D)$ that may contain a pole and
without $\ln a$; it is absorbed into $\alpha_2$ and $\alpha_3$. The finiteness of $\beta_2$ and the pole of $\beta_3$ are therefore
confirmed within Paper~I's treatment of the outer operators; a double counting, if present, would shift poles and finite parts by
counterterm-form amounts. We did not show that the generic reductions \eqI{A35}--\eqI{A38} of Paper~I produce no other local difference
at coincidence, for example a finite $H^4a^D\delta^D(x-x')$; this would require a full distributional analysis. The vanishing
$a^D\square^2$ total of Table~VIII and the finiteness of Table~IX give partial support. Such a mass-type term would give a source
$\mathcal R_1$ that tends to a constant and, by \eqII{31}, a secular $\Phi_1\propto\frac13\ln a$. Heuristically it would conflict with
the shift symmetry that underlies the $k=0$ cancellation, but this has not been demonstrated.
\item \emph{The global checks are numerical.} The tests of the appendix totals, of Tables~IV--V and of Table~X compare sums at two or
three random spacelike points in 30-digit arithmetic. They are not symbolic identities, and they test only sums. Some further checks
were done in specific dimensions: the gauge-breaking Euler--Lagrange check in $D=3$; the completeness of the noninvariant counterterms
among fifteen candidate densities at $D=4,5,6$, with general-$D$ formulas derived for eleven of them; and the pointwise check of the
nonlocal part of \eqI{60}--\eqI{67} at $D=4$.
\item \emph{The location of the $72\gamma_2$ correction is an interpretation.} Reading it as a term $-12x$ dropped from the $\delta$
row of Table~V, or as a change of the constant of the $\gamma_2$ row, gives the same sum; the tests fix only the sum.
\item \emph{The late-time cancellation is exact through $k^{10}$, not proven to all orders.} The $k$ expansion converges at fixed $a$
(an entire integrand on a compact region), but this gives no control over the behavior of the terms beyond $k^{10}$ as $a\to\infty$:
individual rows carry terms growing like $a$ at every even order computed (for example $\pm\frac{8a}{441}\hat k^6$ in the rows
$(aa')^2(a^2+a'^2)H^2\nabla^2$ and $(aa')^2(a^2+a'^2)\nabla^2\square$), and for $\hat k\gtrsim1$ (modes inside the horizon at $\eta_i$) the
truncated series need not approximate the source. The $a^{-2}$ coefficients of $\mathcal R_1$ and $\Phi_1$ depend on the initial time
and on $\hat k$, i.e.\ they are sensitive to the initial state. The explanation of the $k=0$ cancellation given above is a heuristic.
The operator ordering (scale factors to the left of the derivative operators of Table~X) is assumed; it was verified in the pointwise
test of Table~X\@. Steps done by hand only, without a code check, are the direct light-cone evaluation of the $1/x$ integrals (the code
uses \eqII{35} and agrees with it), the analytic bound \eqref{eq:s6:bound} on the $(aa')^2\nabla^4$ correction, and the heuristic
argument for the $k=0$ cancellation.
\item \emph{Assumptions shared with Paper~I.} Dimensional regularization preserves the shift symmetry and the residual symmetries of
the gauge-fixed action; ultraviolet divergences are polynomial in $H$; coincidence limits are taken in dimensional regularization; and
$D=4$ is set only in terms not multiplied by a pole.
\item \emph{The overall sign rests on the definition \eqI{168}} with $-iM^2=\delta^2(i\Gamma_1)/\delta\Phi\,\delta\Phi$. Several
independent arguments support it: a lemma on second functional derivatives of quadratic forms with derivatives (with two proofs), the
field-equation route, the comparison with the paper's counterterm vertex \eqI{131}, explicit test functions, a lattice Hessian, and
the consistency of the 3-point steps. Once \eqI{168} is adopted, the sign of the nonlocal part is physical.
\end{enumerate}
Of these, points~2 and~5 bear most directly on the late-time result: a finite local term of the type mentioned in point~2, if it
were present, would enter the mode equation directly and could produce a secular $\ln a$, and point~5 limits the proven cancellation to
the orders computed. The late-time conclusion is therefore established through $k^{10}$, given the kernel and the local terms as
derived. The nonlocal
kernel used in the late-time analysis was verified by the global pointwise test of Table~X against a sum that does not depend on the
outer-operator reduction of the appendix.

\begin{acknowledgments}
The author thanks R.~P.~Woodard for the original collaboration on Papers~I and~II\@. The author gratefully acknowledges Claude
Opus~5.5 (Anthropic), which carried out the step-by-step re-derivations and the symbolic and numerical cross-checks reported here,
and prepared the derivation notes and the public repository~\cite{ref:repo}. The complete computation, including the preparation of
this manuscript, used about \TOKOUT{} output tokens (\TOKTOT{} tokens processed in total, the large majority of them cached re-reads
of the working context) over approximately \HOURS{} hours of active computation spread over \NDAYS{} days (October 5, 15 and 16, 2026).
\end{acknowledgments}

\appendix

\section{Verification methods}
\label{app:s0s1s7s8:verification}

This appendix describes the checking tools that carry the most weight in the conclusions: the exact engine for the late-time
Schwinger--Keldysh integrals, the global pointwise tests of the self-mass tables, the momentum expansion beyond \eqII{44} with its
numerical tests, and the reproducibility of all scripts. Throughout, ``30 digits'' means \texttt{mpmath} arithmetic with a working
precision of 30 significant digits.

\subsection{Exact engine for the late-time integrals}
\label{app:s0s1s7s8:engine}

All late-time integrals of Paper~II, with or without the simplification \eqII{44}, are generalizations of the integrals
$S_K[f]=-\frac{i}{16\pi^2}\int d^4x'\,a'^K\{f(\frac{y_{++}}{4})-f(\frac{y_{+-}}{4})\}$ of Tables~2--4 of Paper~II ($H=1$), with an
additional shift of the power of $a'$ (from the expansion of $u(\eta',k)$) and an additional weight $r^{2m}$, $r\equiv\|\vec x\,'-\vec
x\|$ (from the expansion of the phase). Their reduction to single integrals over $a'$, including the direct light-cone evaluation of the
$1/x$ terms, is given in Sec.~\ref{sec:s6} (Eq.~\eqref{eq:s6:master}). We normalize $a(\eta_i)=1$ and write
$\Delta\eta\equiv\eta-\eta'=1/a'-1/a$ and $X_0\equiv\tfrac14aa'\Delta\eta^2$.

\paragraph*{Scaling variable.} With $a'=as$ one has $\Delta\eta=(1-s)/(as)$ and $X_0=(1-s)^2/(4s)$, which is independent of $a$, so
that every integral that occurs reduces to the form \eqref{eq:s6:scaling}, with $G=1$ or $G=\ln X_0$. After binomial expansion
only three elementary integrals remain, $\int_\epsilon^1 s^p\,ds$, $\int_\epsilon^1 s^p\ln s\,ds$ and
\begin{equation}
\int_\epsilon^1\!ds\,s^p\ln(1-s)=C_p-\sum_{n=1}^{q}\frac1n\int_\epsilon^1\!ds\,s^{p+n}+\sum_{n>q}\frac{\epsilon^{p+n+1}}{n(p+n+1)},
\qquad q\equiv\max(0,-p-1),\qquad \epsilon\equiv\frac1a ,
\label{eq:s0s1s7s8:elem}
\end{equation}
which follows from $\ln(1-s)=-\sum_ns^n/n$: the terms with $n\le q$ are kept as divergent or logarithmic integrals, and those with
$n>q$ are integrated term by term. The constants $C_p=-\sum_{n>q}1/[n(n+p+1)]$ are given in closed form in Eq.~\eqref{eq:s6:Cp}.
The result of every integral is thus an exact expansion in $\epsilon=1/a$ whose coefficients are polynomials in $\ln a$ with
rational coefficients and the constants $\ln2$ and $\pi^2$; the only source of $\pi^2$ is $C_{-1}=-\pi^2/6$. Worked examples are
given in Eq.~\eqref{eq:s6:examples}, and the next order of every entry of Tables~2--4 of Paper~II in our
Table~\ref{tab:s6:tab234}.

\paragraph*{Implementation.} The helper module \texttt{engine.py} implements Eqs.~\eqref{eq:s6:scaling}, \eqref{eq:s0s1s7s8:elem} and \eqref{eq:s6:Cp} in
exact rational arithmetic in \texttt{sympy}, including the shift of the power of $a'$ and the weight $r^{2m}$, the operator
$\Dop\equiv-(a^2\partial_a^2+4a\partial_a)$, which is $\square/H^2$ acting on functions of $\eta$ alone, and truncation in powers of
$a$. A first attempt, which used generic symbolic integration followed by an expansion in $1/a$, failed because the resulting
polylogarithms could not be expanded asymptotically; it was replaced by this closed-form engine. With it all 12 entries of Tables~2--4
and all 19 entries of Table~5 of Paper~II are reproduced exactly (in the paper convention), and the next order of Tables~2--4 is
obtained. Paper~II's Tables~2--5 were never used as input.

\subsection{Global pointwise tests of the self-mass tables}
\label{app:s0s1s7s8:global}

The final tables of Paper~I are obtained after long chains of manipulations: extraction of derivatives, reduction to external
operators, and separation of divergences. A global pointwise test compares the end product directly with a reference that does not
depend on these manipulations. All quantities are in the paper convention, at $D=4$ and for $x\neq x'$.

\paragraph*{Reference side.} The normal part of the 3-point contribution is the sum of the seven parts $P_1,\dots,P_7$ defined in
Eqs.~\eqI{A1}--\eqI{A7} of the appendix of Paper~I, each of the generic form $\partial\partial'[(aa')^{D-2}F\,\partial\partial'A]$ with
$F\in\{A,B,\Delta C\}$. They are evaluated by brute-force differentiation in explicit coordinates, with naive (pointwise) inner
derivatives. For the test of Table~X the nonlocal logarithm term of \eqI{67},
$-\frac{\kappa^2H^2}{8\pi^2}\,2H(aa')^2(a\partial_0+a'\partial'_0)\nabla^2A$, is added. The actual $D=4$ propagator functions are
used: $A=h\{4/y-2\ln(y/4)-1\}$, $B=C=4h/y$ and $\Delta C=4(C-A)$, with $h\equiv H^2/16\pi^2$.

\paragraph*{Table side.} The external operators of the table act pointwise on their coefficient functions, with the scale factors
standing to the left of the derivative operators.

\paragraph*{Appendix totals and Tables~IV--V.} The ten external operators acting on the totals of the appendix Tables~XI--XX reproduce
the direct sum $\sum_kP_k$ exactly (to 30 digits), once the five appendix misprints C9--C13 of our Table~\ref{tab:s0s1s7s8:errata} are
corrected. This was tested with the generic powers $A=y^{-7/3}$, $B=y^{-5/4}$, $\Delta C=y^{2/7}$ at $D=5$ and at $D=4$ (two random
spacelike points each), which checks the reduction algebra independently of the propagators, and with the actual $D=4$ propagators
(three random points). Evaluated at $D=4$, the printed Tables~IV and~V miss the direct sum by residuals of order $10^{-2}$ to $10$; the residual vanishes
to 30 digits only when the sign of the $\gamma_2$ entry of Table~IV is corrected (C15) and a term equivalent to $72\gamma_2$
acting on a constant is added (S3; our Table~\ref{tab:s4b:hyp}).

\paragraph*{Table~X.} Three random rational spacelike pairs of points (seed 29) were used. The two terms of \eqI{159} cancel
pointwise, so \eqI{159} does not enter the test, and at $D=4$ the nonlocal parts of \eqI{152}--\eqI{153} reproduce $(4/y)^2$ and
$(4/y)^3$ pointwise, so that the rows of Table~X that come from Table~VII can be compared pointwise with the reference side. The
residuals are listed in our Table~\ref{tab:s5:global}: for the printed Table~X and for Table~X with the $(aa')^2\nabla^4$ correction
alone they range from about $10^{-2}$ to $10^2$ (with $|{\rm LHS}|$ between $0.1$ and $9$), and they vanish to 30 digits ($|{\rm residual}|\lesssim4\times10^{-31}$) only with both the
$(aa')^2\nabla^4$ (S2) and the $72\gamma_2$ ($-54\to+18$, S3) corrections.

The tests are numerical, at sample points, and they test sums only (Sec.~\ref{sec:s0s1s7s8:limitations}). Their strength is that the
reference side is independent of the outer-operator reduction of the appendix and of the separation of divergences, and that random
points make an accidental agreement to 30 digits implausible.

\subsection{Momentum expansion and quadrature checks}
\label{app:s0s1s7s8:kexp}

Beyond \eqII{44} the full $k$ dependence of the source $\mathcal R_1$ of Eq.~\eqref{eq:s0s1s7s8:mode} is kept: the exact
$u(\eta',k)/u(0,k)=(1-iw)e^{iw}$ with $w\equiv\hat k/a'$, the angular average $\sin kr/kr$ of the phase, the $k^2/a^2$ part of
$\square$, and all rows containing $\nabla^2$ (Sec.~\ref{sec:s6}, Eq.~\eqref{eq:s6:kser}). Each term of the
expansion in $k$ is an integral of the type \eqref{eq:s6:scaling}. The integrand is an entire function of $k$ and the
integration region is compact ($\eta'\in[\eta_i,\eta]$, $r\le\Delta\eta$), so the expansion converges for every $k$ at fixed $a$. This
convergence gives no control over the late-time behavior of the terms beyond the order computed, since individual rows contain terms
growing like $a$ at every even order computed (Sec.~\ref{sec:s6:beyond}). The
script \texttt{kexp.py} expands every row of the corrected Table~X exactly through $k^N$ and prints the part of the source that does
not fall off as $a\to\infty$. For $N=2,4,6,10$ this non-falling part vanishes at every order $k^0,\dots,k^N$. The same holds for the
printed table (scripts \texttt{analysis2.py} and \texttt{extract.py}), whose difference from the corrected one consists of falling
terms only. Hand computations reproduce selected entries: the growing terms $\mp\tfrac{4}{15}\hat k^2a$ of the rows
$(aa')^2(a^2+a'^2)H^2\nabla^2$ and $(aa')^2(a^2+a'^2)\nabla^2\square$, the order-$\hat k^2$ entry of the row $(aa')^4H^4$, and the
order-$\hat k^4$ entry of the $(aa')^2\nabla^4$ row, which comes from its $24x^2\ln x$ term (the same in the printed and in the
corrected row). For the correction of the
$(aa')^2\nabla^4$ row an analytic bound valid for every $k$ shows that its contribution to $\mathcal R_1$ falls at least like
$\hat k^4/a$, independently of the expansion.

Two numerical tests compare the engine and the expansion with direct quadrature in \texttt{mpmath}.
\begin{itemize}
\item \texttt{numcheck.py}: the exact expansion through $a^{-6}$ versus direct quadrature over $\eta'$ and $z\equiv r/\Delta\eta$, at
$a=50$ and $a=400$, for seven integrals, including $\ln^2x$, $x\ln x$, $x^2\ln x$, an $r^2$ weight and a power of $a'$ shifted as by
the $u(\eta',k)$ correction. The largest difference is $1.0\times10^{-19}$, on values between $10^{-2}$ and $10^2$.
\item \texttt{numcheck\_k.py}: the series truncated at $k^8$ versus quadrature with the exact $u(\eta',k)$ and $\sin kr/kr$, at $a=200$.
The differences are between $10^{-10}$ and $2\times10^{-8}$ at $\hat k=0.3$, and between $5\times10^{-6}$ and $9.3\times10^{-4}$ at
$\hat k=1$, consistent with the expected truncation error of order $\hat k^9$.
\end{itemize}

\subsection{Reproducibility}
\label{app:s0s1s7s8:repro}

The repository~\cite{ref:repo} contains the English derivation notes (nine chapters, \texttt{notes/chapters/}, and a summary chapter
listing every finding with the labels S1--S4, B1--B6 and C1--C23) and, for each chapter,
a folder \texttt{code/NN\_\ldots/} with standalone Python scripts (Python~3.10 or later, \texttt{sympy}, \texttt{mpmath}),
a README describing what each script checks, and the expected output of every script. Each script runs from its own folder with
\texttt{python3 <script>.py}, without network access, in at most a few minutes. All scripts were re-run for the public version and
reproduce the outputs recorded in the original working logs; the only difference is the order in which one script prints two tied
entries, now fixed by a sort key. Our Table~\ref{tab:s0s1s7s8:code} lists the folders and the main checks they contain.

\begin{table}[tb]
\caption{Code folders of the repository~\cite{ref:repo} and the main checks they contain.}
\label{tab:s0s1s7s8:code}
\renewcommand{\arraystretch}{1.25}
\small
\begin{ruledtabular}
\begin{tabular}{@{}>{\raggedright\arraybackslash}p{0.26\linewidth}>{\raggedright\arraybackslash}p{0.70\linewidth}@{}}
Folder & Main checks\\
\hline
\texttt{01\_feynman\_rules} & expansion \eqI{6} against the exact metric; vertices of Tables~I--II\\
\texttt{02\_four\_point} & derivative rule and overall sign of the 4-point part (test functions, lattice Hessian, calibration with
\eqI{131}); coincidence limits \eqI{34}--\eqI{40}; \eqI{41}--\eqI{44}\\
\texttt{03\_three\_point\_setup} & 3-point rule \eqI{45}--\eqI{46}; the nine terms \eqI{47}--\eqI{49} and their sign; \eqI{52};
first coefficient of \eqI{14}\\
\texttt{04\_three\_point\_local\_log} & local terms \eqI{54}--\eqI{59}; logarithm terms \eqI{60}--\eqI{67} and their cancellation
against \eqI{44} in general $D$\\
\texttt{05\_normal\_terms\_tableIII} & derivative-extraction identities \eqI{68}--\eqI{80}; \eqI{81}--\eqI{93} and Table~III\\
\texttt{06\_tables\_IV\_V} & \eqI{94}--\eqI{125}; global test of the appendix totals; Table~IV in general $D$; Table~V and the
$72\gamma_2$ hypothesis tests\\
\texttt{07\_desitter\_counterterms} & de Sitter isometries and the residual symmetries of the gauge; \eqI{136}--\eqI{145} and the
completeness of the noninvariant counterterms\\
\texttt{08\_renormalization\_tableX} & \eqI{146}--\eqI{159} and Tables~VI--IX; $\beta_i$ of \eqI{161}--\eqI{165}; assembly of Table~X;
global pointwise test of Table~X\\
\texttt{09\_schwinger\_keldysh} & exact late-time engine; Tables~2--5 of Paper~II; momentum expansion through $k^{10}$; quadrature
checks; late-time coefficients of the source\\
\end{tabular}
\end{ruledtabular}
\end{table}

\begin{thebibliography}{99}

\bibitem{ref:Woodard2004}
R. P. Woodard, ``Quantum Effects during Inflation,'' in Quantum Field Theory under the Influence of External Conditions
(Rinton Press, Princeton, 2004) ed. K. A. Milton, pp. 325--330, astro-ph/0310757.

\bibitem{ref:Mukhanov1981}
V. F. Mukhanov and G. V. Chibisov, JETP Lett. 33, 532 (1981).

\bibitem{ref:Starobinsky1979}
A. A. Starobinski, JETP Lett. 30, 682 (1979).

\bibitem{ref:Weinberg2005}
S. Weinberg, Phys. Rev. D 72, 043514 (2005); 74, 023508 (2006).

\bibitem{ref:Miao2006}
S. P. Miao and R. P. Woodard, Classical Quantum Gravity 23, 1721 (2006); Phys. Rev. D 74, 024021 (2006).

\bibitem{ref:Mukhanov1992}
V. Mukhanov, H. Feldman, and R. Brandenberger, Phys. Rep. 215, 203 (1992).

\bibitem{ref:Liddle1993}
A. R. Liddle and D. H. Lyth, Phys. Rep. 231, 1 (1993).

\bibitem{KW07}
E. O. Kahya and R. P. Woodard, ``Quantum gravity corrections to the one loop scalar self-mass during inflation,''
Phys. Rev. D 76, 124005 (2007), arXiv:0709.0536.

\bibitem{ref:Tsamis1994}
N. C. Tsamis and R. P. Woodard, Commun. Math. Phys. 162, 217 (1994).

\bibitem{ref:Woodard2006}
R. P. Woodard, in Deserfest: A Celebration of the Life and Works of Stanley Deser, edited by J. T. Liu, M. J. Duff,
K. S. Stelle, and R. P. Woodard (World Scientific, Hackensack, 2006), pp. 339--351.

\bibitem{ref:Bogoliubov1957}
N. N. Bogoliubov and O. Parasiuk, Acta Math. 97, 227 (1957); K. Hepp, Commun. Math. Phys. 2, 301 (1966); W. Zimmermann,
Commun. Math. Phys. 11, 1 (1968); 15, 208 (1969); Lectures on Elementary Particles and Quantum Field Theory, edited by
S. Deser, M. Grisaru, and H. Pendleton (MIT Press, Cambridge, 1971), Vol. I.

\bibitem{KW08}
E. O. Kahya and R. P. Woodard, ``Scalar Field Equations from Quantum Gravity during Inflation,'' Phys. Rev. D 77, 084012 (2008),
arXiv:0710.5282.

\bibitem{ref:Schwinger1961}
J. Schwinger, J. Math. Phys. 2, 407 (1961).

\bibitem{ref:Keldysh1965}
L. V. Keldysh, Sov. Phys. JETP 20, 1018 (1965).

\bibitem{ref:Jordan1986}
R. D. Jordan, Phys. Rev. D 33, 444 (1986).

\bibitem{ref:Calzetta1987}
E. Calzetta and B. L. Hu, Phys. Rev. D 35, 495 (1987).

\bibitem{ref:repo}
E. O. Kahya, \emph{claude-audits-quantum-gravity}: derivation notes and verification code,
\url{https://github.com/eokahya/claude-audits-quantum-gravity}.

\bibitem{ref:Allen1987}
B. Allen and A. Folacci, Phys. Rev. D 35, 3771 (1987).

\bibitem{ref:Allen1985}
B. Allen, Phys. Rev. D 32, 3136 (1985).

\bibitem{ref:Onemli2002}
V. K. Onemli and R. P. Woodard, Classical Quantum Gravity 19, 4607 (2002).

\bibitem{ref:Onemli2004}
V. K. Onemli and R. P. Woodard, Phys. Rev. D 70, 107301 (2004).

\bibitem{ref:Birrell1982}
N. D. Birrell and P. C. W. Davies, Quantum Fields in Curved Space (Cambridge University Press, Cambridge, England, 1982).

\bibitem{ref:Candelas1975}
P. Candelas and D. J. Raine, Phys. Rev. D 12, 965 (1975).

\bibitem{ref:Dowker1976}
J. S. Dowker and R. Critchley, Phys. Rev. D 13, 3224 (1976).

\bibitem{ref:Tsamis2006}
N. C. Tsamis and R. P. Woodard, Ann. Phys. (N.Y.) 321, 875 (2006).

\bibitem{ref:Prokopec2007}
T. Prokopec, N. C. Tsamis, and R. P. Woodard, Classical Quantum Gravity 24, 201 (2007).

\end{thebibliography}

\end{document}
